\documentclass[11pt,a4paper]{article}
\usepackage[margin=27mm]{geometry}
\usepackage[T1]{fontenc}
\usepackage[utf8]{inputenc}
\usepackage{lmodern,amsmath,amssymb,amsthm,mathtools,booktabs,longtable,array}
\providecommand{\mathscr}[1]{\mathcal{#1}}
\usepackage[dvipsnames]{xcolor}
\usepackage{microtype}
\PassOptionsToPackage{hyphens}{url}
\usepackage[hidelinks]{hyperref}
\Urlmuskip=0mu plus 1mu\relax
\usepackage{enumitem}
\setlist{itemsep=3pt,topsep=5pt}
\setlength{\parskip}{4pt}
\setlength{\parindent}{0pt}
\setlength{\emergencystretch}{3em}
\newtheorem{theorem}{Theorem}[section]
\newtheorem{lemma}[theorem]{Lemma}
\newtheorem{proposition}[theorem]{Proposition}
\newtheorem{corollary}[theorem]{Corollary}
\newtheorem{definition}[theorem]{Definition}
\DeclareMathOperator{\CH}{CH}
\DeclareMathOperator{\Pic}{Pic}
\DeclareMathOperator{\NS}{NS}
\DeclareMathOperator{\End}{End}
\DeclareMathOperator{\Hom}{Hom}
\DeclareMathOperator{\Alb}{Alb}
\providecommand{\tightlist}{\setlength{\itemsep}{0pt}\setlength{\parskip}{0pt}}
\title{JAC1964: primitive incidence polarizations,\newline pencil heights, and the arithmetic complement}
\author{Run JAC1964--20261121\\GPT-6 Astra Pro; collaborative mathematical investigation}
\date{21 November 2026\\Working manuscript: continuously checked findings}
\begin{document}
\maketitle
\begin{abstract}
Accepting the supplied 032 results, we prove that the correctly signed
primitive incidence self-duality of an odd-dimensional abelian variety is
a rational polarization in every characteristic. An explicit quadric
pencil refutes the proposed uncorrected H9 height identity, even when
all fibres are geometrically integral and the restricted cycles are
rationally zero. We give the admissible local correction and a precise
no-correction criterion. Algebraic Lefschetz correspondences prove the
restriction and fixed-part assertions of the sixth theorem. An
all-primes monodromy argument also determines the variable incidence
sign under explicit Lefschetz monodromy hypotheses. The arithmetic
investigation constructs an actual orthogonal Chow Picard summand for
CM varieties and isolates its complementary height problem. For a
nonzero complementary Fermat-quartic Fourier--Ceresa class, an exact
finite Hodge-lattice calculation and a reproducible theta certificate
prove that its normalized height is greater than $1/100$.
Its central CM orbit is a positive rational plane. Algebraically
trivial cycles in the complementary motive lie in the height radical,
reducing the remaining semipositivity problem to a Griffiths quotient.
The full geometric height on an abelian fourfold pencil is semipositive;
the same holds in higher dimension for geometric-generic axes of the
specified degree. We also prove the primitive polarization on actual
codimension-two algebraic representatives in the stated cases.
The unrestricted seventh theorem, with the customary meaning of
$\tau$-equivalence, is false even for CM abelian threefolds.
Over finite fields and their algebraic closures, the accepted Tate
consequence and Kahn's theorem instead prove rational Chow equal to
numerical equivalence for all varieties of geometric abelian type.
The statement-level audit is in \texttt{CHECK\_EVENTS.tsv}.
Every assertion below distinguishes its hypotheses, check status, and
remaining gate. No general arithmetic Hodge-index theorem is claimed.
\end{abstract}
\tableofcontents
\newpage
\section{Scope, inputs, and conventions}
All new deductions are labelled by stable finding IDs. ``Proved'' means
proved from the explicitly accepted and cited inputs, with the
whole-proof reading recorded in the check log; it does not mean formal
verification. The model effort setting is not exposed in this session.
The original archive images were not supplied, so the transcription
audit is mathematical, not palaeographic.

The incidence interpretation of the higher Picard variety is specified
below. We work with rational cycles and rational homomorphisms of
abelian varieties, clearing denominators on actual families. We do not
tensor geometric point groups with $\mathbf Q$ to define varieties.
No integral universal cycle or principal polarization is asserted.

The attached 032, Theorem 1.1 and Corollaries 8.3--8.4, the supplied
van Hoften--Sempliner and Milne reports, and Alp\"oge's advertised
Riemann-hypothesis and Dirichlet-GRH statements are accepted as inputs,
as requested. Their full proofs and the omitted Lean dependency tree
were not independently certified here. The bridge sketch and H9 are
questions to test, not input theorems. The exact scopes of the analytic
inputs are recorded in the arithmetic section.

\part{Checked findings and cited background}
% JAC1964-20261121. /root/programme_audit; GPT-6 Astra Pro.
% Overview only; root owns integration and the master proof text.

\section{The seven-theorem programme: interpretations and scope}
\label{sec:programme-overview}

The operative programme in cote~44 is on scans~2, 4 and~6.
The letter to Serre dated 12~August~1964 occupies scans~8--9.
Scan~1 is a modern cover; scans~3, 5 and~7 concern graded
completion or root-system material and are not further claims in
this programme. The supplied sources are TeX transcriptions;
the manuscript images were unavailable. The following is a
mathematical audit, not a palaeographic verification. Uncertain
readings and deleted clauses remain part of the source record.

Throughout the table, $P^p(X)$ means Saito's higher Picard variety
of algebraically trivial cycles modulo incidence equivalence.
This is the explicit interpretation under which the programme is
tested. A universal regular algebraic representative is identified
with it only in the cases where the comparison has been proved
separately. Polarizations and duality statements are rational;
no integral universal cycle or principal polarization is asserted.

\begingroup
\small
\setlength{\tabcolsep}{4pt}
\renewcommand{\arraystretch}{1.15}
\begin{longtable}{@{}p{0.08\linewidth}p{0.19\linewidth}p{0.68\linewidth}@{}}
\toprule
Claim & Exact archive locus & Interpretation and result\\
\midrule
\endfirsthead
\toprule
Claim & Exact archive locus & Interpretation and result\\
\midrule
\endhead
Th.~1 & Cote~44, scan~2; clarified by scan~8 &
The family map from $\operatorname{Alb}(S)$ has the incidence
kernel, as a statement on geometric points. Saito's construction
supplies this reading; the uncertain archive notation is not
silently replaced by universal regular equivalence.\\

Th.~2 & Cote~44, scan~2; scan~8, Betti bound &
$\dim P^p(X)\le b_{2p-1}(X)/2$. This is established background for
smooth projective varieties over algebraically closed fields in
every characteristic: Saito, Proposition~3.4, pp.~109--110.\\

Th.~3 & Cote~44, scan~2; scan~8, complementary indices &
The incidence quotient is an abelian variety and $P^p(X)$ is
canonically dual, up to isogeny, to $P^{n+1-p}(X)$.
The index $n+1-p$ is forced by the pairing and is explicit in
scan~8. Saito, Theorem~3.6 and Proposition~4.4, supplies the
all-characteristic construction and duality.\\

Th.~4 & Cote~44, scan~2; scan~8, primitive margin; scan~9, formula for $D$ &
The correct signed class on the primitive middle factor is
$(-1)^m\Delta^*d/2$ for $\dim X=2m-1$.
JAC-T4 proves ampleness for abelian $X$ in every characteristic
under 032, Corollary~8.4. The generic-section and complete-intersection
extensions are H9-T6a and H9-T13. The literal unsigned assertion
on the full factor fails already on $E^3$ (JAC-P03).\\

Th.~5 & Cote~44, scan~4, height formula and fibre hypothesis &
The source assumes algebraically trivial restrictions to the generic
section. Its uncorrected formula fails even for geometrically
integral quadric-pencil fibres (H9-C1). H9-T1 gives the corrected
intersection formula; H9-T9 identifies the signed height with a
positive form modulo trace on a polarized algebraically trivial
primitive incidence domain. The larger homological domain requires
the additional hypotheses stated below.\\

Th.~6 & Cote~44, scan~4, three degree ranges &
For $\dim X=n$, restriction is an isogeny when $2p\le n-1$,
its image is the fixed factor when $2p=n$, and it is surjective
when $2p\ge n+1$. JAC-P02 proves these rational assertions when
the normalized Lefschetz lowering operator of $X$ is algebraic,
in particular for abelian $X$ in every characteristic. It also
proves the finite-kernel assertion deleted in the critical clause;
the deletion is not erased from the source.\\

Th.~7 & Cote~44, scan~6 &
With customary $\tau$-equivalence, a nonzero multiple is
algebraically equivalent to zero. The assertion is false in
characteristic zero even for the CM Fermat-quartic Jacobian
(JAC-P05, using Bloch's non-torsion theorem).
It is true for every abelian variety over $\overline{\mathbf F}_p$,
indeed every geometric abelian-type variety, by JAC-P15 using
032, Corollary~8.3, and Kahn's theorem. Divisors in every
characteristic and zero-cycles on connected smooth projective
varieties are classical positive cases.\\
\bottomrule
\end{longtable}
\endgroup

\paragraph{The height domains are part of the assertion.}
Replacing the algebraic-triviality hypothesis of Th.~5 by
homological triviality changes the problem. JAC-P04 exhibits that
distinction on an abelian fourfold. For a genuine Lefschetz pencil
on an even-dimensional abelian variety, JAC-P07 proves condition~A
automatically in every characteristic. H9-T7 then computes the
height on the entire homologically trivial Chow group as an
ambient primitive intersection form plus a variable-axis form,
and identifies the remaining radical with the radical of the axis
form. This is positive in the surface-axis case
H9-T8 and for the geometric generic ordered axes of degree
$d\ge3$ in JAC-P09, using 032 for the ambient term. The latter
scope is the original pencil field with a geometric generic axis,
not every closed axis or every separable extension of that field.
H9-T12 proves the precise finite-extension statements: the whole
result persists under purely inseparable extensions; arbitrary
finite extensions have an orthogonal restriction--norm
decomposition on the admissible domain; the algebraically trivial
primitive domain remains positive modulo the unchanged fixed trace.

\paragraph{The primitive and representative qualifications are essential.}
The margin in cote~44, scan~8, explicitly requires the primitive
decomposition. Under $\phi_L(x)=t_x^*L\otimes L^{-1}$, the
elliptic normalized diagonal induces the negative of the ample
polarization. Thus the divisor-incidence sign is $(-1)^m$,
whereas the primitive odd cup-product sign is $(-1)^{m-1}$.
The factor $1/2$ is the diagonal extraction of a normalized
symmetric divisor correspondence. H9-T10 proves the comparison
with the actual regular representative $\operatorname{Ab}^2$
for the stated generic threefold sections; H9-T13 extends that
case to the indicated complete intersections in an abelian variety.
Neither result supplies a universal regular representative in
arbitrary codimension.

\paragraph{The finite-field conclusion uses the stronger accepted input.}
The Hodge-standard assertion of 032, Corollary~8.4, is not by
itself a rational-equivalence theorem. JAC-P15 uses its Tate
consequence, Corollary~8.3, printed pp.~49--50, together with
Kahn, Theorem~1.10, printed p.~982. It proves
$\operatorname{CH}^r_{\mathbf Q}=N^r_{\mathbf Q}$ over finite
fields and their algebraic closures for abelian-type varieties.
This is compatible with the characteristic-zero Ceresa
counterexample and does not assert rational=numerical over
arbitrary transcendental constant-field extensions.

\subsection*{Finding-to-archive atlas}

The IDs below name deductions or calculations made in this run.
The archive column identifies the claim they resolve or the
passage that motivates them; it does not attribute their modern
inputs or their proofs to Grothendieck. Saito's construction,
the SGA monodromy theorems, Bloch's non-torsion theorem, Kimura
nilpotence, and Kahn's Tate--Beilinson theorem remain cited
background.

\begingroup
\small
\setlength{\tabcolsep}{4pt}
\renewcommand{\arraystretch}{1.12}
\begin{longtable}{@{}p{0.12\linewidth}p{0.25\linewidth}p{0.58\linewidth}@{}}
\toprule
ID & Grothendieck locator & Finding and precise role\\
\midrule
\endfirsthead
\toprule
ID & Grothendieck locator & Finding and precise role\\
\midrule
\endhead
JAC-P00 & Cote~44, scans~2 and~8; Th.~6 on scan~4 &
Rational descent of the incidence quotient and its normalized
families, including imperfect pencil fields and inseparable
norm multiplicities.\\
JAC-P01 & Cote~44, scan~2, Ths.~1--3; scan~8, Tate-module comparison &
The relevant odd realization detects incidence homomorphisms;
cohomological algebraic projectors induce actual rational
abelian-variety projectors.\\
JAC-P02 & Cote~44, scan~4, Th.~6; cote~16, scans~40 and~49 &
Restriction, critical fixed trace, and surjectivity under algebraic
normalized lowering. The critical trace proof spreads an actual
cycle and uses the blowup formula.\\
JAC-P03 & Cote~44, scan~8, primitive margin; scan~9, $D$ &
Explicit families on $E^3$ give opposite signs. The calculation
fixes the divisor-incidence sign and shows that quotienting the
kernel alone does not produce positivity.\\
JAC-P04 & Cote~44, scan~4, algebraic-triviality hypothesis &
An explicit primitive cycle on $E^4$ has homologically trivial
but non-algebraically-trivial restriction in the stated high-degree
characteristic-zero setting.\\
JAC-P05 & Cote~44, scan~6, Th.~7 &
Bloch's infinite-order Fermat Ceresa class refutes customary
numerical=$\tau$ even for a CM abelian threefold in characteristic zero.\\
JAC-P06 & Cote~44, scan~6, undefined $\tau$ &
Over $\overline{\mathbf F}_p$, algebraically trivial Chow classes
are torsion, so the numerical--$\tau$ claim is a rational-Chow
comparison. JAC-P15 supplies that comparison in abelian type.\\
JAC-P07 & Cote~44, scan~4, pencil/fibre hypothesis in Th.~5 &
For an abelian $n$-fold, $n$ even, variable rank is
$h^n-\binom{2n}{n}/(n+1)>0$. This proves condition~A and
the required specialization injectivity for genuine Lefschetz
pencils, in every characteristic.\\
JAC-P08 & Cote~44, scan~4, Th.~6; cote~16, scans~40 and~49--50 &
Explicit algebraic lowering for smooth complete intersections,
and stability under products, extend the proved restriction/trace
statements beyond abelian ambient varieties.\\
JAC-P09 & Cote~44, scan~4, Th.~5; enlarged H9 domain &
Orthogonal monodromy and finite orbit of a geometric generic
cycle eliminate the variable algebraic axis classes for ordered
degree-$d$ axes, $d\ge3$, in every characteristic.\\
JAC-P10 & Arithmetic continuation of cote~44, scan~4, height question &
The explicit stable Hodge lattice of the Fermat quartic has
normalized determinant gain $7/4$; hence
$h_F(J)-3h_F(E_{1728})=\log(2)/4$. No node thickness is inferred
from an unweighted intersection matrix.\\
JAC-P11 & Cote~44, scan~6, non-$\tau$ cycles; arithmetic height continuation &
The central CM operation $T$ has $T^2=-1$ and preserves height.
The Ceresa vector and $T$-translate span a positive orthogonal
plane and are independent even modulo algebraic equivalence.\\
JAC-P12 & Same Th.~7 scope question &
For the hyperflex embedding, the ordinary tautological
codimension-two ring is $\mathbf Q\theta^2+\mathbf Q\eta$;
$T\eta$ lies outside it even modulo algebraic equivalence.
Tautological generation does not prove generation of the whole Chow group.\\
JAC-P13 & Cote~44, scan~6; modern motivic continuation &
With explicitly added Kimura background, the exact complement is
$N=h^3(J)^G$. Its whole $\operatorname{CH}^2$ is a
$\mathbf Q(i)$-module carrying the Hermitian height form.\\
JAC-P14 & Same modern continuation &
For $\operatorname{End}^0(J)=M_3(\mathbf Q(i))$, the exact action
on $N$ is $u-vT$ when $\det b=u+vi$. The full abelian-endomorphism
orbit of the Ceresa vector remains its CM plane.\\
JAC-P15 & Cote~44, scan~6, Th.~7 over finite fields &
Accepted full Tate plus Kahn gives rational=numerical, and hence
numerical=$\tau$, for every geometric abelian-type variety over
finite fields or $\overline{\mathbf F}_p$.\\
\bottomrule
\end{longtable}
\endgroup

\paragraph{The relevant cote~16 corrections are mathematical, not palaeographic.}
In no.~7, scans~34, 36, 38, 40, 49 and~50 are the operative
Lefschetz pages; scans~35, 37 and~39 are unrelated printed
material. The lowering operator is normalized along a Lefschetz
string, not the operator with $\mathfrak{sl}_2$ string
coefficients. With that convention the boxed identity on scan~40,
\[
             \Lambda_Y=j^*\Lambda_X^2j_*,
\]
is correct and is used in JAC-P02. Several adjacent indices in
the transcription are type-inconsistent: for example the lower
right target in scan~34 is forced by degree to be
$H^{2n-i}(Y)$, and the first intermediate identity on scan~40
has only the bottom-endpoint sum $0\le i\le n$ in
$L_X\Lambda_X=1-\sum p_i^X$. The full all-endpoint sum there
belongs only to the subsequent product. The detailed
transcription audit retains every uncertainty and distinguishes
these forced mathematical readings from claims about the
unseen manuscript.

% Proof-ready contributions by /root/programme_audit.
% Run JAC1964-20261121; GPT-6 Astra Pro; effort setting not exposed.
% All statements are rational, unless explicitly stated otherwise.

\section{Incidence Picard varieties and the next theorem of the programme}

For a smooth projective variety $V$ of dimension $d$ over an algebraically
closed field, write $A^p(V)=\mathrm{CH}^p(V)_{\mathrm{alg}}$ and
\[
\begin{split}
 I^p(V)=\{z\in A^p(V):{}&\Gamma_*z=0\text{ in }\mathrm{Pic}^0(T)\\
 &\text{for every smooth projective }T\text{ and }
 \Gamma\in\mathrm{CH}^{d-p+1}(V\times T)\}.
\end{split}
\]
Saito's higher Picard variety, denoted $P^p(V)$ here, represents
$A^p(V)/I^p(V)$ on algebraically closed points
\cite[Theorem~(3.6)(i)]{Saito1979}.
This is the incidence construction suggested by cote~44, scans~2 and~9.
It must not be silently substituted for the universal regular algebraic
representative, a different universal problem.  Saito supplies duality
\emph{up to isogeny}, not an assertion of a canonical integral principal
polarization \cite[Proposition~(4.4)]{Saito1979}.
Rational families below mean rational homomorphisms of abelian varieties,
obtained by clearing the denominators of the cycles first.  They do not
mean that one tensors the abstract group of geometric points with
$\mathbf Q$; that operation would lose the construction over
$\overline{\mathbf F}_p$, where all abelian-variety points are torsion.

\begin{lemma}[JAC-P00: rational descent of the incidence construction]
\label{jac:P00}
Let $F$ be any field and let $V/F$ be smooth and projective.  The
incidence Picard variety of $V_{\overline F}$, together with its quotient
map on algebraically trivial cycles, has a model over $F$ in the isogeny
category.  If $T/F$ is abelian and $u:T\to P^p(V)$ is a rational
homomorphism over $F$, there is a cycle
$Z\in\mathrm{CH}^p(T\times V)_{\mathbf Q}$ whose family map is $u$.
\end{lemma}

\begin{proof}
Saito's quotient is canonical: its universal property among Picard
homomorphisms gives a unique isomorphism from each conjugate quotient
to the quotient of $V_{\overline F}$.  These isomorphisms satisfy the
cocycle identity.  The variety, its group law, the finite correspondence
data witnessing that its quotient map is a Picard homomorphism, and
one surjective generating family are defined over a finite extension
of $F$.  The latter can be taken from Saito's Poincar\'e cycle, whose
family map is a positive integer times the identity.  An automorphism
fixing these data fixes the quotient map on that generating family.
Uniqueness of the comparison isomorphism, and surjectivity of the
family map, then force that comparison to be the identity.  This
proves, using finite algebraic data rather than the set of all cycles,
that the canonical descent datum is trivial on an open subgroup of
the absolute Galois group.

Here are the details that also cover an imperfect $F$.  Enlarge a
finite separable extension to a finite Galois extension $L/F$, and
choose $r\ge0$ so that the preceding finite field of definition lies
in $L^{1/p^r}$ in characteristic $p$; take $r=0$ in characteristic
zero.  Raising the coefficients of the variety and its group law to
$p^r$ gives an abelian variety $Q/L$, and relative Frobenius gives an
isogeny $f:P_{\overline F}\to Q_{\overline F}$.
Transport the canonical descent datum by $f$ and its rational inverse.
For $\sigma\in\operatorname{Gal}(L/F)$ this gives
\[
 c_\sigma:{}^\sigma Q\longrightarrow Q
 \quad\text{in the rational category},\qquad
 c_{\sigma\tau}=c_\sigma\,{}^\sigma c_\tau.
\]
These maps are defined over $L$: homomorphisms of abelian varieties
acquire no new elements under purely inseparable extensions, and the
remaining Galois invariance follows from the cocycle identity and the
triviality of the datum on automorphisms fixing $L$.  For the former
standard rigidity fact, see \cite[Example~4.2.14, p.~38]{ConradAV}.

Rational Galois descent is effective by the following explicit
construction.  Let $G=\operatorname{Gal}(L/F)$ and
$R=\operatorname{Res}_{L/F}Q$.  Over $L$ identify
$R_L=\prod_{\sigma\in G}{}^\sigma Q$.  The maps
\[
 i:Q\longrightarrow R_L,\quad
 i(x)=(c_\sigma^{-1}x)_\sigma,
 \qquad
 q:R_L\longrightarrow Q,\quad
 q((x_\sigma))=\frac1{|G|}\sum_\sigma c_\sigma x_\sigma
\]
satisfy $qi=1$.  The cocycle identity makes $e=iq$ Galois invariant,
so $e\in\operatorname{End}^0_F(R)$.  For a positive integer $N$ with
$Ne$ integral, its image is an abelian variety over $F$ whose base
change is isomorphic to $P$ in the isogeny category.  This supplies the
required rational model and its transported quotient map.  Uniqueness
in the rational category follows from the same quotient universal
property.

Finally clear a positive denominator $d$ so that $du$ is a morphism,
pull a normalized rational Poincar\'e cycle back by $du$ over
$\overline F$, and divide the resulting cycle by $d$.
The resulting family is defined over some finite
extension $F'/F$.  Put $d_s=[F':F]_{\rm sep}$ and
$d_i=[F':F]_{\rm insep}$, and let
$\pi:(T\times V)_{F'}\to T\times V$.
Define $Z=(d_sd_i)^{-1}\pi_*Z_{F'}$.  After base change to
$\overline F$, flat base change for Chow groups expresses this as
the average of the $d_s$ conjugate cycles, each with multiplicity
$d_i$.  The multiplicity is the length of the corresponding local
factor of $F'\otimes_F\overline F$.  Every conjugate induces $u$,
because $u$ is defined over $F$ for the compatible quotient descent
datum.  Hence $Z$ induces $u$.  This explains explicitly both the
separable conjugates and the inseparable factor; it does not treat a
purely inseparable extension as if it had extra automorphisms.
\end{proof}

\begin{lemma}[JAC-P01: detection by the odd realization]
\label{jac:P01}
Let $V,W$ be smooth projective varieties over an algebraically closed
field, and let
$\Gamma\in\mathrm{CH}^{\dim V-p+q}(V\times W)_{\mathbf Q}$.
Fix a prime $\ell$ invertible in the field.  If the action of $\Gamma$
from $H^{2p-1}(V,\mathbf Q_\ell)$ to
$H^{2q-1}(W,\mathbf Q_\ell(q-p))$ is zero, then the induced rational
homomorphism $P^p(V)\longrightarrow P^q(W)$ is zero.

More generally, for an abelian variety $T$ and
$Z\in\mathrm{CH}^q(T\times W)_{\mathbf Q}$, the homomorphism
$T\longrightarrow P^q(W)$ defined by $t\mapsto Z_t-Z_0$ depends only
on the K\"unneth component
\[
 [Z]_{1,2q-1}\in H^1(T,\mathbf Q_\ell)\otimes
 H^{2q-1}(W,\mathbf Q_\ell(q)).
\]
In particular, it is zero if that component is zero.
\end{lemma}

\begin{proof}
For the second assertion, test $Z_t-Z_0$ by an arbitrary incidence
correspondence
$\Xi\in\mathrm{CH}^{\dim W-q+1}(W\times S)$.
The composition $\Xi\circ Z$ is a divisor class on $T\times S$.
Its component in $H^1(T)\otimes H^1(S)$ depends only on
$[Z]_{1,2q-1}$, by the K\"unneth formula and integration over $W$.
A divisor correspondence on $T\times S$ gives a homomorphism
$T\to\mathrm{Pic}^0(S)$; its rational Tate-module action is precisely
the mixed degree-one component.  Faithfulness of the rational Tate
module on homomorphisms of abelian varieties therefore says that this
homomorphism is zero when the indicated mixed component is zero.
Consequently every $Z_t-Z_0$ is annihilated by every incidence tester.
It is zero in $A^q(W)/I^q(W)$ by definition, proving the assertion.
The argument applies to rational cycles after clearing denominators.

Every algebraically trivial cycle is a finite sum of differences of
fibres of families parametrized by smooth projective curves.
For one such family
$U\in\mathrm{CH}^{p}(C\times V)$, its only relevant odd component
gives the map $H^1(C)\to H^{2p-1}(V)(p-1)$.
Composition with $\Gamma$ is zero by the hypothesis.  Composing further
with each incidence tester on $W$ gives a divisor correspondence on
$C\times S$ with zero mixed component.  The induced
$\operatorname{Alb}(C)\to\mathrm{Pic}^0(S)$ is zero, by the same
faithfulness argument.  Thus $\Gamma_*$ sends every element of
$A^p(V)$ into $I^q(W)$, and its action on $P^p(V)$ is zero.
\end{proof}

The lemma does not assume the standard conjecture~D.  The use of
cohomology concerns correspondences acting on families of cycles;
it does not identify homological and incidence equivalence of individual
cycles.

\begin{theorem}[JAC-P02: Theorem~6 for incidence Picard varieties]
\label{jac:P02}
Let $X$ be smooth projective of dimension $n$ over an algebraically
closed field, and assume that its normalized Lefschetz lowering
operator is algebraic.  For a smooth hyperplane section $j:Y\hookrightarrow X$,
restriction induces
\[
 j^*:P^p(X)\longrightarrow P^p(Y).
\]
It is an isogeny if $2p\le n-1$, has finite kernel if $2p=n$, and
is surjective if $2p\ge n+1$.  Each assertion means the corresponding
split statement in the rational category of abelian varieties, hence
the asserted statement for the original integral homomorphism.

Suppose further that $Y=Y_\eta$ is the generic member of a pencil with
smooth codimension-two base locus, and put $K=k(t)$.  In the critical
case $n=2p$, the image of $P^p(X)_K$ is the image of the $K/k$-trace
of $P^p(Y_\eta)$, up to isogeny.  Thus the rational form of all clauses
of cote~44, scan~4, Theorem~6, holds whenever the Lefschetz standard
conjecture holds for $X$.  In particular, it holds for abelian $X$
in every characteristic.
\end{theorem}

\begin{proof}
Write $L$ for multiplication by the hyperplane class and $\Lambda$
for the normalized lowering operator: it sends $L^a v$ to
$L^{a-1}v$ for a primitive vector $v$ and $a>0$, and kills $v$.
The standard conjecture supplies an algebraic correspondence for this
operator.  The correspondence
\[
 R=\Lambda_X\circ j_*
\]
has degree zero and induces $P^p(Y)\to P^p(X)$.
Put $r=2p-1$.
For $r<n$, no Lefschetz string terminates in degree $r$, so
\[
 Rj^*=\Lambda_X L_X=\operatorname{id}\quad\text{on }H^r(X).
\]
By Lemma~\ref{jac:P01}, the same identity holds on $P^p(X)$.
If $r<n-1$, weak Lefschetz makes $j^*:H^r(X)\to H^r(Y)$ an
isomorphism, hence $j^*R$ is also the identity on $H^r(Y)$ and on
$P^p(Y)$.  This proves the first assertion and, when $r=n-1$,
the assertion of finite kernel.

For $r\ge n$, the Gysin map
$j_*:H^r(Y)\to H^{r+2}(X)(1)$ is an isomorphism by the dual of
weak Lefschetz.  No Lefschetz string begins in degree $r+2>n$,
so $L_X\Lambda_X=\operatorname{id}$ there.  For $y\in H^r(Y)$,
\[
 j_*j^*Ry=L_X\Lambda_Xj_*y=j_*y.
\]
Injectivity of $j_*$ gives $j^*R=\operatorname{id}$ on $H^r(Y)$;
Lemma~\ref{jac:P01} transfers this identity to $P^p(Y)$.
These are precisely the algebraic inverse mechanisms of cote~16,
scan~33, formula~B(4), and scans~40 and~45, clauses~(iv) and~(iv bis).

Here is a cycle-level proof of the fixed-part assertion, so no inference
from invariant cohomology to algebraicity of a correspondence is left
unstated.  Write $p=m$, $n=2m$, let $B$ be the base locus, and let
\[
 b:\widetilde X=\operatorname{Bl}_B X\longrightarrow X,
 \qquad e:E\hookrightarrow\widetilde X,
 \qquad q:E\longrightarrow B
\]
be the blow-up, exceptional divisor, and its projection.  Let $T/k$
be an abelian variety and
$u:T_K\longrightarrow P^m(Y_\eta)$ a homomorphism.  Saito's
Poincar\'e $m$-cycle supplies, after division by its positive integer
multiplier, a rational cycle on $P^m(Y_\eta)\times Y_\eta$
inducing the identity \cite[\S4.1]{Saito1979}.  Lemma~\ref{jac:P00}
descends its pullback by $u$ to a rational cycle
\[
 Z_\eta\in\mathrm{CH}^m(T_K\times Y_\eta)_{\mathbf Q}
\]
inducing $u$.  In particular the construction is valid even when
$K=k(t)$ is imperfect.  Subtracting $Z_0$ normalizes
the family at the identity of $T$.

Close a cycle representative in $T\times\widetilde X$.
The codimension-two blow-up formula gives
\[
 \overline Z=(1_T\times b)^*Z_X+
 (1_T\times e)_*(1_T\times q)^*Z_B,
\]
where
$Z_X\in\mathrm{CH}^m(T\times X)_{\mathbf Q}$ and
$Z_B\in\mathrm{CH}^{m-1}(T\times B)_{\mathbf Q}$ are defined over
$k$.  Write $k_B:B\hookrightarrow X$ and $i:B_K\hookrightarrow Y_\eta$.
The mixed component relevant to $Z_B$ is
$H^1(T)\otimes H^{n-3}(B)$.
Weak Lefschetz identifies $H^{n-3}(X)$ with $H^{n-3}(B)$.
An algebraic inverse is
\[
 S=\Lambda_X^2\circ(k_B)_*.
\]
Indeed, $(k_B)_*k_B^*=L_X^2$ and $\Lambda_X^2L_X^2$ is the
identity on degree $n-3$.  Define an algebraic rational cycle over $k$
by
\[
 W_X=Z_X+L_X\circ S\circ Z_B\in
 \mathrm{CH}^m(T\times X)_{\mathbf Q}.
\]
For $x\in H^{n-3}(X)$ the projection formula gives
\[
 i_*k_B^*x=i_*i^*j_\eta^*x=L_{Y_\eta}j_\eta^*x
 =j_\eta^*L_Xx.
\]
It follows that $Z_\eta$ and $j_\eta^*W_X$ have the same
$H^1(T)\otimes H^{n-1}(Y_\eta)$ component.  Lemma~\ref{jac:P01}
says they induce the same rational homomorphism from $T_K$ to
$P^m(Y_\eta)$.  The latter factors through $P^m(X)_K$ since $W_X$
is defined over $k$.  When $m=1$ the exceptional term has no mixed
component, since $H^{-1}(B)=0$, and the same conclusion follows.

Every homomorphism from a constant abelian variety has therefore image
in $\operatorname{Im}j_\eta^*$.  Conversely, this image is itself the
image of a homomorphism from the constant variety $P^m(X)_K$.
The universal property of the $K/k$-trace proves the asserted equality
of images up to isogeny, using the rational $K$-model of
Lemma~\ref{jac:P00}.
\end{proof}

\section{An explicit test of the primitive sign}

\begin{proposition}[JAC-P03: the uncorrected correspondence is indefinite]
\label{jac:P03}
There exist normalized algebraic families of codimension-two cycles
on an abelian threefold for which the divisor correspondence of
cote~44, scan~9, is nondegenerate and indefinite.  Neither quotienting
by its kernel nor multiplication by a common sign makes it a
polarization.  The primitive qualification in scan~8 is essential.
\end{proposition}

\begin{proof}
Let $E$ be an elliptic curve, set $X=E_1\times E_2\times E_3$,
write $a_i$ for the pullback of $[0]\in\mathrm{CH}^1(E_i)$, and use
$h=a_1+a_2+a_3$ as polarization.  On $E\times E$ put
\[
 P=[\Delta_E]-[E\times0]-[0\times E].
\]
As a correspondence $P$ is the symmetric degree-one Chow--K\"unneth
projector, hence ${}^tP\circ P=P$.
As a divisor correspondence it induces
\[
 E\longrightarrow\widehat E,
 \qquad t\longmapsto\mathcal O_E(t-0)=-\lambda_{\mathcal O_E(0)}(t),
\]
where the polarization convention is
$\lambda_L(t)=t_t^*L\otimes L^{-1}$.
Thus the sign calibration is fixed, rather than hidden in a convention.

Define two families on $E_t\times X$ by
\[
 z_\pm=P_{t,1}\,(a_2\pm a_3).
\]
Their incidence divisor correspondences are
\[
 D_\pm=\deg_{E_2\times E_3}(a_2\pm a_3)^2\,P
       =\pm2P,
\]
and the mixed incidence correspondence between the two families
is zero because $(a_2+a_3)(a_2-a_3)=0$.
Use independent parameters for the two families, obtaining a family
over $T=E\times E$.  Its symmetric homomorphism $T\to\widehat T$
has diagonal blocks $-2\lambda_{\mathcal O(0)}$ and
$+2\lambda_{\mathcal O(0)}$.  It is an isogeny, but its associated
rational line bundle restricts to a negative-degree line bundle on one
elliptic factor and a positive-degree one on the other.  It cannot be
ample.  Its negative cannot be ample either.  Passing through a finite
quotient does not change this conclusion, since the pullback of an
ample line bundle by an isogeny is ample.

The degree-one parameter components make the primitive issue explicit.
For $\alpha\in H^1(E_1)$, they have images
\[
 u_\pm(\alpha)=\alpha(a_2\pm a_3).
\]
One has $u_+(\alpha)=L_X\alpha$ and $L_Xu_-(\alpha)=0$.
Moreover, for $\alpha,\beta\in H^1(E_1)$,
\[
 \int_Xu_+(\alpha)u_+(\beta)=2\int_E\alpha\beta,
 \qquad
 \int_Xu_-(\alpha)u_-(\beta)=-2\int_E\alpha\beta,
\]
and the cross pairing vanishes.  The lower Lefschetz and primitive
pieces require different signs.  In the divisor convention above,
the primitive $m=2$ correspondence $D_-=-2P$ already gives the
positive homomorphism $2\lambda_{\mathcal O(0)}$.
\end{proof}

\section{Primitive middle cycles need not reach the incidence Jacobian}

\begin{proposition}[JAC-P04: an abelian-fourfold domain test for H9]
\label{jac:P04}
Let $X=E_1\times E_2\times E_3\times E_4$ over $\mathbf C$,
$a_i=[0]$ pulled back from $E_i$, and $h=a_1+a_2+a_3+a_4$.
For a sufficiently high-degree Lefschetz pencil relative to $h$, let
$Y_\eta$ be its geometric generic member.  The codimension-two cycle
\[
 \alpha=(a_1-a_2)(a_3-a_4)
\]
satisfies $h\alpha=0$ and $\deg\alpha^2=4$.  Its restriction to
$Y_\eta$ is homologically trivial, but no nonzero multiple of that
restriction is algebraically equivalent to zero.
\end{proposition}

\begin{proof}
The relations $a_i^2=0$ give
\[
 h\alpha=0,
 \qquad
 \alpha^2=(-2a_1a_2)(-2a_3a_4)=4a_1a_2a_3a_4.
\]
In particular $\alpha$ has nonzero cohomology class.  Since the Gysin
map $H^4(Y_\eta)\to H^6(X)$ is an isomorphism by weak Lefschetz
and duality, the equality $j_*j^*\alpha=d\,h\alpha=0$ gives
$j^*\alpha=0$ in cohomology, for the degree-$d$ pencil.

Apply SGA~7~II, Expos\'e~XX, Theorem~4.1
\cite[printed p.~355]{SGA7II}.  Its hard-Lefschetz hypothesis holds
for $X$.  Its algebraic-retraction hypothesis~B is supplied by the
algebraic Lefschetz operator, as noted in Expos\'e~XX, \S3.4.
For sufficiently high degree the pencil satisfies condition~A by
Expos\'e~XVIII, Corollary~6.4.  In characteristic zero its variable
middle cohomology has level greater than one for sufficiently high
degree, by Expos\'e~XX, Remark~4.3 (printed p.~356).
The cited theorem says that a middle-dimensional cycle on $X$ with
algebraically trivial generic restriction has zero cohomology class
on $X$.  Applied to any nonzero multiple of $\alpha$, this contradicts
$\deg\alpha^2=4$.  Therefore no such multiple is algebraically
trivial on $Y_\eta$.
\end{proof}

This example identifies the enlargement between cote~44, scan~4,
Theorem~5, which explicitly requires algebraic triviality on $Y$,
and H9, which assumes only homological triviality.  The latter condition
does not provide a point of the incidence Picard variety.  It provides
a class in the Griffiths group in this example.  A general intermediate
Jacobian or normal function is a different object and does not remove
this distinction.

\begin{proposition}[JAC-P05: the customary reading of Theorem~7 is false]
\label{jac:P05}
If $\tau$-equivalence means that a nonzero integral multiple is
algebraically equivalent to zero, the unrestricted equality of numerical
and $\tau$-equivalence in cote~44, scan~6, fails already for complex
abelian threefolds of CM type, including one defined over a number field.
\end{proposition}

\begin{proof}
Take the Fermat quartic
\[
 C:\quad T_0^4+T_1^4=T_2^4
\]
and embed it in its Jacobian $J=J(C)$ using a rational base point.
Bloch's Theorem~(4.1), printed p.~104 of \cite{Bloch1984L}, states
that its Ceresa cycle $\Xi=[C]-[-1]_*[C]$ has infinite order modulo
algebraic equivalence.  This is the rational non-torsion assertion,
stronger than merely $\Xi\not\sim_{\rm alg}0$.
Bloch's Lemma~(3.1), printed p.~99, gives an isogeny
\[
 J\longrightarrow E\times E\times E',\qquad
 E:y^2=x^3-x,\quad E':y^2=x^3+x.
\]
Both elliptic curves have CM by $\mathbf Q(i)$, as is seen from
$(x,y)\mapsto(-x,iy)$; thus $J$ is of CM type.
On $H^4(J,\mathbf Q)$ the involution $[-1]^*$
acts as $(-1)^4=1$.  Thus this cycle is homologically trivial, hence
numerically trivial by compatibility of cycle classes and intersection.
Its infinite order modulo algebraic equivalence means that it is not
$\tau$-equivalent to zero.  This gives the required counterexample.
\end{proof}

Ceresa's original theorem also gives a very general genus-three
Jacobian example.  The explicit Bloch example has the advantage that
both the infinite-order statement and the CM isogeny have directly
verified primary-source locators.  A theorem making Hodge classes
algebraic says nothing that would force a homologically trivial
cycle such as $\Xi$ to become algebraically trivial.

For divisors the statement is correct in every characteristic:
SGA~6, Expos\'e~XIII, Theorem~4.6 identifies numerical triviality
with $\tau$-triviality; definitions are in \S\S4.1 and~4.4.
The supplied archive page does not itself define $\tau$.  Accordingly,
the preceding refutation is explicitly tied to this customary reading,
not to a guessed replacement of $\tau$ by homological equivalence.

\begin{lemma}[JAC-P06: the finite-field meaning of the last claim]
\label{jac:P06}
Over $k=\overline{\mathbf F}_p$, every algebraically trivial Chow class
is torsion.  Hence $\tau$-triviality means torsion in the Chow group,
and numerical equivalence equal to $\tau$-equivalence would identify
the rational Chow group with cycles modulo numerical equivalence.
For abelian varieties, and more generally geometric abelian-type
varieties, JAC-P15 proves that comparison using the accepted Tate
consequence of 032 and Kahn's additional established theorem.
\end{lemma}

\begin{proof}
An algebraically trivial cycle is a finite sum of images, under cycle
correspondences, of degree-zero zero-cycles on smooth projective curves.
For a curve $C/k$, its degree-zero Chow group is $J(C)(k)$.  Every
point and the curve are defined over a finite field, and each finite
group $J(C)(\mathbf F_q)$ is torsion.  Thus $J(C)(k)$ is torsion,
and so is each of its images under a correspondence.  A finite sum of
torsion elements is torsion.  If a nonzero multiple of a Chow class
is algebraically trivial, another nonzero multiple is therefore zero
in the Chow group.  The converse is immediate.
\end{proof}


% JAC1964-20261121. /root/programme_audit; GPT-6 Astra Pro.
% Accepted input: 032 Corollary 8.3/Milne's full Tate consequence.
% Additional established input: Kahn 2003 Theorem 1.10.

\subsection*{JAC-P15. The numerical--tau assertion over finite fields}

\begin{theorem}[JAC-P15]\label{jac:P15}
Accept the Tate consequence of 032, Corollary~8.3, printed
pp.~49--50, as also formulated in the supplied Milne report,
Theorem~3.1, printed p.~6. Let $k$ be a finite field
or an algebraic closure of a finite field. Let $X/k$ be smooth and
projective, and suppose that its geometric rational Chow motive is
of abelian type: it belongs to the thick rigid tensor category
generated by abelian varieties and Artin motives.
Then, for every $r$, the natural map
\[
 \operatorname{CH}^r(X)_{\mathbf Q}
        \longrightarrow N^r(X)_{\mathbf Q}                 \tag{JAC.P15.1}
\]
is an isomorphism, where $N^r$ denotes cycles modulo numerical
equivalence. In particular, for an integral cycle $z$ on $X$,
\[
 z\equiv_{\rm num}0
 \quad\Longleftrightarrow\quad
 [z]\text{ is torsion in }\operatorname{CH}^r(X)
 \quad\Longleftrightarrow\quad
 z\equiv_\tau0.                                           \tag{JAC.P15.2}
\]
Here $\tau$ has the customary meaning that a nonzero multiple is
algebraically equivalent to zero.

Consequently cote~44's Th.~7, with that meaning of $\tau$, holds
for every abelian variety over $\overline{\mathbf F}_p$, in every
dimension and codimension. The conclusion also holds for smooth
projective varieties whose geometric Chow motive is of abelian
type, including products of curves.
\end{theorem}

\begin{proof}
We separate the accepted arithmetic input from the additional
established motivic theorem.

\emph{Tate passes to abelian-type summands.}
The accepted 032, Corollary~8.3, printed pp.~49--50, says that for every abelian variety
$A/\mathbf F_q$, every $r$ and every $\ell\ne p$, the cycle-class
map
\[
 \operatorname{CH}^r(A)\otimes\mathbf Q_\ell
   \longrightarrow H^{2r}(A_{\overline{\mathbf F}_q},
                         \mathbf Q_\ell(r))^{G_q}
\]
is surjective. This applies to every product of abelian varieties
and after every finite extension of the finite field.

Every construction of an object of the thick rigid category uses
finitely many sums, tensor products, duals, projectors, and splitting
correspondences. After a finite extension $k'/k$, its Artin factors
split. Tensor products of abelian motives are motives of abelian
products; their duals are Tate twists of such motives, by Poincare
duality. Thus, after enlarging $k'$ if necessary, $h(X_{k'})$ is
a direct summand of a finite direct sum
\[
              \bigoplus_j h(A_j)(t_j)                    \tag{JAC.P15.3}
\]
with each $A_j/k'$ abelian. If the initial abelian-type presentation
was given only over the algebraic closure, the same conclusion
holds: the varieties, projectors, correspondences and the finitely
many rational equivalences giving their splitting identities all
descend to a finite field extension.

A Galois-invariant cohomology class on $X_{k'}$ maps, by the
inclusion in (JAC.P15.3), to invariant classes of the appropriate
degrees and Tate twists on the $A_j$. The accepted Tate statement
represents them by $\mathbf Q_\ell$-linear combinations of cycles.
Applying the algebraic retraction in (JAC.P15.3) gives a cycle
representing the original class on $X_{k'}$. Thus $X_{k'}$ satisfies
Tate in every codimension. If the class originally came from a
$G_k$-invariant class on $X$, its cycle representative descends
rationally by the normalized norm
\[
           [k':k]^{-1}\pi_*,\qquad \pi:X_{k'}\to X.
\]
Indeed pullback of this norm is the Galois average and fixes that
class. No assertion that every tensor product of previously known
Tate classes spans all tensor invariants is used; Tate is invoked
on the abelian product itself.

\emph{Kahn's theorem upgrades the equivalence relation.}
Kahn, \emph{Equivalences rationnelle et numerique sur certaines
varietes de type abelien sur un corps fini}, Ann.\ Sci.\ Ecole
Norm.\ Sup.\ (4) 36 (2003), 977--1002, Definition~1 and Theorem~1,
printed p.~978, and Theorem~1.10, printed p.~982, proves that
rational and numerical equivalence agree with rational
coefficients for an abelian-type variety satisfying Tate.
All of these hypotheses have just been verified for $X_{k'}$.
We therefore obtain (JAC.P15.1) over $k'$.

For precision, the mechanism in that source is stronger than
mere equality of homological and numerical equivalence. The
nilpotence theorem for Kimura motives lifts the numerical simple
decomposition. A non-unit simple factor on which Frobenius is
not $1$ has no morphism from the unit: if its minimal Frobenius
polynomial is $P$ with $P(1)\ne0$, nilpotence gives
$P(F)^a=0$, whereas $Ff=f$ for a morphism $f$ from the unit.
The unit factors lift to actual unit motives. Tate and Frobenius
semisimplicity identify exactly these factors. The precise source
loci are Theorem~1.4, Proposition~1.5, Lemma~1.6 and
Proposition~1.7, printed p.~981, and Lemma~1.9 and
Theorem~1.10, printed p.~982.

Suppose now that $k$ is finite and $z$ is numerically trivial on
$X$. Its pullback to $k'$ is numerically trivial: for each
complementary cycle $w$ on $X_{k'}$, the projection formula gives
\[
 [k':k]\,\deg_{k'}(\pi^*z\cdot w)
       =\deg_k(z\cdot\pi_*w)=0.
\]
Hence $\pi^*z=0$ in rational Chow, and
$\pi_*\pi^*z=[k':k]z$ gives $z=0$ there as well.
This proves (JAC.P15.1) even when the abelian-type presentation
only becomes available over $k'$.

If $k=\overline{\mathbf F}_p$, choose a finite field of definition
for $X$, $z$, and its abelian-type presentation. Numerical
triviality of $z$ over $k$ implies numerical triviality over that
finite field, since its test cycles form a subset of the geometric
test cycles. The finite-field conclusion makes $z$ zero in
rational Chow there, hence also over $k$. The reverse implication
in (JAC.P15.1) follows from rational equivalence implying numerical
equivalence. Surjectivity of (JAC.P15.1) is built into the
definition of the numerical quotient.

Finally, an element of an abelian group vanishes after tensoring
with $\mathbf Q$ exactly when it is torsion. Thus numerical
triviality implies Chow torsion. Chow torsion implies
$\tau$-triviality because a rationally trivial multiple is
algebraically trivial. Conversely a $\tau$-trivial cycle is
numerically trivial: all intersection numbers of its algebraically
trivial multiple vanish, and integer intersection numbers have
no nonzero torsion. This proves (JAC.P15.2).
\end{proof}

\paragraph{A finite-field consequence for higher $K$-theory.}
For $X$ as in JAC-P15 over a finite field,
\[
                    K_i(X)\otimes\mathbf Q=0
                    \qquad(i>0).                         \tag{JAC.P15.4}
\]
For $X$ already of abelian type over that field, this is Kahn's
Corollary~2.2, printed p.~982, applied to the Tate result above.
In the geometric abelian-type case, apply that corollary after
the finite extension giving (JAC.P15.3), then use the restriction
and transfer identity $\mathrm{cor}\,\mathrm{res}=[k':k]$ on
$K_i(X)$. This is a consequence through Kahn's cited theorem,
not a new proof of the general Parshin conjecture.

\paragraph{What changes in the programme.}
JAC-P06 correctly identifies the finite-field meaning of $\tau$,
but does not itself prove the numerical comparison. JAC-P15 now
supplies that comparison using the \emph{Tate} consequence of the
accepted input and Kahn's additional established theorem. It
does not follow from the Hodge standard statement of Corollary~8.4
alone. It does not apply to every smooth projective variety in
positive characteristic without the abelian-type hypothesis, to
arbitrary transcendental extensions of $\overline{\mathbf F}_p$,
or to the number-field Fermat Jacobian of JAC-P05. The latter
retains its non-torsion Ceresa obstruction.

\paragraph{Primary source verification.}
The published 26-page Kahn article was directly read from
\path{https://www.numdam.org/item/ASENS_2003_4_36_6_977_0.pdf}.
The author's 35-page \texttt{cell4.pdf} has the same theorem numbers
but different pagination: Definition~1 p.~2, Theorem~1 p.~3,
the nilpotence argument pp.~6--7, and Theorem~1.10 p.~8.
The older arXiv version uses different numbering and is not used
for these exact locators. The conclusion above is a deduction
in this run from accepted 032/Milne and established Kahn; no
historical originality or formal verification is asserted.


% JAC1964-20261121. /root/programme_audit; GPT-6 Astra Pro.
% Accepted input: numerical=homological on all abelian varieties,032Cor8.4.
% Additional established input: Jannsen semisimplicity.

\subsection*{JAC-P16. Lefschetz inverses for varieties of geometric abelian type}

\begin{theorem}[JAC-P16]\label{jac:P16}
Let $k$ be any field and let $X/k$ be smooth projective of dimension
$d$. Suppose that $h(X_{\bar k})$ belongs to the thick rigid
category of rational Chow motives generated by abelian varieties
and Artin motives. Accept the equality of numerical and
$\ell$-adic homological equivalence on every abelian variety over
an algebraically closed field, for every invertible prime $\ell$,
as stated in 032, Corollary~8.4, printed p.~50.

Then numerical and homological equivalence agree on $X^a$ in
every codimension, for every $a\ge1$. In particular this assertion
holds on $X\times X$, not only for cycles on $X$.
The Kunneth components of the diagonal of $X$ are algebraic.
For every ample class $h$ defined over $k$ and every $0\le i\le d$,
the inverse of
\[
 L_h^{d-i}:H^i(X_{\bar k},\mathbf Q_\ell)
       \longrightarrow
 H^{2d-i}(X_{\bar k},\mathbf Q_\ell(d-i))                 \tag{JAC.P16.1}
\]
is induced by a rational algebraic correspondence defined over $k$.
One can choose these correspondences to work for every
$\ell\ne\operatorname{char}k$. The normalized lowering operator
$\Lambda_X$ is algebraic as well; thus $B(X)$ holds.

If $k$ is algebraically closed, JAC-P02 consequently proves all
three restriction/fixed-trace assertions of cote~44, Th.~6, for
every such $X$ and a genuine Lefschetz pencil on $X$, in every
characteristic. This includes its deleted critical finite-kernel
clause. No Hodge-standard assertion for $X$'s own intersection
form is included in this theorem.
\end{theorem}

\begin{proof}
We first work over $\bar k$ and fix one invertible prime $\ell$.
Let $\mathcal S$ be the category whose objects are finite direct
sums of Tate twists of motives of abelian varieties. Products and
duals remain of this form after adding direct summands. Artin
motives are sums of units over $\bar k$, and the unit is already
an abelian-type summand. Therefore every Chow motive of abelian
type is a direct summand of an object of $\mathcal S$.

\emph{Numerical and homological ideals agree on the whole category.}
For $S,T\in\mathcal S$, a morphism is a matrix of algebraic
cycle correspondences on products of abelian varieties, in the
codimensions determined by the Tate twists. Those products are
themselves abelian varieties. The accepted equality in 032
therefore identifies the numerically zero and homologically zero
matrices. Equivalently, the categorical numerical ideal
\[
 \mathcal N(S,T)
  =\{f:\operatorname{tr}(gf)=0\text{ for every }g:T\to S\}
                                                               \tag{JAC.P16.2}
\]
is exactly the kernel of realization. This is not an assumption
about just endomorphisms: choosing $g$ with one matrix component
tests each component of $f$ against all complementary cycles.
Under Poincare duality, the trace is their intersection number.
Jannsen's trace identity is Lemma~1, printed p.~448, of
\emph{Motives, numerical equivalence, and semi-simplicity},
Invent.\ Math.\ 107 (1992), 447--452.

Here is the direct-summand step explicitly. Suppose $M$ is a
summand of $S$ and $N$ a summand of $T$, with inclusions $u_M,u_N$
and retractions $v_M,v_N$. For $f:M\to N$ extend it by zero to
\[
                 F=u_N f v_M:S\longrightarrow T.
\]
For every $G:T\to S$, cyclicity of the trace gives
\[
 \operatorname{tr}_S(GF)
     =\operatorname{tr}_M(v_MG u_N f).                    \tag{JAC.P16.3}
\]
Every $g:N\to M$ occurs as $v_MG u_N$, by taking
$G=u_Mg v_N$. Thus $f$ is numerically zero if and only if $F$ is.
The ambient equality of ideals makes this equivalent to $H(F)=0$,
which is in turn equivalent to $H(f)=0$, since the inclusions and
retractions split after realization. This proves the equality of
the two ideals on all abelian-type summands and all their Hom
spaces, including twists.

For a variety $W$ of abelian type, codimension-$r$ cycles are
morphisms from the appropriate Tate motive to $h(W)$; testing
their categorical numerical ideal is exactly testing intersection
with complementary cycles on $W$. Hence the previous paragraph
proves numerical=homological on $W$ in every codimension.
Apply this to $W=X^a$, whose motive is of abelian type for every
$a$. In particular, all codimensions on $X\times X$ have the
required equality. This supplies the full correspondence-level
statement used below.

\emph{Algebraic Kunneth projectors.}
Choose a Chow splitting $h(X)\xrightarrow{u}S\xrightarrow{v}h(X)$
with $vu=1$. Abelian varieties have algebraic Kunneth projectors,
with their functorial grading; a direct sum with Tate twists has
the induced degree projectors $\varepsilon_i$. The correspondences
\[
                         \pi_i=v\varepsilon_i u          \tag{JAC.P16.4}
\]
act as the degree-$i$ projections on $H^*(X)$. They are therefore
orthogonal idempotents modulo homological equivalence, and also
modulo numerical equivalence by the equality just proved. We
need only algebraicity of their cohomological classes, not a claim
that the particular Chow representatives in (JAC.P16.4) are
orthogonal idempotents before passing to that equivalence relation.

\emph{Semisimplicity gives the inverse, without a positivity argument.}
The pseudo-abelian category of homological abelian-type motives
is now fully faithfully identified with the full numerical
subcategory generated by the same objects: the equality of ideals
was proved on all Hom spaces, and idempotent completion preserves
that equality. Jannsen's Theorem~1, printed p.~448, makes the
numerical category semisimple abelian. Its full subcategory closed
under direct sums and summands is also semisimple abelian.

Using (JAC.P16.4), (JAC.P16.1) is a morphism between the two
appropriate degree motives, with its indicated Tate twist.
Hard Lefschetz makes its realization an isomorphism. In a
semisimple category every short exact sequence splits, so any
additive realization is exact. Its faithfulness forces both kernel
and cokernel of this morphism to vanish. The morphism is therefore
invertible in the numerical category.

The inverse is a rational correspondence class. Lift it to a
cycle on $X\times X$ with the prescribed degree supports. The
two inverse identities initially hold numerically; the proved
equality on $X\times X$ makes them hold homologically. Thus the
lift induces precisely the hard-Lefschetz inverse, rather than
just some map with a numerically inverse class.
The construction is independent of the selected prime at the
numerical-motive level. The same Kunneth classes work in every
realization, and 032 gives the equality of ideals for every
invertible prime. The lifted inverse therefore works for all of
them.

For clarity, the normalized lowering operator can be reconstructed
by finite algebraic compositions. For $r\le d$ let
$I_r=(L^{d-r})^{-1}:H^{2d-r}(X)\to H^r(X)$, with the degree
supports understood. The primitive projector on $H^r(X)$ is
\[
 e_r=1-L I_{r-2}L^{d-r+1}\quad(r\ge2),
 \qquad e_0=e_1=1,                                     \tag{JAC.P16.5}
\]
where $I_{r-2}=(L^{d-r+2})^{-1}$.
It kills $LH^{r-2}$ and is the identity on
$\ker L^{d-r+1}$. For $0\le a\le d-r$, the projector onto
$L^aP^r\subset H^{r+2a}(X)$ is
\[
             L^a e_r I_r L^{d-r-a}\pi_{r+2a}.           \tag{JAC.P16.6}
\]
On that string, the lowering map for $a\ge1$ is obtained by
replacing the first $L^a$ in (JAC.P16.6) by $L^{a-1}$; it is
zero for $a=0$. Summing these finitely many algebraic
correspondences gives the normalized $\Lambda_X$ used in the
archive and JAC-P02.

\emph{Descent to the original field.}
Each correspondence constructed over $\bar k$ is defined over a
finite extension $L/k$. Its desired cohomological action is
canonical: a Kunneth projection, a hard-Lefschetz inverse for
the $k$-defined class $h$, or the resulting normalized lowering
operator. It is therefore Galois invariant. If
$\pi:(X\times X)_L\to X\times X$, replace a representing
correspondence $\gamma$ by
\[
                         [L:k]^{-1}\pi_*\gamma.          \tag{JAC.P16.7}
\]
After geometric base change, this is the average of the conjugate
correspondences, with their inseparable multiplicities. Each has
the same canonical action. The number of embeddings is the
separable degree and each multiplicity is the inseparable degree,
so their sum is $[L:k]$; (JAC.P16.7) has exactly the required
action. This deals with imperfect $k$ without pretending every
finite extension is Galois or separable.

The numerical/homological equality also descends. A numerically
trivial $k$-cycle remains numerically trivial after any finite
extension, by the projection formula against the norm of an
arbitrary complementary cycle. Every complementary cycle over
$\bar k$ is defined over a finite extension. Hence a numerically
trivial cycle is geometrically numerically trivial and therefore
homologically zero by the geometric result. The reverse implication
always follows from intersection compatibility. This applies
to $X^a$ in every codimension.

We have proved $B(X)$. Over an algebraically closed field, insert
the normalized algebraic $\Lambda_X$ into JAC-P02. Its proof
already establishes the below-middle isogeny, the critical image
and fixed trace, the critical finite kernel, and above-middle
surjectivity. No change to its cycle-spread or imperfect-pencil-field
argument is needed.
\end{proof}

\paragraph{Why the Hodge-standard form is still a separate question.}
The proof establishes algebraic inverses and nondegeneracy of the
pairings on the relevant numerical summands. It does not identify
the geometric primitive form
\[
       (x,y)\longmapsto(-1)^r\deg_X(xyh^{d-2r})           \tag{JAC.P16.8}
\]
with a positive form transported from an abelian variety. Even
if a positive form $Q_0$ is supplied on an abstract summand,
another nondegenerate symmetric form has the shape
$Q(x,y)=Q_0(ax,y)$ with $a$ invertible and $Q_0$-self-adjoint.
These two properties alone do not make its eigenvalues positive;
$a=-1$ already demonstrates that logical gap. Here one must prove
positivity for the particular comparison operator determined by
(JAC.P16.8). The Chow-motive presentation and Jannsen
semisimplicity supply no such assertion. This is not a
counterexample to Hodge standard for abelian-type varieties;
it records exactly why that stronger conclusion has not been
claimed from this proof.

\paragraph{Primary source and archive scope.}
Jannsen's original six-page article was directly read at
\url{https://epub.uni-regensburg.de/26642/1/jannsen10.pdf}:
Theorem~1 and the trace identity are on printed p.~448;
the proof of semisimplicity continues on p.~449.
The functorial abelian Kunneth decomposition is recalled as
Deninger--Murre's theorem in Ancona--Enright-Ward--Huber,
Theorem~4.2.2, printed p.~824, in the source checked for JAC-P14.
The accepted numerical/homological input is explicitly the first
assertion of 032, Corollary~8.4, printed p.~50.
This provides a precise extension of cote~44, scan~4, Th.~6,
and of the inverse-correspondence discussion in cote~16,
scan~49. Neither a universal regular representative nor
Hodge-standard positivity for every such $X$ is inferred.


\section{The primitive incidence polarization in every characteristic}
\label{sec:T4}

All varieties in this section are over an algebraically closed field $k$.
We work with rational correspondences and with $\ell$-adic cohomology,
$\ell\ne\operatorname{char}k$; Tate twists in displayed K\"unneth components
are suppressed when no confusion is possible. A rational polarization means
an element of $\operatorname{NS}(B)_{\mathbf Q}$ a positive integral multiple
of which is the class of an ample line bundle. It does not mean a principal
polarization or an integral universal cycle.

Write $P^p(X)=\operatorname{Pic}^p(X)$ for Saito's higher Picard variety:
\[
 A^p(X)/K^p(X)\simeq P^p(X)(k),\qquad
 A^p(X)=\operatorname{CH}^p(X)_{\rm alg}.
\]
Here $K^p$ is the incidence kernel, tested by divisor-valued
correspondences. The construction, rational Poincar\'e cycles and the
canonical duality isogeny are the cited results
\cite[Theorem~3.6, \S4.1, Proposition~4.4]{Saito1979}.
The notation is the precise incidence interpretation of cote~44,
scan~2, Theorems~1--3, and scan~8. It must not be replaced without proof
by a universal regular algebraic representative.

\begin{lemma}[JAC-T4.1: primitive incidence summand]
\label{lem:primitive-summand}
Let $A$ be an abelian variety of dimension $n=2m-1$, and let $h$ be a
rational ample divisor class. Let $e_h$ denote the cohomological projector
of $H^n(A)$ onto
\[
 H^n(A)_{\rm prim}=\ker\bigl(L_h:H^n(A)\longrightarrow H^{n+2}(A)\bigr)
\]
along $L_hH^{n-2}(A)$, extended by zero in other degrees.
It is represented by a rational algebraic correspondence. It induces an
idempotent on $P^m(A)$ in the isogeny category. Its image $B=P^m(A)_{\rm prim}$
has a nondegenerate symmetric rational incidence correspondence
$d_B\in\operatorname{NS}(B\times B)_{\mathbf Q}$, rigidified along the two
axes. For a rational Poincar\'e family $z\in\operatorname{CH}^m(B\times A)_{\mathbf Q}$
inducing the inclusion $B\longrightarrow P^m(A)$, it is
\begin{equation}\label{eq:incidence}
 d_B=\operatorname{pr}_{12*}(z_{13}\cdot z_{23})_{11}.
\end{equation}
The subscript $11$ removes the pullbacks of divisor classes on the two
factors. The family may be chosen with cohomology entirely in
$H^1(B)\otimes H^n(A)_{\rm prim}$.
\end{lemma}

\begin{proof}
The Lefschetz standard conjecture and the algebraic K\"unneth projectors
hold for abelian varieties; for an exact classical locator see
\cite[Theorem~2A11]{Kleiman68}. Thus the primitive projector is algebraic.
The incidence-detection lemma Lemma~\ref{jac:P01} says that a correspondence
whose relevant odd cohomological action is zero acts as zero on higher
Picard varieties. In particular, its action depends only on the
homological class of that component. Applying this to $e_h^2-e_h$ gives
an actual idempotent in $\operatorname{End}^0(P^m(A))$.

The Poincar\'e pairing on $H^n(A)$ makes its primitive summand orthogonal
to $L_hH^{n-2}(A)$: if $L_hu=0$, then
$\int_Au\,L_hv=\int_A(L_hu)\,v=0$.
Hard Lefschetz makes this an orthogonal direct sum of nondegenerate
spaces. Consequently $e_h$ is self-transpose in degree $n$. Saito's
transpose compatibility for incidence duality
\cite[Proposition~4.4, especially (4.4.1)--(4.4.2)]{Saito1979}, together
with incidence detection, says that the induced idempotent is
self-adjoint for the canonical rational duality of $P^m(A)$.
Writing the resulting isogeny decomposition as
$P^m(A)\sim B\times B'$, this duality is block diagonal. Since the full
duality is an isogeny, each block is an isogeny. In particular the block
on $B$ is nondegenerate.

Choose Saito's Poincar\'e cycle and divide it by its positive integer
$k_m$ so that it induces the identity, then pull it back to $B\times A$.
Apply the degree-one K\"unneth projector on the parameter factor and
$e_h$ on the $A$ factor. These operations preserve the required map to
the incidence quotient and give the asserted mixed cohomological
component. Subtracting constant restrictions rigidifies the family at
the origin without changing this component. Composition of
correspondences gives \eqref{eq:incidence}. Equivalently, the
unnormalized universal cycle would multiply this correspondence by
$k_m^2$, a positive rational square. Symmetry follows by interchanging
the two parameter factors. The block argument proves its
nondegeneracy. This also proves independence of the chosen normalized
family, in the N\'eron--Severi group, from Saito's incidence identities.
\end{proof}

\begin{lemma}[JAC-T4.2: transfer of the curve test to an abelian product]
\label{lem:curve-test}
Let $C$ be a smooth projective pointed curve of positive genus $b$, and
let $A$ have dimension $2m-1$. Suppose a rational algebraic class
$\xi\in H^1(C)\otimes H^{2m-1}(A)$ satisfies
$(1\otimes L_h)\xi=0$. Let $J=J(C)$, let
$a:C\longrightarrow J$ be its Abel map, and let $\theta$ be its canonical
principal polarization. There is a rational algebraic class
$w\in H^1(J)\otimes H^{2m-1}(A)$ with
$(a\times1)^*w=\xi$. For $H=\theta+h$ on $J\times A$,
\begin{equation}\label{eq:curve-transfer}
 H^bw=0,\qquad
 \int_{J\times A}w^2H^{b-1}
 =(b-1)!\int_{C\times A}\xi^2.
\end{equation}
Hence 032, Corollary~8.4, gives
\[
 (-1)^m\int_{C\times A}\xi^2\ge0,
\]
with equality if and only if $\xi=0$.
\end{lemma}

\begin{proof}
The Abel map identifies $h^1(C)$ and $h^1(J)$ by rational algebraic
correspondences. One can use its graph, the divisor correspondence of
the Poincar\'e bundle, and the inverse Lefschetz correspondence on $J$
to obtain an algebraic inverse to $a^*:H^1(J)\to H^1(C)$.
More explicitly, $w=(b-1)!\,(L_\theta^{b-1})^{-1}(a\times1)_*\xi$ on the indicated mixed component; the inverse is algebraic by the Lefschetz standard conjecture for $J$. Indeed, $a_*a^*x=x\theta^{b-1}/(b-1)!$ for $x\in H^1(J)$. Apply this inverse to the first factor of $\xi$ and project to
$H^1(J)\otimes H^{2m-1}(A)$; the latter projector is algebraic because
both varieties in the product are abelian. This constructs $w$ as an
algebraic class, not merely as a cohomological tensor. It is unique as a
cohomological tensor, and $h\,w=0$.

The dimension of $J\times A$ is $b+2m-1$, and $w$ has codimension $m$.
Every term in $(\theta+h)^bw$ with a positive power of $h$ vanishes.
The remaining term $\theta^bw$ also vanishes because its degree on $J$
is $2b+1$. Thus $w$ is primitive for the product ample class.
Likewise $w^2H^{b-1}=w^2\theta^{b-1}$.
The canonical polarization of the Jacobian gives
\[
 a_*[C]=\frac{\theta^{b-1}}{(b-1)!}
 \quad\hbox{in cohomology}.
\]
Only its pairing with $H^2(J)$ is needed here. In this form it follows
immediately from the identification of the canonical alternating
pairings on $H^1(C)$ and $H^1(J)$: in a symplectic basis
$\theta=\sum_{i=1}^b x_i\wedge y_i$, expansion of
$\theta^{b-1}$ gives $(b-1)!$ times the complementary wedges.
The identification of the pairings is Jacobian autoduality
\cite[Chapter~III, Theorem~6.8 and Summary~6.11]{MilneAV}.
Projection formula now gives \eqref{eq:curve-transfer}.

Apply the Hodge standard conjecture in the exact form of 032,
Corollary~8.4, to the abelian variety $J\times A$, codimension $m$ and
ample class $H$. Its degree exponent is
$(b+2m-1)-2m=b-1$, and its sign is $(-1)^m$.
The factor $(b-1)!$ is positive. The equality assertion follows from
positive definiteness and the injectivity of $a^*$ in degree one.
\end{proof}

\begin{theorem}[JAC-T4.3: cote 44, Theorem 4, on the primitive incidence quotient]
\label{thm:T4}
Accept 032, Corollary~8.4, as in the task specification. For every
abelian variety $A$ of odd dimension $2m-1$ over an algebraically closed
field of any characteristic, and every rational ample class $h$, the
symmetric homomorphism on $P^m(A)_{\rm prim}$ defined by
\[
 (-1)^m d_B
\]
is a rational polarization. Equivalently, with
$\phi_L(x)=t_x^*L\otimes L^{-1}$, the rational divisor class
\begin{equation}\label{eq:T4-divisor}
 \ell_B=\frac{(-1)^m}{2}\,\Delta_B^*d_B
 \quad\in\operatorname{NS}(B)_{\mathbf Q}
\end{equation}
is ample. After multiplication by a positive integer it defines an
actual polarization in every characteristic, including characteristic
$2$.
\end{theorem}

\begin{proof}
Use Lemma~\ref{lem:primitive-summand} to choose a normalized universal
family $z$ and the nondegenerate symmetric incidence correspondence
$d_B$. For a symmetric rigidified divisorial correspondence, the
class $\Delta^*d_B/2$ is the rational N\'eron--Severi class inducing
that correspondence. Indeed, for a line bundle $L$ the associated
biextension class is
$m_B^*c_1(L)-p_1^*c_1(L)-p_2^*c_1(L)$, whose diagonal pullback is
$[2]^*c_1(L)-2c_1(L)=2c_1(L)$.
Thus \eqref{eq:T4-divisor} induces $(-1)^md_B$, which is still
nondegenerate.

Let $C\to B$ be the normalization of any integral curve in $B$.
It is nonconstant, so $C$ has positive genus. Translate it so that a
chosen point maps to the origin; translations do not change
N\'eron--Severi classes. Pulling back $z$ gives $z_C$ on $C\times A$,
whose cohomology is in
$H^1(C)\otimes H^{2m-1}(A)_{\rm prim}$.
Refined pullback and projection formula for the incidence product give
\begin{equation}\label{eq:degree-test}
 \deg_C(\ell_B)=\frac{(-1)^m}{2}
       \deg_{C\times A}(z_C\cdot z_C).
\end{equation}
There are no omitted degree-zero or degree-two parameter components:
the family was projected to the degree-one component before this
calculation. Lemma~\ref{lem:curve-test} makes the right side
nonnegative. Hence $\ell_B$ is nef.

A nef line bundle on an abelian variety whose associated symmetric
homomorphism is an isogeny is ample. To recall a short precise version
of this criterion, clear denominators and call the resulting bundle
$L$. Nondegeneracy gives $\chi(L)\ne0$ and
$\deg\phi_L=\chi(L)^2$, whereas nefness and Riemann--Roch give
$\chi(L)>0$. For an ample $H_0$, the polynomial
$\chi(L+tH_0)=(c_1(L)+tc_1(H_0))^{\dim B}/(\dim B)!$
is positive for all $t\ge0$. Mumford's index theorem therefore gives
index zero; its effective nondegenerate criterion then gives
ampleness \cite[Chapter~III, \S16, Index Theorem, and Chapter~II,
\S6, Application~1]{MumfordAV}. This argument is valid in arbitrary
characteristic. Applying it to \eqref{eq:T4-divisor} proves the theorem.
The zero-dimensional $B$ case is immediate. Rational coefficients
justify all divisions even when the characteristic divides a chosen
integral denominator; no assertion about an integral principal
polarization follows.
\end{proof}

\paragraph{Archive meaning and the exact input.}
This is the primitive reading required by the marginal instruction in
cote~44, scan~8, and the incidence product in scan~9. Cote~16,
scans~34, 36, 38, 40 supplies the Lefschetz language; scans~49--50 places
positivity in the ring of correspondences. The passage from cycle
positivity to an ample divisor is the curve test above; it is not the
unproved assertion that a cohomologically nondegenerate correspondence
is automatically a polarization. The input 032 is used only on
$J(C)\times A$ in Lemma~\ref{lem:curve-test}. No Hodge conjecture for a
non-CM characteristic-zero abelian variety and no new algebraicity
claim for such Hodge classes is inferred.

\paragraph{Sign calibration.}
For an elliptic curve, the normalized divisor
$P=\Delta-E\times\{0\}-\{0\}\times E$ induces
$x\mapsto\mathcal O_E(x-0)=-\phi_{\mathcal O_E(0)}(x)$.
Its self-incidence is $P$, and $\deg\Delta^*P=-2$.
Thus the sign in \eqref{eq:T4-divisor} is indeed $(-1)^m$ under the
stated convention: for $m=1$ it gives the positive elliptic
polarization. The frequently written sign $(-1)^{m-1}$ belongs instead
to the alternating primitive cohomological pairing before the
incidence-to-divisor identification. Confusing these two conventions
would reverse the conclusion already for a curve.

\begin{corollary}[JAC-T4.5: descent over an arbitrary field]
Let $A/F$ be an abelian variety of odd dimension and let the ample
class $h$ be defined over $F$. In the rational descent category of
JAC-P00 the primitive incidence Picard variety and its signed
polarization are defined over $F$. A positive integral multiple on an
abelian representative of this isogeny class is an actual polarization
over $F$.
\end{corollary}

\begin{proof}
Over $\overline F$, the primitive cohomological projector is invariant
under every automorphism over $F$, because it is characterized by the
Lefschetz decomposition for $h$. Its induced endomorphism on the
incidence Picard variety is independent of the algebraic representative
of that projector, by JAC-P01. Thus it commutes with the rational
descent datum of JAC-P00. Galois descent and invariance of homomorphisms
under purely inseparable extensions, as used there, give a rational
idempotent over $F$ and hence its image over $F$ in the isogeny category.
JAC-P00 supplies a normalized rational Poincar\'e family for this image
over $F$. Its self-incidence and diagonal pullback therefore give the
rational divisor class of JAC-T4.3 over $F$. That class is ample after
base change to $\overline F$, and ampleness descends under a field
extension (Stacks Project, Tag 0D2P, Lemma 35.25.6).
Clearing a positive denominator gives an ample line bundle
and its polarization over $F$. This proves the assertion even when
$F$ is imperfect; it makes no claim that a preferred integral model of
the higher Picard variety or a principal polarization descends.
\end{proof}


\subsection{The actual algebraic representative in the threefold case}

\begin{theorem}[JAC-T4.4]
\label{thm:T44}
Let $A$ be an abelian threefold over an algebraically closed field of
arbitrary characteristic, let $h$ be an ample rational divisor class,
and accept 032, Corollary~8.4. Let $\operatorname{Ab}^2(A)$ be the
universal regular algebraic representative for algebraically trivial
codimension-two cycles. The algebraic projector onto
$H^3(A)_{\rm prim}$ induces an idempotent on
$\operatorname{Ab}^2(A)$ in the isogeny category. Write $B_{\rm reg}$
for its image. The canonical quotient homomorphism induces an isogeny
\[
 q:B_{\rm reg}\longrightarrow P^2(A)_{\rm prim}.
                                                        \tag{JAC.T44.1}
\]
The normalized primitive incidence divisor on $B_{\rm reg}$ is a
rational polarization; in the divisor convention of JAC-T4.3 it is
$\Delta^*d_{\rm reg}/2$. This assertion needs neither a universal
integral cycle nor surjectivity of the second Bloch map onto all of
$H^3(A)$.
\end{theorem}

\begin{proof}
We use the following precise known results. The codimension-two
algebraic representative exists, and a positive-degree miniversal
cycle exists for any surjective regular homomorphism:
Achter--Casalaina-Martin--Vial, \emph{Decomposition of the diagonal,
intermediate Jacobians, and universal codimension-2 cycles in positive
characteristic}, Ast\'erisque 455 (2025), \S1.1 and \S1.3, printed
pp.~3--6. Their Proposition~8.1, printed pp.~41--42, states that a
homologically zero correspondence induces zero on an algebraic
representative when the target codimension is $1$, $2$, or the
dimension. It is valid in every characteristic. We use only
codimension $2$.

Let $R=\operatorname{Ab}^2(A)$, with universal regular homomorphism
$\varphi:\operatorname{CH}^2(A)_{\rm alg}\to R(k)$.
Represent the primitive cohomological projector by a rational
correspondence $e$ supported in degree $3$. Proposition~8.1 applied
to $e^2-e$ shows that its action $e_R$ on $R$ is an actual rational
idempotent. Let $\iota:B_{\rm reg}\to R$ denote its inclusion.
Pull a miniversal family on $R\times A$ back by $\iota$ and divide
by its positive degree. Subtract its fibre over the origin. The
result is a rational family whose induced morphism is $\iota$.
Applying $e$ on the $A$ factor preserves this morphism.

We can furthermore arrange that the cohomology of this family is
entirely in $H^1(B_{\rm reg})\otimes H^3(A)_{\rm prim}$ without
altering its induced morphism. Here is an explicit justification
which does not assume that regular homomorphisms are incidence
homomorphisms. Put $b=\dim B_{\rm reg}$ and choose a rational
polynomial $Q$ with
\[
 Q(2)=1,\qquad Q(2^j)=0\quad(0\le j\le2b,\ j\ne1).
\]
Apply $Q(([2]\times1)^*)$ to the family. On cohomology this is the
degree-one projector on the parameter factor. On induced morphisms
to $R$ it acts as $Q(2)$, since the original morphism is a group
homomorphism and pullback by $[2]$ replaces it by $2\iota$.
Thus the induced morphism remains $\iota$. If $b=0$ the theorem is
immediate, so this construction may be made with $b>0$.
Denote the resulting normalized family by $z$.

Form
\[
 d_{\rm reg}=\operatorname{pr}_{12*}(z_{13}z_{23})_{11},
 \qquad \ell_{\rm reg}=\tfrac12\Delta^*d_{\rm reg}.
                                                        \tag{JAC.T44.2}
\]
For the normalization $f:C\to B_{\rm reg}$ of any integral curve,
translate so that a chosen point maps to zero and put
$z_C=(f\times1)^*z$. Its cohomological class is nonzero. Indeed,
if that class were zero, Proposition~8.1, applied with source $C$,
source codimension $1$ and target codimension $2$, would make the
induced homomorphism
$J(C)=\operatorname{Ab}^1(C)\to R$ zero. But this homomorphism
composed with the pointed Abel map is $\iota f$, which is
nonconstant. This is the pointed-family consequence of
Proposition~8.1; no assertion that an absolute constant family is
determined by its cohomology is needed.

JAC-T4.2 with $m=2$ now gives
\[
 \deg_C\ell_{\rm reg}
       =\tfrac12\deg_{C\times A}(z_C^2)>0.
                                                        \tag{JAC.T44.3}
\]
In particular $\ell_{\rm reg}$ is nef. It is also nondegenerate.
Otherwise the reduced identity component of the kernel of its
symmetric homomorphism would be a positive-dimensional abelian
subvariety $K$. The restriction of $\ell_{\rm reg}$ to $K$ has
zero associated homomorphism, hence is algebraically trivial. A
curve in $K$ would then have degree zero, contrary to
(JAC.T44.3). The nef, nondegenerate ampleness criterion used in
JAC-T4.3 therefore proves that $\ell_{\rm reg}$ is ample.

It remains to prove the comparison (JAC.T44.1), rather than assume
it. The incidence quotient map is a regular homomorphism. The
universal property of $R$ supplies a surjection
$R\to P^2(A)$. It commutes with correspondences, so its primitive
restriction $q$ is surjective in the isogeny category. Saito's
incidence identities, Proposition~4.4, identify
\[
 d_{\rm reg}=(q\times q)^*d_{P,\rm prim}
       \quad\hbox{in }\operatorname{NS}(B_{\rm reg}\times
                                      B_{\rm reg})_{\mathbf Q}.
                                                        \tag{JAC.T44.4}
\]
One can also check this identity directly: subtract the pullback
of a normalized Poincar\'e family on $P^2(A)$ from $z$. The
difference has zero incidence family map; testing it against either
family makes the mixed divisor component zero. Thus replacing one
family at a time leaves the $11$ incidence component unchanged.
If $q$ had positive-dimensional kernel, (JAC.T44.4) would make
$\ell_{\rm reg}$ numerically trivial on a curve in that kernel,
contradicting (JAC.T44.3). Its kernel is finite, proving the
isogeny and completing the proof.
\end{proof}

\paragraph{Comparison with existing autoduality results.}
The ACMV monograph proves substantially more integral information
under other hypotheses; those hypotheses must be kept separate.
Theorem~12.12(1), printed p.~68, assumes that the rational second
Bloch map is an isomorphism onto $H^3$. Its parts (2)--(3) use
universal $\operatorname{CH}_0$-triviality to construct a purely
inseparable isogeny or an isomorphism. Proposition~10.3, pp.~50--51,
establishes its standard torsion assumption in characteristic zero
or under the stated decomposition of the diagonal. Proposition~13.1
and Corollary~13.3, pp.~71--75, treat polarization by specialization.
None of these extra hypotheses was asserted here for an abelian
threefold. JAC-T4.4 instead uses 032 on products with curve Jacobians
and proves the rational \emph{primitive} assertion directly.

The monograph records its cup-product sign in Definition~12.1,
Remark~12.2 and Example~12.5, pp.~63--64. Its positivity assertion
also requires the appropriate primitive qualification or vanishing
of lower odd Hodge pieces. The raw full middle incidence form can
have both signs, as JAC-P03 shows. Comparisons of signed formulas
must use the specified divisor and alternating-pairing conventions;
JAC-T4.3 fixes ours by the elliptic-curve calculation.

\paragraph{Archive location.}
This strengthens the codimension-two interpretation of cote~44,
scan~2, Theorems~1--4, with the primitive qualification of scan~8.
It proves an isogeny comparison in this case; it does not identify
the incidence quotient with the universal regular representative in
every codimension or assert a principal polarization.


% Research contribution of /root/h9_height, JAC1964-20261121.
% This is a section fragment, not a standalone document.
% Sources are identified in H9_sources.md. All claims below have an explicit
% logical status. No formal proof certificate is asserted.

\section{The pencil formula: a counterexample and its local repair}

The text of cote 44, scan 4, assumes that the restrictions of the two
cycles are algebraically equivalent to zero. H9 replaces this hypothesis by
homological triviality. These are different hypotheses. Moreover, geometric
integrality of the fibres does not remove higher-codimension vertical
corrections. The following counterexample satisfies even the stronger
condition that the restrictions are rationally equivalent to zero.

\subsection*{H9-C1. A quadric counterexample, with both local contributions}

\textbf{Claim.} For every integer $m\geq2$, over every algebraically closed
field $k$, there is a
Lefschetz pencil on a smooth $2m$-dimensional quadric $X$, all of whose fibres
are geometrically integral, and a codimension-$m$ cycle $Z$ such that
\[
 Z|_{Y}=0\quad\text{in }\operatorname{CH}^{m}(Y),\qquad
 (-1)^m\deg(Z^2)=2.
\]
Here $Y$ is the generic fibre over $K=k(t)$. Consequently every
bilinear pairing on its Chow group gives height zero to the restriction,
whereas the uncorrected expression in H9(a) gives $2$. The claim also
refutes the literal formula of cote 44, scan 4, under just its recorded
geometrically integral-fibre hypothesis.

\textbf{Proof.}
Write
\[
 X=\{uv+\textstyle\sum_{i=1}^{m} w_i z_i=0\}
       \subset\mathbf P^{2m+1}.
\]
Use the pencil of hyperplanes $a u+b v=0$, with parameter $[a:b]$.
Its base is the smooth quadric
\[
 B=\{u=v=0,\ \textstyle\sum w_i z_i=0\}\simeq Q^{2m-2}.
\]
The members with $ab\ne0$ are smooth odd-dimensional quadrics. This remains
true in characteristic two: after eliminating $v$, the equation is
$-t u^2+\sum w_i z_i=0$; the only possible projective singular point
from its partial derivatives would have all $w_i,z_i$ zero, and then
the equation forces $u=0$ since $t\ne0$. The two
remaining members, $u=0$ and $v=0$, are cones over $B$, each with precisely
one ordinary quadratic singularity, at the $v$- or $u$-vertex respectively.
For $m\geq2$ the base $B$ is a smooth integral positive-dimensional quadric,
so its cones are integral. The axis is transverse, and the total space
obtained by blowing up $B$ is smooth. At each singular point the local
equation is $\sum w_i z_i=0$ in $2m$ variables; its polar form is
nondegenerate also in characteristic two. Thus this is a Lefschetz pencil in
the precise sense of SGA 7 II, XVII, 2.2. It lies in the generic orbit of
pencils under the orthogonal group: equivalently its pencil line meets the
dual quadric transversely in two points.

Put
\[
 L_+=\mathbf P\langle w_1,\ldots,w_m\rangle,\qquad
 L_-=\mathbf P\langle w_1,\ldots,w_{m-1},z_m\rangle
       \subset B.
\]
These are maximal $(m-1)$-planes in the two different families on $B$.
Let $P_{\infty,+},P_{\infty,-}$ be their cones with the $u$-vertex and
$P_{0,+},P_{0,-}$ their cones with the $v$-vertex. They are $m$-planes on
$X$. Denote their Chow classes by the same symbols and set
\[
 Z=P_{\infty,+}-P_{\infty,-}.
\]
The two planes in this difference belong to the two different families on
$X$. Replacing the vertex while retaining $L_+$ interchanges these families;
the same holds for $L_-$. This assertion can be checked without a
classification theorem. Quotient by
$\langle w_1,\ldots,w_{m-1}\rangle$ and use the rank-one matrix description
of the quadric surface $uv+w_mz_m=0$, with matrix
$\left(\begin{smallmatrix}u&w_m\\-z_m&v\end{smallmatrix}\right)$.
The two row lines are fibres of one ruling and the two column lines are
fibres of the other. Varying the row (respectively column) gives a
$\mathbf P^1$-family joining $P_{\infty,+}$ to $P_{0,-}$ (respectively
$P_{\infty,-}$ to $P_{0,+}$), after adjoining the common subspace.
Consequently
\[
 [P_{0,+}]-[P_{0,-}]=-Z.
\]

We check the intersection number directly. For an $m$-plane $P\subset X$,
the normal-bundle sequence is
\[
 0\longrightarrow N_{P/X}\longrightarrow
 \mathcal O_{\mathbf P^m}(1)^{\oplus(m+1)}
 \longrightarrow\mathcal O_{\mathbf P^m}(2)\longrightarrow0.
\]
The last map is surjective since $X$ is smooth. Therefore
\[
 [P]^2=[h^m]\frac{(1+h)^{m+1}}{1+2h}
       =\frac{1+(-1)^m}{2}.
\]
Here $[h^m]$ denotes extraction of the coefficient of $h^m$. If $\alpha$
and $\beta$ are the two middle plane classes, a linear section degenerating
to two planes gives $h^m=\alpha+\beta$. Since $\deg(h^m\alpha)=1$,
\[
 \alpha\beta=\frac{1-(-1)^m}{2},\qquad
 (\alpha-\beta)^2=2(-1)^m.
\]
For the degeneration $h^m=\alpha+\beta$, intersect $X$ with the
$\mathbf P^{m+1}$ on which all $z_i$ vanish. Its equation is $uv=0$, and
its components are $P_{\infty,+}$ and $P_{0,+}$; the row--column families
above identify $[P_{0,+}]=\beta$. Thus every Chow relation used in this
calculation has an explicit geometric realization.

On the generic fibre, take the affine parameterization $v=-t u$. Its
equation is $-t u^2+\sum w_i z_i=0$. The restrictions of the two planes
appearing in $Z$ are $L_+$ and $L_-$. They are rationally equivalent over
$K$, not merely after passing to $\overline K$. Indeed, put
$W=\langle w_1,\ldots,w_{m-1}\rangle$ and consider the smooth $K$-conic
\[
 C_t=\{-t u^2+w_mz_m=0\}\subset\mathbf P^2_K.
\]
It is parameterized by
\[
 [r:s]\longmapsto[u:w_m:z_m]=[rs:r^2:t s^2].
\]
The family of projective $(m-1)$-planes
$\mathbf P(W\oplus\ell)$, for the isotropic lines $\ell$ represented by
$C_t$, lies in $Y$. At $[r:s]=[1:0]$ and $[0:1]$ its members are $L_+$
and $L_-$. Since these are $K$-rational points of $C_t\simeq\mathbf P^1_K$,
the two fibre cycles are rationally equivalent. Thus $Z|_Y=0$ as asserted.

We now compute the correction at each singular fibre. Write
$q:\widetilde X=\operatorname{Bl}_{B}X\to X$,
$j:E\hookrightarrow\widetilde X$, and $p:E\to B$. Set
\[
 t_X=q^*Z,\quad \delta_B=[L_+]-[L_-],\quad
 e=j_*p^*\delta_B,
\]
and let $u_\infty,u_0$ be the differences of the strict transforms of
$P_{\infty,\pm}$ and $P_{0,\pm}$, respectively. Then
\[
 u_\infty=t_X-e,\qquad u_0=-t_X-e.                 \tag{H9.1}
\]
For completeness, these identities follow directly from the incidence
model $\widetilde X\subset X\times\mathbf P^1$. On
$P_{\infty,\pm}\times\mathbf P^1$, the incidence equation is $a u=0$;
its divisor is the union of the plane in the fibre $a=0$ and
$L_\pm\times\mathbf P^1$. This is precisely the refined pullback of
$P_{\infty,\pm}$. On $P_{0,\pm}\times\mathbf P^1$ the equation is
$b v=0$, yielding the analogous identity. Taking differences proves
(H9.1), including the coefficients and signs.

The projection formula gives $t_X e=0$ because $q_*e=0$. Since
$\mathcal O_E(E)=\mathcal O_E(-1)$ and $p_*c_1(\mathcal O_E(1))=1$,
\[
 e^2=-\delta_B^2=-2(-1)^{m-1}=2(-1)^m,
 \qquad t_X^2=2(-1)^m.
\]
It follows that
\[
 \begin{array}{c|cc}
  &u_\infty&u_0\\\hline
 u_\infty&4(-1)^m&0\\
 u_0&0&4(-1)^m
 \end{array},\qquad
 t_Xu_\infty=2(-1)^m,\quad t_Xu_0=-2(-1)^m.
                                                        \tag{H9.2}
\]
In particular $t_X=(u_\infty-u_0)/2$ is itself a rational vertical cycle.
Its admissible correction is
\[
 t_X-\tfrac12u_\infty+\tfrac12u_0=0.
\]
Each local Schur-complement term is
$(2(-1)^m)^2/(4(-1)^m)=(-1)^m$. Thus each singular fibre contributes $1$
after multiplication by the H9 sign $(-1)^m$, and the two contributions
cancel the uncorrected value $2$. This proves the counterexample and the
claimed local refinement.\hfill$\square$

This is a mathematical insufficiency of the stated hypothesis, not evidence
that the word ``int\`egres'' was transcribed incorrectly. Without images of
the original manuscript no palaeographic alteration is justified.

\subsection*{H9-T1. The corrected intersection formula}

Let $f:\mathcal X\to C$ be a projective morphism from a smooth projective
$2m$-fold to a smooth projective curve over an algebraically closed field,
with smooth generic fibre $Y$. Use rational Chow groups. Call
$x\in\operatorname{CH}^m(\mathcal X)_{\mathbf Q}$ admissible if
$x|_Y$ is numerically trivial and
\[
 \deg(xv)=0
 \quad\text{for every codimension-$m$ cycle $v$ supported in fibres.}
                                                        \tag{H9.3}
\]
The geometric intersection-height normalization used here is
\[
 h_m(x|_Y,y|_Y)=(-1)^m\deg(xy)                    \tag{H9.4}
\]
for admissible extensions. Its independence of extensions follows from
the localization exact sequence: two extensions of the same Chow class
differ, modulo rational equivalence, by a vertical cycle, and (H9.3)
annihilates this difference. This is the signed degree of the pairing in
Kahn, \emph{Refined height pairing}, Definition 2.2, Proposition 2.8 and
equation (2.9). The terminology ``geometric intersection height'' specifies
the construction; comparison with a particular Bloch biextension or a
N\'eron--Tate pairing is a separate compatibility statement.

Suppose $x_0,y_0$ extend numerically trivial generic-fibre cycles and that
vertical corrections are to be made inside a finite-dimensional space
$W$ spanned by numerical classes of vertical cycles. Let
$R=\operatorname{rad}(W)$ for the intersection form. Assume $x_0$ and $y_0$
are orthogonal to $R$. Choose representatives $v_1,\ldots,v_r$ for a basis
of $W/R$, put $M_{ij}=\deg(v_iv_j)$ and
\[
 a_i=\deg(x_0v_i),\qquad b_i=\deg(y_0v_i).
\]
Then $M$ is invertible, the corrected classes are
\[
 x=x_0-\sum_i(M^{-1}a)_i v_i,\qquad
 y=y_0-\sum_i(M^{-1}b)_i v_i,
\]
and direct bilinear expansion gives their signed intersection. If the
corrected classes are admissible (in particular, if $W$ contains every
vertical numerical class), this is the height:
\[
 h_m(x|_Y,y|_Y)=(-1)^m\left(\deg(x_0y_0)-a^{\!t}M^{-1}b\right).
                                                        \tag{H9.5}
\]
The expression is independent of the chosen representatives modulo $R$.
For a smaller $W$, the same expression still computes
$(-1)^m\deg(xy)$, but is not called a height until orthogonality to the
remaining vertical classes has been verified.
If $W$ contains all the vertical numerical classes, the corrected cycles
satisfy (H9.3). Conversely, existence of a correction in $W$ requires
$x_0\perp R$, because adding an element of $W$ cannot change its pairing
with $R$. The finite-dimensional linear algebra proves this criterion is
also sufficient. Hence (H9.5) states both the formula and its exact
existence gate.

For a pencil with blowup map $q$ and $x_0=q^*Z$, $y_0=q^*Z'$, the
projection formula gives $\deg(x_0y_0)=\deg(ZZ')$. Fibres at different
closed points are disjoint, so the matrix can be organized into local
blocks, after separating the radical. Formula (H9.5) is the requested
local variant, including higher-dimensional vertical cycles; it is not
restricted to irreducible components of fibres. For $m=1$ it recovers the
usual intersection-matrix correction in Shioda, \emph{On the Mordell--Weil
lattices}, Theorem 8.6, equations (8.11)--(8.15), printed p.~228. The
original printed page was rendered and read in this run.

\subsection*{H9-T2. A cohomological hypothesis that removes all corrections}

\textbf{Claim.} Let a Lefschetz pencil on a smooth projective $2m$-fold
$X$ have generic fibre $Y$ and smooth blowup $q:\widetilde X\to X$.
Fix $\ell\ne\operatorname{char}k$. Suppose at each singular fibre $Y_s$
the specialization map
\[
 H^{2m}(Y_s,\mathbf Q_\ell(m))\longrightarrow
 H^{2m}(Y_{\bar\eta_s},\mathbf Q_\ell(m))       \tag{H9.6}
\]
is injective. If $Z,Z'$ restrict to homologically trivial cycles on $Y$,
then $q^*Z,q^*Z'$ are admissible and
\[
 h_m(Z|_Y,Z'|_Y)=(-1)^m\deg(ZZ').               \tag{H9.7}
\]
This statement holds in every characteristic subject to the explicit
injectivity hypothesis (H9.6).

\textbf{Proof.} On the smooth locus of the pencil the fibre classes of
$q^*Z$ form a section of the locally constant sheaf
$R^{2m}f_*\mathbf Q_\ell(m)$. The section is zero at the geometric generic
point, hence zero on this locus. At a singular fibre, compatibility of the
class of a global cycle with specialization and (H9.6) imply
$i_s^*\operatorname{cl}(q^*Z)=0$ in the ordinary cohomology of $Y_s$,
where $i_s:Y_s\hookrightarrow\widetilde X$. For any $m$-dimensional cycle $V$
inside any fibre, the compatibility of the cycle class with cap products
and the trace gives
\[
 \deg(q^*Z\cdot V)=
 \operatorname{Tr}_{V}(i_s^*\operatorname{cl}(q^*Z)\cap[V])=0.
\]
Thus $q^*Z$ is admissible; similarly for $Z'$. Formula (H9.7) follows
from (H9.4) and the projection formula.\hfill$\square$

In the ordinary nodal Picard--Lefschetz setting, a sufficient condition
for (H9.6) is that every local vanishing cycle is nonzero in the middle
cohomology of the nearby smooth fibre. Indeed, the relevant exact sequence
is
\[
 H^{2m-1}(Y_{\bar\eta_s})\longrightarrow\mathbf Q_\ell(-m)
 \longrightarrow H^{2m}(Y_s)
 \longrightarrow H^{2m}(Y_{\bar\eta_s})\longrightarrow0.
\]
The first arrow is pairing with the vanishing cycle. Poincar\'e duality
makes it surjective when that cycle is nonzero, proving injectivity of the
next specialization map. This identifies the additional local condition
missing in H9; the ordinary quadric pencil has zero middle cohomology and
therefore zero vanishing cycles. The precise SGA framework is SGA 7 II,
XVIII, 5.3.2 (printed p.~287), 5.3.4--5.3.5 (p.~288): condition $(A)$ says
that the direct-image sheaves agree with the extensions of their
restrictions to the smooth locus. Proposition 6.2 (p.~319) identifies
$(A)$ with nontrivial inertia action on middle cohomology at every
singular fibre. These original printed pages were independently read in
the programme audit. Using condition $(A)$, rather than an unqualified
quadratic-monodromy convention, makes (H9.6) available in every
characteristic covered by that condition.

Assuming the attached 032 Corollary 8.4, if $X$ is an abelian $2m$-fold,
the right side of (H9.7) is positive definite on its primitive algebraic
cohomology. If $h$ is the ample hyperplane class defining the pencil, put
\[
\begin{split}
 C_{\rm prim}&=\{Z\in\operatorname{CH}^m(X)_{\mathbf Q}:
                          h\operatorname{cl}(Z)=0\},\\
 R&=\operatorname{im}(C_{\rm prim}\to\operatorname{CH}^m(Y)_{\mathbf Q}),\\
 R_0&=\operatorname{im}(\operatorname{CH}^m_{\rm hom}(X)_{\mathbf Q}
                         \to\operatorname{CH}^m(Y)_{\mathbf Q}).
\end{split}
\]
Here $\ker(j^*:H^{2m}(X)\to H^{2m}(Y))=\ker L_h$: indeed
$j_*j^*=L_h$, and $j_*:H^{2m}(Y)\to H^{2m+2}(X)$ is injective by
Poincar\'e duality and weak Lefschetz in degree $2m-2$. Thus every element
of $C_{\rm prim}$ satisfies the restriction hypothesis of H9-T2.
For pencils satisfying (H9.6), $h_m$ on $R$ has radical exactly $R_0$:
the inclusion of $R_0$ in the radical follows because a homologically
trivial ambient class has zero intersection with every algebraic class;
the opposite inclusion follows from (H9.7) and positive definiteness of
the primitive form. In particular the restriction map induces an
isometry
\[
 C_{\rm prim}/\operatorname{CH}^m_{\rm hom}(X)_{\mathbf Q}
       \xrightarrow{\sim}R/R_0
\]
for the signed primitive intersection form on the left and the geometric
height on the right. For injectivity, if a restriction belongs to $R_0$,
subtract a homologically trivial ambient lift; the resulting restriction
is zero, so its height is zero and positivity forces its ambient
cohomology class to vanish. This proves the assertion without pretending
that a cohomology class has a canonical Chow representative. The argument
does not turn homologically trivial restrictions into points of an
algebraic intermediate Jacobian.

\subsection*{H9-T3. The precise quotient statement behind H9(b)}

Here is the valid linear-algebra implication, separating all geometric
inputs. Let $V$ be a rational vector space of cycle classes equipped with
a symmetric form $B$, and let
\[
 a:V\longrightarrow
 (J(K)/\operatorname{Tr}_{K/k}J(k))\otimes\mathbf Q
\]
be a defined homomorphism. Suppose $J$ is polarized and a height comparison
has established
\[
 B(v,w)=\langle a(v),a(w)\rangle_{\rm NT}.
                                                        \tag{H9.8}
\]
Then $B$ is positive semidefinite and
\[
 \operatorname{rad}B=\ker a.
\]
Indeed the right side is positive definite on the Lang--N\'eron space
after quotienting torsion, so $B(v,v)=0$ is equivalent to $a(v)=0$; the
bilinear identity proves the radical statement. Thus $B$ is positive
definite on $V$ if and only if $a$ is injective; it is always positive
definite on $V/\ker a$.

Applied to a pencil, (H9.8) must use the corrected form (H9.5). Replacing
it by the uncorrected primitive intersection form additionally requires
that every local correction vanish. The quadric example has nonzero
primitive classes with zero generic restriction, so injectivity on the
original primitive space fails. This does not by itself refute a bare
logical equivalence of two positivity assertions (both are true for a
quadric); it refutes the claimed identification and its asserted proof.

There is also a distinct domain obstruction. H9 assumes only homological
triviality, whereas Grothendieck's $P^m$ is constructed from algebraically
trivial cycles. SGA 7 II, XX, Theorem 4.1 and Remark 4.3 provide examples
where these restrictions are not algebraically trivial, even with
$X=E^4$: with $a_i$ the divisor class from the $i$th elliptic factor and
$h=\sum a_i$, the class
\[
 \alpha=(a_1-a_2)(a_3-a_4)
\]
satisfies $h\alpha=0$ and $\alpha^2=4$. On sufficiently high-degree
Lefschetz hyperplane sections, its restriction is homologically trivial
but not algebraically trivial. The original locators are printed
pp.~355--356; the full source hypotheses and their verification are
supplied by the independent programme audit. In particular, the regular
homomorphism from algebraically trivial cycles defining $P^m$ cannot be
applied to this cycle. One can ask for other, additional maps or
projections on a larger group, but H9 has not specified or proved such a
construction. Thus its displayed ``image'' in the algebraic $P^m$ of
cote 44 is not defined on the stated domain by the construction cited.

The repaired route must retain three pieces: the admissible height on
homologically trivial cycles; the local vertical form; and, only on the
algebraically representable part, comparison with an abelian
N\'eron--Tate height. The supplied CM-Hodge and standard-conjecture theorems
give positivity of the primitive intersection form for abelian $X$; they
do not remove either of the two explicit counterexamples above.

\subsection*{H9-T4. Polarization of the fixed primitive part of an ample section}

\textbf{Claim.} Assume 032, Corollary 8.4. Let $X$ be an abelian variety
of dimension $2m$ over an algebraically closed field of any characteristic,
let $j:Y\hookrightarrow X$ be a smooth ample divisor, and write
$h=[Y]$, $n=2m-1$. The fixed primitive summand of the incidence higher
Picard variety $P^m(Y)$ has a rational polarization given by its
$(-1)^m$-signed incidence correspondence. Its cohomological subspace is
\[
 F=j^*\ker\bigl(L_h^2:H^{n}(X)\longrightarrow H^{n+4}(X)\bigr)
       \subset H^n(Y).
\]
This proves the relevant part of Theorem 4 even though $Y$ need not be
abelian. The residual summand is the variable part, on which $j_*$ is zero.

\textbf{Proof.}
Hard Lefschetz on $X$ gives an isomorphism
$L_h:H^n(X)\to H^{n+2}(X)$; its inverse $\beta$ is algebraic because
the Lefschetz standard conjecture holds for abelian varieties. Let $e$
be the algebraic primitive projector in degree $n$ on $X$, with image
$\ker L_h^2$, and extend it by zero in other degrees. Define
\[
 r=\beta j_*,\qquad \Pi_F=j^*e\beta j_*.
\]
Since $j_*j^*=L_h$, one has $rj^*=1$ and $\Pi_F^2=\Pi_F$ on
cohomology. The correspondence $\Pi_F$ is algebraic. The incidence
detection lemma from JAC-T4.1 therefore induces an idempotent on
$P^m(Y)$ in the isogeny category.

The decomposition
\[
 H^n(Y)=j^*H^n(X)\oplus\ker j_*
\]
is orthogonal for the middle pairing on $Y$. On its first summand the
pairing is $\int_X uvh$. For $u\in\ker L_h^2$ and
$v=L_hw\in L_hH^{n-2}(X)$ one has
$\int_Xuvh=\int_X(L_h^2u)w=0$. Thus $\Pi_F$ is self-adjoint for
the pairing on $Y$. Saito's transpose compatibility and nondegenerate
incidence duality imply that its image $B_F$ is a nondegenerate
self-dual isogeny factor of $P^m(Y)$. Moreover it is primitive for $Y$:
if $u\in\ker L_h^2$, then
$j_*L_{j^*h}j^*u=L_h^2u=0$, and
$j_*:H^{n+2}(Y)\to H^{n+4}(X)$ is injective by dual weak Lefschetz.

It remains to verify the sign on all curves in $B_F$. Pull a normalized
rational Poincar\'e family back to a smooth pointed curve $C$ mapping to
$B_F$, and apply the parameter degree-one projector. Its class is an
algebraic tensor
\[
 \xi\in H^1(C)\otimes F.
\]
Set $\eta=(1\otimes r)\xi\in H^1(C)\otimes H^n(X)$.
It is algebraic, $j^*\eta=\xi$, and $h^2\eta=0$.
If $b$ is the genus of $C$ and $\theta$ the canonical principal
polarization on $J(C)$, the algebraic inverse to the Abel pullback
lifts $\eta$ to an algebraic class
$w\in H^1(J(C))\otimes H^n(X)$. It satisfies $h^2w=0$.
For $H=\theta+h$ on the abelian variety $J(C)\times X$,
\[
 H^{b+1}w=0.
\]
Indeed terms with $h$-exponent at least two vanish, while terms with
$h$-exponent zero or one have excessive degree on $J(C)$.
Consequently $w$ is primitive of codimension $m$. In its quadratic form
only one term of $H^b$ survives, and the Abel curve-class formula yields
\[
 \begin{aligned}
 \int_{J(C)\times X}w^2H^b
 &=b\int_{J(C)\times X}w^2\theta^{b-1}h\\
 &=b!\int_{C\times X}\eta^2h
 =b!\int_{C\times Y}\xi^2.
 \end{aligned}                                                \tag{H9.9}
\]
The factor $b!$ is positive. By 032, Corollary 8.4, the last intersection
has sign $(-1)^m$ and is zero only if $\xi=0$. The same curve test and
nef-plus-nondegenerate criterion as in JAC-T4.3 therefore prove that
the signed incidence class on $B_F$ is ample. This finishes the proof.
\hfill$\square$

The projectors used here implement the hyperplane formula of cote 16,
scan 40; the cohomological fixed/variable decomposition belongs to
SGA 7 II, XVIII, \S5. The input 032 is applied only to the genuinely
abelian product $J(C)\times X$, never to $J(C)\times Y$.

\subsection*{H9-T5. A finite-field Newton witness eliminates the remaining part}

\textbf{Claim.} Let a Lefschetz pencil of ample sections of an abelian
$2m$-fold be defined over a finite field of characteristic $p$. Write
$U$ for its smooth locus and $\mathcal V_\ell$ for its variable
$R^{2m-1}$-local system. Assume that $\mathcal V_\ell$ is nonzero and
remains absolutely irreducible under every finite-index restriction of
geometric monodromy, for a fixed prime $\ell\ne p$. Let
$B_{\rm var}$ be the variable incidence factor of $P^m(Y)$ over the
generic field, after a finite extension making the factor and universal
cycle defined. If one smooth closed fibre has a Frobenius eigenvalue on
its variable cohomology with normalized $p$-adic slope outside
$[m-1,m]$, then $B_{\rm var}=0$. In that case H9-T4 proves the primitive
Theorem 4 for the full incidence quotient of the generic fibre.

\textbf{Proof.}
The algebraic correspondence $1-j^*L_h^{-1}j_*$ induces the variable
idempotent on the incidence quotient. It is self-adjoint, so its image
$B_{\rm var}$ has nondegenerate incidence duality. A normalized Poincar\'e
cycle gives, for every $\ell\ne p$, an algebraic correspondence map
\[
 V_\ell B_{\rm var}\longrightarrow\mathcal V_\ell(m)
                                                        \tag{H9.10}
\]
over a finite cover of $U$. This map is injective: a vector in its
kernel pairs to zero with every other vector under the pullback of the
middle pairing, whereas that pullback is the realization of the
nondegenerate incidence isogeny. A choice of polarization on
$B_{\rm var}$ identifies $V_\ell B_{\rm var}$ with
$H^1(B_{\rm var},\mathbf Q_\ell)(1)$. Thus (H9.10) is equivalently
\[
 H^1(B_{\rm var},\mathbf Q_\ell)(-(m-1))
       \lhook\joinrel\longrightarrow\mathcal V_\ell.       \tag{H9.11}
\]
If $B_{\rm var}$ is nonzero, the irreducibility hypothesis, which survives
the finite field-of-definition extension, forces (H9.11) to be an
isomorphism.

At every point of the finite cover lying above a smooth fibre, the
right side is unramified. Therefore the left side is unramified and
the N\'eron--Ogg--Shafarevich criterion gives good reduction of
$B_{\rm var}$ there. Smooth proper base change specializes (H9.11).
If the residue field has cardinality $q$, the Frobenius eigenvalues on
$H^1$ of this reduction are algebraic integers $\gamma$ for which
$q/\gamma$ is also an algebraic integer. Hence, at every embedding into
$\overline{\mathbf Q}_p$,
\[
 0\leq\frac{v_p(\gamma)}{v_p(q)}\leq1.
\]
The Tate twist $-(m-1)$ multiplies geometric-Frobenius eigenvalues by
$q^{m-1}$. Thus all slopes on the variable part must belong to
$[m-1,m]$. Passing to a finite extension of a residue field raises
each eigenvalue and $q$ to the same positive power and leaves these
normalized slopes unchanged. A single smooth original fibre with a
slope outside the interval therefore contradicts $B_{\rm var}\ne0$.
This proves the claim.\hfill$\square$

The geometric irreducibility hypothesis is the standard monodromy
condition studied in SGA 7 II, XVIII, 6.7.2; it is stated explicitly
here so that no exceptional pencil is silently included. The good
reduction criterion and the integrality/duality of Frobenius on abelian
varieties are background inputs. This proof does not assert that every
high-degree pencil has the required Newton witness. Producing that
witness is the residual concrete gate of this method; purity of the
complex absolute values alone is insufficient to do so.


% Whole-proof self-check and independent programme-audit reading completed.
% Root integration review remains a separate check, recorded in its own ledger.

\subsection*{H9-T6. The variable incidence sign from integral monodromy}

The Newton test in H9-T5 is a useful way to prove that the variable
incidence factor is zero. The following argument instead determines its
sign when it is nonzero. It uses Grothendieck's integral monodromy pairing,
not a sign assigned to one $\ell$-adic scalar.

\textbf{Hypotheses.} Let $Y$ be the generic fibre of an odd-dimensional
Lefschetz pencil of fibre dimension $n=2m-1$. Suppose its middle
cohomology has an algebraic self-adjoint variable projector, and let
$B$ be the corresponding incidence factor of $P^m(Y)$. Assume:
\begin{enumerate}
\item for every $\ell\ne\operatorname{char}k$, the variable local system
      $\mathcal V_\ell$ is nonzero and absolutely irreducible after every
      finite-index restriction of geometric monodromy;
\item at one ordinary quadratic singular fibre the vanishing cycle in
      $\mathcal V_\ell$ is nonzero for every such $\ell$;
\item $B$ has its nondegenerate symmetric incidence duality and a
      normalized rational Poincar\'e cycle, as in Saito's construction.
\end{enumerate}
For the ample-section pencils considered in H9-T4, the algebraic variable
projector is $1-j^*L_h^{-1}j_*$ on the relevant odd degree. The first two
hypotheses are the precise Picard--Lefschetz monodromy hypotheses; their
standard generic-pencil instances are in SGA 7 II, XVIII, 6.7.2.

\textbf{Claim.} Under these hypotheses the rational divisor class
\[
 \ell_m=\frac{(-1)^m}{2}\Delta_B^*d_B
\]
is a rational polarization of $B$, in every characteristic. Consequently,
under the same monodromy hypotheses, H9-T4 and H9-T6 together prove the
primitive Theorem 4 for the full incidence quotient of a generic ample
section of an abelian variety, accepting 032 Corollary 8.4 for its fixed
part.

\paragraph{H9-L1: a rational square detected at almost all primes.}
If $r\in\mathbf Q^*$ is a square in $\mathbf Q_\ell$ for all but
finitely many primes $\ell$, then $r$ is a square in $\mathbf Q$.
In particular $r>0$.

\emph{Proof.} Otherwise the quadratic field $\mathbf Q(\sqrt r)$
defines a nonprincipal quadratic Dirichlet character $\chi$ of some
conductor. At primes outside the conductor, local squareness implies
$\chi(\ell)=1$, so only finitely many Euler factors of $L(s,\chi)$
differ from those of $\zeta(s)$, for real $s>1$. The quotient
$L(s,\chi)/\zeta(s)$ therefore has a finite positive nonzero limit as
$s\downarrow1$. On the other hand, a nonprincipal periodic character
has bounded partial sums and zero sum over a period. Summation by parts
shows that $\sum_{n\ge1}\chi(n)n^{-s}$ stays bounded as $s\downarrow1$
(indeed converges at $s=1$), whereas $\zeta(s)\to+\infty$.
This is a contradiction. No Riemann hypothesis is used.\hfill$\square$

\paragraph{H9-L2: the rank-one monodromy coefficient.}
Let an abelian variety have semistable reduction over a strictly
henselian trait, and suppose the toric rank is one. Let $L_0$ be a
genuine polarization and let $E_0$ be the alternating form on the Tate
module defined by $c_1(L_0)$. Choose a primitive generator $v_0$ of the
toric cocharacter lattice, and choose the same generator of
$\mathbf Q_\ell(1)$ in the monodromy operator and in the values of $E_0$.
Then the tame monodromy logarithm has the form
\[
 N_\ell(x)=-q_0 E_0(x,v_0)v_0,\qquad q_0\in\mathbf Q_{>0},
                                                        \tag{H9.12}
\]
where the rational number $q_0$ is independent of $\ell$.

\emph{Derivation from the cited theorem.} SGA 7 I, IX, 9.2, printed
pp.~419--421, identifies
the monodromy operator with the scalar extension of the integral
pairing between the two toric character lattices. Theorem 10.4(a),
printed p.~444, asserts that this integral pairing is common to all
primes; 10.4(b) asserts positivity after composing with a genuine
polarization. To specify the lattice conventions, use the underlined
$M,M'$ of IX, 5.5 and 10.4: $M$ is the character lattice of the torus
of the dual abelian variety, whereas $M'$ is the character lattice
of the original torus. (The ordinary $M,M'$ used temporarily in
IX, 9.2 occur in the opposite order.) In rank one choose generators
$a\in M$, $b\in M'$ and the dual toric cocharacter $v_0$, and write
$u=u(a,b)$ and $\lambda_t(a)=s b$ for the polarization map.
Positivity in 10.4(b) is exactly $us>0$.
For a rational lift $y$ of the corresponding quotient generator,
the canonical identifications give $N(y)=u v_0$ and
$E_0(v_0,y)=s$, so
\[
 q_0=u/s=(us)/s^2>0.
\]
This proves rationality and positivity simultaneously; it never interprets
one $\ell$-adic square as a real square. The sign in (H9.12) is the minus
sign from alternating the arguments of $E_0$. Here $E_0$ is defined
by $c_1(L_0)$, or equivalently by the biextension
$m^*L_0\otimes p_1^*L_0^{-1}\otimes p_2^*L_0^{-1}$ of
VIII, 4.12.1, rather than by an unlabelled Cartier-duality pairing.
The anticommutative comparison in VIII, 2.3.11 must be respected.
Our convention agrees with the curve
case: a one-node Jacobian has $N(y)=v_0$ with
$E_0(v_0,y)=1$. The compatibility of this curve normalization with
the general monodromy pairing is explicitly SGA 7 I, IX, 12.5 and 12.7.
Under a ramified extension of traits of index $e$, the monodromy pairing
is multiplied by $e>0$ (IX, 10.3.5, printed p.~443).

\paragraph{Proof of H9-T6.}
If $B=0$ there is nothing to prove, so suppose $B\ne0$. Work over one
finite extension of the generic field on which $B$, the normalized
Poincar\'e cycle and a genuine polarization $L_0$ are defined. Let
\[
 f_\ell:V_\ell B\longrightarrow\mathcal V_\ell(m)
\]
be the map of realizations given by this cycle. As in H9-T5,
nondegenerate incidence duality makes $f_\ell$ injective. The
irreducibility hypothesis makes it an isomorphism, for every
$\ell\ne\operatorname{char}k$.

Every geometric endomorphism of $B$ is defined over a finite extension.
Tate-module faithfulness gives
\[
 \operatorname{End}^0(B_{\overline K})\otimes\mathbf Q_\ell
   \lhook\joinrel\longrightarrow
 \operatorname{End}_{G'}(V_\ell B)
\]
after a sufficiently small finite-index monodromy subgroup $G'$.
Absolute irreducibility makes the right side $\mathbf Q_\ell$.
It follows that $\operatorname{End}^0(B_{\overline K})=\mathbf Q$.
Consequently the nondegenerate symmetric class $\ell_m$ is
\[
 \ell_m=c L_0\quad\text{in }\operatorname{NS}(B)_{\mathbf Q},
 \qquad c\in\mathbf Q^* .                              \tag{H9.13}
\]
The task is to determine the sign of this single rational number.

Let $Q_\ell$ denote the middle cup-product pairing on
$\mathcal V_\ell(m)$, with values in $\mathbf Q_\ell(1)$.
The incidence formula gives the following exact sign:
\[
 E_{\ell_m}=(-1)^{m-1}f_\ell^*Q_\ell=cE_0.            \tag{H9.14}
\]
To verify it, write the mixed class of the universal cycle as
$z=\sum_i e_i\otimes u_i$, with $e_i\in H^1(B)$ and
$u_i\in H^n(Y)(m)$. In the product $z_{13}z_{23}$ the odd-degree
class $u_i$ crosses the degree-one parameter class $e_j$, producing
a minus sign. Thus
\[
 d_B=-\sum_{i,j}Q_\ell(u_i,u_j)e_i\boxtimes e_j.
\]
Pulling back to the diagonal and dividing by two shows that the
alternating form of $\Delta^*d_B/2$ is $-f_\ell^*Q_\ell$.
Multiplying by $(-1)^m$ proves (H9.14).

At the chosen singular fibre, take a finite trait extension that makes
the abelian variety semistable and the toric lattice split. Let its
ramification index relative to the original pencil parameter be $e>0$.
The nonzero Picard--Lefschetz operator has rank one. Hence the toric rank
of $B$ is one: by IX, 9.2 and the nondegeneracy in 10.4, rank of monodromy
equals rank of the toric character lattice.

SGA 7 II, XV, 3.3.3--3.3.6, printed pp.~191--192, supplies a vanishing generator
$\delta_\ell$, normalized up to sign by the primitive quotient of the
exceptional quadric. With the same tame-character convention as in
(H9.12), its formula, in relative dimension $2m-1$, is
\[
 N_\ell(x)=e(-1)^m Q_\ell(x,\delta_\ell)\delta_\ell.
                                                        \tag{H9.15}
\]
Here twists are suppressed only after choosing a generator of
$\mathbf Q_\ell(1)$. More intrinsically, XV defines the untwisted
vanishing generator in $H^{2m-1}(Y)(m-1)$; tensoring it with that
generator puts it in $H^{2m-1}(Y)(m)$, and dividing both cup-product
values and the tame character by that same generator gives (H9.15).
Thus no unrecorded unit is introduced into the coefficient.

Put $v_\ell=f_\ell^{-1}\delta_\ell$. Since $f_\ell$ intertwines
monodromy, $v_\ell$ spans the toric line, so
$v_\ell=a_\ell v_0$ for some $a_\ell\in\mathbf Q_\ell^*$.
Combining (H9.14) and (H9.15) gives
\[
 N_\ell(x)=-e c a_\ell^2 E_0(x,v_0)v_0.
\]
Comparison with (H9.12), using $E_0$ nondegenerate, proves
\[
 a_\ell^2=\frac{q_0}{e c}\quad\text{for every }\ell\ne
 \operatorname{char}k.                                \tag{H9.16}
\]
The right side is the same nonzero rational number for all these primes.
H9-L1 therefore says it is a rational square and in particular positive.
Since $q_0>0$ and $e>0$, one obtains $c>0$. Equation (H9.13) now proves
that $\ell_m$ is a rational polarization.

For the final assertion, the fixed primitive and variable incidence
summands are orthogonal and their dualities are block diagonal.
H9-T4 polarizes the former and the argument just completed polarizes
the latter. A product of the two rational polarizations is ample,
and ampleness descends under the isogeny giving the decomposition.
This proves the claimed extension of Theorem 4.\hfill$\square$

\paragraph{Scope.}
The result concerns the geometric generic fibre under explicit monodromy
hypotheses. It does not assert positivity on every special fibre of a
pencil, where new incidence factors can appear. Nor does it identify
the height of a merely homologically trivial generic cycle with a
N\'eron--Tate height on $P^m$: the $E^4$ domain obstruction remains.
The original H9(b) therefore does not follow from this extension.


% Follow-up to H9-T6, independently checked by /root/programme_audit.
% Run JAC1964-20261121; GPT-6 Astra Pro; effort setting not exposed.

\subsection*{H9-T6a. A concrete generic-pencil polarization theorem}

\begin{theorem}[H9-T6a]
Let $X$ be smooth, projective and connected of even dimension $2m$
over an algebraically closed field $k$ of arbitrary characteristic.
Assume that the Lefschetz lowering operator of $X$ is algebraic, and
let $Y_\eta$ be the geometric generic fibre of a Lefschetz pencil of
hyperplane sections for a projective embedding of $X$.
The variable incidence factor of $P^m(Y_\eta)$ is polarized by its
$(-1)^m$-signed incidence correspondence.

If $X$ is abelian, then, accepting 032 Corollary~8.4, the same statement
holds for the entire primitive factor of $P^m(Y_\eta)$.
If $X$ is a smooth complete intersection in projective space, it holds
for all of $P^m(Y_\eta)$, without requiring 032.

For any starting projective embedding, a general pencil of degree-$d$
hypersurface sections is a Lefschetz pencil for every $d\ge2$.
Consequently these assertions give concrete families in every
characteristic, including characteristic two.
\end{theorem}

\begin{proof}
Write $n=\dim X=2m$ and $h$ for the hyperplane class of the pencil.
Hard Lefschetz gives a nondegenerate ambient summand
$j^*H^{n-1}(X)$ in $H^{n-1}(Y_\eta)$.  Its orthogonal complement is
the variable space $E_\ell=\ker j_*$.  The algebraic correspondence
\[
 1-j^*(L_h:H^{n-1}(X)\longrightarrow H^{n+1}(X))^{-1}j_*
\]
is the self-adjoint projector onto $E_\ell$ on the relevant degree.
Lemma~\ref{jac:P01} and Saito's transpose compatibility give a
nondegenerate incidence factor $B$ and a normalized rational
Poincar\'e cycle.  If $B=0$, its polarization assertion is empty.
Suppose $B\ne0$.

The incidence formula makes the realization map
$V_\ell B\to E_\ell(m)$ injective for every invertible prime $\ell$,
so $E_\ell\ne0$ for every such prime.  SGA~7~II, XVIII,
Theorem~6.6, printed pp.~323--325, says that the local vanishing
cycles span $E_\ell$ and are conjugate up to sign.  Thus every local
vanishing cycle is nonzero, for every $\ell$.  In particular the
critical set is nonempty and one can choose the same critical point
for all primes.  Corollary~6.7.2, printed p.~326, says that every
finite-index subgroup of geometric monodromy acts absolutely
irreducibly on $E_\ell$.  Its hypotheses are satisfied: $X$ is smooth
projective, the pencil is Lefschetz, hard Lefschetz holds for $X$, and
$n$ is even.  The exception in Theorem~6.6 is the conjunction
``$n$ odd and characteristic two'', which is absent here.

These are exactly the hypotheses used in H9-T6.  Its all-primes
monodromy argument therefore proves that the signed variable class
is a rational polarization.  A separate condition~A assumption was
unnecessary: nonzero $B$ itself supplied the nonzero variable
cohomology and hence the nonzero local vanishing cycles.

If $X$ is abelian, H9-T4 polarizes the fixed primitive summand by 032
on $J(C)\times X$.  The primitive fixed and variable projectors are
orthogonal, so their polarizations combine and descend under the
isogeny decomposition.  If $X$ is a smooth complete intersection of
even dimension, $H^{n-1}(X)=0$, so there is no fixed odd summand.
Moreover $H^{n+1}(Y_\eta)=0$ by dual weak Lefschetz, so all of its
middle odd cohomology is primitive.  Algebraicity of the lowering
operator for a smooth complete intersection is supplied explicitly
in JAC-P08 below.  Thus the variable conclusion is the entire
assertion in this case.

Finally, SGA~7~II, XVII, Theorem~2.5.1, printed p.~217, states that
every degree-$d$ multiple of a projective embedding, for $d\ge2$,
is a Lefschetz embedding.  The theorem has no characteristic
restriction.  The term ``quadratic ordinary'' includes the usual
characteristic-two distinction, but in our even ambient dimension
the local quadratic singularities are nondegenerate: XVII, p.~219
explicitly notes this when the ambient dimension is even or the
characteristic is not two.  This verifies that the odd-fibre
Picard--Lefschetz formula used in H9-T6 applies here as stated.
\end{proof}

\paragraph{The finite-index step can be checked without a density theorem.}
The proof of XVIII~6.7.2, printed p.~327, is short.  With the
alternating middle pairing, each local inertia generator acts by
$T_i=1+N_i$, where $N_i^2=0$ and
$\operatorname{im}N_i=\mathbf Q_\ell\delta_i$.
For any finite-index subgroup $G'$ a positive power $T_i^{a_i}$
belongs to its image, and
$T_i^{a_i}-1=a_iN_i$.  Since $a_i$ is nonzero in $\mathbf Q_\ell$,
a subspace stable under $G'$ is stable under every $T_i$.
The inertia groups topologically generate the tame monodromy group,
so it is stable under the full group.  Given a nonzero vector, the
nondegenerate pairing and the spanning vanishing cycles produce
some $N_i$ with nonzero image; conjugacy then gives every vanishing
cycle, hence the entire space.  The argument survives scalar
extension to $\overline{\mathbf Q}_\ell$, proving absolute
irreducibility.  No orthogonal-monodromy assertion is used here.

\begin{proposition}[JAC-P07: condition A is automatic for abelian pencils]
\label{jac:P07}
Let $X$ be an abelian variety of even dimension $n=2m\ge2$ over an
algebraically closed field, and let $D$ be a genuine Lefschetz pencil
for a projective embedding of $X$.  Let $L$ be its ample section
line bundle, $h=c_1(L)$, and $Y$ a smooth member.  For every prime
$\ell$ invertible in the field, the variable middle cohomology has
dimension
\[
 \dim E_\ell=\deg h^n-C_n,
 \qquad C_n=\frac1{n+1}\binom{2n}{n}.
                                                        \tag{JAC.P07.1}
\]
This dimension is strictly positive.  Every local vanishing cycle
is nonzero and the pencil satisfies condition~A of SGA~7~II,
XVIII~5.3, in every characteristic.  In particular the
specialization map $H^n(Y_s,\mathbf Q_\ell)\to
H^n(Y_{\overline\eta_s},\mathbf Q_\ell)$ at every singular fibre
is injective.
\end{proposition}

\begin{proof}
The tangent bundle of $X$ is trivial.  The Euler-characteristic
formula SGA~7~II, XVIII~6.4.2.3, printed p.~321, therefore gives
\[
 \chi(Y,\mathbf Q_\ell)=(-1)^{n-1}\deg h^n=-\deg h^n.
\]
Weak Lefschetz, Poincar\'e duality on $Y$, and
$b_i(X)=\binom{2n}{i}$ give
\[
 b_{n-1}(Y)=\deg h^n+
       2\sum_{i=0}^{n-2}(-1)^i\binom{2n}{i}
 =\deg h^n+2\binom{2n-1}{n-2}.
\]
The ambient subspace has dimension $b_{n-1}(X)=\binom{2n}{n-1}$.
Subtracting and using Pascal's identity yields (JAC.P07.1).
All identities are uniform in $\ell$.

Riemann--Roch on $X$ gives $\deg h^n=n!\chi(L)$, and for an
ample line bundle on an abelian variety $\chi(L)=h^0(L)$ is a
positive integer.  If $n\ge4$, then $n!>C_n$.  One may check this
without an estimate: $3!=6>C_3=5$, and
\[
 \frac{C_{r+1}}{C_r}=\frac{2(2r+1)}{r+2}<4\le r+1
 \qquad(r\ge3),
\]
so the inequality propagates.  If $n=2$, a genuine pencil requires
two independent sections of $L$, hence $h^0(L)\ge2$ and
$\deg h^2=2h^0(L)\ge4>C_2=2$.  This proves strict positivity
in every even dimension under consideration.

By XVIII~6.6.1--6.6.2, the local vanishing cycles span $E_\ell$
and are conjugate up to sign.  Since $E_\ell\ne0$, none of these
cycles is zero, and there is at least one critical fibre.
Proposition~6.2, printed p.~319, identifies this condition with
condition~A.  The excluded parity/characteristic combination in
XVIII~6.6 is odd ambient dimension in characteristic two; our
ambient dimension is even.  Thus the conclusion includes
characteristic two.

For clarity, the specialization consequence can also be seen from
the vanishing-cycle sequence at an ordinary quadratic fibre.
Its potentially nonzero connecting portion is
\[
 H^{n-1}(Y_{\overline\eta_s})\longrightarrow
 \mathbf Q_\ell(-m)\longrightarrow H^n(Y_s)
 \longrightarrow H^n(Y_{\overline\eta_s})\longrightarrow0.
\]
The first arrow is pairing with the nonzero local vanishing cycle,
up to the fixed nonzero normalization, and is surjective.
Consequently the final specialization map is injective (indeed an
isomorphism).  These are exactly the local hypotheses used in
H9-T2; integrality of the singular fibres alone would not suffice.
\end{proof}

\begin{proposition}[JAC-P08: further ambient varieties for Theorem~6]
\label{jac:P08}
The rational hyperplane and trace conclusions of JAC-P02 apply to
smooth complete intersections and to products of smooth complete
intersections, curves, and abelian varieties, in every characteristic.
For an even-dimensional smooth complete intersection $X$, the
critical incidence variety $P^m(Y_\eta)$ has zero $K/k$-trace up to
isogeny.  Products with nonzero ambient odd cohomology give further
instances of the nonzero fixed-part equality in JAC-P02.
\end{proposition}

\begin{proof}
For a smooth complete intersection $X$ of dimension $n$, with
hyperplane class $h$ and $e=\deg h^n>0$, weak Lefschetz and duality
show that its cohomology outside the middle degree consists of the
one-dimensional spaces generated by $h^r$ in degree $2r$.
The normalized lowering operator is the algebraic correspondence
\[
 \Lambda_X=\frac1e\sum_{r=1}^{n}
           h^{n-r}\times h^{r-1}
 \quad\in\mathrm{CH}^{n-1}(X\times X)_{\mathbf Q}.
\]
It sends $h^r$ to $h^{r-1}$.  It kills the middle primitive
cohomology: this is automatic by parity in odd dimension, and in
even dimension follows from its orthogonality to $h^{n/2}$.
The same elementary construction applies to curves.  For abelian
varieties algebraicity of the lowering operator is the known
Lefschetz standard theorem used throughout this paper.

For completeness, algebraic Lefschetz operators are stable under
products.  Replace each normalized lowering operator by the usual
$\mathfrak{sl}_2$ lowering operator, which acts on a primitive string
by $L^av\mapsto a(d-i-a+1)L^{a-1}v$.  The degree projectors and
primitive-string projectors are algebraic, so this replacement is
algebraic.  Explicitly, put $P=1-L\Lambda$. On primitive vectors the
operator
$P(\Lambda^rL^r-\Lambda^{r+1}L^{r+1})P$ selects precisely strings
of length $r+1$; conjugating these selectors by $L^a$ and
$\Lambda^a$ selects each level. This also constructs the needed
grading projectors directly from the two algebraic operators.
On a product the sums of the raising, lowering, and
grading operators give the tensor-product $\mathfrak{sl}_2$ action.
Its Casimir and grading have finite integer spectra.  Rational
polynomial interpolation in these commuting operators rescales the
nonzero lowering coefficients to one, and gives zero on primitive
vectors; hence the normalized product lowering operator is again
algebraic.  Applying JAC-P02 proves the assertions.

In the even complete-intersection case all odd cohomology of $X$
vanishes.  Lemma~\ref{jac:P01}, applied to the identity correspondence,
then gives $P^p(X)=0$ for every $p$.  The critical trace equality
of JAC-P02 therefore gives zero trace for $P^m(Y_\eta)$, even when
that incidence variety itself is nonzero.  For example, this is a
substantive assertion for generic curve sections of smooth surfaces
and for odd-dimensional generic complete intersections carrying a
nonzero middle incidence Picard variety.
\end{proof}

For a degree-$d$ hypersurface section relative to a starting class
$h$, formula (JAC.P07.1) reads
$\dim E_\ell=d^n\deg h^n-C_n$.  The dimension calculation is not
needed for the variable-polarization argument in H9-T6a, which can
start from $B\ne0$.  It is needed to discharge condition~A uniformly
in the geometric height theorem, including when $B=0$.


% Further H9 contribution: full generic height group, not only ambient cycles.
% New deduction in this run; no historical originality or formal check claimed.

\subsection*{H9-T7. The whole geometric height group and the axis contribution}

Let $X$ be an abelian $2m$-fold over an algebraically closed field $k$,
let $h$ define a Lefschetz pencil with smooth axis
$\kappa:B\hookrightarrow X$, and let
\[
 q:\widetilde X=\operatorname{Bl}_B X\longrightarrow X,
 \qquad f:\widetilde X\longrightarrow\mathbf P^1
\]
be its blowup and pencil morphisms. Write $e:E\hookrightarrow\widetilde X$
and $p:E\to B$ for the exceptional divisor maps. Put $K=k(\mathbf P^1)$
and let $Y/K$ be the generic fibre, with maps
$i:B_K\hookrightarrow Y$ and $j:Y\hookrightarrow X_K$.
Fix one $\ell\ne\operatorname{char}k$ and use its cycle-class map.
Assume the specialization injectivity of H9-T2 at every singular fibre;
in particular SGA 7 condition $(A)$ suffices.
For the genuine abelian Lefschetz pencils in this statement the
assumption is automatic by JAC-P07. It is kept visible below to
identify the precise geometric input to the proof.

Let $\operatorname{Alg}^r(T)$ denote the image of $\operatorname{CH}^r(T)_{\mathbf Q}$
in $H^{2r}(T,\mathbf Q_\ell(r))$. Define the rational vector spaces
\[
 \begin{split}
 A_X&=\ker\bigl(L_h:\operatorname{Alg}^m(X)\longrightarrow H^{2m+2}(X)(m+1)\bigr),\\
 A_B&=\ker\bigl(\kappa_*:\operatorname{Alg}^{m-1}(B)
                       \longrightarrow H^{2m+2}(X)(m+1)\bigr),\\
 \mathcal H&=\operatorname{CH}^m(Y)_{\rm hom,\mathbf Q},\\
 \mathcal H_0&=\operatorname{im}\bigl(
       \operatorname{CH}^m(\widetilde X)_{\rm hom,\mathbf Q}
           \longrightarrow\operatorname{CH}^m(Y)_{\mathbf Q}\bigr).
 \end{split}                                                   \tag{H9.17}
\]
The Tate twists in the first two target spaces only record the degree of
the Lefschetz or Gysin operation.

\textbf{Claim.} Every element of $\mathcal H$ has an admissible extension.
The signed geometric height descends to $\mathcal H/\mathcal H_0$, and
there is a canonical isomorphism of rational spaces with bilinear forms
\[
 \begin{split}
 \Theta:A_X\oplus A_B&\xrightarrow{\ \sim\ }
                          \mathcal H/\mathcal H_0,\\
 ([Z],[w])&\longmapsto[j^*Z+i_*w],\\
 Q((z,w),(z',w'))&=(-1)^m\deg_X(zz')
                    +(-1)^{m-1}\deg_B(ww').
 \end{split}                                                   \tag{H9.18}
\]
Every class in $A_B$ is primitive on $B$ for $h_B=\kappa^*h$.
No positivity conjecture for $B$ is used to prove this isomorphism.
Accepting 032 Corollary 8.4, semipositivity of the height on the entire
group $\mathcal H$ is therefore equivalent to semipositivity of
$(-1)^{m-1}\deg_B(ww')$ on this specified subspace $A_B$.
If that form is positive definite, the radical of the full height is
exactly $\mathcal H_0$.

\paragraph{Proof: admissibility and the blowup decomposition.}
Any cycle on $Y$ spreads to a cycle on $\widetilde X$ by taking its closure.
For a homologically trivial generic cycle, the argument of H9-T2 applies
to the cohomology class of any such extension, not only to classes pulled
back from $X$. On smooth fibres its restriction is zero; injectivity of
specialization forces its restriction to every singular fibre to be zero
as well. Pairing with every vertical cycle by cap product shows that the
extension is admissible. A homologically trivial global extension pairs
to zero with every global cycle, so $\mathcal H_0$ lies in the height
radical.

The Chow blowup formula for a smooth centre of codimension two
(Fulton, \emph{Intersection Theory}, Proposition 6.7(e)) gives
\[
 \operatorname{CH}^m(\widetilde X)_{\mathbf Q}
   =q^*\operatorname{CH}^m(X)_{\mathbf Q}
      \oplus e_*p^*\operatorname{CH}^{m-1}(B)_{\mathbf Q}.        \tag{H9.19}
\]
Its cohomological analogue has the same two summands. On the generic
fibre, $e_*p^*w$ restricts to $i_*w$. Hence every $\zeta\in\mathcal H$
has an expression
\[
 \zeta=j^*Z+i_*w.
\]
The Gysin map $j_*:H^{2m}(Y)\to H^{2m+2}(X)$ is an isomorphism by
Poincar\'e duality and weak Lefschetz in degree $2m-2$. Thus the
homological triviality of $\zeta$ is exactly
\[
 h[Z]+\kappa_*[w]=0.                                           \tag{H9.20}
\]

Hard Lefschetz on $X$ gives an isomorphism
$L_h^2:H^{2m-2}(X)\to H^{2m+2}(X)$ whose inverse is algebraic because
$X$ is abelian. Choose a rational algebraic correspondence $\beta$
inducing this inverse, and set, at the Chow level,
\[
 a=\beta_*\kappa_*w,\qquad Z_0=Z+h a,
 \qquad w_0=w-\kappa^*a.                                      \tag{H9.21}
\]
Here $a$ has codimension $m-1$ on $X$. Since the axis is the
intersection of two divisors of class $h$, its class is $h^2$, and the
projection formula gives $\kappa_*\kappa^*a=h^2a$.
Equations (H9.20)--(H9.21) give
\[
 h[Z_0]=0,\qquad \kappa_*[w_0]=0.
\]
Moreover $B_K$ is a divisor of class $j^*h$ on $Y$, so the projection
formula gives the exact Chow identity $i_*\kappa^*a=j^*(h a)$.
Consequently
\[
 \zeta=j^*Z_0+i_*w_0.                                         \tag{H9.22}
\]
This proves surjectivity in (H9.18). It also proves well-definedness:
changing a representative of $[Z_0]$ or $[w_0]$ changes the global
blowup lift by a homologically trivial cycle, whose restriction lies
in $\mathcal H_0$.

To check that $A_B$ is primitive, use
$j_*i_*[w_0]=\kappa_*[w_0]=0$ and the injectivity of $j_*$ just noted.
Then $i_*[w_0]=0$, whence
$h_B[w_0]=i^*i_*[w_0]=0$ by self-intersection for the Cartier divisor
$B_K\subset Y$.

\paragraph{Proof: all vertical cohomology classes.}
We need an exact kernel statement, not merely a positivity argument.
Let $F=f^*[\mathrm{point}]=q^*h-[E]$. Under the specialization
injectivity hypothesis, every vertical cycle $V$ of codimension $m$
has cohomology class
\[
 [V]=F\,q^*u=q^*(h u)-e_*p^*\kappa^*u
 \quad\text{for some }u\in H^{2m-2}(X,\mathbf Q_\ell(m-1)).     \tag{H9.23}
\]
This assertion is cohomological; $u$ need not be asserted algebraic.
It suffices to prove it for $V$ supported in one fibre $Y_s$.
The cycle defines a linear functional on $H^{2m}(Y_s)(m)$ by cap
product and trace. Extend that functional along the injective map to
$H^{2m}(Y_{\bar\eta_s})(m)$. Poincar\'e duality on the smooth nearby
fibre identifies the extension with pairing against a class in
$H^{2m-2}(Y_{\bar\eta_s})(m-1)$. Weak Lefschetz identifies this last
space with $j^*H^{2m-2}(X)(m-1)$, so the class is $j^*u$ for some $u$.

For any global $\alpha\in H^{2m}(\widetilde X)(m)$, compatibility of
restriction with specialization now gives
\[
 \int_{\widetilde X}\alpha[V]
   =\int_{Y_{\bar\eta_s}}\alpha|_{Y_{\bar\eta_s}}\,j^*u
   =\int_{\widetilde X}\alpha Fq^*u.
\]
Perfect Poincar\'e duality on $\widetilde X$ proves the first equality
of classes in (H9.23); intersecting with the exceptional divisor gives
the second. For a smooth fibre the same argument starts directly with
its Poincar\'e duality. Summing proves the assertion for every vertical
cycle.

Suppose an adapted pair $(Z_0,w_0)$ maps to zero in
$\mathcal H/\mathcal H_0$. Its lift
$x=q^*Z_0+e_*p^*w_0$ then differs from a homologically trivial global
cycle by a vertical cycle, using localization for the restriction to
the generic fibre. By (H9.23) and the direct cohomological blowup
decomposition,
\[
 [Z_0]=h u,\qquad [w_0]=-\kappa^*u
\]
for some $u$. Since $h[Z_0]=0$, one has $h^2u=0$.
Hard Lefschetz on $X$ gives $u=0$. Both components vanish, proving
injectivity of $\Theta$. This proof does not assume definiteness of
either intersection form.

\paragraph{Proof: the pairing and its radical.}
For two adapted lifts, the mixed intersection is zero by
$q_*e_*p^*w=0$. The self-intersection formula for the exceptional
divisor, $\mathcal O_E(E)=\mathcal O_E(-1)$ and
$p_*c_1(\mathcal O_E(1))=1$, gives
\[
 \deg_{\widetilde X}
  (q^*Z_0+e_*p^*w_0)(q^*Z'_0+e_*p^*w'_0)
   =\deg_X(Z_0Z'_0)-\deg_B(w_0w'_0).                          \tag{H9.24}
\]
Both lifts are admissible, so multiplying by $(-1)^m$ is exactly the
height. This proves the isometry (H9.18). The map is canonical because
its definition uses only restriction and Gysin maps; the choice of
$\beta$ was used solely to prove its surjectivity.

The first form in (H9.18) is positive definite by 032, Corollary 8.4.
Hence the stated semipositivity equivalence follows by restricting to
or summing the two orthogonal summands. If the second form is also
positive definite, the quotient form is positive definite and its
preimage radical is exactly $\mathcal H_0$. More generally, without
that assumption, the radical modulo $\mathcal H_0$ is precisely the
radical of the displayed axis form.\hfill$\square$

\subsection*{H9-T8. Complete positivity for the threefold pencil fibre}

\textbf{Corollary.} In H9-T7 let $m=2$. Then $Y$ is a threefold and
$B$ is a smooth projective surface. Accepting 032 Corollary 8.4, the
geometric height on all of $\operatorname{CH}^2(Y)_{\rm hom,\mathbf Q}$
is positive semidefinite, and its radical is exactly
\[
 \operatorname{im}\bigl(\operatorname{CH}^2(\widetilde X)_{
       \rm hom,\mathbf Q}\longrightarrow
       \operatorname{CH}^2(Y)_{\mathbf Q}\bigr).
\]
The quotient is canonically the orthogonal direct sum in (H9.18), with
positive form $\deg_X(zz')-\deg_B(ww')$.

\emph{Proof.} A class $w\in A_B$ is a divisor class on $B$ and is
primitive for the ample class $h_B$. The surface Hodge index theorem
(SGA 6, XIII, Theorem 7.1, printed p.~662), valid in every
characteristic, makes $-\deg_B(ww')$ positive definite modulo numerical
equivalence. SGA 6, XIII, Theorem 4.6, printed pp.~645--646, identifies
numerical triviality of a divisor with algebraic triviality after a
nonzero multiple; such a divisor has zero rational cycle class.
Conversely a zero cycle class is numerically trivial. Thus this is
positive definiteness on precisely the cohomological divisor classes
used in $A_B$. Apply H9-T7. There is
no need to place these generic cycles in an algebraic intermediate
Jacobian.\hfill$\square$

\paragraph{Scope and the remaining higher-dimensional question.}
These results concern cycles defined over the original pencil field
$K=k(\mathbf P^1)$. After an arbitrary finite cover of the base, the
total space need no longer be the blowup of $X$, so no such extension is
silently claimed. For $m>2$ the remaining positivity statement is the
Hodge-standard sign on the explicit algebraic subspace $A_B$ of the
axis, rather than an unspecified height conjecture on all cycles of
$Y$. If a geometric generic axis has $A_B=0$, H9-T7 immediately removes
this gate; proving that vanishing requires an additional monodromy or
Hodge-type argument and cannot be inferred for an arbitrary fixed axis.


% New deduction in run JAC1964-20261121, independently source-checked.
% /root/programme_audit; GPT-6 Astra Pro; effort setting not exposed.

\subsection*{JAC-P09. A geometric generic axis has no variable algebraic classes}

\begin{proposition}[JAC-P09]
\label{jac:P09}
Let $X\subset\mathbf P^N_k$ be an abelian variety of dimension
$2m\ge4$ over an algebraically closed field $k$ of arbitrary
characteristic. Fix $d\ge3$. Let $(H_1,H_2)$ be a geometric generic
ordered pair of degree-$d$ hypersurfaces, and put
\[
 Y_1=X\cap H_1,\qquad B=X\cap H_1\cap H_2,
 \qquad \kappa:B\hookrightarrow X.
\]
All varieties hereafter are taken over the algebraic closure $k'$ of
the function field of this parameter space. Then for every invertible
prime $\ell$,
\[
 \ker\!\left(\kappa_*:\operatorname{Alg}^{m-1}(B)
             \longrightarrow H^{2m+2}(X,\mathbf Q_\ell(m+1))\right)=0.
                                                        \tag{JAC.P09.1}
\]
Equivalently, the variable middle cohomology of this geometric generic
axis contains no nonzero rational algebraic cycle class.

The pencil generated by $(H_1,H_2)$ is Lefschetz. Accepting 032,
Corollary~8.4, its signed geometric height on the entire group
$\operatorname{CH}^{m}(Y/K)_{\mathrm{hom},\mathbf Q}$, where
$K=k'(\mathbf P^1)$ and $Y/K$ is its generic fibre, is positive
semidefinite. Its radical is exactly
\[
 \operatorname{im}\!\left(
  \operatorname{CH}^{m}(\operatorname{Bl}_B X)_{
       \mathrm{hom},\mathbf Q}
     \longrightarrow\operatorname{CH}^{m}(Y)_{\mathbf Q}\right),
                                                        \tag{JAC.P09.2}
\]
and the resulting quotient, with its height form, is canonically
isometric to
\[
 \left(\ker\!\left(L_h:\operatorname{Alg}^{m}(X)
                        \longrightarrow H^{2m+2}(X)(m+1)\right),
       \ (z,z')\longmapsto(-1)^m\deg_X(zz')\right),
                                                        \tag{JAC.P09.3}
\]
where $h$ is the degree-$d$ section class.
\end{proposition}

\begin{proof}
We first identify the parameter field precisely. Choose generic
coefficients for $H_1$, let $k_1$ be an algebraic closure of their
function field inside $k'$, and regard $Y_1$ as a smooth projective
variety over $k_1$, with the projective embedding induced from
$X\subset\mathbf P^N$. The independent coefficients of $H_2$ are
still generic over $k_1$. Thus $B/k'$ is the geometric generic member
of the universal family of smooth degree-$d$ hypersurface sections
of this fixed embedded $Y_1/k_1$. This is a statement about its
geometric generic fibre, not about every closed parameter point in
a Zariski open subset.

The dimension of $Y_1$ is $2m-1\ge3$, which is odd.
Katz--Pandharipande, \emph{Inequalities related to Lefschetz pencils
and integrals of Chern classes}, Corollaries~2 and~3, author PDF
p.~3, show that its variable middle local system has rank $N_d>8$
and geometric monodromy group the full orthogonal group $O(N_d)$.
The only exceptional case in those corollaries is a linearly
embedded $\mathbf P^3$ with $d=3$. It cannot occur here:
weak Lefschetz gives
$H^1(Y_1,\mathbf Q_\ell)\simeq H^1(X,\mathbf Q_\ell)\ne0$,
whereas $H^1(\mathbf P^3,\mathbf Q_\ell)=0$.
The hypotheses in the paper's introduction are only that $k_1$ is
algebraically closed and $\ell$ is invertible. Corollary~3 imposes
no restriction on the characteristic. In particular this use of
orthogonal monodromy is valid in characteristic two.

Write $i:B\hookrightarrow Y_1$ and $j_1:Y_1\hookrightarrow X$.
Hard Lefschetz on $Y_1$ makes the ambient subspace
$i^*H^{2m-2}(Y_1)$ nondegenerate for the middle pairing on $B$.
Consequently its orthogonal complement
\[
 E_B=\ker\!\left(i_*:H^{2m-2}(B)
                         \longrightarrow H^{2m}(Y_1)(1)\right)
\]
identifies equivariantly with the variable quotient local system
used by Katz--Pandharipande. The projection onto $E_B$ commutes with
geometric monodromy; it need not be asserted algebraic for this step.

Every algebraic cycle on the geometric generic fibre is defined
over a finite algebraic extension of the generic parameter field.
Its cohomology class therefore has a finite orbit under the
geometric fundamental group. Purely inseparable field extensions
do not change this assertion: their absolute Galois groups and
prime-to-characteristic \^{e}tale cohomology are unchanged, and the
separable part of the field of definition has finite degree.
The variable projection of the class also has a finite orbit.
Its stabilizer has finite index, so the Zariski closure of that
stabilizer contains the identity component $SO(N_d)$ of $O(N_d)$.
The standard representation of $SO(N_d)$, with $N_d>8$, has no
nonzero fixed vector. The variable projection is therefore zero.

Finally,
\[
 (j_1)_*:H^{2m}(Y_1)\longrightarrow H^{2m+2}(X)(1)
\]
is an isomorphism, by Poincar\'e duality and weak Lefschetz in degree
$2m-2$. Since $\kappa_*=(j_1)_*i_*$, its kernel on
$H^{2m-2}(B)$ is precisely $E_B$. The preceding paragraph proves
(JAC.P09.1), including for rational algebraic combinations of cycles.

A geometric generic ordered independent pair determines a geometric
generic line in the projective parameter space of degree-$d$
sections. SGA~7~II, XVII, Theorem~2.5.1, printed p.~217, ensures
that this is a Lefschetz pencil for $d\ge2$. Its axis is smooth.
JAC-P07 supplies condition~A, so H9-T7 applies to all homologically
trivial cycles over its original pencil field $K$. In the
orthogonal decomposition of H9-T7, formula (JAC.P09.1) says that
the axis summand $A_B$ is zero. Only the ambient primitive summand
remains. Its form is positive definite by 032, Corollary~8.4;
H9-T7 then gives (JAC.P09.2)--(JAC.P09.3) and the asserted
semipositivity.
\end{proof}

\paragraph{Scope of the genericity and of the height assertion.}
The assertion about $B$ is geometric generic. It does not assert
vanishing of $A_B$ for every sufficiently general closed point over
the original field, nor for a fixed special axis. One may equivalently
describe the pencil as the geometric generic pencil in the
Grassmannian, after adjoining a generic ordered frame and taking an
algebraic closure. The result concerns cycles defined over the
original pencil function field $k'(\mathbf P^1)$; it makes no claim
after an arbitrary finite covering of the pencil base, whose total
space need not be a blowup of $X$.

\paragraph{Primary source for the new monodromy input.}
The author-hosted PDF is
\url{https://web.math.princeton.edu/~nmk/Inequalities15.pdf}.
Its introduction, p.~1, defines the variable quotient and geometric
monodromy over an algebraically closed field with one invertible
prime. Theorem~0, p.~2, cites Katz's theorem that $N_d>8$ and
$d\ge3$ force full orthogonal monodromy in even fibre dimension.
Corollaries~2--3, p.~3, remove the rank hypothesis with the unique
linear cubic-surface exception. Those pages were inspected in the
original author PDF in this run. The characteristic restriction
appearing later in the separate Theorem~4 is not a hypothesis of
Corollary~3.


% Further H9 deduction in run JAC1964-20261121.
% /root/h9_height; GPT-6 Astra Pro; effort setting not exposed.
% Classical results used below are explicitly cited; no formal verification
% or historical originality is asserted.

\subsection*{H9-T9. The algebraically trivial domain: incidence and trace}

Let $k$ be algebraically closed, let $C/k$ be a smooth projective connected
curve, and put $K=k(C)$. Let $Y/K$ be smooth projective of dimension
$2m-1$. Write $H_V$ for the degree of the unsigned geometric intersection
height on a smooth projective $K$-variety $V$; thus the signed middle
height on $Y$ is
\[
 h_m(x,y)=(-1)^m H_Y(x,y).
\]
We use the pairing of Kahn, \emph{Refined height pairing}, Theorem 4.14,
printed p.~31. Its restriction to algebraically trivial cycles exists in
every characteristic by Theorem 5.6(a), printed pp.~33--35. On a model
with admissible extensions it is their intersection degree, so this
normalization agrees with H9-T1.

Let $P=P^m(Y)$ be the rationally descended incidence Picard variety, and
let $a$ denote its quotient map on algebraically trivial cycles. Choose
an abelian direct summand $B\subset P$ in the isogeny category, and a
rational family $U\in\operatorname{CH}^m(B\times Y)_{\mathbf Q}$
whose family map is the inclusion $B\to P$. The existence over $K$ is
the rational Poincar\'e-family descent in JAC-P00. Subtract $U_0$ so
that the family is pointed. Let
\[
 D=\bigl(\operatorname{pr}_{12*}(U_{13}U_{23})\bigr)_{11},
 \qquad \ell=\frac{(-1)^m}{2}\Delta_B^*D
       \in\operatorname{NS}(B)_{\mathbf Q},                 \tag{H9.25}
\]
where the subscript removes the two axis restrictions. Suppose $\ell$
is a rational polarization. For the primitive summand of a generic
Lefschetz section of an even-dimensional abelian variety this hypothesis
is supplied by H9-T4, H9-T6, H9-T6a and JAC-P07, accepting 032.

Define
\[
 \mathcal A_B=
 \{z\in\operatorname{CH}^{m}(Y)_{\mathrm{alg},\mathbf Q}:
                      a(z)\in B(K)\otimes\mathbf Q\}.
\]
One may also use the group of $K$-cycles which become algebraically
trivial after a finite extension: admissibility and the curve-family
arguments below are obtained over that extension and descended by
transfer, with rational coefficients. The quotient map and the
subspace are understood in the rational category; they are independent
of clearing the denominators of the chosen isogenies.

\textbf{Claim.} The height on $\mathcal A_B$ is positive semidefinite,
and its radical is exactly
\[
 \{z\in\mathcal A_B:
      a(z)\in\operatorname{im}((\operatorname{Tr}_{K/k}B)(k)
                                  \longrightarrow B(K))
                                  \otimes\mathbf Q\}.       \tag{H9.26}
\]
More precisely, if $L$ is a symmetric ample line bundle representing
$N\ell$ for some positive integer $N$, and
\[
 \phi_L(b)=t_b^*L\otimes L^{-1}\in\operatorname{Pic}^0(B),
\]
then
\[
 h_m(z,z')=\frac1N\,
 H_B\bigl(\phi_L(a(z')),\,[a(z)]-[0]\bigr).                 \tag{H9.27}
\]
The right side uses the divisor--Albanese pairing, so the notation is
extended bilinearly to rational points. No assertion is made that an
arbitrary homologically trivial cycle has a point $a(z)$.
Explicitly, clear a positive denominator $r$ so that $rb$ is an actual
$K$-point; the notation $[b]-[0]$ in that pairing means the Albanese
class of $r^{-1}([rb]-[0])$. Its value is independent of $r$ by the
Albanese factorization. The same clearing of positive denominators
turns every rational inclusion, projector and family map used here
into an actual homomorphism before applying it to points.
The induced map from $\mathcal A_B$ modulo the subgroup in (H9.26) to
$\operatorname{LN}(B,K/k)\otimes\mathbf Q$ is an isomorphism with
bilinear forms: surjectivity is supplied by the pointed family $U$.

\paragraph{Proof: the required correspondence adjointness.}
Kahn's Proposition 3.8, printed pp.~19--20, and Theorem 4.14(b) state
adjointness for correspondences of degree zero. The degree shift needed
here reduces to exactly that statement. Given
$\Gamma\in\operatorname{CH}^{m}(T\times Y)_{\mathbf Q}$, a degree-zero
zero-cycle $\delta$ on $T$, and an algebraically trivial middle cycle
$\eta$ on $Y$, put $W=T\times\mathbf P^{m-1}$ and
\[
 \Gamma'=\Gamma\boxtimes H^{m-1}
   \in\operatorname{CH}^{2m-1}(W\times Y)_{\mathbf Q},
\]
where $H$ is the hyperplane class. This is a degree-zero
correspondence in Kahn's contravariant convention. Its transpose action
on $\eta$ is $\Gamma^t_*\eta\boxtimes H^{m-1}$, while its pushforward
action on $\delta\boxtimes[\mathbf P^{m-1}]$ is $\Gamma_*\delta$.
Adjointness, followed by adjointness for the projection $W\to T$, gives
\[
 H_Y(\Gamma_*\delta,\eta)
    =H_T(\delta,\Gamma^t_*\eta).                            \tag{H9.28}
\]
The projection formula uses only
$\deg_{\mathbf P^{m-1}}H^{m-1}=1$. Every cycle paired here is
algebraically trivial, so Kahn's admissibility theorem applies. This
argument also includes $m=1$, with $\mathbf P^0$.

\paragraph{Proof: the incidence kernel is killed by the height.}
Let $\eta$ have zero incidence image. An algebraically trivial cycle
$z$ is a finite sum of curve-family differences
$\Gamma_*\delta$, after a finite extension if necessary. By definition
of incidence equivalence, $\Gamma^t_*\eta$ is zero in
$\operatorname{Pic}^0(T)_{\mathbf Q}$ for every such tester. If this
vanishing is initially over $\overline K$, it already holds rationally
over the finite field of definition: the kernel of restriction on a
Picard group is killed by transfer after a suitable finite extension.
Equation (H9.28) therefore proves $H_Y(z,\eta)=0$. Finite base change
multiplies the degree height by its positive field degree, by Kahn's
Theorem 4.14(b) and Proposition 3.9(iv), so the assertion descends to
$K$.

Consequently the height on algebraically trivial cycles factors through
the incidence quotient. In particular, if $b=a(z)\in B(K)\otimes
\mathbf Q$, then $z$ may be replaced in the pairing by
$U_*([b]-[0])$. The difference has incidence image zero. The same
applies to $z'$.

Apply (H9.28) with $T=B$ and $\Gamma=U$. The composition $U^t\circ U$
is the divisor correspondence before rigidification in (H9.25).
Pullbacks of divisors from either factor act trivially on degree-zero
zero-cycles, so its action is that of $D$. The symmetric rigidified
divisor $D$ is the biextension divisor
\[
 D=\frac{(-1)^m}{N}
       (\mu^*L-p_1^*L-p_2^*L)
       \quad\text{in }\operatorname{Pic}(B\times B)_{\mathbf Q}
       \text{ after rigidification}.                        \tag{H9.29}
\]
Indeed symmetric divisor correspondences rigidified on the axes are
identified with symmetric homomorphisms $B\to\widehat B$, and
diagonal pullback is twice the associated N\'eron--Severi class.
Alternatively, the difference in (H9.29) has zero N\'eron--Severi class,
and every line bundle in $\operatorname{Pic}^0(B\times B)$ is a tensor
product of pullbacks from the factors; rigidification removes them.
The restriction of the displayed biextension at its first argument $b$
is precisely $\phi_L(b)$ under the stated translation convention.
This proves (H9.27), including its sign.

\paragraph{Proof: an explicit polarization--intersection identity.}
Let $d=\dim B>0$, and define the homomorphism
\[
 T_L:\operatorname{Pic}^0(B)\longrightarrow B,
 \qquad \alpha\longmapsto
        \operatorname{alb}_B(L^{d-1}\alpha).
\]
Intersection with a fixed cycle followed by the Albanese map defines
a homomorphism of abelian varieties. We claim the identity
\[
 T_L\phi_L=-c_L\,1_B,
 \qquad c_L=(d-1)!\chi(L)>0.                                 \tag{H9.30}
\]
We give an algebraic proof that fixes the sign in every characteristic.
First, the Albanese image of an intersection containing at least two
algebraically trivial divisor factors is zero. To see this, vary those
factors independently in copies of $\operatorname{Pic}^0(B)$, using
the Poincar\'e line bundles. Intersection followed by Albanese gives a
morphism from the product of these abelian varieties to $B$, taking
the origin to zero. Such a morphism is a homomorphism by rigidity, and
therefore is the sum of its restrictions to the factors. Each such
restriction is zero because another factor is the trivial line bundle.
The morphism is zero. This also proves the assertion when several
factors coincide, by restriction to a diagonal.

Now write $\alpha=\phi_L(b)$, so $t_b^*L=L+\alpha$ in additive
notation. Expanding the $d$th power and applying Albanese gives
\[
 \operatorname{alb}_B((t_b^*L)^d-L^d)=d\,T_L(\alpha).
\]
On the other hand, pullback by $t_b$ translates every point of the
zero-cycle $L^d$ by $-b$. Its Albanese sum therefore changes by
$-\deg(L^d)b=-d!\chi(L)b$, by Riemann--Roch for abelian varieties.
Dividing by $d$ proves (H9.30). Both sides are homomorphisms, so
checking on geometric points proves the identity over $K$ as well.

\paragraph{Proof: positivity and the exact radical.}
Kahn's Theorem 6.2, printed pp.~38--39, shows that the divisor--zero-cycle
height factors through the Albanese map and, after taking degree, through
the two Lang--N\'eron quotients. Its Theorem 6.6, printed pp.~41--42,
shows that
\[
 q_L(\alpha)=H_B(\alpha,L^{d-1}\alpha)
             =H_B(\alpha,T_L\alpha)
\]
is negative definite on
$\operatorname{LN}(\operatorname{Pic}^0(B),K/k)\otimes\mathbf Q$.
There is no conjectural hypothesis here, because the base $C$ has
dimension one. Since $\phi_L$ is an isogeny, it induces an isomorphism
of rational Lang--N\'eron groups. Equation (H9.30) gives
\[
 H_B(\phi_L(b),[b]-[0])
       =-\frac{1}{c_L}\,q_L(\phi_L(b)).                     \tag{H9.31}
\]
Thus this is positive definite modulo the trace, with exactly the trace
as its rational zero space. Equations (H9.27) and (H9.31) prove
semipositivity and (H9.26). If $B=0$, the incidence-kernel argument
already gives the same conclusion.\hfill$\square$

\paragraph{What this repairs, and what it does not assert.}
This is a precise trace-kernel theorem on the algebraically trivial
domain where the incidence quotient is defined. It proves the height
consequence of a primitive incidence polarization in every characteristic.
The proof identifies the pairing with a divisor--Albanese height and
derives its positivity directly; it does not invoke an unspecified
comparison theorem for a higher Abel--Jacobi map. It supplies no map
from every homologically trivial cycle to $P^m(Y)$, and therefore does
not repair the overlarge domain in literal H9(b). The E4 example in
JAC-P04 and the quadric example in H9-C1 retain their stated force.

\paragraph{Source and novelty record.}
The primary source used for the height properties is
\url{https://arxiv.org/pdf/2009.00533}, version 7, directly read in this
run: Proposition 3.8, pp.~19--20; Theorem 4.14, p.~31; Theorem 5.6(a),
pp.~33--35; Theorem 6.2, pp.~38--39; and Theorem 6.6, pp.~41--42.
The padding argument, the explicit proof of (H9.30), and their assembly
with the incidence quotient are deductions written in this run. No
historical originality, integral universal cycle, or formal proof check
is claimed.


% H9-T10: actual codimension-two algebraic representatives.
% New deduction in run JAC1964-20261121, by /root/h9_height.
% GPT-6 Astra Pro; effort setting not exposed.

\subsection*{H9-T10. The actual algebraic representative of the generic threefold}

\begin{theorem}[H9-T10]
Let $X/k$ be a smooth projective fourfold over an algebraically closed
field of arbitrary characteristic. Suppose that the Lefschetz lowering
operator of $X$ is algebraic, and let $Y/K$, $K=k(\mathbf P^1)$, be the
generic fibre of a genuine Lefschetz pencil of hyperplane sections of
$X$. Then the variable factor of the actual universal regular algebraic
representative $\operatorname{Ab}^2(Y)$ is polarized by the rational
incidence divisor $\Delta^*d/2$. Its canonical map to the variable
incidence Picard factor is an isogeny.

If $X$ is abelian and 032, Corollary 8.4, is accepted, the same assertions
hold for the whole primitive factor:
\[
 \operatorname{Ab}^2(Y)_{\mathrm{prim}}
          \longrightarrow P^2(Y)_{\mathrm{prim}}
       \quad\text{is an isogeny,}                           \tag{H9.32}
\]
and the normalized primitive incidence divisor on its source is a
rational polarization. If $X$ is a smooth complete intersection in
projective space, all of $H^3(Y)$ is variable and primitive, so the
assertion holds for the entire $\operatorname{Ab}^2(Y)$, without 032.
The assertion concerns the generic fibre and its displayed abelian
varieties after base change; it does not assert the same conclusion
for every special fibre.
\end{theorem}

\paragraph{Existence over the imperfect pencil field.}
The symbol $\operatorname{Ab}^2(Y)$ here denotes the actual universal
regular algebraic representative over $K$, even when $K=k(t)$ is
imperfect. Its existence is the arbitrary-field case of
Achter--Casalaina-Martin--Vial, \emph{A functorial approach to regular
homomorphisms}, \emph{Algebraic Geometry} 10 (2023), Theorem 2,
printed p.~89, proved as Theorem 6.1, printed p.~114. Take
$S=\Lambda=\operatorname{Spec}K$ in that theorem. The authors explicitly
state the new arbitrary-field codimension-two conclusion immediately
before Theorem 2. No perfectness assumption or descent from a geometric
object is needed for this existence assertion. We do not assert an
integral universal base-change theorem along purely inseparable
extensions.

\paragraph{The codimension-two realization input.}
We first prove the exact consequence of the known Bloch-map results
needed here. Let $V/F$ be smooth projective over any field, put
$R=\operatorname{Ab}^2(V)$, and let $B$ be an abelian factor of $R$ in
the rational category. Let $U\in\operatorname{CH}^2(B\times V)_{\mathbf Q}$
be a pointed family whose induced regular homomorphism is the inclusion
$\iota:B\to R$. Then, for every prime $\ell\ne\operatorname{char}F$,
its realization map is injective:
\[
 f_\ell:V_\ell B\lhook\joinrel\longrightarrow
                   H^3(V_{\overline F},\mathbf Q_\ell(2)).  \tag{H9.33}
\]

The sources are Achter--Casalaina-Martin--Vial, Ast\'erisque 455 (2025),
\S1.3 and Lemmas 1.3--1.4, printed pp.~6--7, and the proof of
Proposition 8.1, printed pp.~41--42. The first gives a miniversal cycle
over an arbitrary field for every surjective regular homomorphism.
Lemma 1.3 says that $b\mapsto[b]-[0]$ identifies the torsion of an
abelian variety with the torsion of its algebraically trivial zero-cycle
group. Lemma 1.4 consequently gives the compatibility on Tate modules
of a pointed family with its regular homomorphism. The proof of
Proposition 8.1 uses the injectivity of the codimension-two Bloch map
$V_\ell\lambda^2$, in every characteristic with $\ell$ invertible.

For completeness, these statements give (H9.33) directly. Clear a
positive denominator of $U$ and then tensor all Tate modules with
$\mathbf Q_\ell$. The family gives
\[
 V_\ell B\xrightarrow{V_\ell w_U}
 V_\ell A^2(V_{F^{\rm sep}})
 \xrightarrow{V_\ell\lambda^2}
 H^3(V_{\overline F},\mathbf Q_\ell(2)).
\]
Writing $\varphi$ for the universal regular homomorphism,
Lemma 1.4 gives
$(V_\ell\varphi)(V_\ell w_U)=V_\ell\iota$.
Since $\iota$ is an inclusion up to isogeny, $V_\ell w_U$ is
injective. The second arrow is injective by the cited Bloch-map
result. Their composite is the cohomological correspondence map
by the naturality used in Proposition 8.1, proving (H9.33).
The same input is also stated explicitly in the 2023 paper:
Theorem 8.1, printed p.~119, gives injectivity of the second
$\ell$-adic Bloch map over a separably closed field, and Proposition
8.2, printed p.~120, gives the miniversal realization injection at
primes away from its degree. The rational argument just given also
covers primes dividing that degree, since its positive integer is
invertible in $\mathbf Q_\ell$.
This proves injectivity for a chosen normalized miniversal realization;
it does not assert that $V_\ell\varphi$ is injective, or that
$V_\ell\lambda^2$ is surjective onto all of $H^3$.

\paragraph{Projectors and normalized families.}
Put $h=[Y]$. On $H^3(Y)$ the algebraic correspondence
\[
 \Pi_V=1-j^*(L_h:H^3(X)\longrightarrow H^5(X))^{-1}j_*
                                                               \tag{H9.34}
\]
is the self-adjoint variable projector. Choose the inverse supported
from degree $5$ to degree $3$, using the algebraic K\"unneth
projectors supplied by the Lefschetz standard theorem. On the other
degrees the subtracted correspondence is zero, so (H9.34) is still
an idempotent on the whole cohomology. Proposition 8.1 applied to
its idempotence defect gives an actual idempotent on
$R=\operatorname{Ab}^2(Y)$ in the isogeny category. Denote its image
by $B_V$.

For abelian $X$, let $e$ be the primitive projector of $H^3(X)$ onto
$\ker L_h^2$, and put
\[
 \Pi_F=j^*eL_h^{-1}j_*.
\]
Exactly as in H9-T4, this is an algebraic self-adjoint idempotent,
orthogonal to $\Pi_V$. It gives a factor $B_F$ of $R$. Their sum
acts on $H^3(Y)$ as the primitive projector. This defines the factor
$B_{\rm prim}$ in (H9.32), with
$B_{\rm prim}\sim B_F\times B_V$. The action on the algebraic
representative depends only on the degree-three action: to verify
this precise point, take a normalized miniversal family on $R\times Y$
and apply the parameter degree-one projector constructed below. If a
correspondence acts as zero on $H^3(Y)$, its composite with this family
has zero total cohomology class; Proposition 8.1 with source $R$ and
source codimension $\dim R$ makes the induced homomorphism $R\to R$
zero. Thus no off-middle action of the chosen representatives changes
the factor just defined.

Pull a miniversal family back to any of these factors, divide by its
positive degree, subtract its origin fibre, and apply the relevant
projector on $Y$. Its regular family map remains the factor inclusion.
As in JAC-T4.4, a rational polynomial $Q$ with $Q(2)=1$ and
$Q(2^a)=0$ for $0\le a\le2\dim B$, $a\ne1$, applied to
$([2]\times1)^*$, preserves this family map and leaves its cohomology
entirely in $H^1(B)\otimes H^3(Y)$. It has the required fixed or
variable image. This construction is valid in characteristic two:
the relevant cohomology uses $\ell\ne2$, the eigenvalues are $2^a$,
and the projector and family are rational correspondences.

\paragraph{The variable factor.}
If $B_V=0$, its assertion is empty. Otherwise (H9.33) gives an injection
\[
 V_\ell B_V\lhook\joinrel\longrightarrow E_\ell(2),
 \qquad E_\ell=\ker(j_*:H^3(Y)\to H^5(X)),                  \tag{H9.35}
\]
for every invertible prime. In particular $E_\ell\ne0$. SGA 7 II,
XVIII, 6.6 and 6.7.2, printed pp.~323--327, give spanning and conjugacy
of vanishing cycles and absolute irreducibility after every finite-index
monodromy restriction. Their parity hypotheses hold because $\dim X=4$;
the characteristic-two exception with odd ambient dimension is absent.
As in H9-T6a, all local vanishing cycles are nonzero and (H9.35) is an
isomorphism, even after the finite extension needed to define the family.

The cup product on $E_\ell$ is nondegenerate, because it is the
orthogonal complement of the nondegenerate ambient summand. Hence the
self-incidence divisor of the normalized family has a nondegenerate
alternating realization. Its associated symmetric homomorphism is
therefore an isogeny. This establishes the nondegeneracy here directly;
it was not borrowed from the incidence quotient.

The remaining sign argument is precisely H9-T6, now applicable to
$B_V$ and this family. To make clear which hypotheses have been
supplied, absolute irreducibility gives
$\operatorname{End}^0(B_{V,\overline K})=\mathbf Q$. If $L_0$ is a
genuine polarization, write
$\ell_V=\tfrac12\Delta^*d_V=cL_0$, $c\in\mathbf Q^*$.
The incidence Koszul sign gives
$E_{\ell_V}=-f_\ell^*Q_\ell=cE_0$. At one quadratic singular fibre,
after a common semistable extension of ramification degree $e>0$,
the comparison with the integral positive toric monodromy pairing of
H9-L2 gives
\[
 a_\ell^2=\frac{q_0}{e c}\quad(\ell\ne\operatorname{char}k),
 \qquad q_0\in\mathbf Q_{>0}
\]
with the same rational $q_0$ for every prime. H9-L1 forces $c>0$.
Thus $\ell_V$ is a rational polarization. All ingredients in this
argument are properties of the normalized family, its injective
realization and the Lefschetz monodromy; none requires that $B_V$
already be identified with an incidence Picard variety.

\paragraph{The fixed primitive factor for an abelian fourfold.}
Work over $\overline K$ to test ampleness. Let $C_0\to B_F$ be the
normalization of a nonconstant curve, pointed at the origin after
translation, and let $\xi$ be the class of its induced family.
The induced homomorphism $J(C_0)\to B_F$ is nonzero, so its rational
Tate-module map is nonzero by faithfulness. Composing with the
injection (H9.33) for $B_F$ remains nonzero. By naturality this
composite is the realization of the pulled-back family. Thus
$\xi\ne0$. This argument uses the base change of the existing
realization injection; it requires no claim that a universal regular
representative commutes integrally with purely inseparable base change.
The class lies in $H^1(C_0)\otimes F$, with $F$ as in H9-T4.

H9-T4 transports it, by the algebraic Lefschetz inverse on $X$ and the
Abel map of $C_0$, to a nonzero primitive codimension-two algebraic
class on the abelian variety $J(C_0)\times X$. Equation (H9.9) and
032 give
\[
 \deg_{C_0}\ell_F
       =\tfrac12\int_{C_0\times Y}\xi^2>0.
\]
Translation does not alter this degree, because it acts trivially on
the N\'eron--Severi group. Thus $\ell_F$ is nef and has positive degree
on every curve. If its associated symmetric homomorphism had a
positive-dimensional connected reduced kernel, the restriction to
that abelian subvariety would be algebraically trivial and a curve
inside it would have degree zero. This is impossible. The nef and
nondegenerate ampleness criterion proves that $\ell_F$ is a rational
polarization, exactly as in JAC-T4.4.

The fixed and variable cohomological subspaces are orthogonal. The
mixed incidence divisor between the corresponding families therefore
has zero mixed cohomology class. Faithfulness of the Tate module for
homomorphisms of abelian varieties makes the corresponding homomorphism
$B_F\to\widehat B_V$ zero. Thus their incidence forms are orthogonal
already in the rational N\'eron--Severi group. The product of the two
rational polarizations descends through the isogeny to polarize
$B_{\rm prim}$.

\paragraph{Comparison with the incidence quotient and the complete-intersection case.}
The incidence quotient map is regular, so universality gives a
surjection $R\to P^2(Y)$, compatible with algebraic correspondences.
Its variable and primitive restrictions are therefore surjective.
As in (JAC.T44.4), Saito's incidence identities give
\[
 d_B=(q\times q)^*d_{P,B}
       \quad\text{in }\operatorname{NS}(B\times B)_{\mathbf Q}.
\]
This can also be checked by replacing the normalized miniversal family
by the pullback of a normalized Poincar\'e family one factor at a time:
their difference has zero incidence map, so testing it against either
family kills the mixed divisor component. Since $\Delta^*d_B/2$ is
ample, the kernel of $q$ cannot contain a curve. Its kernel is finite,
which proves the isogeny claims, allowing a finite inseparable kernel
in positive characteristic.

If $X$ is a smooth four-dimensional complete intersection, weak
Lefschetz gives $H^3(X)=H^7(X)=0$. Hence $E_\ell=H^3(Y)$ and
$H^5(Y)=0$, so this cohomology is wholly primitive. The variable
projector acts as the identity on the algebraic representative by the
degree-three detection proved above. Its lowering operator is algebraic
by the explicit construction in JAC-P08. The variable proof alone
therefore gives the full assertion without 032.\hfill$\square$

\paragraph{Scope and status.}
This is a deduction written in this run from the cited codimension-two
Bloch injectivity, the source-checked SGA monodromy statements, and 032
where explicitly used. It supplies the actual algebraic representative
in the first higher-dimensional generic-pencil test. It asserts a
rational primitive polarization and an isogeny comparison, not an
integral principal polarization. No standard torsion assumption for
$\operatorname{Ab}^2$, no surjectivity of its realization onto all of
$H^3$, and no map on arbitrary homologically trivial cycles has been
assumed.

The primary 2023 existence and realization source is the publisher PDF
\url{https://api.algebraicgeometry.nl/Article/19846/2023-1-003.pdf}.
Theorem 2 (p.~89), Theorem 6.1 (p.~114), Theorem 8.1 (p.~119), and
Proposition 8.2 (p.~120) were directly read in this run. All additional
data used in the monodromy proof are algebraic and of finite type, so
a common finite field extension suffices to define them; the
finite-index irreducibility and common finite semistable extension
arguments are exactly those of H9-T6.


% Optional precise replacement for the overlarge-domain reduction in H9(b).
% Deduction in run JAC1964-20261121. /root/h9_height.
% This is conditional exactly where stated. Root whole-read the final proof,
% including the accepted032 upgrade, and reported PASS on 2026-11-21.

\subsection*{H9-T11. A conditional reduction using an elliptic product}

The invalid map from all homologically trivial hyperplane restrictions
to an incidence Jacobian is unnecessary for the following more limited
reduction. It identifies separately the positivity input and the remaining
numerical-versus-homological gate.

\begin{proposition}[H9-T11]
Let $X/k$ be smooth projective of dimension $2m$ over an algebraically
closed field, let $h$ be an ample class, and let $E/k$ be an elliptic
curve with origin $0$. Suppose that the primitive projector for the
middle odd cohomology of $V=X\times E$, for the ample class
$h+[0]_E$, is algebraic. Suppose the primitive incidence assertion of
Theorem 4 holds for $V$: its signed incidence class is a rational
polarization on $P^{m+1}(V)_{\rm prim}$.

Then the form
\[
 q_X(Z)=(-1)^m\deg_X(Z^2)
\]
is nonnegative on all rational algebraic cycles $Z$ with
$h\,\operatorname{cl}(Z)=0$. On this primitive domain its zero space
is exactly the subgroup of numerically trivial cycles. Consequently
it is positive definite on the quotient of that domain by numerical
equivalence. If numerical and homological equivalence coincide for
middle cycles on $X$, it is positive definite on the primitive
algebraic cohomology classes: this is the middle Hodge-standard
assertion for $(X,h)$.
\end{proposition}

\begin{proof}
On $E\times E$ use the normalized Poincar\'e divisor
\[
 P_E=\Delta_E-E\times\{0\}-\{0\}\times E.
\]
As a rational Chow correspondence it is the degree-one projector:
$P_E^t\circ P_E=P_E$. Its diagonal pullback has degree $-2$,
because the diagonal normal bundle is $T_E$ of degree zero and the
two axis terms each restrict to $[0]$. Thus
\[
 \tfrac12\Delta_E^*P_E=-[0]
       \quad\text{in }\operatorname{NS}(E)_{\mathbf Q}.       \tag{H9.36}
\]

For a codimension-$m$ cycle $Z$ on $X$, form the pointed family
\[
 z_Z=Z\boxtimes P_E
    \in\operatorname{CH}^{m+1}(E\times(X\times E))_{\mathbf Q},
\]
with the first copy of $E$ the parameter. Every fibre difference is
algebraically trivial. Its relevant class is in
$H^1(E)\otimes H^{2m}(X)\otimes H^1(E)$.
If $h\operatorname{cl}(Z)=0$, then this class is primitive on $V$:
the $h$ term kills its $X$ component, and $[0]_E H^1(E)=0$.
The algebraic primitive projector and incidence detection (JAC-P01)
therefore show that its family homomorphism factors through
\[
 a_Z:E\longrightarrow B:=P^{m+1}(V)_{\rm prim}
\]
in the rational category. Clear positive denominators whenever an
actual homomorphism is needed.

Direct composition of the two family correspondences gives the
self-incidence divisor
\[
 D_Z=\deg_X(Z^2)P_E.
\]
If $\ell_B$ is the normalized signed incidence class on $B$, the
incidence identity and (H9.36) give
\[
 a_Z^*\ell_B
   =\frac{(-1)^{m+1}}2\Delta_E^*D_Z
   =(-1)^m\deg_X(Z^2)[0].                                  \tag{H9.37}
\]
The pullback of a rational polarization is nef, and its degree on
the elliptic curve is zero if and only if $a_Z=0$. Indeed a nonzero
homomorphism from an elliptic curve has finite kernel and its image
is an elliptic subvariety on which an ample line bundle has positive
degree. Thus $q_X(Z)\ge0$, and equality is equivalent to $a_Z=0$.

If $a_Z=0$, test the family against $z_{Z'}$ for an arbitrary rational
codimension-$m$ cycle $Z'$ on $X$. Incidence equivalence says that
the resulting divisor correspondence on the two parameter elliptic
curves induces the zero homomorphism. That correspondence is
\[
 \deg_X(ZZ')P_E.
\]
Since $P_E$ induces a nonzero isomorphism $E\to\operatorname{Pic}^0(E)$,
one has $\deg_X(ZZ')=0$ for every $Z'$. This is precisely numerical
triviality of $Z$. Conversely a numerically trivial $Z$ has $Z^2=0$.
We have proved the exact zero-space assertion. Positivity of the
induced rational bilinear form on the numerical quotient follows by
polarization. If middle numerical and homological equivalence coincide,
this quotient is exactly the space of primitive algebraic cohomology
classes, which proves the final assertion.
\end{proof}

\paragraph{Exact inputs and the lower-degree consequence.}
The algebraic-projector hypothesis holds, for example, if the
Lefschetz standard conjecture holds for $X$, since it then holds for
$X\times E$. The middle equality of numerical and homological
equivalence is a separate input; it has not been inferred merely from
positivity on an incidence quotient.

Here is a precise way to recover another degree from middle degree.
Let $X_0$ have dimension $n$, let $1\le p\le n/2$, and choose a smooth
complete intersection $i:W\hookrightarrow X_0$ of $n-2p$ divisors
in a sufficiently ample multiple of a chosen class $h_0$, with
$\dim W=2p$. Write their common class as $h=d h_0$.
For a primitive class $z\in\operatorname{Alg}^p(X_0)$, meaning
$h^{n-2p+1}z=0$, its restriction $i^*z$ is primitive on $W$.
Indeed $i_*L_h i^*z=h^{n-2p+1}z=0$, and
\[
 i_*:H^{2p+2}(W)\longrightarrow H^{2n-2p+2}(X_0)
\]
is injective by dual weak Lefschetz. Weak Lefschetz also makes
$i^*:H^{2p}(X_0)\to H^{2p}(W)$ injective. Finally
\[
 \deg_W((i^*z)^2)=\deg_{X_0}(z^2h^{n-2p}).
\]
Thus the middle assertion for $W$ implies the degree-$p$ assertion
for $X_0$. H9-T11 supplies that middle assertion from primitive
Theorem 4 for the odd variety $W\times E$, together with its displayed
algebraic-projector hypothesis and the middle numerical/homological
equality for $W$. These are the precise hypotheses for this reduction.
They must not be replaced by the assertion that a merely homologically
trivial hyperplane restriction defines a point of $P^p$.

\paragraph{Exactly what accepted 032 adds when $X$ is abelian.}
For an abelian $2m$-fold $X$, 032, Corollary 8.4, asserts that
\[
 (-1)^m\deg_X(z^2)>0
 \quad\text{for every nonzero primitive rational algebraic
 cohomology class }z.
\]
Consequently a primitive cycle which is numerically trivial has zero
cohomology class: its self-intersection is zero, and the displayed
definiteness rules out a nonzero class. Thus 032 explicitly removes
the middle numerical/homological zero-space gate on the primitive
domain of H9-T11. The strict conclusion here is an application of
the accepted theorem; it is not being inferred from semipositivity
on a quotient alone.

One can also derive the whole equality of numerical and homological
equivalence on an abelian variety from exactly the inputs already
accepted. Here are the linear-algebra details. Algebraicity of the
Lefschetz operators makes the primitive decomposition of an algebraic
class algebraic. For a class $a\in\operatorname{Alg}^r(X)$ with
$r\le n/2$, write
\[
 a=\sum_j h^j a_j,
 \qquad a_j\in\operatorname{Alg}^{r-j}(X)_{\rm prim}.
\]
Set the algebraic complementary class
\[
 b=\sum_j(-1)^{r-j}h^{n-2r+j}a_j.
\]
Different primitive strings are orthogonal: if $j<k$, the exponent
$n-2r+j+k$ in the mixed intersection exceeds
$n-2(r-j)$, so it kills $a_j$ by primitivity; interchange the factors
when $k<j$. Hence
\[
 \int_X ab
   =\sum_j(-1)^{r-j}\int_Xa_j^2h^{n-2(r-j)}>0
       \quad\text{if }a\ne0,                               \tag{H9.38}
\]
by 032. A numerically trivial class cannot have this nonzero algebraic
test. For $r>n/2$, put $s=n-r$ and write
$a=h^{n-2s}\sum_j h^j a_j$ using the algebraic inverse of hard
Lefschetz, with $a_j$ primitive of codimension $s-j$. The complementary
algebraic class $b=\sum_j(-1)^{s-j}h^j a_j$ gives
$\int_Xab=\sum_j(-1)^{s-j}\int_Xa_j^2h^{n-2s+2j}>0$
when $a\ne0$, by the same orthogonality and 032. The reverse
implication, homological triviality implies numerical triviality,
is immediate from compatibility of cycle classes and intersection.
This proves the needed equality for abelian varieties rather than
silently inserting it as another hypothesis.

For a general $X$ or for a fixed complete-intersection axis in an
abelian variety, 032 has not been assumed for that variety. Neither
an embedding into an abelian variety nor an unspecified relation of
motives replaces (H9.38). The separate numerical/homological gate
therefore remains unless an actual transfer preserving the relevant
forms is proved.

\paragraph{Relation to the archive and to the height statements.}
Cote 44, scan 9, writes the self-incidence divisor of a pointed family
and asks for positivity on the abelian quotient by its kernel. The
family $z_Z$ is an explicit instance of that construction, with the
primitive qualification recorded in the margin of scan 8. Scan 8
also expressly makes the Hodge reduction dependent on further
unproved assertions. The present argument identifies a precise
conditional numerical reduction; it does not identify those
unwritten assertions historically.

Theorem 5 in scan 4 starts with \emph{algebraically} trivial generic
restrictions, whereas literal H9 starts with homologically trivial
restrictions. H9-T9 proves the height/trace consequence on the former
domain. H9-T11 uses an elliptic family which is algebraically trivial
by construction, so it never crosses that invalid domain step.
It uses polarization directly, not a N\'eron--Tate height on a trace
quotient. In particular, when $X\times E$ is taken constant over a
function field, its constant elliptic family is killed by the trace
quotient and cannot be recovered from that quotient's height.

No assertion here concerns the whole complementary arithmetic Chow
summand $\mathrm{CH}(N)$. Positivity of the separately constructed
Fermat class and its rational span, and positivity of a geometric
incidence quotient, remain different statements with their stated
domains.

\paragraph{Status.}
This is a conditional deduction written in this run, with its complete
intersection and elliptic Poincar\'e calculations included. It is not
claimed to reproduce the unpublished 1964 argument, and it is not an
unconditional proof of the Hodge standard conjecture for arbitrary
varieties. It concerns the incidence polarization itself; it does not
identify a function-field N\'eron--Tate height on a trace quotient with
every primitive middle cycle.


% H9-T12: exact finite-base-change statements and the remaining new-cycle gate.
% Run JAC1964-20261121; /root/h9_height; GPT-6 Astra Pro.

\subsection*{H9-T12. Finite extensions: what persists and the exact new summand}

Let $K=k(\mathbf P^1)$ and let $Y/K$ be the generic fibre of a genuine
Lefschetz pencil on an abelian $2m$-fold $X/k$, with $k$ algebraically
closed. Keep the notation of H9-T7. Let $L/K$ be a finite extension of
degree $d$, and let $C/k$ be the smooth projective curve with function
field $L$. Write
\[
 \mathcal H_K=\operatorname{CH}^m(Y)_{\rm hom,\mathbf Q},
 \qquad
 \mathcal H_L^{\rm adm}
 =\operatorname{CH}^m(Y_L)_{\rm hom,\mathbf Q}
       \cap\operatorname{CH}^m(Y_L)_{(0),\mathbf Q},
\]
where the second factor denotes Kahn's saturated admissible height
domain. No equality with all homologically trivial cycles is built into
this notation. H9-T7 and JAC-P07 show that all of $\mathcal H_K$ is
admissible.

\paragraph{(i) Purely inseparable extensions give the whole theorem again.}
If $L/K$ is purely inseparable, restriction gives an isomorphism
\[
 \mathrm{res}:\operatorname{CH}^m(Y)_{\mathbf Q}
             \xrightarrow{\ \sim\ }
             \operatorname{CH}^m(Y_L)_{\mathbf Q},
 \qquad \mathrm{res}^{-1}=d^{-1}\mathrm{cor}.                \tag{H9.39}
\]
It identifies the two homologically trivial subgroups. Every such
cycle over $L$ is admissible, and
\[
 h_L(\mathrm{res}\,x,\mathrm{res}\,y)=d\,h_K(x,y).           \tag{H9.40}
\]
Thus H9-T7's full isomorphism and intersection-form decomposition hold
over $L$ after multiplying the displayed form by $d$. In the positive
cases H9-T8 and JAC-P09, the full height remains positive semidefinite
and its radical is exactly $\mathrm{res}(\mathcal H_0)$.

\emph{Proof.} The finite flat map $\pi:Y_L\to Y$ is a universal
homeomorphism. The push--pull identities on Chow groups are
$\pi_*\pi^*=d$ and $\pi^*\pi_*=d$. The second can be checked on an
integral cycle: there is one component above its image, and its residue
field degree times the flat-pullback multiplicity is $d$. Both
identities respect rational equivalence. This proves (H9.39).
Prime-to-characteristic geometric \'etale cohomology is unchanged by
purely inseparable base extension, compatibly with cycle classes, so
the homologically trivial groups correspond as well. Pullback preserves
the admissible domain, and the finite-base-change identity for heights
gives (H9.40). These are Kahn, Proposition 3.9(i), (iv), printed
pp.~20--23, and its extension in Theorem 4.14. Every cycle over $L$
is a pullback rationally by (H9.39), so all the assertions follow.
No smooth blowup model after the inseparable base change is needed.
\hfill$\square$

\paragraph{(ii) An orthogonal restriction--norm decomposition for every finite extension.}
For arbitrary finite $L/K$, restriction and corestriction preserve the
groups just specified and satisfy
\[
 \mathrm{cor}\,\mathrm{res}=d,
 \qquad
 h_L(\mathrm{res}\,x,z)=h_K(x,\mathrm{cor}\,z).               \tag{H9.41}
\]
Consequently
\[
 \begin{split}
 \mathcal H_L^{\rm adm}
   &=\mathrm{res}(\mathcal H_K)\ \mathbin{\perp}\
                           \ker(\mathrm{cor}),\\
 z&=d^{-1}\mathrm{res}(\mathrm{cor}\,z)
       +\bigl(z-d^{-1}\mathrm{res}(\mathrm{cor}\,z)\bigr).
 \end{split}                                                \tag{H9.42}
\]
The first form is exactly $d$ times the old form. In particular, when
the old height is positive semidefinite, positivity of the whole
admissible group over $L$ is equivalent to positivity on the explicit
new norm-zero subgroup $\ker(\mathrm{cor})$. Its radical is
\[
 \mathrm{res}(\operatorname{rad}h_K)
       \oplus\operatorname{rad}
          \bigl(h_L|_{\ker(\mathrm{cor})}\bigr).              \tag{H9.43}
\]
If $L/K$ is Galois with group $G$, the first summand is exactly the
$G$-invariants, and its orthogonal projector is
$|G|^{-1}\sum_{\sigma\in G}\sigma$.

\emph{Proof.} Pushforward of a geometrically homologically trivial cycle
is homologically trivial, and the admissibility statements are the
pullback and pushforward assertions in Kahn, Proposition 3.9(i)--(ii),
extended by Proposition 4.12(c). The first identity in (H9.41) is
finite flat push--pull. The second is exactly Proposition 3.9(iv),
equation (3.13), after taking the degree on the two base curves;
Theorem 4.14(b) extends it to the rational height domains. Taking
degrees of equation (3.14) gives (H9.40) for arbitrary finite extension.

The operator $R=d^{-1}\mathrm{res}\,\mathrm{cor}$ is an idempotent,
and (H9.41) makes it self-adjoint. Its image is the old restriction
space, and its kernel is the norm-zero space. This proves the direct
sum and its orthogonality, without using any positivity. The form on
its image is (H9.40). Restriction is injective rationally because
$\mathrm{cor}\,\mathrm{res}=d$. Formula (H9.43) follows by testing a
vector separately against the two orthogonal summands. In the Galois
case, $\mathrm{res}\,\mathrm{cor}=\sum_{\sigma\in G}\sigma$ on
cycles and hence on Chow groups, proving the last assertion.
\hfill$\square$

\paragraph{(iii) All homologically trivial cycles are admissible for covers unramified above the critical values.}
Suppose $L/K$ is separable and $C\to\mathbf P^1$ is unramified over
every singular-fibre parameter of the pencil. Then
\[
 \mathcal H_L^{\rm adm}
       =\operatorname{CH}^m(Y_L)_{\rm hom,\mathbf Q}.          \tag{H9.44}
\]
Thus (H9.42)--(H9.43) describe the entire height group in this case.

\emph{Proof.} The base-changed total space
$M=\widetilde X\times_{\mathbf P^1}C$ is smooth. Above a smooth
fibre this follows from smoothness of the original morphism and the
smoothness of $C/k$. Above a critical value the base change is \'etale,
so smoothness of $\widetilde X$ is preserved. Its singular fibres are
the same nodal fibres after an unramified local extension. Their
specialization maps in degree $2m$ are injective by JAC-P07.

Any cycle over $Y_L$ has a closure on $M$. If its generic class is zero,
the class restricts to zero on every smooth fibre by smooth proper base
change, and then on every singular fibre by that local specialization
injectivity. The cap-product argument of H9-T2 makes the closure
orthogonal to every vertical cycle, hence admissible. This proves
(H9.44). It does not identify the new norm-zero group with cycles on
$X$ or on the original axis.\hfill$\square$

\paragraph{(iv) The complete algebraically trivial primitive domain over every finite extension.}
Let $B=P^m(Y)_{\rm prim}$ and accept 032. H9-T4, H9-T6a and JAC-P07
provide its primitive polarization. For every finite $L/K$, this
polarization remains ample after base change. The incidence variety
over $L$ is this base-changed isogeny class: both are the rational
descent of the same geometric incidence quotient, by JAC-P00.
H9-T9 therefore applies
to the entire algebraically trivial primitive domain over $L$: its
signed height is positive semidefinite and its radical is the inverse
image of the $L/k$-trace.

In fact this trace does not acquire a new variable factor after any
finite extension. Write $B\sim B_F\times B_V$ for the fixed primitive
and variable factors. JAC-P02 identifies $B_F$, up to isogeny, with
the primitive projection of the image of the constant
$P^m(X)_K$. It is the original trace factor. If $B_V\ne0$, the proof
of H9-T6a identifies every $V_\ell B_V$ with the variable local system,
which remains absolutely irreducible after every finite-index
restriction and has a nonzero local monodromy operator. Hence this
representation has no invariant vector after any finite extension.
A homomorphism from a constant abelian variety to $(B_V)_L$ would
induce a map from a trivial Tate-module representation into those
invariants, and is therefore zero by faithfulness. Thus
\[
 \operatorname{Tr}_{L/k}(B_V)_L=0,
 \qquad
 \operatorname{im}(\operatorname{Tr}_{L/k}B_L)
        \sim (B_F)_L.                                      \tag{H9.45}
\]
Purely inseparable extensions do not alter the absolute Galois
representation; a nonzero unipotent operator also remains nonzero
under a finite ramification power in characteristic-zero coefficient
fields. These observations cover all finite extensions in (H9.45).
In codimension two, the proof of H9-T10 applies directly to $Y_L$:
the fixed-part argument is unchanged, and the variable argument uses
exactly the finite-index irreducibility just established. It gives the
same statement with the actual
$\operatorname{Ab}^2(Y_L)_{\rm prim}$ in place of $P^2$, through the
proved isogeny comparison. This does not assume an integral purely
inseparable base-change theorem for universal regular representatives.

\paragraph{The exact remaining limitation.}
For a separable cover, new norm-zero homologically trivial cycles can
appear. They need not be algebraically trivial and so need not lie in
the domain in (iv). The old blowup calculation and its axis form do not
calculate their self-heights. Parts (ii)--(iii) isolate this issue as a
specific summand; they do not assert its positivity. For covers ramified
over critical values, (ii) is stated on the admissible subgroup, while
(i) already handles the purely inseparable case completely. Proving
admissibility and positivity for all remaining cycles would require
additional local/model or Hodge-standard input, which is not supplied
by replacing the base curve in H9-T7's blowup formula.

\paragraph{Source and novelty scope.}
Restriction--corestriction, invariance of rational Chow groups under a
finite flat universal homeomorphism, and Kahn's base-change formulas
are background results used with explicit proofs or exact locators
above. Their combination with the pencil height decomposition and
the finite-index monodromy theorem is the deduction recorded here.
No formal verification or historical originality is asserted.


% H9-T13: generic odd complete intersections in abelian varieties.
% Run JAC1964-20261121; /root/h9_height; GPT-6 Astra Pro.

\subsection*{H9-T13. Generic odd complete intersections in an abelian variety}

\begin{theorem}[H9-T13]
Let $A/k$ be an abelian variety of dimension
$N=2m+c-1$, where $m\ge1$ and $c\ge1$, over an algebraically closed
field of arbitrary characteristic. Let $h$ be an ample divisor class.
Choose a smooth successive complete intersection
\[
 X=A\cap H_1\cap\cdots\cap H_{c-1},\qquad [H_i]=h,
 \qquad \dim X=2m,
\]
with the convention $X=A$ when $c=1$. Let $Y/K$ be the generic member
of a genuine Lefschetz pencil of sections of $X$ in the restricted
class $h$, where $K=k(\mathbf P^1)$. Thus
$\dim Y=2m-1$ and $Y$ has total codimension $c$ in $A$.

Accept 032, Corollary 8.4. The normalized signed incidence class
\[
 \frac{(-1)^m}{2}\Delta^*d
       \quad\text{on }P^m(Y)_{\rm prim}
\]
is a rational polarization. The first $c-1$ smooth sections need not
be generic; genericity is required only of the last member through
the stated genuine Lefschetz-pencil hypothesis.

For a starting projective embedding with class $h_0$, taking
$h=d h_0$ with $d\ge2$ supplies such terminal Lefschetz pencils in
every characteristic. In particular the conclusion holds for a
geometric generic complete intersection of $c$ degree-$d$ sections
of an embedded abelian variety, for every $d\ge2$.

If $m=2$, the actual primitive regular algebraic representative
$\operatorname{Ab}^2(Y)_{\rm prim}$ has the same polarization, and
its canonical map to $P^2(Y)_{\rm prim}$ is an isogeny.
\end{theorem}

\paragraph{The fixed and variable projectors.}
Put $n=2m-1$, let $i:X\hookrightarrow A$, $j:Y\hookrightarrow X$,
and $k=i\circ j:Y\hookrightarrow A$. Weak Lefschetz identifies
$H^n(X)$ with $H^n(A)$ and embeds this space in $H^n(Y)$.
On the ambient subspace the middle pairing is
\[
 (u,v)\longmapsto\int_Auvh^c.
\]
It is nondegenerate, because hard Lefschetz on the $N$-dimensional
abelian variety gives the isomorphism
\[
 L_h^c:H^n(A)\xrightarrow{\ \sim\ }H^{n+2c}(A).
\]
Its inverse is algebraic. Let
\[
 r=(L_h^c)^{-1}k_*:H^n(Y)\longrightarrow H^n(A),
 \qquad r k^*=1.
\]
Choose this inverse supported in the indicated degrees, using the
algebraic K\"unneth projectors of $A$. It gives an algebraic
correspondence $r$. Hence
\[
 \Pi_V=1-k^*r
\]
is the algebraic self-adjoint variable projector on $H^n(Y)$.
This is exactly the variable space for the terminal pencil: the
ambient space of that pencil is $j^*H^n(X)=k^*H^n(A)$.
No separate standard-conjecture hypothesis for $X$ is needed for this
projector.

Let $e_A$ be the algebraic primitive projector on $H^n(A)$, whose
image is $\ker L_h^{c+1}$. Put
\[
 F=k^*\ker L_h^{c+1},\qquad \Pi_F=k^*e_A r.
\]
The ambient primitive decomposition is orthogonal for
$\int_Auvh^c$: if $L_h^{c+1}u=0$ and $v=L_hw$, then
$\int_Auvh^c=\int_A(L_h^{c+1}u)w=0$.
Thus $\Pi_F$ is self-adjoint and orthogonal to $\Pi_V$.

We check precisely that these give the full primitive part on $Y$.
The Gysin map
\[
 k_*:H^{n+2}(Y)\longrightarrow H^{n+2+2c}(A)
\]
is injective by dual weak Lefschetz in degree $n-2$ (and the assertion
is immediate when the source is zero, as for $n=1$).
For an ambient class $k^*u$, its product with $h_Y$ pushes to
$h^{c+1}u$. Therefore it is primitive on $Y$ exactly when
$u\in\ker L_h^{c+1}$. A variable class $v\in\ker j_*$ is primitive
as well: $j_*(h_Yv)=h_Xj_*v=0$, and
$j_*:H^{n+2}(Y)\to H^{n+4}(X)$ is injective by the same dual weak
Lefschetz argument. Consequently
\[
 H^n(Y)_{\rm prim}=F\ \mathbin{\perp}\ E,
 \qquad E=\ker j_*.
                                                               \tag{H9.46}
\]
Incidence detection and Saito's transpose compatibility turn these
algebraic projectors into orthogonal nondegenerate factors of
$P^m(Y)$ in the rational category, as in JAC-T4.1 and H9-T4.

\paragraph{The fixed-part curve test with all exponents.}
Let $C_0$ be a smooth pointed parameter curve of genus $b\ge1$, and
pull a normalized family on the fixed factor back to it. After the
parameter degree-one projection its class is
\[
 \xi\in H^1(C_0)\otimes F.
\]
Set $\eta=(1\otimes r)\xi$. This is algebraic, satisfies
$k^*\eta=\xi$, and obeys $h^{c+1}\eta=0$.
The algebraic inverse to Abel pullback lifts it to an algebraic class
\[
 w\in H^1(J(C_0))\otimes H^n(A),
       \qquad h^{c+1}w=0,
\]
on the abelian variety $J(C_0)\times A$ of dimension
$b+2m+c-1$. The class $w$ has codimension $m$.
Let $\theta$ be the canonical principal polarization of $J(C_0)$
and put $H=\theta+h$. Then
\[
 H^{b+c}w=0.
\]
Indeed every term with $h$-exponent at least $c+1$ vanishes, while
every remaining term has $\theta$-exponent at least $b$ and hence
exceeds the top degree of $J(C_0)$ after multiplication by its
degree-one component. This is exactly the primitive exponent
$(b+2m+c-1)-2m+1=b+c$.

In the quadratic form, degree on the two factors permits only the term
with $h$-exponent $c$ and $\theta$-exponent $b-1$. The Abel curve-class
formula therefore gives
\[
 \begin{aligned}
 \int_{J(C_0)\times A}w^2H^{b+c-1}
  &=\binom{b+c-1}{c}
       \int_{J(C_0)\times A}w^2\theta^{b-1}h^c\\
  &=\frac{(b+c-1)!}{c!}
       \int_{C_0\times A}\eta^2h^c\\
  &=\frac{(b+c-1)!}{c!}
       \int_{C_0\times Y}\xi^2.
 \end{aligned}
 \tag{H9.47}
\]
The coefficient is positive. By 032, the right-hand intersection has
sign $(-1)^m$ and is zero only if $\xi=0$. This proves nefness of the
signed incidence divisor on the fixed factor. Its nondegeneracy was
already supplied by the self-adjoint-factor argument, so the
nef-plus-nondegenerate criterion makes it a rational polarization.
Curves of genus zero have constant maps to an abelian variety and
contribute degree zero; they cause no exception to the curve test.

\paragraph{The variable factor and the terminal pencil hypothesis.}
The terminal ambient variety $X$ is smooth projective of even
dimension $2m$, and its pencil is genuinely Lefschetz. Hard Lefschetz
holds for $X$. SGA 7 II, XVIII, 6.6 and 6.7.2, printed pp.~323--327,
therefore give spanning and conjugacy of the vanishing cycles and
absolute irreducibility after every finite-index monodromy restriction.
Their characteristic-two exception is restricted to odd ambient
dimension, so it is absent here.

If the variable incidence factor is nonzero, its normalized family
gives an injective realization into $E_\ell$ for every invertible
prime. Thus $E_\ell\ne0$, every local vanishing cycle is nonzero,
and the hypotheses of H9-T6 are satisfied, exactly as in H9-T6a.
That theorem proves positivity of the variable factor. If it is zero,
there is nothing to prove on that factor. Combining its polarization
with the fixed polarization through (H9.46) proves the main assertion.

SGA 7 II, XVII, Theorem 2.5.1, printed p.~217, says that a degree-$d$
multiple of a projective embedding is Lefschetz for $d\ge2$. Apply
this to the embedding of the already chosen smooth $X$ induced from
the starting embedding of $A$. It supplies the last-pencil hypothesis
in the stated degree-$d$ construction. Successive Bertini sections
supply the earlier smooth chain. No assertion that an arbitrary
nodal morphism has the global monodromy of such a pencil is used.

\paragraph{The actual representative when $m=2$.}
Existence over the possibly imperfect $K$ and injective realizations
of normalized miniversal families are exactly the source-checked
ACMV results used in H9-T10. They make the same projectors act on
$\operatorname{Ab}^2(Y)$ and give the variable injection required
for H9-T6's sign argument. On the fixed factor, a nonconstant
geometric curve has a nonzero pulled-back family class because its
Tate-module map remains nonzero after that injective realization.
Equation (H9.47), now with $m=2$, makes the incidence degree strictly
positive. The argument excluding a positive-dimensional kernel and
the canonical comparison with the incidence quotient are verbatim
the explicit arguments of H9-T10. They prove the stated rational
polarization and isogeny, without any integral purely inseparable
base-change assertion.\hfill$\square$

\paragraph{Why this does not close the arbitrary fixed-axis gate.}
Let $\kappa:W\hookrightarrow X_0$ be the codimension-two axis in
an abelian variety $X_0$ from H9-T7. The odd variety $W\times E$ is
cut out in $X_0\times E$ by two pullbacks of divisors from $X_0$.
These divisor classes have degree zero along every elliptic fibre
$\{x\}\times E$ and are not ample. Degree on that fibre is constant
in a connected algebraic deformation of the line bundle, so one
cannot make them ample while retaining those divisor classes.

There is also an exact obstruction to using the ambient transfer
above on the remaining axis classes. For
$w\in A_W=\ker(\kappa_*:\operatorname{Alg}^{m-1}(W)
\to H^{2m+2}(X_0))$, the elliptic family $w\boxtimes P_E$
on $E\times(W\times E)$ pushes to
\[
 (\kappa_*w)\boxtimes P_E=0
       \quad\text{in the cohomology of }E\times X_0\times E.
\]
Applying ambient algebraic Lefschetz inverses to this zero class
still gives zero. The transfer in (H9.47) was injective on fixed
cohomology; this attempted transfer kills precisely the unresolved
axis summand. A new injective correspondence or a theorem that lifts
these special classes through a degeneration would be required.
Generic positivity does not supply either assertion. The separate
classical characteristic-zero Hodge--Riemann argument for special
fibres is unaffected by this positive-characteristic limitation.

\paragraph{Status and archive scope.}
This is a deduction written in this run from 032, algebraic Lefschetz
operators on abelian varieties, and the explicitly cited terminal-pencil
monodromy theorems. It extends the primitive incidence construction of
cote 44, scans 8--9, to the indicated generic odd complete intersections.
It neither asserts polarization on every special member in positive
characteristic nor upgrades the original H9 homological domain to an
incidence Jacobian. No historical originality or formal verification
is claimed.


% H9-T14: constant finite-field varieties after transcendental base extension.
% Run JAC1964-20261121; /root/h9_height; GPT-6 Astra Pro.
% Input P15 uses the Tate consequence032 Cor8.3 and Kahn2003 Th1.10.

\subsection*{H9-T14. Constant finite-field varieties and the full primitive height domain}

\begin{lemma}[Spreading algebraic equivalence]\label{h9:spread}
Let $k_0$ be an algebraically closed field and let $X_0/k_0$ be
smooth projective. Fix a codimension $r$ and a prime $\ell$ invertible
in $k_0$, and use rational $\ell$-adic cycle classes. These classes are
compatible with smooth proper base change and are unchanged under
extension of algebraically closed fields.
Suppose
\[
 \operatorname{CH}^r(X_0)_{\rm hom,\mathbf Q}=0.
                                                               \tag{H9.48}
\]
Then, for every extension field $F/k_0$, including imperfect fields,
\[
 \operatorname{CH}^r((X_0)_F)_{\rm hom,\mathbf Q}
       =\operatorname{CH}^r((X_0)_F)_{\rm alg,\mathbf Q}.       \tag{H9.49}
\]
Homological triviality on the left is geometric. Algebraic triviality
on the right means algebraic equivalence already over $F$, with
rational coefficients, not merely after passage to $\overline F$.
\end{lemma}

\begin{proof}
Only the inclusion from left to right requires proof. Let $z$ be a
rational cycle class on $(X_0)_F$ whose geometric class is zero.
Choose a cycle representative and a finitely generated subfield
$F_0/k_0$ of $F$ carrying that cycle. Since $k_0$ is perfect,
$F_0/k_0$ is separably generated. It has a smooth geometrically
integral affine model $S/k_0$ with $k_0(S)=F_0$.

Taking closures of the finitely many cycle components and shrinking
$S$ if necessary gives a relative cycle
\[
 Z\in\operatorname{CH}^r(X_0\times S)_{\mathbf Q}
\]
whose generic fibre pulls back to $z$, and whose components are flat
over $S$. All fibres are interpreted by refined Gysin pullback; on
this flat locus these are the usual cycle-theoretic fibres. The
global cycle class of $Z$ has an image under the edge restriction map
\[
 H^{2r}(X_0\times S,\mathbf Q_\ell(r))
  \longrightarrow
 H^0\!\left(S,R^{2r}\operatorname{pr}_{S*}\mathbf Q_\ell(r)\right).
                                                               \tag{H9.50}
\]
Here the displayed sheaf is constant, with fibre
$H^{2r}(X_0,\mathbf Q_\ell(r))$, because the family is the product
family. The image of the class at a geometric generic point is zero:
its pullback to $\overline F$ is the zero class of $z$, and extension
between algebraically closed fields preserves this cohomology.
Since $S$ is connected, the section in (H9.50) is zero everywhere.
Compatibility of cycle classes with Gysin restriction therefore
gives $\operatorname{cl}(Z_s)=0$ for every geometric fibre.

Choose a closed point $s\in S(k_0)$; such a point exists because
$k_0$ is algebraically closed and $S$ is a nonempty scheme of finite
type. Hypothesis (H9.48) gives $Z_s=0$ in
$\operatorname{CH}^r(X_0)_{\mathbf Q}$.

There are now two \emph{$F$-rational} points of the smooth
geometrically integral parameter space $S_F$: the point $s_F$ and
the point $t$ induced by $F_0=k_0(S)\hookrightarrow F$. Although
$t$ maps to the generic point of $S$, it is a rational closed point
of $S_F$. By construction,
\[
 (Z_F)_t=z,\qquad (Z_F)_{s_F}=(Z_s)_F=0.
                                                               \tag{H9.51}
\]
Thus (H9.51) is already a smooth connected parameter-space witness
for algebraic triviality over $F$.

To verify the curve formulation explicitly, if $\dim S>0$, choose
a smooth geometrically integral curve in the quasi-projective smooth
variety $S_F$ through $t$ and $s_F$. The required Bertini statement
is Achter--Casalaina-Martin--Vial, \emph{Parameter spaces for
algebraic equivalence}, Theorem A.1(1),(4), author manuscript p.~17.
Pulling $Z_F$ back to that curve preserves (H9.51). If $F$ is
algebraically closed, normalize its projective compactification and
close the cycle family there; the resulting smooth projective
parameter curve has the same two fibres. Over an imperfect $F$, a
smooth quasi-projective parameter curve suffices for algebraic
equivalence. If a smooth projective curve is desired, Lemma 3.8(2)
of the same source, p.~6, supplies it after multiplication by a power
of the characteristic, harmless with rational coefficients. If
$\dim S=0$, then $F_0=k_0$ and (H9.48) already says $z=0$.
This proves (H9.49).
\end{proof}

\paragraph{The numerical variant, without a numerical/homological conjecture.}
Suppose instead that
$\operatorname{CH}^r(X_0)_{\rm num,\mathbf Q}=0$.
Then for every extension $F/k_0$ one has
\[
 \operatorname{CH}^r((X_0)_F)_{\rm num,\mathbf Q}
   =\operatorname{CH}^r((X_0)_F)_{\rm hom,\mathbf Q}
   =\operatorname{CH}^r((X_0)_F)_{\rm alg,\mathbf Q}.         \tag{H9.49a}
\]
Here numerical triviality may even be tested only against cycles
defined over $F$; the proof shows it implies geometric numerical
triviality as well. To see this, spread a numerically trivial $z$
as above. For \emph{each} complementary cycle $w$ on $X_0$, form
\[
 \operatorname{pr}_{S*}
       \bigl(Z\cdot\operatorname{pr}_{X_0}^*w\bigr)
       =d_w[S]\quad\text{in }\operatorname{CH}^0(S)_{\mathbf Q}.
\]
The intersection product exists because $X_0\times S$ is smooth;
no proper-intersection assumption on the chosen representatives is
needed. The projection is proper, and $S$ is integral, so a
codimension-zero class has precisely this form. Compatibility of
proper pushforward with refined Gysin restriction says that every
fibre degree $\deg(Z_s\cdot w)$ equals $d_w$. The same assertion
at the generic point is flat base change. Pulling that generic
intersection to $F$ preserves its degree, and numerical triviality
of $z$ gives $d_w=0$.

Thus the closed $k_0$-fibre $Z_s$ is numerically trivial, since the
argument applies to every complementary $w$ on $X_0$. It is zero
in rational Chow by the new hypothesis. Equation (H9.51) again
proves that $z$ is algebraically trivial over $F$. Finally,
algebraic triviality implies geometric homological triviality,
which implies numerical triviality. These inclusions prove all
equalities in (H9.49a). This argument uses intersection degrees,
not a standard conjecture identifying numerical and homological
equivalence.

\begin{theorem}[H9-T14]\label{h9:T14}
Put $k_0=\overline{\mathbf F}_p$. Accept the Tate consequence of
032, Corollary 8.3, and the established Kahn input in JAC-P15.

\begin{enumerate}
\item
Let $X_0/k_0$ be smooth projective with geometric rational Chow
motive of abelian type. For every extension field $F/k_0$ and every
codimension $r$ one has the three equalities (H9.49a). In particular
this applies to every
abelian variety over $k_0$ and to every product of curves over $k_0$.
For an abelian variety isogenous over $F$ to $(A_0)_F$, with
$A_0/k_0$ abelian, the same equalities hold. Consequently an integral
cycle on any of these varieties is numerically trivial if and only
if a nonzero multiple is algebraically trivial: this is customary
$\tau$-triviality, including after transcendental extension.

\item
Let $k/k_0$ be algebraically closed, let $C/k$ be a smooth projective
connected curve, and put $K=k(C)$. Let $A_0/k_0$ be an abelian
variety of dimension $2m-1$, and write $A=(A_0)_k$. Choose an ample
class on $A$ and let
\[
 B=P^m(A)_{\rm prim},\qquad
 \ell_B=\frac{(-1)^m}{2}\Delta^*d_B.
\]
Accept also 032, Corollary 8.4. Then $\ell_B$ is a rational
polarization by JAC-T4. Let
\[
 a_K:\operatorname{CH}^m(A_K)_{\rm hom,\mathbf Q}
          \longrightarrow P^m(A)_K(K)_{\mathbf Q}
\]
be the incidence map, which is now defined on the entire displayed
homological group by part (1). Its primitive domain is
\[
 \mathcal H_{\rm prim}
   :=a_K^{-1}(B(K)_{\mathbf Q}).
                                                               \tag{H9.52}
\]
Every class in this domain is admissible for the geometric height.
The normalized signed middle height
$h_m=(-1)^m H_{A_K}$ is positive semidefinite on
$\mathcal H_{\rm prim}$, with exact radical
\[
 \operatorname{rad}(h_m|_{\mathcal H_{\rm prim}})
      =a_K^{-1}(B(k)_{\mathbf Q}).                           \tag{H9.53}
\]
It induces a positive definite form on the quotient
\[
 \frac{\mathcal H_{\rm prim}}
       {a_K^{-1}(B(k)_{\mathbf Q})}
      \ \simeq\
 \frac{B(K)_{\mathbf Q}}{B(k)_{\mathbf Q}}
      \ \simeq\
 \operatorname{Hom}_k(J(C),B)\otimes\mathbf Q.               \tag{H9.54}
\]
The same assertion holds on the image of an actual primitive middle
Chow projector for $A$, with the radical obtained by intersection
with (H9.53). A normalized projected universal family makes its
quotient in (H9.54) surjective as well.
\end{enumerate}
\end{theorem}

\begin{proof}
JAC-P15 proves rational equivalence equals numerical equivalence on
$X_0$ with rational coefficients. Thus the numerical variant of the
spreading lemma applies in every codimension and proves part (1).
For an integral class, membership in the rational span of the
algebraically trivial subgroup means that a nonzero integer multiple
belongs to that subgroup; clear denominators and any remaining
torsion. The converse follows from numerical intersection numbers
having no integer torsion. This proves the $\tau$ formulation.

If $f:(A_0)_F\to A'$ is an isogeny and $z$ is numerically trivial
on $A'$, the projection formula shows that $f^*z$ is numerically
trivial and hence algebraically trivial by the result just proved.
Proper pushforward preserves algebraic
equivalence, and $f_*f^*z=\deg(f)z$. Dividing by $\deg(f)>0$
proves the isogeny variant, including inseparable isogenies.

For part (2), take $F=K$ in part (1). Every geometrically
homologically trivial class is algebraically trivial over $K$ with
rational coefficients. Kahn's Theorem 5.6 puts all of these classes
in the saturated admissible height domain. The primitive incidence
polarization supplied by JAC-T4 remains ample over $K$.
H9-T9 therefore gives positive semidefiniteness and identifies the
radical with the inverse image of the $K/k$-trace of $B_K$.

Since $B_K$ is constant and $k$ is algebraically closed, its trace is
$B$, with the identity trace morphism. This assertion concerns the
trace \emph{variety}; its group of trace points in $B(K)$ is
$B(k)$, which need not equal $B(K)$. Thus H9-T9 gives exactly
(H9.53), not the assertion that every height is zero.

For completeness the first arrow of (H9.54) is surjective.
A normalized rational universal incidence family on $B\times A$
induces the identity of $B$ in the isogeny category. For $b\in B(K)$
its pointed fibre difference is an algebraically trivial cycle with
incidence image $b$, after clearing the single normalization
denominator. Applying the primitive projector to this family gives
the same surjectivity on the primitive Chow image. This is the same
normalized-family construction used in H9-T9, and no surjectivity
from arbitrary homological cycles is being assumed independently.

To prove the last isomorphism of (H9.54), choose $c_0\in C(k)$.
A point $b\in B(K)$ extends uniquely to a morphism $f_b:C\to B$:
properness extends it over every discrete valuation ring on the
smooth curve. The map $f_b-f_b(c_0)$ factors uniquely through the
Abel map $C\to J(C)$. Hence
\[
 B(K)=B(k)\oplus\operatorname{Hom}_k(J(C),B)
\]
as abelian groups after the chosen base point, and the quotient
isomorphism is independent of the constant translation. Tensoring
with $\mathbf Q$ gives (H9.54). In particular a nonzero homomorphism
$J(C)\to B$ is a positive direction in this quotient.

One can also see its height normalization directly. For an integral
point $b$ represented by $f_b:C\to B$, pull the normalized universal
family back to $C\times A$. All fibres are homologically trivial,
so its closure on this smooth product model is admissible. The
definition of the incidence divisor gives
\[
 h_m(z_b,z_b)
      =(-1)^m\deg_C f_b^*\Delta^*d_B
      =2\deg_C f_b^*\ell_B.                               \tag{H9.55}
\]
Rigidification removes constant-family terms; equivalently their
pairings vanish by the incidence-kernel/trace argument in H9-T9.
For rational points, clear denominators first and divide the two
sides by their square. Formula (H9.55) is positive for a nonconstant
$f_b$ and zero for a constant map, because $\ell_B$ is ample.
It is consistent with (H9.53)--(H9.54).
\end{proof}

\paragraph{Why passage through an algebraic closure causes no hidden descent assumption.}
The proof of the spreading lemma already constructs its two parameter
points over $F$ itself. It therefore bypasses any assumption that
integral algebraic equivalence descends from $\overline F$.
Here is the alternative rational norm argument, to specify its scope.
An algebraic-equivalence witness over $\overline F$ is defined over
a finite extension $L/F$. Achter--Casalaina-Martin--Vial, Lemma 4.14,
author pp.~14--15, applied with rational coefficients, replaces that
witness by a family $Z$ on a smooth geometrically integral curve
$T/F$, with two $L$-points $t_0,t_1$. This statement includes
inseparable $L/F$ because the characteristic is inverted. Proper
pushforward of the two fibres gives
\[
 [L:F]z
  =Z_{\beta},\qquad
 \beta=(t_1)_*[\operatorname{Spec}L]
            -(t_0)_*[\operatorname{Spec}L],\qquad
 \deg_F\beta=0.
\]
Lemma 3.9 and the proof of Lemma 4.15 in that source show that
$Z_\beta$ is algebraically trivial over $F$. Division by $[L:F]$
proves rational descent. This is precisely a degree-zero cycle
norm argument; it does not assert integral descent over imperfect
fields or treat an inseparable parameter curve as smooth over $F$.

\paragraph{Other constant varieties and exact remaining limits.}
The same height conclusion applies to a constant odd-dimensional
$X_0/k_0$ of abelian type whenever an algebraic primitive factor and
its signed incidence polarization have been established: use part (1)
to supply the homological-to-algebraic domain equality and H9-T9 to
supply the height and trace statement. No polarization for an
unspecified abelian-type variety is inferred from the phrase
``abelian type'' alone.

The equalities in (H9.49a) do not assert rational equivalence equals
numerical equivalence after transcendental extension. Even for
abelian varieties, the theorem concerns varieties isogenous
to constants from $\overline{\mathbf F}_p$. It does not apply to a
general abelian variety with transcendental moduli, nor does it
annul every algebraically trivial cycle after extension of constants.
For example a constant elliptic curve has nonconstant points over
its own function field. Their degree-zero divisors are homologically
and algebraically trivial, while their signed geometric self-height
is positive by (H9.55). The radical includes the incidence kernel
and the preimage of constant trace points, not all incidence points.

This gives a valid full homological domain for the corrected H9(b)
height/trace statement in this constant finite-field situation. It
does not repair the invalid homological-to-incidence map for arbitrary
pencils, and it has no implication for the whole arithmetic
$\operatorname{CH}(N)$ of the number-field Fermat Jacobian.

\paragraph{Sources and status.}
The new deduction is the combination of JAC-P15 with the displayed
spreading lemma and the already proved primitive height theorem.
It uses the \emph{Tate} consequence032 Corollary 8.3 through
Kahn2003, Theorem 1.10, published p.~982, and uses the separate
Hodge-standard consequence032 Corollary 8.4 for the primitive
incidence polarization. The parameter-space source was directly read
at \url{https://www.math.uni-bielefeld.de/~vial/abtriv.pdf}:
Definition 3.1 and Proposition 3.5 pp.~4--5, Lemma 3.8 p.~6,
Lemma 3.9 pp.~6--7, Theorem 4.11 and Lemmas 4.14--4.15 pp.~14--16,
and Theorem A.1 pp.~17--18. These are author-manuscript page numbers.
No historical originality or formal verification is asserted.
\section{Checked arithmetic reductions and explicit remaining hypotheses}
The following deductions include conditional implications with substantive
hypotheses. A proved reduction is distinguished from a proof of its
remaining height inequality. The statement-level checking events are
recorded in the accompanying log.
\section{Arithmetic sibling of cote 44: corrected theorem and remaining
gate}\label{arithmetic-sibling-of-cote-44-corrected-theorem-and-remaining-gate}

Run: \texttt{JAC1964-20261121}. Author/checker:
\path{/root/arithmetic_sign}, GPT-6 Astra Pro; effort setting not
exposed. Date: 2026-11-21 UTC. These are working proof sections for
incorporation into the parent manuscript. They are not an independent
review of the parent manuscript. No Lean build was attempted here.
Accepted attached results are used at their stated scope.

\subsection{A. Conventions and exact
background}\label{a.-conventions-and-exact-background}

Let \(K\) be a number field, \(A/K\) an abelian variety of dimension
\(g\), and \(\ell=c_1(L)\) for a symmetric ample line bundle. All Chow
groups and Chow motives have rational coefficients. Write \(L\) also for
multiplication by \(\ell\). The raw Beilinson--Bloch pairing, normalized
by \([K:\mathbf Q]\), is denoted \(h_K\). Thus \(h_K\) is invariant
under finite extension of the field of definition. This normalization is
distinct from the choice whether the bilinear Néron--Tate form is the
polarization of \(\widehat h\) with or without a factor two; below the
Néron--Tate convention is Künnemann's.

\textbf{Background AR-B1 (existence and functoriality).} Künnemann,
\emph{Height pairings for algebraic cycles on abelian varieties}, Ann.
Sci. ENS 34 (2001), Theorem 1.6 and Corollary 1.7, give the pairing on
every abelian variety over a number field. Section 1.3 records
compatibility with correspondences; §1.4, equation (7), gives the
extension-of-field scaling before normalization. The pairing uses
extensions perpendicular to vertical cycles; an arbitrary closure on a
model is insufficient. The correspondence projection formula is also
Künnemann 1996, §5, equation (18), from Lemma 4.1.

\textbf{Background AR-B2 (Picard comparison and sign).} Künnemann,
\emph{Higher Picard varieties and the height pairing}, Amer. J. Math.
118 (1996), Theorem 8.2, equation (23), identifies the height of
algebraically trivial cycles with the Néron--Tate biextension pairing on
higher Picard varieties, with the positive comparison integer \(k_X^p\).
His Theorem 6.1 calibrates the raw curve height as \textbf{minus} the
positive Jacobian pairing. Let \(B^p(A_K)\) be algebraically trivial
cycles modulo geometric incidence equivalence. Theorem 12.1(ii) gives

\[
 x\ne0,\quad L^{g+2-2p}x=0
 \quad\Longrightarrow\quad
 (-1)^p h_K(x,L^{g+1-2p}x)>0.
\]

The old regular-model hypothesis in the 1996 statement can be met after
finite extension using AR-B1; the positive base-change factor does not
affect the conclusion. This is a known theorem, not a new consequence of
GRH or of the attached CM Hodge theorem. For \(g=2m-1,p=m\), the
positive candidate is

\[
q_A(x,y)=(-1)^m h_K(x,y),\qquad Lx=Ly=0.
\tag{AR.1}
\]

Consequently ``the BB height equals a positive Néron--Tate height''
requires a specified signed duality. With the parent's convention
\(\phi_L(a)=t_a^*L\otimes L^{-1}\), the signed incidence is
\((-1)^m D\). Its Hodge cup-product counterpart carries \((-1)^{m-1}\);
the incidence dictionary supplies the extra minus. The \(m=1\)
Faltings--Hriljac case is an indispensable check.

\textbf{Background AR-B3 (motivic Lefschetz).} Künnemann, \emph{A
Lefschetz decomposition for Chow motives of abelian schemes}, Invent.
Math. 113 (1993), Theorem 5.1, Theorem 5.2, and Corollary 5.3, all on
printed p.96, supply actual Chow projectors, motivic hard Lefschetz, and
transposition compatibility. For \(g=2m-1\), the primitive projector
\(\pi_P\) on \(h^g(A)\) satisfies \({}^t\pi_P=\pi_P\); the primitive
motive \(P\) has its perfect middle Poincaré pairing. These three
statements were checked directly in the rendered source page, not
inferred from a cohomological decomposition.

\textbf{Background AR-B4 (degree one is faithful).} The functor from
abelian varieties up to isogeny to degree-one Chow motives is fully
faithful (contravariantly). In particular,

\[
\operatorname{End}_{\mathrm{Chow}}(h^1B)=\operatorname{End}(B)_{\mathbf Q}^{\mathrm{op}},
\]

and the Betti realization of this endomorphism algebra is faithful. A
directly inspected modern statement is Laga--Shnidman, \emph{Ceresa
cycles of bielliptic Picard curves}, arXiv:2312.12965, Proposition 4.5,
citing Kings, \emph{Higher regulators, Hilbert modular surfaces, and
special values of L-functions}, Duke Math. J. 92 (1998), Proposition
2.2.1. The contravariant/opposite convention changes no invertibility
assertion below.

\subsection{B. Exact removal of the irrelevant Beauville
grades}\label{b.-exact-removal-of-the-irrelevant-beauville-grades}

For \([n]^*x=n^{2p-s}x\), say \(x\) has Beauville grade \(s\). The
well-known height selection rule is

\[
h_K(\operatorname{CH}^p_{(s)},\operatorname{CH}^{g+1-p}_{(t)})=0
\quad\text{unless}\quad s+t=2.
\tag{AR.2}
\]

An explicit literature locator is Zhang, \emph{Gross--Schoen cycles and
dualising sheaves}, Invent. Math. 179 (2010), §5.3, printed p.68. For
completeness its proof is one line: the projection formula gives

\[
h_K([n]^*x,[n]^*y)=n^{2g}h_K(x,y),
\]

whereas the eigenvalue computation gives \(n^{2(g+1)-s-t}\) times the
original pairing. Take \(n>1\).

\textbf{Proposition AR-1 (primitive grade reduction).} Put

\[
V^p=\ker\bigl(L^{g+2-2p}:\operatorname{CH}^p(A)_{\mathrm{hom}}
\longrightarrow\operatorname{CH}^{g+2-p}(A)\bigr),\qquad 2p\le g+1.
\]

Then \(V^p_{(s)}=0\) for \(s\ge2\). Every \(V^p_{(s)}\) with \(s\le0\)
belongs to the radical of the signed height on \(V^p\). Thus
semipositivity on the full primitive group is equivalent to
semipositivity on grade one. Strict positivity on the full group is
equivalent to strict positivity on grade one together with the vanishing
of the primitive grades \(s\le0\).

\emph{Proof.} A grade-\(s\) class of codimension \(p\) is a Chow class
of \(h^i(A)\) with \(i=2p-s\). For \(s\ge2\), put \(e=g+2-2p\). Then
\(i\le g\) and \(e\le g-i\). Motivic hard Lefschetz makes \(L^{g-i}\) an
isomorphism. If \(L^e x=0\), multiply by \(L^{g-i-e}=L^{s-2}\) to obtain
\(L^{g-i}x=0\), hence \(x=0\). Consequently every surviving primitive
grade is at most one. Multiplication by a symmetric \(\ell\) has grade
zero. Equation (AR.2) therefore pairs a primitive grade \(s\le0\) only
with grade \(2-s\ge2\), which has just been proved zero. This proves the
radical assertion and both equivalences. No assertion that all negative
Beauville grades vanish was used. \(\square\)

For \(g=2m-1,p=m\), the grade-one component of \(V^m\) is exactly

\[
V^m_{(1)}=\operatorname{CH}^m(P).
\tag{AR.3}
\]

Indeed the motivic Lefschetz decomposition of \(h^g(A)\) is a direct sum
of the primitive middle summand and longer Lefschetz strings. On the
middle-degree term of each longer string, \(L\) has a left inverse; on
\(P\), it is zero. This checks that (AR.3) is a Chow-level assertion,
rather than merely a statement about realization.

\textbf{Corollary AR-1a (the threefold has no additional lower-grade
gate).} If \(g=3,p=2\), then \(\operatorname{CH}^2_{(s)}(A)=0\) for
every \(s<0\), and \(\operatorname{CH}^2_{(0)}(A)_{\mathrm{hom}}=0\).
Consequently the full primitive homologically trivial group equals
\(\operatorname{CH}^2(P)\), and strict positivity has no separate
lower-grade vanishing condition.

\emph{Proof.} Fourier gives \[
\mathcal F:\operatorname{CH}^2_{(s)}(A)
 \simeq\operatorname{CH}^{1+s}_{(s)}(\widehat A).
\] For \(s<-1\) the target has negative codimension. For \(s=-1\) it is
\(\operatorname{CH}^0_{(-1)}=0\). For \(s=0\), a homologically trivial
class maps to a homologically trivial divisor of grade zero. Such a
divisor lies in \(\operatorname{Pic}^0\), which has grade one, so it is
zero. Fourier is invertible and preserves homological triviality. AR-1
removes all primitive grades at least two; only grade one remains.
\(\square\)

\subsection{C. A CM Chow summand with a strictly positive
height}\label{c.-a-cm-chow-summand-with-a-strictly-positive-height}

Assume the attached 032 Theorem 1.1, i.e.~the rational Hodge conjecture
for products of CM abelian varieties. Its Lemma 8.1 supplies effective
CM realizations; in weight one it is the usual realization by a CM
abelian variety. The statement below uses those accepted inputs,
classical motivic Lefschetz, and the established Picard height theorem.
It does not use GRH.

\textbf{Theorem AR-2 (orthogonal Picard summand and exact remaining
arithmetic gate).} Let \(A/K\) be CM, \(\dim A=2m-1\), with symmetric
ample \(L\). After a finite number-field extension \(K'/K\), there are a
CM abelian variety \(B/K'\), an actual self-transpose Chow idempotent
\(e\) of the primitive middle motive \(P\), and an orthogonal
decomposition

\[
P=M\oplus N,\qquad M\simeq h^1(B)(1-m),
\tag{AR.4}
\]

with the following properties.

\begin{enumerate}
\def\labelenumi{\arabic{enumi}.}
\tightlist
\item
  \(H(M)\) is the maximal rational sub-Hodge structure of \(H(P)\)
  having only types \((m,m-1)\) and \((m-1,m)\). Thus \(H(N)\) has no
  nonzero such sub-Hodge structure.
\item
  \(\operatorname{CH}^m(M)\simeq\operatorname{Pic}^0(B)(K')_{\mathbf Q}\).
  The restriction of \(q_A=(-1)^m h_{K'}\) to this full Chow summand is
  a Néron--Tate form for a rational polarization, with the positive
  comparison normalization of AR-B2. It is strictly positive on every
  nonzero class.
\item
  The two Chow summands are height orthogonal. Therefore the
  semipositivity problem on the full primitive homologically trivial
  Chow group over \(K'\) reduces \textbf{exactly} to
\end{enumerate}

\[
(-1)^m h_{K'}(z,z)\ge0
\qquad(z\in\operatorname{CH}^m(N)).
\tag{AR-GATE}
\]

For strict positivity one must additionally prove the vanishing of the
primitive grades \(s\le0\) and strict positivity on the nonzero classes
of \(\operatorname{CH}^m(N)\).

\emph{Proof.} Let \(W\subset H(P)\) be the sum of all rational sub-Hodge
structures of the indicated types. Polarizable Hodge structures are
semisimple; this sum is a sub-Hodge structure and has an orthogonal
complement. Because \(A\) is CM, so is \(W\). Its twist \(W(m-1)\) is
effective of weight one. The degree-one case of 032 Lemma 8.1, followed
by Poincaré complete reducibility, gives a CM \(B\) and an isomorphism

\[
H^1(B)(1-m)\simeq W.
\]

A complex CM abelian variety admits a number-field model (Milne, Abelian
Varieties, version 1.10, Proposition 16.10, printed p.65). After
enlarging \(K\), both \(B\) and chosen polarizations are defined over
the same number field. Write \(b=\dim B\). The inclusion of this
weight-one twist into \(H^g(A)\) is a Hodge class of codimension
\(b+m-1\) on \(B\times A\). Theorem 1.1 applied to \(B\times A\) makes
it an algebraic correspondence. This is exactly the codimension
bookkeeping in the proof of 032 Corollary 8.2: \(n=1,r_0=m-1\).

The correspondence can be chosen over \(\overline{\mathbf Q}\). Indeed
spread each of its finitely many cycle components over a finitely
generated extension of \(\overline{\mathbf Q}\), or equivalently
represent it by a complex point of the Hilbert scheme of the fixed
variety \(B\times A\). Take the connected component of that finite-type
Hilbert scheme containing this point. Because the scheme is defined over
the algebraically closed field \(\overline{\mathbf Q}\), that component
contains a \(\overline{\mathbf Q}\)-point. The universal family is flat,
and its Betti cycle class on the fixed complex variety is locally
constant on the analytification; connectedness therefore preserves the
class. Specializing each component in this way preserves the original
rational Hodge correspondence. Finitely many algebraic coefficients and
cycle equations require only one finite extension \(K'\). Sandwich the
resulting correspondence between the actual degree-one and primitive
Chow projectors. We obtain a Chow morphism

\[
u:M_0=h^1(B)(1-m)\longrightarrow P
\]

whose realization is the chosen injection onto \(W\).

There is an algebraic, perfect alternating form on \(M_0\) given by a
polarization of \(B\) and motivic hard Lefschetz. On \(P\), use
\((-1)^{m-1}\) times the middle Poincaré form. Its restriction to the
types occurring in \(W\) is a Hodge polarization. Both forms are perfect
at the Chow-motive level: for \(B\) this is motivic hard Lefschetz; for
\(P\) it is AR-B3. They define a Chow adjoint \(u^\dagger\). The
realization of

\[
a=u^\dagger u\in\operatorname{End}_{\mathrm{Chow}}(M_0)
\]

is positive and invertible: for the positive Hodge Hermitian forms,
\((av,v)=(uv,uv)>0\) for nonzero \(v\). By AR-B4 this is an actual
rational isogeny endomorphism of \(B\), so \(a^{-1}\) exists in the Chow
endomorphism algebra. This step does \textbf{not} infer Chow
invertibility from a general cohomological isomorphism; it uses full
faithfulness specifically in degree one.

Set

\[
v=a^{-1}u^\dagger,\qquad e=uv=u a^{-1}u^\dagger.
\]

Then \(vu=1\), \(e^2=e\), and \(e^\dagger=e\), by direct composition.
Since the form on \(P\) is a scalar multiple of its Poincaré pairing,
\(e^\dagger=e\) says precisely \({}^t e=e\) within \(P\). Its
realization is the orthogonal projector onto \(W\). The complementary
motive \(N=(P,1-e)\) therefore has the stated realization. No assertion
that a realization-zero complementary motive is zero is made; such a
conclusion would need a separate nilpotence theorem. This avoids a
hidden lift from homological to Chow motives.

The retraction identifies \(\operatorname{CH}^m(M)\) with
\(\operatorname{CH}^1(h^1B)=\operatorname{Pic}^0(B)(K')_{\mathbf Q}\).
Such classes are algebraically trivial. This identification is also
injective modulo incidence equivalence: if \(u_*x\) is incidence
equivalent to zero, apply the correspondence \(v\); then \(x\) is
incidence equivalent to zero as a divisor on \(B\), hence \(x=0\) in
\(\operatorname{Pic}^0(B)_{\mathbf Q}\). The latter assertion follows
either from the identity correspondence in the definition of incidence
equivalence, or from the ordinary Picard universal property. Thus
AR-B2's primitive strict positivity applies to every nonzero class of
this Chow summand, not only to an unspecified quotient. Its Theorem 8.2
identifies the form with the pullback of the signed primitive Picard
polarization, giving the stated Néron--Tate description.

For \(x=e_*x\) and \(y=(1-e)_*y\), the height projection formula gives

\[
h_{K'}(x,y)=h_{K'}(x,{}^t e(1-e)_*y)=0.
\]

Combining this with AR-1 and (AR.3) proves the exact gate and the
strict-positivity assertion. Finite base change is harmless for classes
descending to \(K\): normalized heights are unchanged, and pullback on
rational Chow groups is injective since pushforward composed with
pullback multiplies by \([K':K]\). \(\square\)

\textbf{Scope caution.} The decomposition is asserted over \(K'\). For a
fixed initial \(K\), positivity on all of \(N(K')\) is a sufficient
condition for positivity on descended classes, not a claim of logical
equivalence with positivity over \(K\) alone. The exact equivalence in
the theorem concerns the group over \(K'\). The theorem is a derived
consequence of stated inputs; no claim of literature novelty is made.

\subsection{D. A small explicit complement which the reduction does not
erase}\label{d.-a-small-explicit-complement-which-the-reduction-does-not-erase}

Let \(E/K\) have an automorphism \(\rho\) of order six acting on
\(H^{1,0}(E)\) by a primitive sixth root \(\zeta\), and assume this
automorphism is defined over \(K\). For example take a \(j=0\) elliptic
curve after extending its number field. On \(A=E^3\), use the product
polarization and let \(\rho_\Delta=(\rho,\rho,\rho)\). Let
\(\Pi=\pi_1\otimes\pi_1\otimes\pi_1\) be the degree-one tensor Chow
projector. Put

\[
e_3=\Pi\,\frac16\sum_{r=0}^5(-1)^r\rho_\Delta^{*r}.
\tag{AR.5}
\]

\textbf{Proposition AR-3.} Equation (AR.5) is an actual rational Chow
idempotent, self-transpose on the middle motive, and its motive \(T\) is
primitive. Its Hodge realization has dimension two and types
\((3,0),(0,3)\). Its entire \(\operatorname{CH}^2(T)\) lies in Beauville
grade one and is homologically trivial. Consequently its height is a
concrete surviving case of (AR-GATE), with positive-sign candidate
\(h_K(z,z)\).

\emph{Proof.} The cyclic group character projector is idempotent because
the coefficient character \(r\mapsto(-1)^r\) is a rational character of
the order-six group. It commutes with \(\Pi\). Transposition replaces
group elements by their inverses and fixes the character projector;
\(\pi_1\) of an elliptic curve is self-transpose. In the tensor
\(H^1(E)^{\otimes3}\), the eigenvalues are \(\zeta^3=-1\), \(\zeta\),
\(\zeta^{-1}\), and \(\zeta^{-3}=-1\), according to the number of
holomorphic factors. The two all-holomorphic/all-antiholomorphic tensors
are exactly the \(-1\)-eigenspace. Multiplication by each factor's
divisor class kills the factor \(h^1(E)\), since it would land in
\(h^3(E)=0\); this is exact in the Chow-motive grading. Hence
\(L\Pi=0\), so \(T\) is primitive. Finally \(T\subset h^3(A)\), and
\(\operatorname{CH}^2(h^3(A))=\operatorname{CH}^2_{(1)}(A)\). It has
zero ordinary degree-four cycle realization. \(\square\)

This is a specified motive and a specified height problem. It is not an
assertion that \(\operatorname{CH}^2(T)\) is known to contain a nonzero
class, nor an assertion that all its heights have been determined. No
general height-to-\(L'\) identity for this group follows from the
supplied papers.

\subsubsection{D.2 The complementary gate contains a nonzero
Fermat-quartic
cycle}\label{d.2-the-complementary-gate-contains-a-nonzero-fermat-quartic-cycle}

\textbf{Theorem AR-7 (a nonzero CM complementary class).} Let
\(C:x^4+y^4=z^4\), let \(J=\operatorname{Jac}(C)\), and choose a
degree-one base class after a finite number-field extension. Write \[
\kappa=[C]-[-1]_*[C]\in\operatorname{CH}^2(J)_{\rm hom,\mathbf Q}.
\] Let \(\pi_P\) be the primitive middle Chow projector for a symmetric
theta divisor, and let \(e\) be the Picard-summand projector of AR-2.
Then \[
z_N=(\pi_P-e)_*\kappa\in\operatorname{CH}^2(N)
\tag{AR.9}
\] is nonzero modulo algebraic equivalence, even over
\(\overline{\mathbf Q}\). Thus (AR-GATE) concerns a group containing a
specified nonzero class; AR-13 subsequently proves that this class has
strictly positive height, while the sign on the entire group remains
unresolved. In this example \(H(N)\) has precisely types
\((3,0),(0,3)\), each of dimension one.

\emph{Proof.} Bloch, \emph{Algebraic cycles and values of L-functions},
J. reine angew. Math. 350 (1984), Lemma (3.1), p.99, gives an isogeny
\(J\to E^2\times E'\), of degree a power of two, where \(E:y^2=x^3-x\)
and \(E':y^2=x^3+x\). Thus \(J\) is CM; over \(\overline{\mathbf Q}\) it
is isogenous to the cube of a \(j=1728\) elliptic curve. Bloch's Theorem
(4.1), p.104, says that this Ceresa cycle has infinite order modulo
algebraic equivalence. Its class modulo algebraic equivalence is
independent of the chosen embedding (Theorem (3.4) and following text,
p.102).

We first check that \(\kappa\) has grade one, so the Chow projector
\(\pi_3\) does not discard its nonzero Griffiths class. It is odd under
\([-1]^*\), hence only odd Beauville grades can occur. The Fourier
isomorphism \[
\mathcal F:\operatorname{CH}^2_{(s)}(J)
 \simeq\operatorname{CH}^{1+s}_{(s)}(\widehat J)
\] kills every odd \(s\ge3\), by dimension, and every odd \(s\le-3\), by
negative codimension. For \(s=-1\) the target is
\(\operatorname{CH}^0_{(-1)}(\widehat J)=0\), since
\(\operatorname{CH}^0\) is spanned by the fundamental class and has
grade zero. Therefore
\(\kappa\in\operatorname{CH}^2_{(1)}(J)=\operatorname{CH}^2(h^3J)\).

Motivic Lefschetz gives \(h^3(J)=P\oplus Lh^1(J)\). The codimension-two
Chow group of the second summand is an image under \(L\) of
\(\operatorname{Pic}^0(J)_{\mathbf Q}\); all these classes are
algebraically trivial. The same is true of \(\operatorname{CH}^2(M)\),
by AR-2. Consequently \[
\kappa-z_N=(1-\pi_P)_*\kappa+e_*\kappa
\] is algebraically trivial. Bloch's theorem now forces \(z_N\) to be
nonzero modulo algebraic equivalence. This argument is unchanged after
finite extension.

For the final Hodge-type assertion, the rational CM decomposition of
\(H^3(E^3)\) splits off its two extreme types \((3,0),(0,3)\). One can
see rationality explicitly by applying the scalar CM endomorphism
\(1+i\): its eigenvalues on the extreme types satisfy \(T^2+4T+8\),
while those on \((2,1),(1,2)\) satisfy \(T^2-4T+8\). These coprime
rational polynomials give the two rational Hodge substructures. The full
level-one part has dimension \(18\). The Lefschetz summand \(LH^1(J)\)
is six-dimensional and lies entirely in that part. The extreme
two-dimensional substructure is already primitive, by Hodge type. Thus
the primitive middle piece has a twelve-dimensional maximal level-one
part and the indicated two-dimensional complement. \(\square\)

This proof refutes a CM-wide assertion that homological triviality is
algebraic triviality up to torsion: the rational Griffiths quotient is
nonzero already for this CM Jacobian. It does not refute the arithmetic
Hodge-index sign conjecture. The new statement here is the precise
placement of the known nonzero cycle in the orthogonal complementary
Chow summand, using AR-2; Bloch's nonvanishing theorem is an accepted
established input.

The Beauville Fourier grading used above is the standard Fourier
decomposition (Beauville, \emph{Sur l'anneau de Chow d'une variété
abélienne}, Math. Ann. 273 (1986), 647--651). An explicitly read
statement with the grading convention is Shen--Vial, \emph{The Fourier
transform for certain hyperkähler fourfolds}, Introduction A, p.2,
displayed (1) and property (a). Its Fourier inversion is also recalled
in Künnemann 1993, §2, printed p.88.

\subsubsection{D.3 A canonical nonzero complementary cycle with an exact
local height
formula}\label{d.3-a-canonical-nonzero-complementary-cycle-with-an-exact-local-height-formula}

\textbf{Theorem AR-8 (Fermat symmetry and Gross--Schoen reduction).} For
the Fermat quartic choose the hyperflex \(P_0=(1:0:1)\) and its Abel
embedding \(f(x)=[x-P_0]\). Let \(X_1\) be the grade-one part of
\(f_*[C]\), and set \[
\eta=\mathcal F(X_1)\in\operatorname{CH}^2(J).
\] One can choose the splitting in AR-2 so that
\(\eta\in\operatorname{CH}^2(N)\). This class is nonzero modulo
algebraic equivalence. Let \(K'\) contain the relevant definitions and
admit a semistable model of \(C\), and put \[
\Omega_{K'}=\frac{\omega_a^2}{[K':\mathbf Q]},
\qquad
\Phi_{K'}=\frac{1}{[K':\mathbf Q]}
 \sum_v\varphi(C_v)\deg(v).
\] Then the raw normalized middle height has the exact value \[
h_{K'}(\eta,\eta)
 =\frac16 h_{K'}(\Delta_{P_0},\Delta_{P_0})
 =\frac7{24}\Omega_{K'}-\frac16\Phi_{K'}.
\tag{AR.10}
\] Thus this specified nonzero complementary class has nonnegative
height exactly when \(7\Omega_{K'}\ge4\Phi_{K'}\), and strictly positive
height exactly when the inequality is strict. There is no positive
Picard contribution to subtract for this choice of splitting.

\emph{Proof.} The tangent \(x=z\) meets \(C\) at \(P_0\) with
multiplicity four, so \(K_C\sim4P_0\). Let \(G=\mu_4^2\) act by the two
coordinate scalings on the affine equation \(x^4+y^4=1\). Each
\(\gamma P_0\) satisfies \(4\gamma P_0\sim K_C\), so \(\gamma P_0-P_0\)
is four-torsion in \(J\). Translation by a torsion point acts
identically on rational Chow groups: if \(na=0\), then
\([n]\circ t_a=[n]\); since \([n]_*\) is invertible on rational Chow
groups by the Chow--Künneth eigenvalue decomposition,
\([n]_*t_{a*}=[n]_*\) implies \(t_{a*}=1\). Therefore the class
\(f_*[C]\), its grade-one part, and its Fourier transform \(\eta\) are
\(G\)-invariant. The last assertion uses the \(G\)-invariance of the
symmetric principal theta class; the Fourier kernel is consequently
\(G\)-equivariant.

The holomorphic differentials have \(G\)-characters \[
(1,1),\quad(1,2),\quad(2,1)\qquad\text{in }(\mathbf Z/4)^2.
\] Their exterior triple has trivial character. In a mixed exterior
triple of type \((2,1)\), canceling a holomorphic and antiholomorphic
character leaves one of the three nonzero characters above; using three
distinct indices gives \((0,2),(2,0),(2,2)\), again nonzero. Conjugation
handles type \((1,2)\). Hence \(H^3(J)^G\) consists exactly of the two
extreme Hodge types. In particular it has no overlap with \(W=H(M)\).

Let \(p_G=|G|^{-1}\sum_{\gamma\in G}\gamma_*\). This is an exact
self-transpose Chow idempotent. It commutes with the primitive projector
because \(G\) preserves theta. In AR-2 replace the constructed inclusion
by \[
u'=(1-p_G)u.
\] Its realization is unchanged on \(W\); therefore the same adjoint
argument gives a split inclusion and a self-transpose projector \(e'\).
Now \(p_Gu'=0\) holds exactly, and transposition gives
\((u')^\dagger p_G=0\), hence \(e'p_G=0\). Since \(p_G\eta=\eta\), we
get \(e'\eta=0\). Zhang's Theorem 1.5.6 says \(\eta\) is primitive;
therefore \(\eta\in\operatorname{CH}^2(N)\) for this splitting.

By AR-7, the Ceresa class is \(2X_1\) and is nonzero modulo algebraic
equivalence. Fourier is an invertible algebraic correspondence and
preserves algebraic equivalence, so \(\eta\) is likewise nonzero. Zhang,
\emph{Gross--Schoen cycles and dualising sheaves} (2010), Theorem 1.5.6,
printed p.14, gives the first equality in (AR.10) in genus three. His
Theorem 1.3.1, printed p.5, gives the second because \(P_0-K_C/4=0\) in
the rational Picard group. Dividing every global intersection by
\([K':\mathbf Q]\) gives the displayed normalization. \(\square\)

This identity supplies the local terms used in the subsequently
completed proof of (AR.10) in AR-13. Good finite places contribute zero
to \(\Phi_{K'}\). For the displayed integral plane quartic, every odd
prime is a smooth-reduction prime by the partial derivatives. Thus only
places over \(2\) and the archimedean places can contribute to
\(\Phi_{K'}\). The admissible self-intersection \(\Omega_{K'}\) also
includes arithmetic metric and model information. AR-13 combines a local
Hodge-lattice estimate at \(2\) with certified archimedean bounds to
settle the height of this cycle. Positivity of this specified class does
not establish positivity of the whole group \(\operatorname{CH}^2(N)\).

\subsubsection{D.4 The positive finite local term and the archimedean
reduction}\label{d.4-the-positive-finite-local-term-and-the-archimedean-reduction}

\textbf{Proposition AR-9 (local arithmetic progress for the Fermat
cycle).} Work over a number field \(K'\) on which \(C\) is semistable.
For each finite place set \[
b_v=\frac1{18}\operatorname{ord}_v(\chi'_{18})-\lambda(\Gamma_v),
\] and put \[
D_2=\frac1{[K':\mathbf Q]}\sum_{v\mid2}\delta(\Gamma_v)\log N(v),\qquad
A_\infty=\frac1{[K':\mathbf Q]}\sum_{v\mid\infty}
\left(-\frac1{18}\log\|\chi'_{18}\|_{\mathrm{Hdg},v}-\lambda(C_v)\right),
\] using de Jong's archimedean multiplicity convention. Then \[
h_{K'}(\eta,\eta)
 =\frac72\left(
  \frac1{[K':\mathbf Q]}\sum_{v\mid2}b_v\log N(v)+A_\infty
 \right)
 \ge\frac{D_2}{6}+\frac72 A_\infty,\qquad D_2>0.
\tag{AR.12}
\] In particular \(A_\infty>-D_2/21\) is a sufficient condition for
strict positivity of this specified nonzero complementary class.
Moreover the graph contribution in (AR.10) is explicitly \[
\varphi(\Gamma_v)=\frac43\,\delta(\Gamma_v)\qquad(v\mid2).
\tag{AR.13}
\]

\emph{Proof.} Van Bommel--Docking--Dokchitser--Lercier--Lorenzo García,
\emph{Reduction of Plane Quartics and Cayley Octads}, Table 7.1 in the
read preprint (published version Table 9), identifies the Fermat curve's
stable reduction at \(2\) as type \((0eee)\): one genus-zero central
component with three elliptic tails. The source cites Liu,
\emph{Algebraic Geometry and Arithmetic Curves}, Example 10.4.39, for
this characteristic-two calculation. We use the stated stable type.
Programme JAC-P10 and the independent h9 review subsequently read Liu's
original Example 10.4.39, pp.548--549, and use its local coordinates to
prove the Hodge-lattice estimate entering AR-13. The dual graph is a
tree, all three edges separate genus one from genus two, and all have
strictly positive thickness. The curve has smooth nonhyperelliptic
reduction at every odd prime, so \(b_v=0\) there.

De Jong, \emph{On the height of Gross--Schoen cycles in genus three}
(2018), Theorem 8.2, expresses the canonical Gross--Schoen height as
\(21\) times the sum of these finite and archimedean terms. Divide by
\([K':\mathbf Q]\) and by \(6\), using AR-8, to obtain the equality in
(AR.12). De Jong's Lemma 11.2 applies to every genus-three metrized tree
and proves \(b_v\ge\delta(\Gamma_v)/21\). This proves the inequality and
\(D_2>0\).

For (AR.13), Zhang 2010, Proposition 4.4.1, printed pp.58--59, gives the
separating-edge contribution \(2i(g-i)\ell_e/g\). Here \(g=3,i=1\) for
all edges, and there is no cyclic contribution, so summing gives
\(4\delta(\Gamma_v)/3\). One may also check the accompanying invariant
from de Jong, Example 5.4: \(\lambda(\Gamma_v)=2\delta(\Gamma_v)/7\).
The inequality \(b_v\ge\delta_v/21\) is then
\(\operatorname{ord}_v(\chi'_{18})\ge6\delta_v\), exactly the tree bound
used in his proof. \(\square\)

This determines the sign of the finite local contribution in the
\(\chi'_{18}\)-formula; AR-12 and AR-13 subsequently bound the
archimedean term and prove the global sign for this class. Neither
(AR.12) nor (AR.13) identifies the bad-prime logarithmic coefficient of
the different automorphic comparison (AR.11); those two local
decompositions must not be conflated.

\textbf{Subsequent resolution for this class, AR-12 and AR-13.} The
companion FERMAT\_THETA\_GATE.md reduces the archimedean term to an
explicit logarithmic theta integral at the exact Fermat period matrix.
FERMAT\_THETA\_CERTIFICATE.md combines a new finite Hodge-lattice
estimate JAC-P10 with a deterministic, enclosed theta computation to
prove \(h_K(\eta,\eta)>1/100\). The computation uses a torsion-grid
lower bound for the Frobenius-theta mean and a cell-Jensen upper bound
for the ordinary-theta mean; it does not use an \(L'\)-height
comparison. The final certificate and its independent-review status are
in fermat\_theta\_certify.py/.json and CHECK\_EVENTS.tsv. This settles
the sign on \(\mathbf Q\eta\), while positivity on the full
complementary Chow group remains open.

\textbf{Corollary AR-14 (a positive arithmetic space containing a
Griffiths class).} For the Fermat threefold and the splitting chosen in
AR-8, the height is strictly positive on \[
\operatorname{CH}^2(M)(K')_{\mathbf Q}\oplus\mathbf Q\eta.
\] Indeed AR-2 gives strict positivity on the first summand and
orthogonality to \(\eta\in\operatorname{CH}^2(N)\); AR-13 gives
positivity on the second. Thus
\(h(x+a\eta,x+a\eta)=h(x,x)+a^2h(\eta,\eta)\), which is positive unless
both terms vanish. No assertion that \(\eta\) spans
\(\operatorname{CH}^2(N)\) is needed or made.

\subsection{E. Mixed-characteristic and constant-family
checks}\label{e.-mixed-characteristic-and-constant-family-checks}

The correct arithmetic analogue of an intersection formula has the form

\[
h_K(z,z')=\frac1{[K:\mathbf Q]}\widehat{\deg}
 \bigl(\widehat z^{\mathrm{adm}}\cdot\widehat z'^{\mathrm{adm}}\bigr),
\tag{AR.6}
\]

with vertical-orthogonal extensions and the compatible Green-current
data. Künnemann 2001, §1.2--1.4, is the exact construction used here.
Thus a raw closure intersection does not become a height merely by
inserting \((-1)^m\). The parent's quadric counterexample to the
uncorrected pencil formula is consistent with this distinction: integral
fibers do not remove all higher-codimension local correction terms.

\textbf{Constant-family check.} For the constant family
\(B\times C\to C\) over a field \(k\), and a constant section \(s_P\)
with \(P\in B(k)\), the geometric canonical height over \(k(C)\) is
zero. Indeed the pullback of a line bundle from \(B\) along the constant
map is a constant one-dimensional vector space tensored with
\(\mathcal O_C\), hence has degree zero; the same holds after every
multiplication by \(n\). If \(k=K\) is a number field and \(P\) is
nontorsion, its arithmetic Néron--Tate height on \(B/K\) is strictly
positive for an ample polarization. Therefore the arithmetic value
cannot be recovered from that geometric degree alone. The missing term
is genuine arithmetic metric/finite-place information. This is a
counterexample to any claimed direct replacement of a geometric height
by the number-field height with all these terms omitted, not a
counterexample to a correctly normalized biextension formula.

For ordinary divisors, Faltings--Hriljac gives \(h_K(D,D)<0\) when the
Jacobian class of \(D\) is nontorsion. For a primitive codimension-two
class on a threefold the proposed positive sign is \(+h_K\). These two
low-dimensional checks prevent an unsigned ``BB is NT-positive''
statement.

\subsection{F. What GRH actually
gives}\label{f.-what-grh-actually-gives}

\textbf{Lemma AR-4 (central sign).} Let \(F\) be holomorphic in a
neighborhood of the real interval \([c,\infty)\), real on that interval,
positive for sufficiently large real arguments, and nonzero for every
real \(s>c\). If \(F\) vanishes to order \(r\ge1\) at \(c\), then
\(F^{(r)}(c)>0\). In particular \(F'(c)\ge0\), with strict inequality
exactly when \(r=1\).

\emph{Proof.} Continuity and the absence of a real zero on
\((c,\infty)\) give \(F(s)>0\) there. In the expansion
\(F(c+t)=a_rt^r+O(t^{r+1})\), \(a_r\) is real and nonzero. Positivity
for small \(t>0\) forces \(a_r>0\), and \(F^{(r)}(c)=r!a_r\).
\(\square\)

If \(F=G L\), with \(G(c)>0\) real and \(L(c)=0\), then
\(F'(c)=G(c)L'(c)\). Thus GRH supplies the lemma's zero-free interval
for an appropriately normalized real self-dual completed \(L\)-function,
provided poles and the positivity normalization have also been handled.
A complex non-self-dual Hecke \(L\)-function has no ordered derivative
sign until an appropriate real combination has been specified.

\textbf{Corollary AR-5 (conditional height sign).} If a proven formula
gives \(q_A(z,z)=C L'(c)\), \(C>0\), with the hypotheses of AR-4, then
\(q_A(z,z)\ge0\). It is strictly positive if the central zero is simple.
If the central zero has higher order, the formula gives zero; it does
not prove \(z=0\). To deduce positivity for an entire cycle space one
needs a spanning family and control of cross terms, or a formula for
every vector in that space. A regulator determinant alone is
insufficient: the real symmetric matrix
\(\operatorname{diag}(-1,-1,1,\ldots,1)\) has positive determinant and
is not positive definite.

The attached Alpöge \texttt{Statements.lean} quantifies over Dirichlet
characters \(\chi:\operatorname{DirichletCharacter}(\mathbf C,N)\). Its
\texttt{Main.lean} proves exactly those advertised Dirichlet statements
from imported results; it contains no theorem for arbitrary number-field
Hecke characters. We accept the cited theorem; this is a scope check,
not a re-verification of its proof. The Masullo index labels 14.1.2.1
(printed p.529) a \textbf{Conjecture}, and 14.1.2.3 (p.531) a
proposition conditional on that conjecture and 14.1.2.2. The row on
p.522 whose opening reads ``13.3.1.2 (Divisor of L(s, chi))'' includes
unfinished statement text (the index parses its number as 13.3.1).
Accordingly any use of broad Hecke GRH is explicitly \textbf{conditional
(Hecke-GRH/Masullo)}. No such input supplies the height identity in
AR-5.

\subsubsection{F.2 The concrete CM derivative and the newer
triple-product
comparison}\label{f.2-the-concrete-cm-derivative-and-the-newer-triple-product-comparison}

Bloch 1984, Proposition (5.2), printed p.107, proves that the Hecke
\(L\)-function of the Fermat-quartic rank-two motive of types
\((3,0),(0,3)\) has a simple zero at \(s=2\). This is a theorem about
order of vanishing. The proposition does not state an arithmetic height
identity or an ordered sign for the derivative. The introduction,
printed p.95, explicitly presents the height formula as a possible
further assertion. Thus the supplied Dirichlet-GRH theorem, Bloch's
nonvanishing, and (AR.10) are three distinct inputs.

\textbf{New analytic deduction AR-11.} The companion proof
\path{FERMAT_LDERIVATIVE_SIGN.md} proves unconditionally \[
L'(M,2)>\frac7{20}>0
\] for this particular Fermat rank-two motive, using the known character
\(\psi^3\mu\), its conductor \(64\), its root number \(-1\), and an
elementary Mellin-transform tail estimate. A supplementary certified
interval is \(1.5071685<L'(M,2)<1.5071697\). Thus the analytic sign is
no longer a gate for this example. AR-13 independently proves the
local-height inequality and the positive height of the specified class.
A geometric identity comparing that height with this central derivative
remains a separate unresolved requirement.

A stronger next comparison is available in Yuan--Zhang--Zhang,
\emph{Triple product L-series and Gross--Kudla--Schoen cycles},
manuscript dated 14 November 2023, with appendix by Yifeng Liu.
Proposition 1.3.2, printed p.7, expresses the trilinear cycle height as
a coefficient \(\mathcal L(\pi)\) times a product of local invariant
pairings. Theorem 1.3.3, printed p.8, states \[
\mathcal L(\pi)=
\frac{16\zeta_F(2)^2}{L(1,\pi_L,\mathrm{ad})}
 L'(1/2,\pi_L)+\sum_{p\in S}c_p\log p,
\qquad c_p\in\overline{\mathbf Q},
\tag{AR.11}
\] where \(S\) consists of rational primes below ramification places of
\(\pi\). The exact generalized Gross--Kudla formula holds in their
everywhere-unramified case. This is stronger than the older September
2012 manuscript, Theorem 1.2.4, whose hypotheses included at least two
finite nonsplit quaternion places.

The Fermat quartic is \(X_0(64)\); for example the explicit modular-form
embedding and the relation \(x^4+y^4=z^4\) are given in Fang-Ting Tu,
\emph{Arithmetic of Quaternion Orders and Its Applications} (Waseda
dissertation, February 2012), §5.3.2, p.59. Hence the relevant modular
elliptic constituents are unramified away from \(2\). If one carries out
the precise trilinear-cycle identification and the local test-vector
normalizations for this example, the possible correction in (AR.11) is
supported at \(2\). Its shape is \(c_2\log2\); its vanishing or sign is
not supplied by an \(L\)-function zero-free theorem. The whole
identification is retained as a task, rather than treating a general
modular parametrization as an automatic equality with \(\eta\).

The newer Gao--Zhang manuscript, \emph{Heights of Ceresa and
Gross--Schoen cycles}, Theorems 1.1 and 1.4, constructs canonical adelic
height objects and proves a lower bound on a suitable Zariski-open part
of moduli. This offers genuine geometry for the Bost route below. A
bound with an additive constant on an unspecified open subset does not
determine the sign at this fixed Fermat point. The specific formula
(AR.10) is subsequently evaluated from below in AR-13, which settles
this fixed class. A completed automorphic comparison including the term
in (AR.11) remains separate and unresolved.

AR-13 subsequently proves the required local-height inequality for the
specified Fermat class directly. The full complementary-group inequality
and the automorphic cycle/test-vector comparison, including \(c_2\),
remain distinct unresolved tasks.

\subsection{G. Testing the two suggested Alpöge
mechanisms}\label{g.-testing-the-two-suggested-alpuxf6ge-mechanisms}

\subsubsection{G.1 Sign-only Weil/explicit-formula
route}\label{g.1-sign-only-weilexplicit-formula-route}

The note's footnote 3, printed p.3, explains the sign in the
finite-field trace calculation using concentration in middle cohomology
and the opposing zero-mode parity, with a bound from repeated Parseval.
Step 5 (pp.3--4) changes to a von Mangoldt/Guinand--Weil variance setup.
Neither statement identifies the cycle intersection or arithmetic height
form with the analytic form.

A minimal algebra check isolates the missing information. Let \(F=pI\)
on \(\mathbf Q^2\), \(p>0\). Both \(b_+=I\) and
\(b_\pm=\operatorname{diag}(1,-1)\) satisfy \(F^tbF=p^2b\); they have
the same Frobenius eigenvalues and traces and different signatures. Thus
eigenvalue purity and trace parity do not determine a compatible form's
signature. This is an information-loss test, not a proposed geometric
counterexample. A valid geometric transfer still needs an identification
of the form and its polarization.

\subsubsection{G.2 Bost slope route}\label{g.2-bost-slope-route}

Bost, \emph{Algebraic leaves of algebraic foliations over number
fields}, Publ. Math. IHES 93 (2001), §4.1.4, Proposition 4.5, equation
(4.11), gives for an injective \(K\)-linear map between hermitian
bundles

\[
\widehat{\deg}_{n}(E)
\le\operatorname{rk}(E)\bigl(\widehat\mu_{\max}(F)+h(\varphi)\bigr).
\tag{AR.7}
\]

Proposition 4.6, equation (4.14), is its filtered version. The note's
step 13, p.4, names this inequality. The attached Lean map points to
\path{Formalism/LatticeSlopes/LatticeSlopesProofs/SlopeInequality.lean}
and \path{Formalism/LatticeSlopes/Statement.lean}; these
implementation files were \textbf{not attached or read}, and no claim is
made to have run them.

\textbf{Proposition AR-6 (precise slope gate).} Suppose that a geometric
construction associates to a cycle \(z\) a hermitian line
\(\overline{\mathcal B}_z\) with
\(\widehat{\deg}_n\overline{\mathcal B}_z=q_A(z,z)\), and an injective
map \(\varphi:\mathcal B_z^\vee\to F_z\) satisfying

\[
\widehat\mu_{\max}(F_z)+h(\varphi)\le0.
\tag{AR.8}
\]

Then \(q_A(z,z)\ge0\); a strict upper bound in (AR.8) gives strict
positivity.

\emph{Proof.} Apply (AR.7) to the rank-one source
\(\overline{\mathcal B}_z^\vee\), whose normalized degree is
\(-q_A(z,z)\). \(\square\)

The hypotheses in AR-6 are substantive: the line must be the cycle's
canonical arithmetic biextension/height object, and the map and norm
bounds must be proved. Merely declaring a trivial line to have metric
\(e^{-q_A(z,z)}\) is circular. To see why the inequality by itself
yields no sign, let \(E_a=(\mathcal O_K,\|1\|_\sigma=e^{-a})\). Then
\(\widehat{\deg}_n(E_a)=\widehat\mu_{\max}(E_a)=a\), and the identity
map has height zero. Equation (AR.7) becomes \(a\le a\) for every
positive or negative real \(a\).

This explicitly tests the next tool: a Bost proof of (AR-GATE) would
require (AR.8) for the height line of every class in the complementary
motive \(N\), including local and archimedean terms. Neither the
four-page note nor the supplied Lean map provides this bridge. On the
Picard summand the Poincaré biextension already gives the canonical
object and the result is known independently.

\textbf{Construction supplied by AR-10.} Root's
work/root/AR10\_adelic.tex, independently whole-read here, specializes
Gu 2026, Theorem 1.3 and Proposition 5.3, to the exact middle signed
height and constructs \(\overline{\mathcal B}_z\) as a metrized Deligne
pairing. This resolves the construction of the height line for a chosen
rational representative, with representative-independent arithmetic
degree. The resulting line is integrable adelic. One applies the adelic
rank-one slope inequality, or a justified model passage to Bost's
hermitian version; clearing rational tensor powers by itself is not that
passage. The injection and its norm/slope bound remain open.

\subsection{H. Sources and reading
record}\label{h.-sources-and-reading-record}

\begin{itemize}
\tightlist
\item
  Attached 032: Theorem 1.1; Lemma 8.1; Corollary 8.2 and its proof.
  Accepted inputs; targeted statement/codimension reading.
\item
  Attached cote 44 scan004: Theorem 5 and its margin. It assumes
  algebraically trivial generic restrictions; its printed intersection
  formula is handled separately by the parent's H9 counterexample and
  correction.
\item
  Künnemann 1993:
  \url{https://gdz.sub.uni-goettingen.de/download/pdf/PPN356556735_0113/LOG_0011.pdf}
  . Direct PDF downloaded, rendered printed pp.96--98; Theorems 5.1,5.2,
  Corollary 5.3 directly read.
\item
  Künnemann 1996: DOI \url{https://doi.org/10.1353/ajm.1996.0033} . Primary
  paper's indexed PDF at
  \url{https://citeseerx.ist.psu.edu/document?doi=24752c170be79846ab861e9c80a8977cac60035f&repid=rep1&type=pdf}
  . Web retrieval exposed full extracted text after direct-download
  failures. Lemma 4.1, equation (18), Theorems 6.1,8.2,9.1,12.1, and
  Remark 12.2 were targeted-read; not a full proof audit.
\item
  Künnemann 2001:
  \url{https://www.numdam.org/item/10.1016/s0012-9593(01)01068-0.pdf} .
  §1.2--1.4, Theorem 1.6, Corollary 1.7 directly read via extracted
  primary PDF.
\item
  Milne, Abelian Varieties v1.10:
  \url{https://www.jmilne.org/math/CourseNotes/AV110.pdf} . Proposition 16.10,
  printed p.65, model over a number field; statement read, proof omitted
  in this source.
\item
  Zhang 2010:
  \url{https://web.math.princeton.edu/~shouwu/publications/gross-schoen.pdf}
  . §5.3, p.68, height selection; directly read.
\item
  Laga--Shnidman: \url{https://arxiv.org/pdf/2312.12965} . Proposition 4.5,
  Proposition 5.9 and Proposition 5.10 directly read; Kings Prop.2.2.1
  is cited through this source, not directly read.
\item
  Bost 2001:
  \url{https://www.numdam.org/item/PMIHES_2001__93__161_0.pdf} . §4.1.4,
  Propositions 4.5--4.6, equations (4.11),(4.14), directly read.
\item
  Gu 2026: \url{https://arxiv.org/abs/2609.33223} . Local primary v2 read at
  \path{work/references/Gu2026.pdf} and its extracted text: Theorem 1.3, Proposition 5.3, §2.2
  and Corollary 1.6; root's complete AR-10 proof independently checked
  against these statements.
\item
  Additional Fermat derivative sources, calculation and
  application-scope checks: see \path{FERMAT_LDERIVATIVE_SIGN.md} and
  KNOWN\_STATEMENT\_COMPARISONS.md.
\item
  Attached Alpöge note: pp.3--4, footnote3, steps5 and13; attached
  Main/Statements; Lean-map §4 and §5. No formal build and no claim to
  inspect omitted implementation files.
\item
  Attached Masullo index: rows 13.3.1.2,14.1.2.1,14.1.2.3. The full
  paper is not attached; broad Hecke GRH is retained as an explicit
  hypothesis.
\end{itemize}

\subsection{I. Self-check summary}\label{i.-self-check-summary}

AR-1: exponent \(g-i-e=s-2\) checked; grade zero of symmetric \(\ell\)
checked; no unproved Beauville vanishing inserted. AR-2: codimension
\(b+m-1\), twist \(1-m\), Chow degree \(m+(1-m)=1\), exact inverse
specifically in \(h^1\), self-transpose projector, incidence injection,
field extension scope, and height orthogonality each checked. AR-3:
idempotent group algebra character, transpose/inverse, eigenvalue
multiplicities \(1,3,3,1\), primitive tensor action, and grade
\(2p-i=1\) checked. AR-4: poles, real-valuedness, positivity
normalization, order of zero, and derivative/leading-coefficient
distinction checked. AR-6: dual-source sign and rank-one rescaling test
checked. AR-7: Fourier odd-grade elimination, primitive versus Lefschetz
summand, algebraic triviality of the discarded terms, and the two CM
characteristic polynomials checked; root independently confirmed this
argument. AR-1a: all negative grades and the homologically trivial
grade-zero divisor image checked. AR-8: hyperflex canonical class,
torsion translation using invertible multiplication, all twenty exterior
character triples, exact averaged-projector annihilation, and the
constants 7/24 and 1/6 checked. The character enumeration independently
returned only the all-holomorphic and all-antiholomorphic triples as
invariant.

Remaining parent-facing review request: independently check AR-2's
motivic-adjoint sign and field-descent argument before promoting it in
the final paper. It is a complete written argument from stated
background, not a formal certificate and not a claimed solution of
(AR-GATE).
\section{The canonical adelic line in the remaining slope problem}
\label{sec:AR10}

\begin{theorem}[AR-10: an explicit adelic representative of the complementary height]
Let $A/K$ have dimension $g=2m-1$, let $L$ be a symmetric ample line
bundle rigidified at the identity, and equip $L$ with its canonical
adelic metric. Put $d=m-1$ and
\[
 \overline Q=\mu^*\overline L-p_1^*\overline L-p_2^*\overline L
 \quad\text{on }A\times A,
\]
where $\mu$ is addition. If $z,w\in\operatorname{CH}^m_{(1)}(A)$
are primitive in the Chow sense $Lz=Lw=0$, then
\begin{equation}\label{eq:AR10}
 (-1)^m h_K(z,w)=
 \frac{(-1)^{m-1}}{(2m-1)!\,h^0(A,L)\,[K:\mathbf Q]}
 \widehat{\deg}\bigl(\overline Q^{\,2m-1}\cdot(z\times w)\bigr).
\end{equation}
Here $h_K$ is the normalized raw Beilinson--Bloch pairing of the
arithmetic section. The arithmetic intersection on the right uses
rational cycle representatives and canonical metrics.

For a chosen rational representative $z=\sum_i a_i[V_i]$, define the
integrable adelic rational line on $\operatorname{Spec}K$
\[
 \overline{\mathcal B}_z=
 \frac{(-1)^{m-1}}{(2m-1)!h^0(A,L)}
 \sum_{i,j}a_i a_j
 \bigl\langle\underbrace{\overline Q,\ldots,\overline Q}_{2m-1}
 \bigr\rangle_{(V_i\times V_j)/K}.
\]
Its normalized arithmetic degree is $(-1)^m h_K(z,z)$. Thus for
$z\in\operatorname{CH}^m(N)$ in AR-2, the line required as the source
of the slope criterion AR-6 has an explicit construction from the
cycle and the canonical polarization metric. The residual requirement
of that criterion is an injective map from its dual to a hermitian
bundle satisfying the stated maximal-slope and operator-norm bound.
\end{theorem}

\begin{proof}
We use Gu, \emph{The Beilinson--Bloch heights of generalized Ceresa
cycles}, arXiv:2609.33223v2 (2 October 2026), Theorem 1.3, printed p.~2,
and Proposition 5.3, printed pp.~12--13. These are cited formulas;
the specialization and the slope-line construction are the deduction
here. Gu uses the unnormalized height before division by $[K:\mathbf Q]$,
so we keep this division explicitly on arithmetic intersections.

A grade-one codimension-$m$ class has eigenvalue $(-1)^{2m-1}=-1$
under $[-1]^*$. Since $[-1]$ is an involution, pushforward has the same
action. Hence the normalized generalized Ceresa cycle
$\operatorname{Ce}(z)=(z-[-1]_*z)/2$ is exactly $z$.
Its pure dimension is $d=g-m=m-1$. Gu's Theorem 1.3 therefore gives
\[
 h_{\mathrm{BB},F}(z)=
 \frac{(-1)^d}{(2d+1)!\,[K:\mathbf Q]}
 \widehat{\deg}\bigl(\overline Q^{2d+1}\cdot(z\times z)\bigr).
\]
For the primitive grade-one class $z$, Proposition 5.3 gives
\[
 h_{\mathrm{BB},F}(z)=
 \frac{h^0(A,L)}{(g+1-2m)!}\,h_{\mathrm{BB},\mathrm{Lef}}(z)
 =h^0(A,L)(-1)^m h_K(z,z).
\]
Indeed $g+1-2m=0$, and the Lefschetz operator in that height is the
identity. Combining the two identities proves the diagonal case of
\eqref{eq:AR10}. The kernel $\overline Q$ is invariant under interchange
of the factors, so its arithmetic intersection against $z\times w$
is symmetric in $z,w$. Apply the diagonal identity to $z+w$ and
subtract the identities for $z$ and $w$; division by two proves the
bilinear statement.

The relative dimension of each $V_i\times V_j$ is $2d=2m-2$.
The displayed Deligne pairing therefore has exactly one more line-bundle
argument than the relative dimension. Its arithmetic degree is the
arithmetic intersection of those arguments, including their canonical
metrics. This is the metric Deligne-pairing identity recalled in Gu,
\S2.2, printed p.~5; its extension to rational cycles is multilinearity,
as also stated in Corollary 1.6, printed p.~3. It proves the assertion
about $\overline{\mathcal B}_z$. Finite base change multiplies the degree
by the extension degree, so the normalized equality is unchanged.
The arithmetic degree is independent of the chosen representative by
\eqref{eq:AR10}; no canonical isometry between the lines from all
representatives is needed or asserted.

Finally, use the following adelic rank-one version of the slope
argument AR-6. An integrable adelic rational line is not automatically
an integral hermitian lattice, so this extension is made explicit.
Let $\overline E=\overline{\mathcal B}_z^{\vee}$, let
$\varphi:E\hookrightarrow F$ be an injective linear map of adelic
normed spaces, and give its image $I$ the norms induced from $F$.
With normalized local weights $n_v$, put
$h(\varphi)=\sum_v n_v\log\|\varphi\|_v$.
For a nonzero vector $e\in E$, the local operator bound yields
$\|\varphi(e)\|_v\le\|\varphi\|_v\|e\|_v$.
Negate logarithms and sum to get
\[
 \widehat{\deg}_n E\le
 \widehat{\deg}_n I+h(\varphi)
 \le\widehat\mu_{\max}(F)+h(\varphi).
\]
The product formula makes the degrees independent of the chosen vector.
This proof works for integrable adelic norms and rational powers by
linearity; no passage to an integral lattice is being assumed.
Since $\widehat{\deg}_n E=-(-1)^m h_K(z,z)$, a nonpositive
right side proves the required nonnegative height, and a negative
right side proves strict positivity. The construction has fixed its metric through
$\overline L$; it has not assigned an arbitrary metric using the
unknown height. Therefore the only additional hypotheses of that
argument are the actual injective map and its norm/slope estimate.
\end{proof}

\paragraph{Why this does not already give the sign.}
The adelic kernel $\overline Q$ is integrable, not nef in general.
Even its underlying line bundle has
\[
 \Delta^*Q=2L,\qquad (1,-1)^*Q=-2L
 \quad\text{in }\operatorname{NS}(A)_{\mathbf Q}.
\]
Both identities follow from the theorem of the square and the
rigidification at the identity. The second restricts negatively to
curves. Consequently one cannot replace the remaining arithmetic
inequality by a blanket assertion that intersections of $\overline Q$
are nonnegative. Gu's own Conjecture 1.1, printed p.~1, explicitly leaves
the general Fourier-height positivity problem open. The construction
above supplies the geometric height line for a Bost approach; it does
not supply the norm estimate.
% New deduction, root, JAC1964-20261121. Whole-proof review recorded.
\subsection*{AR-18. A direct comparison of canonical surface heights}

The adelic formula also turns the missing sign into a concrete comparison
between two effective cycles. This gives a second geometric target for
the slope approach; it is not a proof of the comparison.

\begin{proposition}[AR-18]
Let $A/K$ be an abelian threefold over a number field, with a rigidified
symmetric ample line bundle $L$ and its canonical metric. Let $V$ be an
effective rational one-cycle, put
\[
 z=\operatorname{Ce}(V)=\frac{V-[-1]_*V}{2},
 \qquad Lz=0\quad\text{in }\operatorname{CH}^3(A)_{\mathbf Q},
\]
and let $\mu_\pm(x,y)=x\pm y$. Define the effective rational two-cycles
$S_\pm=(\mu_\pm)_*(V\times V)$, with the usual proper-pushforward
multiplicities and zero contribution from a component whose image has
smaller dimension. Their unnormalized canonical intersection heights are
\[
 H_\pm(V)=\frac1{[K:\mathbf Q]}
       \widehat{\deg}(\overline L^3\cdot S_\pm).
\]
Then
\[
 h_K(z,z)=\frac{H_-(V)-H_+(V)}{12h^0(A,L)}.                 \tag{AR.34}
\]
Every primitive class in $\operatorname{CH}^2_{(1)}(A)_{\mathbf Q}$
admits an effective rational representative in this sense. Consequently
semipositivity on that whole space is equivalent to
$H_-(V)\ge H_+(V)$ for every such $V$. The known nonnegativity
$H_\pm(V)\ge0$ does not itself order these two numbers.
\end{proposition}

\begin{proof}
Write $\overline L_i=p_i^*\overline L$ and
$\overline Q=\mu_+^*\overline L-\overline L_1-\overline L_2$.
The theorem of the cube and the canonical rigidifications give the
metric identities
\[
 \mu_\pm^*\overline L=\overline L_1+\overline L_2\pm\overline Q,
 \qquad (1,[-1])^*\overline Q=-\overline Q.
\]
Inversion in either factor therefore makes intersection against an odd
power of $\overline Q$ odd in the corresponding cycle. In particular
\[
 \widehat{\deg}(\overline Q^3\cdot(V\times V))
 =\widehat{\deg}(\overline Q^3\cdot(z\times z)).             \tag{AR.35}
\]
All identities are identities of arithmetic intersections of the chosen
cycles; descent of an arbitrary such intersection to a bare Chow group
is not being assumed.

We need the vanishing of the mixed term. Let $D=L\cdot V$ and let
$u=\operatorname{alb}_A(D)$, where the origin of $A$ defines the
Albanese sum for zero-cycles of any degree. Since $L$ is symmetric,
\[
 \operatorname{alb}_A(Lz)
 =\tfrac12\bigl(\operatorname{alb}_A(D)
                    -\operatorname{alb}_A([-1]_*D)\bigr)=u.
\]
Thus $u=0$ in $A(K)\otimes\mathbf Q$ under the stated primitive
hypothesis. All line bundles and pairings in this paragraph are taken
in the rational adelic Picard group; equivalently, one can first clear
the rational coefficients and then kill the resulting torsion point.

Here is a metric justification for using this ordinary zero-cycle.
Normalize each curve in $V$ and use additivity of the Deligne pairing.
Over the first factor $A$, the relative Deligne pairing of
$\overline Q$ with $\overline L$ along the second curve has underlying
line bundle $Q|_{A\times\{u\}}$. This is the norm description of the
pairing: a divisor for $L$ gives the tensor product of its Poincare
fibres, and the theorem of the cube identifies that product with the
fibre at the Albanese sum. Moreover, pulling back by $[n]$ on the first
factor multiplies this metrized pairing by $n$, since it does so to
$\overline Q$. The rigidification at zero is trivial. Hence its metric
is exactly the canonical Picard metric. Uniqueness follows directly:
the logarithm of a ratio of two such continuous metrics is bounded on
the compact analytic space and satisfies $f([n]x)=n f(x)$, forcing
$f=0$. For $u=0$ in $A(K)\otimes\mathbf Q$ this is the zero class
of the rational adelic Picard group. Applying the
Deligne-pairing projection formula once more gives
\[
 \widehat{\deg}(\overline L_1\overline L_2\overline Q
                   \cdot(V\times V))=0.                   \tag{AR.36}
\]
Also $\overline L_1^2\overline Q$ and
$\overline L_2^2\overline Q$ have zero intersection against
$V\times V$: pair first over the other curve; the degree of its
Poincare restriction is zero, so the projection formula gives the
zero tensor power. These observations apply to singular curves through
normalization and proper projection, and to rational coefficients by
multilinearity.

Expanding the two cubes now gives
\[
 \begin{split}
 [K:\mathbf Q](H_-(V)-H_+(V))
 &=-6\widehat{\deg}\bigl((\overline L_1+\overline L_2)^2
                           \overline Q\cdot(V\times V)\bigr)\\
 &\phantom{={}}-2\widehat{\deg}\bigl(\overline Q^3
                           \cdot(V\times V)\bigr)\\
 &=-2\widehat{\deg}\bigl(\overline Q^3\cdot(z\times z)\bigr).
 \end{split}
\]
The odd part of a one-cycle on a threefold has Beauville grade one:
Fourier sends $\operatorname{CH}^2_{(s)}$ to
$\operatorname{CH}^{1+s}_{(s)}$; negative grades and grades above two
vanish in these degrees, and inversion selects odd $s$.
AR-10 with $m=2$ therefore identifies the last expression with
$12h^0(A,L)[K:\mathbf Q]h_K(z,z)$, proving (AR.34).

Finally represent any grade-one primitive class $z$ by a finite rational
cycle $Z=Z_+-Z_-$ with $Z_+,Z_-$ effective. Set
$V=Z_++[-1]_*Z_-$. It is effective and
$\operatorname{Ce}(V)=\operatorname{Ce}(Z)=z$ in rational Chow groups.
The already proved identity supplies the claimed equivalence.
Nonnegativity of the canonical heights of effective cycles is Zhang,
\emph{Small points and adelic metrics}, Theorem 2.4(b), author manuscript
pp.~11--12; apply it to each component of $S_\pm$ and add. It gives two
separate inequalities, not their difference.\end{proof}

\paragraph{Exact next gate for this route.}
For the complementary Fermat group the comparison is required for all
effective $V$ whose Ceresa class lies in $N$, not only for the Fermat
curve. Neither canonical-height nonnegativity nor the two nef bundles
$\mu_\pm^*\overline L$ produces the missing ordering. The adelic height
line of AR-10 can equally be represented, at the level of arithmetic
degree, by the difference of these two surface Deligne pairings divided
by $12h^0(A,L)$. A Bost proof along this route must provide a comparison
map and its norm bound between those particular pairings.

The exact background is Gu, arXiv:2609.33223v2, Theorem 1.3,
Proposition 5.3 and \S2.2, and Zhang, Theorem 2.4(b),
\url{https://web.math.princeton.edu/~shouwu/publications/smpoint.pdf}.
The relation (AR.34), the primitive mixed-term vanishing, and the
effective representative construction are the deduction here.
\section{A direct positive central derivative for the Fermat-quartic
motive}\label{a-direct-positive-central-derivative-for-the-fermat-quartic-motive}

Run JAC1964-20261121. Author/checker: /root/arithmetic\_sign, GPT-6
Astra Pro; effort setting not exposed. This is an elementary deduction
from the stated known CM identification and functional equation. It
proves no height identity.

\subsection{Proposition AR-11}\label{proposition-ar-11}

Let \(M\) be the rank-two motive of types \((3,0),(0,3)\) in the third
cohomology of the Jacobian of \(x^4+y^4=z^4\), over \(\mathbf Q\).
Normalize its motivic \(L\)-function by its Euler series
\(L(M,s)=\sum_{n\ge1}a_n n^{-s}\), with \(a_1=1\). Then \[
L'(M,2)>\frac7{20}>0.
\tag{AR.14}
\] In particular the analytic sign required in a positively normalized
height formula for this particular motive is unconditional. No GRH is
used.

\subsubsection{Known CM input and
normalization}\label{known-cm-input-and-normalization}

Write \(F=\mathbf Q(i)\), \(\mathfrak q=(1+i)\), and let \(\psi,\psi'\)
be the algebraic Hecke characters of the elliptic curves \(E:y^2=x^3-x\)
and \(E':y^2=x^3+x\). The relevant character is \[
\chi=\psi^2\psi'=\psi^3\mu,
\] where \(\mu\) corresponds to \(F(\sqrt2)=\mathbf Q(\zeta_8)\).
Bertolini--Darmon--Prasanna, \emph{p-adic L-functions and the coniveau
filtration on Chow groups}, §5.2, pp.24--27 of the read author
manuscript, gives this identification, the conductors
\(\mathfrak q^3,\mathfrak q^4\) of \(\psi,\psi'\), the conductor
\(\mathfrak q^2\) of \(\psi^2\), and the root-number formula. For
exponent three with the twist \(\mu\), that root number is \(-1\).

Because \(\psi^2\) is trivial on \(1+\mathfrak q^2\) whereas \(\psi'\)
is nontrivial on \(1+\mathfrak q^3\), the product \(\chi=\psi^2\psi'\)
has conductor exactly \(\mathfrak q^4\). Its infinity type is \((3,0)\).
Thus its normalized CM theta series \(f(z)=\sum a_n e^{2\pi i nz}\) is a
weight-four cusp form of level
\(N=|D_F|N_{F/\mathbf Q}(\mathfrak q^4)=64\), with trivial nebentypus.
Its \(L\)-function is \(L(M,s)\), with center \(2\) and sign \(-1\). The
standard modular description is independently recorded as LMFDB newform
64.4.a.c; the argument below does not use the database's numerical
analytic rank.

For an ideal \(\mathfrak a\) prime to \(\mathfrak q\), we have
\(|\chi(\mathfrak a)|=(N\mathfrak a)^{3/2}\). The number of ideals in
\(\mathbf Z[i]\) of norm \(n\) is at most \(d(n)\), the divisor count.
Consequently \[
|a_n|\le d(n)n^{3/2}\le2n^2.
\tag{AR.15}
\] The final inequality follows by pairing divisors below and above
\(\sqrt n\). The conductor removes every even-norm term. There is no
ideal of norm \(3\). Hence \[
a_1=1,\qquad a_2=a_3=a_4=0.
\tag{AR.16}
\]

\subsubsection{Mellin calculation}\label{mellin-calculation}

Set \(a=1/8=N^{-1/2}\). The Mellin transform is \[
I(s)=\int_0^\infty f(iy)y^{s-1}\,dy
     =(2\pi)^{-s}\Gamma(s)L(M,s).
\] The Fricke relation, with root number \(-1\) and weight four, reads
\[
f\!\left(\frac{i}{Ny}\right)=-N^2y^4 f(iy).
\] Substituting \(y\mapsto1/(Ny)\) in the part below \(a\) gives \[
I(s)=\int_a^\infty f(iy)
 \left(y^{s-1}-N^{2-s}y^{3-s}\right)\,dy.
\] At \(s=2\) the integrand vanishes. Differentiating gives \[
I'(2)=2\int_a^\infty f(iy)y\log(8y)\,dy.
\] Since \(L(M,2)=0\), the Mellin prefactor contributes no derivative
term and \(L'(M,2)=(2\pi)^2I'(2)\).

Let \(E_1(x)=\int_x^\infty e^{-t}\,dt/t\). Integration by parts gives \[
\int_a^\infty e^{-2\pi ny}y\log(y/a)\,dy
 =\frac{e^{-cn}+E_1(cn)}{(2\pi n)^2},
\qquad c=\frac\pi4.
\] The coefficient bound (AR.15) gives absolute convergence on
\(y\ge a\), justifying the exchange of sum and integral. Thus \[
L'(M,2)
 =2\sum_{n\ge1}\frac{a_n}{n^2}
       \left(e^{-cn}+E_1(cn)\right).
\tag{AR.17}
\]

\subsubsection{Elementary lower bound}\label{elementary-lower-bound}

For \(x>0\), \[
0<E_1(x)\le e^{-x}/x.
\] Put \(q=e^{-c}\). Since \(3/4<c<1\), we have \(q<1/2\) and \(q>1/3\):
indeed \(e^{3/4}>1+3/4+(3/4)^2/2>2\), while \(e<3\). The \(n=1\) term of
(AR.17) is therefore greater than \(2/3\).

By (AR.15)--(AR.16), the absolute value of the sum of every remaining
term is bounded by \[
\begin{aligned}
2\sum_{n\ge5}\frac{|a_n|}{n^2}
       \left(e^{-cn}+E_1(cn)\right)
&\le4\sum_{n\ge5}q^n\left(1+\frac1{cn}\right)\\
&<4\cdot\frac{19}{15}\cdot
       \frac{(1/2)^5}{1-1/2}
 =\frac{19}{60}.
\end{aligned}
\] It follows that \(L'(M,2)>2/3-19/60=7/20\). This proves (AR.14).
\(\square\)

\subsection{Optional numerical refinement, independently
checked}\label{optional-numerical-refinement-independently-checked}

The accompanying script enumerates the unique primary Gaussian-integer
generator of each ideal of norm at most \(20\). It obtains exactly \[
a_1=1,\quad a_5=-22,\quad a_9=-27,\quad a_{13}=18,\quad a_{17}=-94,
\] with all other \(a_n\), \(1\le n\le20\), zero. These coefficients
also match the primary database's displayed expansion.

At precision 160 bits, python-flint/Arb evaluates the finite sum in
(AR.17). The tail is bounded mathematically by \[
4\left(1+\frac1{21c}\right)\frac{e^{-21c}}{1-e^{-c}}
 <5.357966\cdot10^{-7}.
\] The resulting enclosing rational interval is \[
1.5071685<L'(M,2)<1.5071697.
\tag{AR.18}
\] The elementary proof of the sign does not depend on this computation.
The script and its output are in fermat\_lderi\_bound.py and
fermat\_lderi\_bound.json.

\subsection{What this resolves and what it
leaves}\label{what-this-resolves-and-what-it-leaves}

For this one rank-two CM motive, the derivative sign is settled without
the supplied Dirichlet-GRH theorem or a hypothetical broader Hecke-GRH
theorem. A proven identity \(h(\eta,\eta)=C\,L'(M,2)\), with \(C>0\),
would now imply strict positivity immediately. No such identity is
established here. The actual known height formula remains (AR.10), with
the stronger finite-local information (AR.12); an automorphic comparison
must still control its ramified contribution and identify the particular
Chow cycle.

The generalized-Heegner formula of Lilienfeldt--Shnidman,
\emph{Derivatives of Rankin--Selberg L-functions and heights of
generalized Heegner cycles}, Theorem 1.1 (arXiv:2408.04375, read author
manuscript dated 8 August 2024), does not directly supply the missing
equality for the BDP Fermat construction. Its hypotheses require an
even-weight form, an unramified Hecke character over an imaginary
quadratic field of odd discriminant, and splitting of every level prime.
The BDP construction uses the weight-three form \(\theta_{\psi^2}\), the
ramified field \(\mathbf Q(i)\), and the ramified character \(\psi'\).
Their §1.5 explicitly lists odd weight and relaxed ramification/Heegner
conditions as further cases. This is an application-scope observation,
not a denial of their proven height formula.

\subsection{Bibliographic locators and read
sources}\label{bibliographic-locators-and-read-sources}

\begin{itemize}
\tightlist
\item
  Massimo Bertolini, Henri Darmon, Kartik Prasanna, with appendix by
  Brian Conrad, \emph{p-adic L-functions and the coniveau filtration on
  Chow groups}, J. reine angew. Math. 731 (2017), 21--86, DOI
  10.1515/crelle-2014-0150. Read author PDF:
  \url{https://sites.lsa.umich.edu/kartikp/wp-content/uploads/sites/1326/2024/07/BDP5-hl.pdf}
  ; §1 p.6, §5.2 pp.24--27, especially the local conductor discussion
  pp.24--25 and displayed root-number formulas p.27.
\item
  LMFDB, newform 64.4.a.c:
  \url{https://www.lmfdb.org/ModularForm/GL2/Q/holomorphic/64/4/a/c/} . CM
  field, conductor, weight, Fricke sign, and first coefficients directly
  read; the database's numerical analytic rank is not a proof input.
\item
  David T.-B. G. Lilienfeldt and Ari Shnidman, \emph{Derivatives of
  Rankin--Selberg L-functions and heights of generalized Heegner
  cycles}, arXiv:2408.04375; read primary PDF
  \url{https://perso.pages.math.cnrs.fr/users/david.lilienfeldt/arXivGZGHC.pdf}
  ; Theorem 1.1 p.2 and §1.5 p.5.
\end{itemize}

Self-check: conductor product on \(1+\mathfrak q^3\), weight/conductor
normalization, Fricke factor \(N^2\), Mellin derivative factor \(2\),
gamma cancellation at the zero, integration-by-parts kernel, coefficient
growth, vanished first terms, geometric-series tail, and
\(2/3-19/60=7/20\) have each been checked. This is a written proof with
a supplementary arithmetic interval check, not a formal proof-assistant
certificate.
% JAC1964-20261121. New finite-place deduction; no reliance on an integer-k
% assertion or on reading node thickness from a stable intersection graph.
% /root/programme_audit; GPT-6 Astra Pro; effort setting not exposed.

\subsection*{JAC-P10. The finite-place Hodge lattice of the Fermat quartic}

\begin{proposition}[JAC-P10]
\label{jac:P10}
Let $C:x^4+y^4+z^4=0$ over $\overline{\mathbf Q}$, let
$J=\operatorname{Jac}(C)$, and let $E$ be any elliptic curve with
$j(E)=1728$. For stable Faltings heights, in any common normalization,
\[
 h_F(J)-3h_F(E)\ =\ \frac14\log 2.                 \tag{JAC.P10.1}
\]
More precisely, for the degree-eight isogeny
$\Phi:J\to E_x\times E_y\times E_z$ induced by the three coordinate
sign quotients, the normalized valuation of the determinant of
pullback on the good-reduction Hodge lattices at every place above
$2$ is exactly $7/4$, when the valuation is normalized by $\nu(2)=1$.
\end{proposition}

\begin{proof}
The statement is unchanged by a finite extension of the number field.
We may therefore work locally over a finite extension of $\mathbf Q_2$
on which the curve is stable, its stable components and nodes are split,
and the three elliptic quotients have good reduction. We may make
further finite extensions in the local calculation: normalized
determinant valuations of maps between good-reduction Hodge lattices
are invariant under such extension.

\paragraph{A smooth chart at an elliptic tail.}
Put $u=x/z$, $v=y/z$. The quotient by $u\mapsto-u$ has affine
equation $w^2+v^4+1=0$, with $w=u^2$.
Liu, \emph{Algebraic Geometry and Arithmetic Curves},
Example~10.4.39, printed pp.~548--549, gives constants
$\xi_8,\alpha,\beta,\gamma$ over a finite extension such that
\[
 \begin{gathered}
 \xi_8v=1+\frac1{t+\alpha},\qquad
 w=\frac{2s+\beta t+\gamma}{(t+\alpha)^2},\\
 s^2+(\beta t+\gamma)s=t^3,\qquad
 \nu(\alpha)=-\frac14,\quad
 \nu(\beta)=\frac12,\quad \nu(\gamma)=0.
                                                        \tag{JAC.P10.2}
 \end{gathered}
\]
The last equation is a smooth good-reduction elliptic model: its
reduction is $s^2+\bar\gamma s=t^3$.
Choose $a=\sqrt\gamma/\alpha$, put $b=\alpha^{-2}$, and substitute
$u=a+bU$. Thus $\nu(a)=1/4$ and $\nu(b)=1/2$.
The curve is described, on a suitable affine open, by the two
integral equations
\[
 \begin{split}
 G={}&s^2+(\beta t+\gamma)s-t^3=0,\\
 F={}&\gamma t^2+(2\gamma\alpha-\beta\alpha^2)t
       -2\alpha^2s
       +2a(t+\alpha)^2U
       +\alpha^{-2}(t+\alpha)^2U^2=0.
 \end{split}                                             \tag{JAC.P10.3}
\]
Indeed $F$ is $\alpha^2$ times
$u^2(t+\alpha)^2-(2s+\beta t+\gamma)$, after using
$a^2=\gamma/\alpha^2$. Every displayed coefficient is integral.
Its reduction is
\[
 \bar G=s^2+\bar\gamma s-t^3,
 \qquad
 \bar F=U^2+\bar\gamma t^2-\overline{\beta\alpha^2}t,
                                                        \tag{JAC.P10.4}
\]
where both $\bar\gamma$ and $\overline{\beta\alpha^2}$ are nonzero.
The Jacobian determinant of $(F,G)$ with respect to $(t,s)$ is a
unit. Hence this chart is smooth over the valuation ring and $dU$
generates its relative differentials.

The reduction of $F$ gives a purely inseparable extension of degree
two of the elliptic function field defined by $\bar G$. It is
nontrivial because the derivative of
$\bar\gamma t^2-\overline{\beta\alpha^2}t$ is nonzero.
It therefore has genus one. This smooth chart supplies the generic
point of one elliptic component $\Gamma_x$ in the stable reduction.
To justify persistence, compare a proper completion of this chart
with the stable model through a common regular resolution. A
component that disappears over a nodal stable model belongs to the
rational exceptional chains resolving nodes or to rational curves
introduced by point blowups. The genus-one divisorial valuation
just constructed cannot be among them. Equivalently, this is the
reduced-special-fibre case of Liu, Proposition~10.4.37, p.~547.
Here
\[
 \nu_{\Gamma_x}(u)=\frac14,\qquad
 \nu_{\Gamma_x}(v)=0,
 \qquad
 \nu_{\Gamma_x}\!\left(\frac{du}{v^3}\right)=\frac12.
                                                        \tag{JAC.P10.5}
\]
The last equality follows directly from $du=\alpha^{-2}dU$.
The cyclic permutation $(x,y,z)\mapsto(y,z,x)$ supplies the other
two elliptic tails. Their distinction follows from their different
valuations of $u,v$. Since these three genus-one components already
account for arithmetic genus three, every other component is
rational and the dual graph is a tree. A stable rational vertex has
valence at least three. The order-three cyclic symmetry permutes
the three positive-genus vertices, so none can be the middle vertex
of a path connecting them. The tree has just one rational branching
vertex and the three genus-one leaves; any additional rational
vertex would have valence at most two. This is the stable star,
also stated in Liu's example.

\paragraph{Three integral stable differentials.}
Write
\[
 \eta_0=\frac{du}{v^3},\qquad
 \eta_1=u\eta_0,\qquad \eta_2=v\eta_0,
 \qquad \kappa=(1+u+v)\eta_0.
                                                        \tag{JAC.P10.6}
\]
These are holomorphic differentials on the smooth plane quartic;
$\eta_0,\eta_1,\eta_2$ form a basis. At $\Gamma_x$ their vertical
valuations are respectively $1/2,3/4,1/2$.
Since $\nu(\xi_8+1)=1/4$ and $\nu((t+\alpha)^{-1})=1/4$,
formula (JAC.P10.2) gives $\nu_{\Gamma_x}(1+v)\ge1/4$.
Together with $\nu_{\Gamma_x}(u)=1/4$ this yields
\[
 \nu_{\Gamma_x}(\kappa)\ge\frac34.
\]
The cyclic permutation permutes $\eta_0,\eta_1,\eta_2$ and fixes
$\kappa$. Thus each $\eta_i$ has valuation at least $1/2$ on
every elliptic tail, and $\kappa$ has valuation at least $3/4$
on every elliptic tail.

We spell out why checking the tails suffices for integrality.
Let $M=H^0(\mathcal C,\omega_{\mathcal C/R})$ for the stable model.
This is a free rank-three lattice and its formation commutes with
residue-field base change. For a compact-type stable curve with
this star graph, restriction gives an injection
\[
 H^0(\mathcal C_s,\omega_{\mathcal C_s})
       \hookrightarrow\bigoplus_{i=x,y,z}
                           H^0(\Gamma_i,\omega_{\Gamma_i}).
                                                        \tag{JAC.P10.7}
\]
Indeed a dualizing differential has opposite residues on the two
branches of each node; the single residue on each elliptic tail
is zero, so all three residues on the rational component are zero,
and that component has no holomorphic differentials.
Suppose a generic holomorphic differential has nonnegative vertical
valuation on all three tails but is not in $M$. Multiplying it by
the smallest positive power of a uniformizer that puts it in $M$
would give a primitive section whose reduction vanishes on all
three tails. Equation (JAC.P10.7) and base change imply that its
reduction is zero, contradicting primitivity. This proves the claim.

Choose constants $c_2,c_3$ with valuations $1/2,3/4$, respectively.
The preceding observations prove that
\[
 f_1=\eta_0/c_2,\qquad f_2=\eta_1/c_2,
 \qquad f_3=\kappa/c_3
 \quad\text{all belong to }M.                            \tag{JAC.P10.8}
\]
In fact these three differentials form a basis of $M$. To verify
this stronger assertion, take $c_3=c_2c_1$ with $\nu(c_1)=1/4$.
Putting $t=s=0$ in the elliptic-coordinate identity gives
\[
 \gamma^2=(\alpha+1)^4-\alpha^4
   =1+4\alpha^3+6\alpha^2+4\alpha.
\]
Every nonconstant term has positive valuation; hence
$\bar\gamma=1$, and also $\overline{\sqrt\gamma}=1$.
At the generic point of $\Gamma_x$ we have the exact expansion
\[
 \begin{split}
 1+u+v={}&(1+\xi_8^{-1})
       +\frac{\sqrt\gamma+\xi_8^{-1}}{\alpha}
       +\alpha^{-2}U\\
 &+\xi_8^{-1}\bigl((t+\alpha)^{-1}-\alpha^{-1}\bigr).
 \end{split}                                             \tag{JAC.P10.8a}
\]
The first term has valuation exactly $1/4$. The second has valuation
strictly greater than $1/4$, because both terms in its numerator
reduce to one; the last two have valuation at least $1/2$.
Thus $(1+u+v)/c_1$ reduces to a nonzero constant $c$.
Choose the nonzero reduction of $\eta_0/c_2$ as a basis of
holomorphic differentials on $\Gamma_x$, and transport it cyclically
to the other tails. Since $\bar v=1$ and $\nu(u)>0$
on $\Gamma_x$, the restrictions of $f_1,f_2,f_3$ to the three tails
are, up to the choice of cyclic ordering, the rows of
\[
 \begin{pmatrix}
 1&0&c\\ 0&1&c\\ 1&1&c
 \end{pmatrix}.
                                                        \tag{JAC.P10.8b}
\]
Its determinant is $c\ne0$ in residue characteristic two.
Restriction in (JAC.P10.7) is an isomorphism, since both sides have
dimension three. The reductions of $f_1,f_2,f_3$ are therefore a
basis of $M/\mathfrak m_R M$, and Nakayama's lemma makes them a
basis of $M$. Their change-of-basis determinant from
$(\eta_0,\eta_1,\eta_2)$ has valuation $-7/4$.
Consequently the lattice spanned by the $\eta_i$ has index in $M$
of normalized valuation exactly $7/4$.

\paragraph{Identification of the pullback lattice and the height.}
An invariant differential on the good model of the first elliptic
quotient is
\[
 \frac{dt}{2s+\beta t+\gamma}.
\]
By (JAC.P10.2), its pullback is a unit multiple of $dv/w=-\eta_1$.
The other two quotient pullbacks are unit multiples of
$-\eta_2$ and $-\eta_0$. Thus the lattice just computed is precisely
the image of the Hodge lattice under
$\Phi^*:H^0(E_x\times E_y\times E_z,\Omega^1)\to H^0(J,\Omega^1)$.
The identification with $M$ follows from the relative Picard
description of a compact-type Jacobian and relative Serre duality.

For completeness, $\deg\Phi=8$. Each quotient map has degree two;
its norm followed by pullback is $[2]$, and the mixed norm--pullback
maps are zero because the three spaces have different characters
under the coordinate-sign group. With principal polarizations,
$\Phi\Phi^t=[2]$ on the product, and hence
$(\deg\Phi)^2=2^6$.

Let $r_v$ be the unnormalized length of the cokernel of the Hodge
pullback over a number field $L$ where all varieties have good
reduction, and let $e_v=\operatorname{ord}_v(2)$.
The local calculation gives $r_v=7e_v/4$ for $v\mid2$.
At every other finite place the isogeny has invertible degree,
so its Hodge pullback is an isomorphism.
The isogeny height formula, with its sign written explicitly, is
\[
 h_F(E_x\times E_y\times E_z)-h_F(J)
 =\frac12\log8-
   \frac1{[L:\mathbf Q]}\sum_{v\mid2} r_v\log N(v).
                                                        \tag{JAC.P10.9}
\]
One may obtain the sign directly: pullback multiplies the squared
archimedean norm of a top differential by $\deg\Phi$, while the
finite lattice inclusion has the displayed cokernel lengths.
It also agrees with R\'emond, \emph{Propri\'et\'es de la hauteur
de Faltings}, Proposition~2.1 and its proof, author PDF pp.~2--3,
and the differential formulation in \S3, p.~4.
As $\sum_{v\mid2}e_v\log N(v)=[L:\mathbf Q]\log2$, we obtain
\[
 h_F(J)-\sum_i h_F(E_i)
 =\left(\frac74-\frac32\right)\log2.
\]
All three elliptic quotients have $j=1728$, so their stable heights
are equal. This proves (JAC.P10.1).
\end{proof}

\paragraph{What this does and does not use.}
This proof computes a Hodge-lattice determinant.
It does not infer a node thickness from the unweighted intersection
graph, and it does not use the claim in Gu\`ardia's Theorem~12 sketch
that every normalized prime-power isogeny height difference has an
integer coefficient of $\log p$. The determinant calculation is
compatible with arbitrary further ramification, since all valuations
were normalized by $\nu(2)=1$.

\paragraph{Primary source locators.}
Liu's item is \emph{Example}~10.4.39, pp.~548--549 of the 2006
paperback, with the elliptic coordinate computation referred to
Example~10.1.28. The full primary book text at these pages was read
through the PDF mirror
\url{https://rexresearch1.com/GeometryLibrary/AlgebraicGeometryArithmeticCurvesLiu.pdf};
the author book page is
\url{https://www.math.u-bordeaux.fr/~qliu/Book/index.html}.
R\'emond's primary author PDF is
\url{https://www-fourier.univ-grenoble-alpes.fr/~remond/2476.pdf}.
The new content here is the smooth tail chart (JAC.P10.3), the
integral basis (JAC.P10.8)--(JAC.P10.8b), and the resulting exact
height difference. The weaker lower-bound version and the exact
basis refinement were independently checked in this run.
\section{A theta-integral reduction for the Fermat
height}\label{a-theta-integral-reduction-for-the-fermat-height}

Run JAC1964-20261121. Author/checker /root/arithmetic\_sign; GPT-6 Astra
Pro, effort setting not exposed.

\subsection{Proposition AR-12}\label{proposition-ar-12}

Let \(C\) be a nonhyperelliptic complex curve of genus three, with
period matrix \(\tau\in\mathfrak H_3\), and let
\(J=\mathbf C^3/(\mathbf Z^3+\tau\mathbf Z^3)\). Use Haar measure
\(\mu\) of mass one. Write \[
T(\tau)=\prod_{\epsilon\ {\rm even}}\theta_\epsilon(0,\tau),
\qquad
I(\tau)=\int_J\log\left|
  \frac{\phi(z,\tau)}{\theta_a(z,\tau)^2}
\right|\,d\mu(z),
\] where \(\phi\) is Frobenius' theta function with the normalization of
de Jong 2020, Theorem 4.1, and \(a\) is any theta characteristic. The
integral is finite and independent of \(a\). The individual archimedean
term in AR-9 is \[
A(C)=-\frac1{18}\log\|\chi'_{18}\|_{\rm Hdg}-\lambda(C)
     =\frac{I(\tau)}7-\frac1{18}\log|T(\tau)|
        -\frac{10}{63}\log2.
\tag{AR.19}
\] Consequently, for the Fermat curve over \(\mathbf Q\), every
archimedean term has the same value. With the normalized averaging
convention of AR-9, the sufficient strict height inequality at this
stage becomes \[
I(\tau_F)>
\frac7{18}\log|T(\tau_F)|+\frac{10}{9}\log2-\frac{D_2}{3}.
\tag{AR.20}
\]

\subsubsection{Proof with all
normalizations}\label{proof-with-all-normalizations}

Put \(H=\log\|H\|(C)\), \(K=\log\|K\|(C)\) in de Jong's notation. His
2020 Theorem A and equation (5), printed pp.1389--1390, give \[
\lambda(C)=-\frac17K+\frac27H
                +\frac{12}{7}\log2-3\log(2\pi).
\] The normalized theta norms from his §5.3 have weight factors
\((\det\operatorname{Im}\tau)^4\) for \(\phi\) and
\((\det\operatorname{Im}\tau)^{1/4}\) for \(\theta_a\). Their Gaussian
factors cancel in the quotient, so \[
K-2H=\frac72\log\det\operatorname{Im}\tau+I(\tau).
\] The equality of the integrals of \(\log\|\theta_a\|\) for different
\(a\) follows from translation invariance and his equation (22);
logarithms of zeros and poles of a nonzero meromorphic function on a
compact complex manifold are locally integrable.

De Jong 2018, §8, uses the primitive algebraic modular form \[
\chi'_{18}=2^{-28}(2\pi i)^{54}T(\tau)
           (dz_1\wedge dz_2\wedge dz_3)^{\otimes18}.
\] His §3, equation (3.1), uses the Hodge norm
\(\|dz_1\wedge dz_2\wedge dz_3\|_{\rm Hdg}
=(\det\operatorname{Im}\tau)^{1/2}\), agreeing with the 2020 convention.
Hence \[
-\frac1{18}\log\|\chi'_{18}\|_{\rm Hdg}
=\frac{14}{9}\log2-3\log(2\pi)
 -\frac1{18}\log|T(\tau)|
 -\frac12\log\det\operatorname{Im}\tau.
\] Subtract the formula for \(\lambda\). The determinant and \(2\pi\)
terms cancel, and \(14/9-12/7=-10/63\). This proves (AR.19). Since \(C\)
is defined over \(\mathbf Q\), all its complex base changes are
isomorphic, so the normalized archimedean average equals this value.
Substituting (AR.19) into \(A_\infty>-D_2/21\) and multiplying by seven
proves (AR.20). \(\square\)

\subsection{An exact period matrix available for
evaluation}\label{an-exact-period-matrix-available-for-evaluation}

The following matrix can be used: \[
\tau_F=
\begin{pmatrix}
\frac12+i&-\frac12&-\frac{1+i}{2}\\
-\frac12&\frac12+i&\frac{1+i}{2}\\
-\frac{1+i}{2}&\frac{1+i}{2}&1+i
\end{pmatrix}.
\tag{AR.21}
\] It follows by specializing Guàrdia 2001, Proposition 5, printed
p.160. Set \(a=1+\sqrt2\). In that proposition the scalar is \[
t_a=\frac{a^2}{2(a^2+1)}
\frac{\mathcal K((a-a^{-1})/(a+a^{-1}))}{\mathcal K(a^{-2})}.
\] Here the arguments are \(1/\sqrt2\) and \((\sqrt2-1)^2\). The
descending Landen identity, DLMF 19.8.11--12, gives \[
\mathcal K(1/\sqrt2)
=(4-2\sqrt2)\mathcal K((\sqrt2-1)^2),
\] and substitution gives \(t_a=1/2\). The general symmetric matrix in
Proposition 5 then gives (AR.21). Its imaginary part has eigenvalues
\(1,1\pm1/\sqrt2\), so it lies in \(\mathfrak H_3\), with determinant
\(1/2\).

The specialized display on Guàrdia's p.161 has an asymmetric entry: the
(1,2) entry contains an extra \(+i\), while the (2,1) entry is \(-1/2\).
This is present in the original scanned page, not just the extracted
text. We therefore derive (AR.21) from the general symmetric Proposition
5 and the exact Landen identity rather than quote that inconsistent
display.

\subsection{Scope of the progress}\label{scope-of-the-progress}

Formula (AR.19) eliminates spectral invariants from the archimedean
quantity and fixes every normalization. It supplies the concrete theta
integral and fixed algebraic period matrix used subsequently in AR-13.
That theorem proves a rigorous lower bound by one-sided torsion-grid and
cell-Jensen estimates with a finite enclosed computation, thereby
closing the sign problem for the specified Fermat class. The earlier
exploratory numerical integrations remain explicitly separate from that
certificate. Positivity on the entire complementary Chow group is not
asserted.

\subsection{Primary sources and
checks}\label{primary-sources-and-checks}

\begin{itemize}
\tightlist
\item
  Robin de Jong, Frobenius' theta function and Arakelov invariants in
  genus three, Pure Appl. Math. Q. 16(5) (2020), 1387--1418. Theorem A
  and equation (5), pp.1389--1390, rendered and visually checked;
  §5.1--5.4 pp.1404--1406 and §4.5 p.1401 read for norms and the
  meromorphic quotient.
  \url{https://archive.intlpress.com/site/pub/files/_fulltext/journals/pamq/2020/0016/0005/PAMQ-2020-0016-0005-a002.pdf}
\item
  Robin de Jong, On the height of Gross--Schoen cycles in genus three
  (2018), §§3,8, for the Hodge metric and primitive algebraic
  modular-form normalization.
  \url{https://pmc.ncbi.nlm.nih.gov/articles/PMC6413647/}
\item
  Jordi Guàrdia, Arakelov computations in genus 3 curves, J. Théor.
  Nombres Bordeaux 13(1) (2001), 157--165, Proposition 5 pp.160--161,
  both pages rendered and visually checked.
  \url{https://www.numdam.org/item/JTNB_2001__13_1_157_0.pdf}
\item
  NIST Digital Library of Mathematical Functions, §19.8(ii), equations
  19.8.11--12. \url{https://dlmf.nist.gov/19.8}
\end{itemize}

Independent review: h9\_height directly read the 2020 formulas and 2018
normalizations and rederived (AR.19), confirming the coefficient
\(-10/63\), determinant cancellation and \(2\pi\) cancellation. No
numerical sign was inferred by that check.
\section{A positive height for the nonzero complementary Fermat
class}\label{a-positive-height-for-the-nonzero-complementary-fermat-class}

Run JAC1964-20261121. Author /root/arithmetic\_sign; GPT-6 Astra Pro;
effort setting not exposed. The finite-place input JAC-P10 is due to
/root/programme\_audit and was independently whole-read here and by
/root/h9\_height.

\subsection{Proposition AR-13}\label{proposition-ar-13}

Let \(\eta=F_{\rm Zhang}(X_1)\) be the canonical Fourier-Ceresa class of
the Fermat quartic specified in AR-8, and use the normalized raw
Beilinson--Bloch height \(h_K\). Then \[
h_K(\eta,\eta)>\frac1{100}>0.
\tag{AR.22}
\] This class is nonzero modulo algebraic equivalence and belongs to the
exact complementary Chow summand constructed in AR-8. The proposition
proves the sign for this class and its rational span. It makes no
dimension claim about the complementary Chow group and does not prove
positivity for all its elements.

The proof below uses the known height formulas in AR-8, AR-9, AR-12 and
the proved local lattice estimate JAC-P10. Its remaining analytic
inequality is certified by finite theta evaluations. It uses no GRH and
no height-to-\(L'\) identity.

\subsection{1. Finite contribution, including the height
normalization}\label{finite-contribution-including-the-height-normalization}

Use the Hodge norm \((i/2)\int\omega\wedge\bar\omega\) throughout. For
\(E:y^2=x^3-x\), the real period of \(dx/(2y)\) is \[
\Omega_E=\int_1^\infty\frac{dx}{\sqrt{x^3-x}}
 =\frac{\Gamma(1/4)^2}{2\sqrt{2\pi}},
\] by the Euler beta integral; its lattice is
\(\Omega_E(\mathbf Z+i\mathbf Z)\). Its discriminant is \(64\). On a
good-reduction model after finite extension, the invariant differential
therefore differs from \(dx/(2y)\) by a factor of valuation \(1/2\),
normalized by \(v(2)=1\). Hence \[
h_E=-\log\Omega_E-\frac12\log2
   =-2\log\Gamma(1/4)+\log2+\frac12\log\pi.
\tag{AR.23}
\] The minus sign in the finite term records that the original
differential is a pole of order \(e_v/2\) in the good Hodge lattice.

JAC-P10 proves \(h_F(J)-3h_E\ge\frac14\log2\). Put \[
c_\infty=-\frac1{18}\log\|\chi'_{18}\|_{\rm Hdg}
 =\frac{37}{18}\log2-3\log(2\pi)-\frac1{18}\log|T(\tau_F)|.
\] Here \(\det\operatorname{Im}\tau_F=1/2\), and every complex embedding
gives the same value. The primitive modular-form degree formula gives \[
F_{18}:=\frac1{18[K:\mathbf Q]}
 \sum_{v\mid2}\operatorname{ord}_v(\chi'_{18})\log N(v)
 =h_F(J)-c_\infty
 \ge3h_E+\frac14\log2-c_\infty.
\tag{AR.24}
\] The tree formula is \(\lambda(\Gamma_v)=2\delta_v/7\), while de
Jong's bound used in AR-9 gives
\(\operatorname{ord}_v(\chi'_{18})\ge6\delta_v\). Thus the entire
normalized finite contribution satisfies \[
B_{\rm fin}=F_{18}-\frac27D_2\ge\frac17F_{18}.
\tag{AR.25}
\] The finite theta computation below gives, with outward enclosures, \[
3h_E+\tfrac14\log2-c_\infty>0.1347786184,
\qquad B_{\rm fin}>0.0192540883.
\tag{AR.26}
\] No integrality assertion about an isogeny height difference and no
inferred node thickness is used.

\subsection{2. An elliptic cover and the exact
metrics}\label{an-elliptic-cover-and-the-exact-metrics}

Use \(\tau_F\) of AR-12, and \(Y=\operatorname{Im}\tau_F\). Set \[
T_0=
\begin{pmatrix}
-1&(-1+i)/2&(1+i)/2\\
-1&(1-i)/2&-(1+i)/2\\
0&0&-1-i
\end{pmatrix}.
\] The columns of \(T_0\) and \(iT_0\) belong to
\(\mathbf Z^3+\tau_F\mathbf Z^3\), so \(T_0\) gives an isogeny
\(E_i^3\to J\). Direct exact multiplication gives \[
T_0^*Y^{-1}T_0=2I_3,\qquad \det T_0=2,
\qquad\deg(T_0)=|\det T_0|^2/\det Y=8.
\] For an explicit lattice check, write \(T_0=A+\tau_F B\) and
\(iT_0=A'+\tau_F B'\), where \[
A=\begin{pmatrix}-1&-1&0\\-1&1&0\\0&0&0\end{pmatrix},
\quad
B=\begin{pmatrix}0&1&0\\0&-1&0\\0&1&-1\end{pmatrix},
\quad
A'=\begin{pmatrix}0&0&-1\\0&0&1\\0&0&2\end{pmatrix},
\quad
B'=\begin{pmatrix}-1&-1&0\\-1&1&0\\0&-1&-1\end{pmatrix}.
\] The theta divisor pulls back to a line with degree two in each
elliptic coordinate. Frobenius' theta function is a section of twice the
theta divisor, hence has degree four in each coordinate.

Let \(W=B^tT_0=\operatorname{diag}(0,-1+i,1+i)\). The functions \[
h_2(u)=e^{\pi i u^tWu}\theta(T_0u,\tau_F),\qquad
h_4(u)=e^{2\pi i u^tWu}\phi(T_0u,\tau_F)
\] are standard theta functions of degrees two and four in each elliptic
coordinate, respectively. Their transformations under
\(u_j\mapsto u_j+1\) and \(u_j\mapsto u_j+i\) verify this directly. The
metrics are exactly \[
\|\theta(T_0u)\|=(\det Y)^{1/4}
 e^{-2\pi|\operatorname{Im}u|^2}|h_2(u)|,\qquad
\|\phi(T_0u)\|=(\det Y)^4
 e^{-4\pi|\operatorname{Im}u|^2}|h_4(u)|.
\tag{AR.27}
\] The identity follows by expanding the Gaussian exponent and using
\(T_0^*Y^{-1}T_0=2I\). Pushforward by the isogeny preserves normalized
Haar measure, so there is no extra factor eight in any mean.

\subsection{3. A finite lower bound for the Frobenius-theta
mean}\label{a-finite-lower-bound-for-the-frobenius-theta-mean}

Let \(g\) be the mean-zero logarithmic theta potential on \(E_i\), with
a zero at the origin. Its distribution relation is \[
\sum_{t\in E_i[n]}g(z+t)=g(nz).
\tag{AR.28}
\] Indeed both sides have the same \(dd^c\), and both have Haar mean
zero. Its maximum is \[
M=\frac12\log2.
\] For completeness, translating the argument does not change a maximum
or mean, so for this calculation use
\(\log\|\theta_3(z,i)\|-\log|\eta(i)|\), whose zero is at \((1+i)/2\).
At fixed imaginary part, the triangle inequality bounds this theta norm
by its value on the imaginary axis. Poisson summation bounds that
Gaussian sum by its value at the origin. Its logarithmic mean is
\(\log|\eta(i)|\), as follows from Jacobi's product and Jensen's circle
formula. Finally \(\theta_3(0,i)/\eta(i)=\sqrt2\), using the standard
theta product and \(\theta_2(i)=\theta_4(i)=2^{-1/4}\theta_3(i)\).
Translating back gives the stated maximum for \(g\).

On an elliptic coordinate, the logarithmic norm of a section of degree
four is a constant plus four translates of \(g\), counted with
multiplicity. Equation (AR.28) therefore gives the one-sided bound \[
\int_{E_i}\log\|s\|
\ge\frac1{n^2}\sum_{t\in E_i[n]}\log\|s(u+t)\|
       -\frac{4M}{n^2}.
\] Iterate this inequality in all three coordinates of the elliptic
cover: \[
K:=\int_J\log\|\phi\|
\ge\frac1{n^6}\sum_{t\in E_i[n]^3}
           \log\|\phi(T_0(u_0+t))\|
       -\frac{6\log2}{n^2}.
\tag{AR.29}
\] The chosen translated grid has certified nonzero values. Zeros of a
restricted section cause no missing positive term in this lower bound.

\subsection{4. A finite upper bound for the ordinary-theta
mean}\label{a-finite-upper-bound-for-the-ordinary-theta-mean}

Partition the real six-dimensional torus \(E_i^3\) into \(n^6\) equal
coordinate boxes \(R\). Jensen's inequality on each box gives \[
H:=\int_J\log\|\theta\|
\le\frac1{2n^6}\sum_R
        \log\left(n^6\int_R\|\theta(T_0u)\|^2\,du\right).
\tag{AR.30}
\] Only upper bounds on these logarithms are needed; zeros do not
require excision.

For \(D=2,4\) use the normalized one-dimensional basis \[
b_{D,r}(x+iy)=
\sum_{k\equiv r\pmod D}
 e^{-\pi(k+Dy)^2/D}e^{2\pi ikx},
\qquad 0\le r<D.
\] Equations (AR.27) express the norms by tensors with \(2^3\) and
\(4^3\) coefficients. For a rectangle \(R_x\times R_y\) on \(E_i\), the
degree-two Gram entry is the absolutely convergent sum \[
\operatorname{av}_{R_x\times R_y}
 b_{2,r}\overline{b_{2,s}}
=\sum_{\substack{k\equiv r(2)\\l\equiv s(2)}}
 e^{-\pi(k-l)^2/4}
 \operatorname{av}_{R_x}e^{2\pi i(k-l)x}
 \operatorname{av}_{R_y}e^{-4\pi(y+(k+l)/4)^2}.
\tag{AR.31}
\] The first average is an exponential times a sinc; the second is a
difference of erf values. Thus every six-dimensional cell mean in
(AR.30) is a finite tensor contraction of rigorously enclosed
one-dimensional integrals.

\subsection{5. Reproducible finite certificate and truncation
bounds}\label{reproducible-finite-certificate-and-truncation-bounds}

The calculation is in fermat\_theta\_certify.py, with output
fermat\_theta\_certify.json. It uses Arb at 160-bit precision for
coefficients, theta values, exponential/erf/sinc factors, and final
logarithms. The tensor contractions use binary64 arithmetic with an
explicit error enclosure described below.

The Frobenius coefficients first come from the one-dimensional space
\(\Gamma_{00}\): in the eight second-order theta basis, impose the value
and six Hessian conditions at zero. A certified invertible
seven-by-seven subsystem fixes the last coefficient to one. Normalize
its absolute scale using de Jong 2020, Theorem 4.1, at the two-torsion
point \(z=(1/2,0,0)\): the absolute value is the product of the sixteen
even thetanulls whose characteristic becomes odd after this translation.
The remaining common unit phase has no effect.

The normalization value itself includes the additional tail ball
\(10^{-45}\sum_p|c_p|\), independently of the coefficient intervals.
This explicit tail was added following h9's implementation review; the
final certificate was rerun after that correction.

The smallest eigenvalue of \(Y\) is \(1-1/\sqrt2>1/4\). Write
\(q=k^tYk\). For the second-order series retain \(q\le100\); the omitted
value and Hessian sums are bounded by \(10^{-45}\). One convenient bound
is to split off \(e^{-0.4\pi q}\) and use \[
\sum_{k\in\mathbf Z^3}e^{-0.1\pi k^tYk}
\le(1+\sqrt{40})^3.
\] The Hessian weights satisfy \(|k_ik_j|\le4q\). For Frobenius
evaluations at \(T_0(j/4)\), \(0\le j_i<4\), the imaginary part has
squared \(Y^{-1}\)-norm at most \(9/8\). The tail is therefore bounded
using \[
-\tfrac12q+2\sqrt{9/8}\sqrt q
\le-0.1q-40+30/\sqrt2\quad(q\ge100).
\] After multiplication by \(\pi\) and summation, this is less than
\(10^{-22}\) times the largest coefficient. The script verifies that
every normalized coefficient has absolute value below one and inserts
the larger radius \(10^{-20}\).

For ordinary thetanulls retain \((n+a')^tY(n+a')\le50\). The same
Gaussian estimate, now splitting off \(e^{-0.8\pi q}\), bounds the
discarded part by less than \(10^{-45}\). The integer enumeration
contains every retained term: for \(u=2n+2a'\), \(u^tYu\le200\) gives
\(|u_{1,2}|\le\sqrt{300}\), \(|u_3|\le20\), using the diagonal of
\(Y^{-1}\). Ordinary-theta evaluations on \(T_0(j/2)\) have imaginary
squared norm at most \(1/2\); the estimate \(-q+\sqrt{2q}\le-0.2q-30\)
for \(q\ge50\) bounds the tail by
\((1+\sqrt{20})^3e^{-30\pi}<10^{-30}\). The product-basis coefficients
are obtained by exact finite Fourier inversion from these \(D^3\)
values. Their scalar divisors are convergent positive one-dimensional
theta sums.

In the normalized one-dimensional bases, the script retains the
indicated nine periods of Fourier indices in each direction, adding
radius \(10^{-100}\). In the degree-two Gram formula it retains
\(|k|,|l|\le12\), adding radius \(10^{-80}\). On \(|y|\le1/2\), the
omitted one-dimensional terms are bounded by Gaussian geometric tails;
these chosen radii are larger than those tails. All enumeration and
phase exponents use integers or exact binary fractions.

For clarity about the finite floating calculations, let \(C\) be a
coefficient tensor and \(B_j\) the three basis arrays. Arb conversion
supplies complex centers and bounds on the sum of absolute real and
imaginary errors. Put \(S=\sum|C|_1\), \(R_C=\sum r_C\),
\(B_j=\max|B_j|_1\), and \(R_j=\max r_{B_j}\). The uniform complex error
used for each contraction is \[
R_C\prod(B_j+R_j)
 +S\left(\prod(B_j+R_j)-\prod B_j\right)
 +\gamma_{4096}S\prod B_j+2^{-1000},
\quad
\gamma_{4096}=\frac{4096\,2^{-53}}{1-4096\,2^{-53}}.
\tag{AR.32}
\] Each contracted entry is evaluated using fewer than 4096 elementary
real arithmetic operations along the expanded expression. The standard
rounding-product bound gives (AR.32); the extra absolute term covers
underflow. The computation has no overflow. Square, sum and square-root
operations used for lower moduli are rounded outward with nextafter.

Final mean logarithms do not rely on the binary64 logarithm. Each
enclosed positive input is binned using its exact binary exponent and a
12-bit dyadic mantissa, downward for lower logarithms and upward for
upper logarithms. Histogram counts are exact integers; Arb evaluates the
resulting integer logarithms and their weighted sum. This deliberately
loses less than \(2^{-11}\) in each mean logarithm.

With \(n=16\), the certificate encloses 16,777,216 grid values and the
same number of cell means. The smallest lower modulus is greater than
\(5.8\cdot10^{-5}\), while the uniform tensor error is below
\(9.3\cdot10^{-12}\). For the cell means the smallest upper input
exceeds \(0.0057\), and the tensor error is below \(1.1\cdot10^{-11}\).
The resulting one-sided bounds are \[
\begin{aligned}
K&>-4.4783736734,\\
H&<-0.6702188747,\\
I=K-2H+\tfrac72\log2&>-0.7119207919,\\
A_\infty&>-0.0159882720.
\end{aligned}
\tag{AR.33}
\] Combining (AR.26), (AR.33), and
\(h_K(\eta,\eta)=\frac72(B_{\rm fin}+A_\infty)\) yields \[
h_K(\eta,\eta)>0.0114303576>\frac1{100}.
\] This proves (AR.22), subject only to the cited height formulas and
the explicitly described finite arithmetic certificate. It is not a
proof-assistant certificate.

\subsection{Reading and independent
checks}\label{reading-and-independent-checks}

The primary sources and exact metric locators are listed in
FERMAT\_THETA\_GATE.md. The finite-place proof is
work/programme/fermat\_hodge\_lattice.tex. Root and h9 independently
checked the AR-12 constants. Agent h9 completed an independent
whole-proof and implementation review: exact lattice/gauge identities,
torsion distribution and maximum, degrees, absence of an isogeny-degree
multiplier, cell Jensen, Gram/DFT phases, Gaussian tails, error
propagation and dyadic logarithms passed after the explicit
normalization-tail fix. The final self-contained certifier was rerun
successfully. This review uses the stated IEEE binary64 arithmetic
assumptions and is not a proof-assistant certificate.
% JAC1964-20261121. /root/programme_audit; GPT-6 Astra Pro.
% Effort setting not exposed. New deductions with their exact scope.

\subsection*{JAC-P11. A positive CM orbit beyond the Picard summand}

\begin{proposition}[JAC-P11]\label{jac:P11}
Let $J$ be the Fermat-quartic Jacobian with symmetric principal theta
class $\theta$, and let $\eta=\mathcal F([C]_{(1)})$ be the
hyperflex-normalized class of AR-8. After a finite extension of the
number field $K$, there is a rational correspondence $T$ on
$\operatorname{CH}^2_{(1)}(J)_{\mathbf Q}$ such that
\[
 T^2=-1,\qquad h_K(Tz,Tw)=h_K(z,w),\qquad
 T\eta\in\operatorname{CH}^2(N).                       \tag{JAC.P11.1}
\]
Here $N$ is the same orthogonal complementary summand chosen in AR-8.
The classes $\eta,T\eta$ are independent even modulo algebraic
equivalence. For $r,s\in\mathbf Q$, not both zero, one has
\[
 h_K(r\eta+sT\eta,r\eta+sT\eta)
   =(r^2+s^2)h_K(\eta,\eta)>\frac{r^2+s^2}{100},        \tag{JAC.P11.2}
\]
where the final bound uses AR-13. Hence the height is strictly
positive on
\[
 \operatorname{CH}^2(M)(K)_{\mathbf Q}
       \oplus\mathbf Q\eta\oplus\mathbf Q T\eta.        \tag{JAC.P11.3}
\]
The last two summands form a one-dimensional $\mathbf Q(i)$-module;
no assertion that they exhaust $\operatorname{CH}^2(N)$ is made.
\end{proposition}

\begin{proof}
Bloch's isogeny gives $J\sim E^3$ geometrically, with $j(E)=1728$.
Thus the centre of $\operatorname{End}^0(J)$ is $\mathbf Q(i)$.
Choose its element $j$ with $j^2=-1$, and enlarge $K$ so it is
defined. Rosati preserves the centre and sends $j$ to $\pm j$.
If $j^\dagger=j$, then
$\operatorname{Tr}(j j^\dagger)=\operatorname{Tr}(-1)=-6$,
contrary to Rosati positivity. Hence $j^\dagger=-j$.
The precise background is Milne, \emph{Abelian Varieties}, v2.0,
Proposition~14.2 and Theorem~14.3, printed pp.~61--62.

Choose a positive integer $n$ such that $a=nj$ is an endomorphism.
It is an isogeny, and
\[
 a^2=[-n^2],\qquad a^\dagger a=[n^2],\qquad
 \deg(a)=n^6,\qquad a^*\theta=n^2\theta
     \text{ in }\operatorname{CH}^1(J)_{\mathbf Q}.     \tag{JAC.P11.4}
\]
For the last assertion the polarization identity first gives
equality in $\operatorname{NS}(J)_{\mathbf Q}$. The difference is
symmetric and algebraically trivial, hence two-torsion, so the
equality holds in the rational Chow group.
Set
\[
 T=n^{-3}a^*\bigm|_{\operatorname{CH}^2_{(1)}(J)_{\mathbf Q}}.
                                                        \tag{JAC.P11.5}
\]
Replacing $(n,a)$ by $(kn,[k]a)$ multiplies this pullback by
$k^{2\cdot2-1}=k^3$, so the definition is independent of clearing
denominators. Since $[-n^2]^*$ acts on this grade by
$(-n^2)^3$, we get $T^2=-1$ exactly in rational Chow groups.
The height projection formula gives
\[
 h_K(a^*z,a^*w)=h_K(z,a_*a^*w)=\deg(a)h_K(z,w)
              =n^6h_K(z,w).
\]
Thus $T$ is an isometry. It preserves primitiveness because
$\theta a^*z=n^{-2}a^*(\theta z)$.

Let $G=\mu_4^2$ and $p_G$ be the averaging projector of AR-8.
Centrality of $a$ makes $T$ commute with $G$, so $T\eta$ is
$G$-invariant. AR-8 constructs its Picard projector $e'$ with
the exact relation $e'p_G=0$. Therefore $e'T\eta=0$; since $T\eta$
is primitive, it belongs to that same $N$.
Symmetry and isometry give
\[
 h_K(\eta,T\eta)=h_K(T\eta,T^2\eta)
               =-h_K(T\eta,\eta)=-h_K(\eta,T\eta).
\]
The mixed term is zero, and both diagonal terms equal
$h_K(\eta,\eta)>1/100$ by AR-13.

Independence modulo algebraic equivalence can also be proved
without using the height. If $r\eta+sT\eta$ is algebraically
trivial, apply $r-sT$. Algebraic correspondences preserve
algebraic equivalence, so $(r^2+s^2)\eta$ is algebraically trivial.
Bloch's infinite-order theorem, transported by Fourier as in AR-8,
forces $r=s=0$. Finally AR-2 gives positivity on
$\operatorname{CH}^2(M)$ and height orthogonality to $N$,
which proves (JAC.P11.3).
\end{proof}

\subsection*{JAC-P12. The exact scope of tautological generation}

\begin{proposition}[JAC-P12]\label{jac:P12}
Let $R(J)$ be the usual tautological ring of the hyperflex-embedded
Fermat quartic: the smallest intersection subring of
$\operatorname{CH}^*(J)_{\mathbf Q}$ containing $[C]$ and stable
under Fourier and all $[n]^*$. Then
\[
 R^2(J)=\mathbf Q\theta^2\oplus\mathbf Q\eta,\qquad
 R^2(J)_{\mathrm{hom}}=\mathbf Q\eta.                  \tag{JAC.P12.1}
\]
The class $T\eta$ from JAC-P11 does not belong to $R^2(J)$ even
after adding an algebraically trivial class. Thus the ordinary
tautological generation theorem cannot identify the full primitive
grade-one Chow group with the Fermat Ceresa line.
\end{proposition}

\begin{proof}
Polishchuk, \emph{Lie symmetries of the Chow group of a Jacobian
and the tautological subring}, arXiv:math/0504448v3,
Theorem~0.2(i), p.~2, proves that $R(J)$ is generated by
$p_n=\mathcal F([C]_{(n-1)})$ and
$q_n=\mathcal F(\theta[C]_{(n)})$.
His Remark~2 following Corollary~0.3, p.~3, says that
$(2g-2)P_0=K_C$ implies $q_n=0$ for all $n\ge1$.
The hypothesis holds for our hyperflex. On p.~4 he gives
$p_n\in\operatorname{CH}^n_{(n-1)}$ and $p_1=-\theta$.
Consequently the only codimension-two monomials in the remaining
generators are $p_1^2$ and $p_2=\eta$ (up to sign if the inverse
Fourier kernel is used). The class $\theta^2$ has nonzero
cohomology, while $\eta$ is homologically trivial and nonzero.
This proves (JAC.P12.1).
If $T\eta$ differed from a tautological class by an algebraically
trivial class, its $\theta^2$ coefficient would vanish by cohomology.
This would contradict the independence of $\eta,T\eta$ modulo
algebraic equivalence proved in JAC-P11.
\end{proof}

\paragraph{The remaining generation gate.}
The two-dimensional Betti realization of $N$ does not bound the
dimension of $\operatorname{CH}^2(N)$: every class in that Chow
group already has zero ordinary cycle class.
Finite-dimensionality of abelian Chow motives supplies nilpotence
for homologically trivial endomorphisms, not a bound on this Chow
group. The CM operation above produces a nonzero positive
one-dimensional $\mathbf Q(i)$-module, but does not prove that
the entire complementary Chow group has that dimension.
A generation theorem must also control a potential Abel--Jacobi
kernel. For comparison, Lilienfeldt--Shnidman,
\emph{Experiments with Ceresa classes of cyclic Fermat quotients},
author manuscript dated 29~November~2021, \S4.1, pp.~5--6,
labels the equality of Griffiths rank and the vanishing order of
the corresponding rank-two motivic $L$-function as Bloch's
conjectural \emph{Son of Recurring Fantasy}.
It is not a Chow-generation theorem.

\paragraph{Primary sources directly read.}
\url{https://www.jmilne.org/math/CourseNotes/AV.pdf},
\url{https://arxiv.org/pdf/math/0504448}, and
\url{https://math.huji.ac.il/~shnidman/Experiments_with_Ceresa_classes.pdf}.
The new deductions are the normalized CM isometry and positive
Griffiths plane in JAC-P11, and the explicit specialization of
the tautological generators in JAC-P12. The cited generation
theorem and Rosati positivity are used as established background.

% JAC1964-20261121. /root/programme_audit; GPT-6 Astra Pro.
% Additional established input: Kimura finite-dimensionality/conservativity.

\subsection*{JAC-P13. The canonical invariant complement}

\begin{proposition}[JAC-P13]\label{jac:P13}
Retain the Fermat notation of AR-8 and JAC-P11. Put
$P=(h^3(J),\pi_P)$, and let $e'$ be AR-8's self-transpose
Picard-summand projector with $e'p_G=p_Ge'=0$.
Then, as actual rational Chow projectors on $h^3(J)$,
\[
 \pi_P-e'=p_G\pi_P=p_G.                               \tag{JAC.P13.1}
\]
In particular the complementary motive is exactly
$N=h^3(J)^G$, independent of the auxiliary inclusion used in AR-2,
and
\[
 \operatorname{CH}^2(N)
       =\operatorname{CH}^2_{(1)}(J)_{\mathbf Q}^{G}.    \tag{JAC.P13.2}
\]
The operation $T$ of JAC-P11 consequently preserves all of
$\operatorname{CH}^2(N)$ and makes it a $\mathbf Q(i)$-module.
Its height is the real part of the complex-valued Hermitian form
\[
 \mathscr H(z,w)=h_K(z,w)+i\,h_K(z,Tw),                 \tag{JAC.P13.3}
\]
where $i$ acts through $T$ and the form is linear in its first
variable. Positivity of $\mathscr H$ on the whole module remains
the full complementary height problem.
\end{proposition}

\begin{proof}
We explicitly add the established Kimura input: a direct summand
of an abelian Chow motive is finite-dimensional; a finite-dimensional
Chow motive with zero Weil realization is zero.
These are Kimura, \emph{Chow groups are finite dimensional, in some
sense}, Math.\ Ann.\ 331 (2005), Example~9.1 and Corollary~7.3.
The exact application is also recalled and used in the directly
read proof of Theorem~A of Laga--Shnidman,
\emph{Vanishing criteria for Ceresa cycles}, author PDF p.~8.
The original Kimura article was not directly read in this run;
the cited standard consequence is carried through that primary
research application.

Set $q=\pi_P-e'-p_G\pi_P$. All three relevant projectors commute:
$e'$ is supported on $P$, $p_G$ commutes with $\pi_P$, and
$e'p_G=p_Ge'=0$. Therefore $q^2=q$.
AR-8's explicit character calculation shows that $p_GH^3(J)$ is
exactly the two-dimensional sum of the extreme Hodge types.
AR-2 identifies $H(e')$ with the remaining maximal level-one
part of $H(P)$. Hence $H(q)=0$.
The motive defined by $q$ is a direct summand of $h(J)$.
Kimura's consequence forces it, and thus the Chow projector $q$,
to be zero. This proves the first equality in (JAC.P13.1).
Also $p_G(1-\pi_P)$ is an idempotent on $h^3(J)$ with zero
realization, because the two extreme types are primitive.
The same argument gives $p_G=p_G\pi_P$.

These arguments can be applied after passing to an algebraic
closure. The resulting rational-equivalence identities descend:
a rational equivalence over the algebraic closure is defined over
a finite extension, and pushforward multiplies its descended class
by the extension degree. Thus no descent ambiguity is hidden here.
The Chow--Kunneth grading gives (JAC.P13.2).
Centrality makes $T$ commute with $p_G$, and JAC-P11 gives
$T^2=-1$ and height isometry. Therefore
$h_K(Tz,w)=-h_K(z,Tw)$.
This verifies conjugate symmetry and
$\mathscr H(Tz,w)=i\mathscr H(z,w)$,
$\mathscr H(z,Tw)=-i\mathscr H(z,w)$.
Finally $\mathscr H(z,z)=h_K(z,z)$, so their positivity questions
are identical. Nothing here bounds the dimension of this
$\mathbf Q(i)$-module.
\end{proof}


\subsection*{JAC-P14. The entire abelian-endomorphism orbit}

\begin{proposition}[JAC-P14]\label{jac:P14}
Choose an isogeny $J\sim E^3$ over $K$, with
$\operatorname{End}^0(E)=F=\mathbf Q(i)$, identifying
$\operatorname{End}^0(J)$ with $M_3(F)$. Choose the sign of $j$
in JAC-P11 so that it corresponds to the scalar matrix $iI_3$.
For $b\in\operatorname{End}^0(J)$ choose $d>0$ with $db$ an
actual endomorphism and define its normalized middle pullback by
\[
 \rho_3(b)=d^{-3}(db)^*\bigm|_{h^3(J)}.
\]
Then $\rho_3(b)$ preserves the exact Chow summand $N$ of JAC-P13.
If $\det_F(b)=u+vi$, with $u,v\in\mathbf Q$, its action is
\[
 \rho_3(b)\bigm|_N=u\,1_N-v\,T.                       \tag{JAC.P14.1}
\]
In particular
\[
 \operatorname{span}_{\mathbf Q}
   \{\rho_3(b)\eta:b\in\operatorname{End}^0(J)\}
        =\mathbf Q\eta\oplus\mathbf Q T\eta.           \tag{JAC.P14.2}
\]
The equality is at the level of rational Chow groups, not merely
after realization. For any $b,c\in\operatorname{End}^0(J)$,
\[
 h_K(\rho_3(b)\eta,\rho_3(c)\eta)
   =\operatorname{Re}
       \bigl(\det_F(b)\overline{\det_F(c)}\bigr)
      h_K(\eta,\eta).                                \tag{JAC.P14.3}
\]
Thus the full matrix endomorphism algebra does not enlarge the
positive CM plane.
\end{proposition}

\begin{proof}
We use the functorial Chow-motive power formula
$h^r(A)=\operatorname{Sym}^r h^1(A)$, with the usual commutativity
constraint on Chow motives. Its realization is the exterior power
of the underlying degree-one cohomology, because realization is
graded-commutative. The exact original reference is Kunnemann,
\emph{On the Chow motive of an abelian scheme}, in \emph{Motives},
Proc.\ Sympos.\ Pure Math.\ 55(1) (1994), 189--205,
Theorem~3.3.1. We directly checked the formulation in
Ancona--Enright-Ward--Huber, \emph{On the motive of a commutative
algebraic group}, Documenta Math.\ 20 (2015), 807--858,
Theorem~4.2.3(3), printed p.~825: the canonical isomorphism respects
the grading. Their Remark~4.2.1, printed p.~824, explains both
transposition to the contravariant convention used here and
Kunnemann's exterior-power notation for this symmetric-power functor.
The scalar normalization in $\rho_3$ is exactly the homogeneous
degree-three normalization in this functor; it is independent of $d$.

Extend the coefficient field of motives to $F$. The CM idempotents
split
\[
 h^1(E)_F=V_+\oplus V_-,
 \qquad j|_{V_+}=i,\quad j|_{V_-}=-i.
\]
Each $V_\pm$ has one-dimensional odd realization.
The motive $\operatorname{Sym}^2 V_\pm$ has zero realization and
is a summand of a tensor power of an abelian motive. By the
Kimura consequence explicitly invoked in JAC-P13 it is zero.
This consequence also holds after the coefficient extension to $F$:
an $F$-coefficient summand is a rational summand after restriction
of scalars, so its finite-dimensionality and the realization-zero
criterion follow from the rational statement.
In particular the interchange map on $V_\pm\otimes V_\pm$ is
exactly $-1$ in the category of Chow motives.

Now
\[
 h^1(J)_F\simeq V_+^{\oplus3}\oplus V_-^{\oplus3}.
\]
In its third symmetric power consider the two extreme summands
\[
 D_+=\operatorname{Sym}^3(V_+^{\oplus3}),\qquad
 D_-=\operatorname{Sym}^3(V_-^{\oplus3}).               \tag{JAC.P14.4}
\]
Since $\operatorname{Sym}^2V_\pm=0$, the direct-sum formula for
symmetric powers identifies $D_\pm$ with $V_\pm^{\otimes3}$:
the only surviving multiplicity pattern is $(1,1,1)$.
The interchange sign just proved shows that a matrix
$b\in M_3(F)$ acts on $D_+$ by $\det_F(b)$ and on $D_-$ by
$\overline{\det_F(b)}$. Indeed expansion of the action on its
three tensor factors is precisely the alternating permutation sum
defining the determinant. This calculation is in Chow motives.

It remains to identify these summands with the particular $N$,
rather than only with a motive having the same realization.
The sum of the projectors onto $D_+$ and $D_-$ is defined over
$\mathbf Q$. For an explicit description, apply $\rho_3(1+j)$.
Its eigenvalues on the two extreme summands are $-2\pm2i$,
while on the two mixed summands they are $2\pm2i$.
The polynomials
\[
 X^2+4X+8,\qquad X^2-4X+8
\]
are coprime over $\mathbf Q$, so their Chinese-remainder
idempotent defines an exact rational Chow projector $p_{\rm ext}$.
The preceding tensor decomposition verifies the polynomial
identity needed for this idempotent exactly, not only in cohomology.
Because $1+j$ is central, $p_{\rm ext}$ commutes with all
$\rho_3(b)$, and in particular with $p_G$.
The two commuting idempotents $p_{\rm ext}$ and $p_G$ have
identical realizations. Therefore
$p_{\rm ext}(1-p_G)$ and $p_G(1-p_{\rm ext})$ are
realization-zero idempotents. Kimura makes both zero.
JAC-P13 now gives $N=(h^3(J),p_G)=(h^3(J),p_{\rm ext})$.

On $D_+$ the operator $T=\rho_3(j)$ is multiplication by
$i^3=-i$, and on $D_-$ it is multiplication by $(-i)^3=i$.
Consequently $u-vT$ has eigenactions $u+vi$ and $u-vi$,
exactly the determinant actions computed above. This proves
(JAC.P14.1) after extension of coefficients to $F$ and hence over
$\mathbf Q$, since extension of a rational coefficient vector
space to $F$ is injective.
The identity endomorphism and $j$ yield $\eta$ and $T\eta$,
while (JAC.P14.1) puts every other image in their span.
This proves (JAC.P14.2). Finally JAC-P11 gives the diagonal
height matrix with equal entries and zero mixed entry, which
implies (JAC.P14.3).
\end{proof}

\paragraph{Interpretation of JAC-P13--P14.}
The remaining question is positivity on an actual
$\mathbf Q(i)$-module of primitive cycles in the exact motive
$h^3(J)^G$. The positive vector $\eta$ generates precisely the
CM plane under all abelian endomorphisms. Further progress on
the entire group therefore needs information on additional
$\mathbf Q(i)$-independent cycles, or a theorem proving there
are none; applying more abelian endomorphisms to $\eta$ cannot
supply that information.


\section{The full complementary Chow group: a rigorous further
reduction}\label{the-full-complementary-chow-group-a-rigorous-further-reduction}

Run \texttt{JAC1964-20261121}. Author/checker:
\path{/root/arithmetic_sign}, GPT-6 Astra Pro; effort setting not
exposed. Date: 2026-11-21 UTC. The positive Fermat height and the
positive CM orbit below are established inputs from AR-13 and programme
JAC-P11. This section does not identify a Selmer rank with a Chow rank.

\subsection{1. Setup and the positive subspace already
obtained}\label{setup-and-the-positive-subspace-already-obtained}

Let \(J\) be the Fermat-quartic Jacobian and use the normalized real
Beilinson--Bloch pairing \(h_F\) of AR-1--AR-13. After a finite
number-field extension containing the required correspondences,
programme JAC-P13 identifies the complementary motive exactly as \[
N=h^3(J)^G,\qquad G=\mu_4^2,
\quad V_F=\operatorname{CH}^2(N_F)
       =\operatorname{CH}^2_{(1)}(J_F)_{\mathbf Q}^{G}.
\tag{AR15.1}
\] The projector is self-transpose. Its Betti realization has only types
\((3,0)\) and \((0,3)\), each once. Every class in \(V_F\) is
homologically trivial.

JAC-P11 and JAC-P13 give an actual rational correspondence \(T\) on this
whole Chow group with \[
T^2=-1,
\qquad h_F(Tz,Tw)=h_F(z,w).
\] Thus \(V_F\) is a \(\mathbf Q(i)\)-module. The complex-valued
Hermitian form, linear in its first variable, is \[
\mathscr H_F(z,w)=h_F(z,w)+i h_F(z,Tw).
\] For the hyperflex-normalized Fourier--Ceresa class \(\eta\), \[
U_F=\mathbf Q\eta\oplus\mathbf Q T\eta,
\qquad h_F(r\eta+sT\eta,r\eta+sT\eta)
  =(r^2+s^2)h_F(\eta,\eta)>
  \frac{r^2+s^2}{100}\quad\bigl((r,s)\ne(0,0)\bigr).
\tag{AR15.2}
\] The two classes are independent even modulo algebraic equivalence.
They form one positive \(\mathbf Q(i)\)-line. The following reduction
concerns all of \(V_F\), without assuming that this line generates it.

\subsection{2. An unconditional radical and an exact Griffiths-quotient
reduction}\label{an-unconditional-radical-and-an-exact-griffiths-quotient-reduction}

\textbf{Theorem AR-15 (algebraically trivial complementary cycles have
zero height against the whole complement).} In the general setup of
AR-2, let \(A/F\) have dimension \(2m-1\), and let
\(N\subset h^{2m-1}(A)\) be its self-transpose complementary Chow
summand. Assume that \(H(N)\) has no nonzero sub-Hodge structure of
types \((m,m-1),(m-1,m)\). Put \[
V_F=\operatorname{CH}^m(N_F),\qquad
\mathcal A_F=V_F\cap\operatorname{CH}^m(A_F)_{\rm alg}.
\] Here algebraic triviality may be interpreted geometrically. Then \[
(-1)^m h_F(a,y)=0
\qquad(a\in\mathcal A_F,\;y\in V_F).
\tag{AR15.3}
\] Consequently the signed height descends to \[
\operatorname{Griff}^m(N_F):=V_F/\mathcal A_F.
\tag{AR15.4}
\] Semipositivity on \(V_F\) is equivalent to semipositivity on this
quotient. Strict positivity on \(V_F\) is equivalent to \textbf{both}
\(\mathcal A_F=0\) and strict positivity on the quotient.

\emph{Proof.} It is enough to work after a finite extension of \(F\):
normalized heights are unchanged, and restriction on rational Chow
groups is injective by the restriction--corestriction identity. Every
algebraically trivial class is a finite sum of classes \[
\Gamma_{0*}D,
\qquad D\in\operatorname{Pic}^0(C)(F)_{\mathbf Q},
\quad\Gamma_0\in\operatorname{CH}^m(C\times A)_{\mathbf Q},
\] where \(C\) is a smooth projective geometrically connected curve. A
finite field extension suffices to define the finitely many curves,
correspondences, and divisors occurring in such an expression. Apply the
projector \(\pi_N\) and put \(\Gamma=\pi_N\circ\Gamma_0\). Since
\(a\in V_F\), these projected expressions still sum to \(a\).

Fix \(y\in V_F\). Its ordinary cycle class is zero because it belongs to
the odd-degree Chow--Künneth component. By codimension bookkeeping, \[
{}^t\Gamma_*y\in\operatorname{CH}^1(C)_{\rm hom,\mathbf Q}
                   =\operatorname{Pic}^0(C)(F)_{\mathbf Q}.
\] After a complex embedding, the cohomological map of \({}^t\Gamma\)
factors through \[
H(N)(m-1)\longrightarrow H^1(C).
\] This morphism is zero. Indeed its transpose is a morphism from
\(H^1(C)\) into \(H(N)(m-1)\), whose image would have only types
\((1,0),(0,1)\); the defining hypothesis on \(N\) excludes such an
image. Equivalently use duality and semisimplicity of polarizable Hodge
structures. Self-transposition of \(\pi_N\) is what places the
transposed correspondence on the same summand.

Functoriality of the complex Abel--Jacobi map therefore gives \[
\operatorname{AJ}_C({}^t\Gamma_*y)=0.
\] For divisors on a curve, Abel's theorem identifies this map with the
injective identification of degree-zero divisor classes with their
Picard class. Hence \({}^t\Gamma_*y=0\) in
\(\operatorname{CH}^1(C_{\mathbf C})_{\mathbf Q}\). This detects zero
already over the number field: the map on Picard points is injective
after field extension, or one can spread the trivialization to a finite
extension and apply the norm. The projection formula for the arithmetic
height gives \[
h_F(\Gamma_*D,y)=h_F(D,{}^t\Gamma_*y)=0.
\] Sum over the finitely many expressions for \(a\). This proves
(AR15.3). A bilinear form whose radical contains \(\mathcal A_F\)
descends to \(V_F/\mathcal A_F\), and the stated positivity equivalences
follow directly. \(\square\)

The background functoriality is recorded explicitly in
Achter--Casalaina-Martin--Vial, \emph{Decomposition of the diagonal,
intermediate Jacobians, and universal codimension-2 cycles in positive
characteristic}, Astérisque 455 (2025), §3.2.4, equation (3.5),
pp.17--18. The arithmetic projection formula is AR-B1. No injectivity of
a higher Abel--Jacobi map is used: the only injectivity is the classical
divisor statement on the auxiliary curve.

For the Fermat motive, AR-15 yields a precise possible zero-direction
mechanism. \textbf{Any nonzero class of \(\mathcal A_F\) would have zero
height and would disprove strict positivity on the full rational Chow
group.} No such nonzero class has been produced here. Its vanishing is
not supplied by the positive value of \(h_F(\eta,\eta)\). The positive
plane (AR15.2) injects into \(\operatorname{Griff}^2(N_F)\), so it
survives this further reduction.

\subsection{3. What would make a Selmer-rank bound
sufficient}\label{what-would-make-a-selmer-rank-bound-sufficient}

Let \(p\) be a rational prime, put \[
W_p=H^3_{\rm et}(N_{\overline F},\mathbf Q_p)(2),
\qquad S_p(F)=H^1_f(F,W_p),
\] and let \[
\alpha_{p,F}:V_F\otimes_{\mathbf Q}\mathbf Q_p
             \longrightarrow S_p(F)
\tag{AR16.1}
\] be the scalar-extended étale Abel--Jacobi map. Its finite local
conditions are valid for the present potentially good-reduction setting.
A primary locator is Nekovář, \emph{p-adic Abel-Jacobi maps and p-adic
heights}, Proposition 2.5, Theorem 3.1(i), and §4.1; CM abelian
varieties acquire good reduction locally after finite extension. Neither
existence of this map nor those local conditions assert injectivity.

\textbf{Proposition AR-16 (a sufficient finite-rank criterion, with its
missing hypothesis exposed).} Fix a number field \(F\) over which
\(N,T,\eta\) are defined. If \[
\dim_{\mathbf Q_p}S_p(F)\le2
\quad\hbox{and}\quad
\ker\alpha_{p,F}=0,
\tag{AR16.2}
\] then \[
V_F=U_F.
\tag{AR16.3}
\] In particular the full complementary height over this \(F\) is
strictly positive, and the complete primitive height is strictly
positive after adjoining the orthogonal Picard block of AR-2.

\emph{Proof.} Extension of scalars preserves linear independence of
rational vectors. The injective map (AR16.1) and the dimension bound in
(AR16.2) imply \(\dim_{\mathbf Q}V_F\le2\), even if
finite-dimensionality had not previously been known. Equation (AR15.2)
already provides two independent vectors in \(V_F\). They therefore form
a basis. Their height matrix is \(h_F(\eta,\eta)I_2\), which is strictly
positive by AR-13. Orthogonality to the Picard block completes the
assertion. \(\square\)

The scalar extension in (AR16.1) is essential. An injective merely
\(\mathbf Q\)-linear map \(V_F\to S_p(F)\) does not bound
\(\dim_{\mathbf Q}V_F\) by \(\dim_{\mathbf Q_p}S_p(F)\): for
\(a\in\mathbf Q_p\setminus\mathbf Q\), the map \((r,s)\mapsto r+as\)
injects \(\mathbf Q^2\) into a one-dimensional \(\mathbf Q_p\)-space.
After scalar extension this map has a nonzero kernel. Thus even an
injectivity claim must be read with its precise coefficients.

A weaker exact criterion is sometimes useful. If a \textbf{rational}
subspace \(R_F\subset\operatorname{rad}(h_F|_{V_F})\) is proved to
satisfy \[
V_F=U_F+R_F,
\tag{AR16.4}
\] then the sum is direct, the height is semipositive, and it is
strictly positive exactly when \(R_F=0\). This follows because a
positive subspace meets the radical trivially. A Selmer-rank bound alone
proves neither the rational decomposition (AR16.4) nor inclusion of an
Abel--Jacobi kernel in the radical of the \textbf{real} pairing.

All these statements are over the \textbf{specified field \(F\)}. A
rank-one theorem over \(\mathbf Q\) cannot simply be applied after
adjoining \(i\), nor after adjoining fields of definition of arbitrary
cycles. Over \(\mathbf Q(i)\), the rank-two Galois representation splits
into two conjugate characters. Its Selmer group is the direct sum of
their Selmer groups, because the Bloch--Kato local conditions respect
direct sums. Conjugation identifies the dimensions of the two character
factors. A rank-one bound for each character therefore gives a
dimension-two bound over \(\mathbf Q(i)\), before any further extension.
It gives no uniform dimension-two bound over every finite extension or
over the direct limit defining \(V_{\overline{\mathbf Q}}\).

\subsection{4. Exact current scope of the relevant CM
theorems}\label{exact-current-scope-of-the-relevant-cm-theorems}

\textbf{Audit AR-17.} The following are primary-source scope checks. The
initially open Katz calculation and conductor-specific rank implication
are now supplied by AR-19 and H9-T15. They give the character Selmer
rank one, hence rank two for the full Fermat representation over
\(K_0=\mathbf Q(i)\). Scalar-extended Chow Abel--Jacobi injectivity in
(AR16.2), and the alternative rational decomposition (AR16.4), remain
unproved.

\subsubsection{4.1 The relevant character really is a cubic power, at
the excluded central
boundary}\label{the-relevant-character-really-is-a-cubic-power-at-the-excluded-central-boundary}

Write \(K_0=\mathbf Q(i)\), \(\mathfrak q=(1+i)\), and use BDP's
characters \(\psi,\psi'\) of \(E:y^2=x^3-x\), \(E':y^2=x^3+x\). BDP
§5.2, author pp.24--25, gives \(\psi'=\psi\mu\), with \(\mu^2=1\) and
conductor \(\mathfrak q^4\). Therefore our motive's character satisfies
the exact identity \[
\chi=\psi^2\psi'=\psi^3\mu=(\psi')^3.
\tag{AR17.1}
\] Its conductor is \(\mathfrak q^4\), its theta series has weight four
and level 64, and AR-11 proves its simple central zero at \(s=2\).

Jonathan Dee, \emph{Selmer groups of Hecke characters and Chow groups of
self products of CM elliptic curves}, arXiv:math/9901155, Theorem 5.1.2,
p.28, treats representations \[
V_{j,k}=V_1(E')^{\otimes(j+k)}(-j)
\quad\text{with}\quad j\ge0,\quad k>j+1.
\] The desired central representation is \(V_{1,2}\): the rational Tate
module is \(H^1(E')(1)\), so its cube twisted by \(-1\) has the required
twist \((2)\). But \(2>1+1\) is false. The theorem's noncentral
critical-value range therefore excludes our case exactly at its
boundary. Theorem 1.0.1's nonvanishing at \(s=k\) is likewise not a
theorem about the cubic central value at \(s=2\).

Dee's Theorems 1.0.2 and 5.1.5, pp.3 and 29, prove a statement about \[
\operatorname{Im}(\mathrm{cl})\cap
F^2H^{2i}(X,\mathbf Q_p(i)).
\] They do not prove injectivity of \(\mathrm{cl}\). There is a
particularly direct check for our summand: the projector of \(N\) kills
every geometric cohomology group except \(H^3\). In the projected
Hochschild--Serre spectral sequence, the only term contributing to total
degree four is \(H^1(F,H^3(N)(2))\). Hence \[
F^2H^4(N,\mathbf Q_p(2))=0
\] already. A theorem saying that the image has zero intersection with
this zero subspace imposes no condition on the kernel in (AR16.1). This
is the precise reason the title's reference to Chow groups does not
solve the present Chow-rank problem.

\subsubsection{4.2 Castella's standing conductor hypothesis and a
separate proof-level
route}\label{castellas-standing-conductor-hypothesis-and-a-separate-proof-level-route}

Castella, \emph{Tamagawa number conjecture for CM modular forms and
Rankin--Selberg convolutions}, author manuscript dated 14 February 2025,
starts §5, printed p.21, with the standing hypothesis \[
(\mathrm{div}):\qquad p\nmid\mathfrak c,\qquad
\mathfrak d_{K_0}\Vert\mathfrak c.
\] Although Theorem 5.2.1, pp.23--24, does not repeat its second clause
in the displayed statement, that omission does not remove the section's
standing restriction. The theorem gives, within its stated scope, \[
\mathcal L_v(\lambda\mathbf N^{-r})(0)\ne0
\quad\Longrightarrow\quad
\dim_\Phi H^1_f(K_0,V_\lambda(r))=1
\tag{AR17.2}
\] for the specified self-dual character of infinity type \((1-2r,0)\).
Theorem C, p.4, and Theorem 5.3.1, p.24, likewise impose exact
divisibility by the different.

For the Fermat character take \(r=2\), \(p=5\), \(K_0=\mathbf Q(i)\),
and the same ideal-valued character \(\lambda=\chi\), written in
Castella's opposite infinity-type convention. The prime splits, is
greater than three, is prime to the class number and conductor, and the
central character is the required quadratic character. However, \[
\mathfrak d_{K_0}=\mathfrak q^2,\qquad
\mathfrak c_\chi=\mathfrak q^4,
\] so the standing exact-conductor condition fails. We therefore do
\textbf{not} cite Theorem 5.2.1 as a theorem that directly covers this
character.

There is a more precise proof-level route. Theorem 4.2.1, printed p.16,
constructs twisted elliptic units and their Coleman-map reciprocity law
for any self-dual character with the specified prime-to-\(p\) conductor;
Theorem 4.3.1, p.17, also permits arbitrary prime-to-\(p\) conductor.
The proof of Theorem 5.2.1 uses that reciprocity law, Lemma 3.4.2's
local conditions, central vanishing, an elliptic-unit Euler-system
bound, and global duality. The independent deduction H9-T15 now checks
those ingredients separately and proves the rank-one implication at the
actual conductor \(\mathfrak q^4\). It uses Rubin's full twisted
elliptic-unit Euler system and strict-dual Selmer bound, with the bad
prime above 2 retained in the excluded-prime ideal, followed by the
explicit Poitou--Tate dimension formula. This is a reconstructed proof,
not a claim that §5 silently withdrew its standing hypothesis.

The analytic premise required by that route was the single \textbf{Katz}
value \[
\mathcal L_v(\chi\mathbf N^{-2})(0)\ne0.
\tag{AR-KATZ}
\] AR-19 now proves (AR-KATZ) directly, by an integral Katz-measure
congruence at 5 and an exact fourth-order theta coefficient. It does not
use the complex derivative to infer p-adic nonvanishing. AR-11's
positive \textbf{complex} derivative is not this premise. Remark 5.2.2,
p.23, explains why the known weight-two implication from a simple
complex zero does not yet carry over to higher weights: it requires both
an Abel--Jacobi injectivity statement and nonzero localization. Remark
1.2.1, p.4, retains the higher-weight implication as an expectation.
Remark 1.2.2 and its footnote 2, p.4, retain injectivity of the
\(p\)-adic Abel--Jacobi map and nondegeneracy of the height as
additional inputs to the proposed full motivic consequence.

\subsubsection{4.3 The BDP calculation does not yet identify the
required Katz value at the ramified
prime}\label{the-bdp-calculation-does-not-yet-identify-the-required-katz-value-at-the-ramified-prime}

BDP, \emph{p-adic L-functions and the coniveau filtration on Chow
groups}, Proposition 5.2.1, author pp.26--27, proves a nonzero local
Abel--Jacobi evaluation for a generalized Heegner-type cycle associated
with the Fermat example; the displayed congruence (5.2.3) can be checked
with \(k=2,p=5\). This is considerably more arithmetic information than
ordinary Betti nonvanishing.

However, their Rankin--Katz factorization Proposition 3.3.1, p.13, is
stated with the conductor of the theta character cyclic and of norm
prime to the discriminant. In the Fermat construction the theta
character is \(\psi^2\), of conductor \(\mathfrak q^2\), at the ramified
prime 2. Their Proposition 3.2.1 also assumes odd discriminant and an
odd auxiliary conductor. Section 5.2, p.27, explicitly notes that its
general exact local formula does not cover this
even-discriminant/even-conductor situation, and leaves the local
2-factor for that comparison uncomputed.

Consequently we have not equated the nonzero BDP evaluation with
(AR-KATZ), and that exact ramified local comparison remains open. AR-19
bypasses it: its direct integral-measure calculation proves the Katz
value is nonzero without the Rankin factorization. Combined with H9-T15
this establishes the stated Selmer-rank result. It does not establish
scalar-extended injectivity in (AR16.1).

\subsubsection{4.4 The generalized Rubin formula retains both the
nonvanishing premise and incompatible
hypotheses}\label{the-generalized-rubin-formula-retains-both-the-nonvanishing-premise-and-incompatible-hypotheses}

Longo--Vigni--Wang, \emph{A generalized Rubin formula for Hecke
characters}, arXiv:2501.03673v1, Assumption 1.1 and Theorem A, p.2,
requires odd imaginary-quadratic discriminant, exact divisibility of the
conductor by the different, and a nonzero Katz value. It gives a
rank-one Selmer space and a formula for its logarithm. Both the
odd-discriminant and exact-conductor hypotheses fail for the present
\(\mathbf Q(i)\) character. The introduction, pp.3--4, also explains the
additional conjectural steps needed to replace the Katz premise by a
nonzero complex derivative. This theorem does not establish (AR-KATZ),
(AR16.2), or a real-height factorization for the Fermat Chow group.

\subsubsection{4.5 A proved p-adic height factorization cannot be
substituted for the real
one}\label{a-proved-p-adic-height-factorization-cannot-be-substituted-for-the-real-one}

Nekovář, \emph{p-adic Abel-Jacobi maps and p-adic heights}, Corollary
4.3 and Remark 4.4(2), printed p.7, distinguish these statements
explicitly. His constructed \(p\)-adic heights factor through \(p\)-adic
Abel--Jacobi maps, with choices of a \(p\)-adic logarithm and
Hodge-filtration splittings. The remark separately asks whether the real
height factors through the complex Abel--Jacobi map in higher
codimension. No result used here identifies these two height pairings or
transfers a radical statement from one to the other. AR-15 obtains its
restricted real radical by a separate curve-correspondence proof.

The fact that the complex intermediate Jacobian of our rank-two Hodge
structure is a one-dimensional complex torus gives no remedy. The Hodge
structure \(H(N)(1)\) has types \((2,-1),(-1,2)\), whereas \(H^1\) of an
elliptic curve has types \((1,0),(0,1)\). A correspondence inducing an
isomorphism between them is excluded by these types. Even a
one-dimensional complex torus has point groups of unbounded rational
rank; its complex dimension does not bound a Chow group's rational
dimension. It supplies neither a number-field Abel--Jacobi morphism to
an elliptic curve nor an arithmetic Néron--Tate comparison for \(V_F\).

\subsubsection{4.6 Infinite-rank results require their own
motive-transfer
argument}\label{infinite-rank-results-require-their-own-motive-transfer-argument}

Lilienfeldt, \emph{Heegner cycles in Griffiths groups of Kuga--Sato
varieties}, arXiv:2107.06731v3, Theorem 1.4, p.4, proves infinite rank
over \(K^{\rm ab}\) for the subgroup generated by Heegner cycles on the
stated Kuga--Sato variety. Its Remark 1.6, on the same page, explicitly
labels injectivity of the étale Abel--Jacobi map for number-field cycles
as conjectural. Remark 4.4, p.13, gives the formula on a fixed newform's
intermediate-Jacobian component.

The theorem as stated does not assert infinite rank for our particular
\(N\). Passing from an infinite-rank subgroup on the entire Kuga--Sato
variety to a specified \(f\)-component requires checking that projection
in the proof, and then transporting it to \(N\) requires an actual
algebraic correspondence. Equality of compatible Galois realizations or
of \(L\)-functions alone is insufficient. The distinction is visible in
Longo--Vigni--Wang §2.4 and Proposition 2.7, pp.6 and 12: their motive
comparisons there are in the category of absolute Hodge cycles. The
existence of separate Chow realizations does not, by itself, turn the
comparison into an isomorphism of Chow motives. We therefore assert
neither a finite-rank generation theorem nor an infinite-rank
counterexample for \(V_{\overline{\mathbf Q}}\) from these sources.

\subsection{5. Current exact remaining
requirements}\label{current-exact-remaining-requirements}

The established result is strict positivity on the orthogonal Picard
block together with the positive \(\mathbf Q(i)\)-line generated by
\(\eta\). The full Fermat complement has now been reduced further to its
Griffiths quotient, with its algebraically trivial subgroup proved to
lie in the height radical.

AR-19 and H9-T15 now give the character Selmer rank one over
\(K_0=\mathbf Q(i)\), hence dimension two for the full rank-two
representation there. Over this specified field, the remaining
hypothesis of the finite-rank criterion is scalar-extended injectivity
(AR16.1). For a further field extension its Selmer dimension must be
established separately; no uniform rank bound is asserted. Thus proving
the analytic input has not identified a Selmer dimension with a Chow
dimension. A different route could prove the real-height inequality
directly on the entire Griffiths quotient, together with
\(\mathcal A_F=0\). A rational generation theorem modulo a proved
real-height radical would yield semipositivity as in (AR16.4), but
strict positivity would still require that radical to vanish.

No negative-height class or nonzero zero-height class in \(V_F\) has
been established here. No general higher Abel--Jacobi kernel has been
declared zero or placed in the real-height radical. These are the
precise limits of the deductions, rather than counterexamples to the
expected arithmetic positivity theorem.

\subsection{Primary locators}\label{primary-locators}

\begin{itemize}
\tightlist
\item
  Castella: \url{https://web.math.ucsb.edu/~castella/Katz.pdf}
  ; local \path{CastellaKatz.pdf/.txt}.
\item
  Dee: \url{https://arxiv.org/pdf/math/9901155} ; local
  \path{Dee1999.pdf/.txt}.
\item
  BDP:
  \url{https://sites.lsa.umich.edu/kartikp/wp-content/uploads/sites/1326/2024/07/BDP5-hl.pdf}
  ; local \path{BDP2017.pdf/.txt}.
\item
  Longo--Vigni--Wang: \url{https://arxiv.org/pdf/2501.03673} ; local
  \newline\path{LongoVigniWang2025.pdf} and its extracted text.
\item
  Nekovář:
  \url{https://webusers.imj-prg.fr/~jan.nekovar/pu/banff.pdf} ;
  primary indexed text read directly, §§2.1--2.5, 3.1--3.2, and
  4.1--4.4. Direct PDF transfer returned HTTP 403; no local PDF is
  claimed.
\item
  Lilienfeldt: \url{https://arxiv.org/pdf/2107.06731} ; local
  \path{LilienfeldtHeegner2024.pdf/.txt} (v3, 5 August 2022; published
  in \emph{Mathematical Research Letters} 31(2), 2024).
\item
  Achter--Casalaina-Martin--Vial: the primary Astérisque 455 (2025)
  volume, local \path{work/references/ACMV2025.pdf/.txt}, §3.2.4,
  equation (3.5), pp.17--18.
\end{itemize}
\section{AR-19. A direct Katz nonvanishing calculation at
5}\label{ar-19.-a-direct-katz-nonvanishing-calculation-at-5}

Run \texttt{JAC1964-20261121}; \path{/root/arithmetic_sign}, GPT-6
Astra Pro; effort setting not exposed. This is a new deduction from the
explicitly cited Katz-measure construction. The complete proof and exact
script have passed independent whole-read checks by the root and h9
agents. It establishes a p-adic analytic input, not injectivity of a
Chow Abel--Jacobi map.

\subsection{Statement and conventions}\label{statement-and-conventions}

Put \(K=\mathbf Q(i)\), \(\mathfrak f=(4)\), and let \(\psi\) be the
ideal-valued Grossencharacter of \[
E:\quad y^2=x^3+x,\qquad \omega=\frac{dx}{2y}.
\] Thus \(\psi\) is the character denoted \(\psi'\) in BDP §5.2, its
conductor is \((4)\), and \(\psi\bar\psi=\mathbf N\). At the split
ordinary prime \(p=5\), define the two ideal-valued characters \[
\rho=\frac{\bar\psi^2}{\psi},\qquad
\rho_0=\frac{\psi^3}{\bar\psi^2}.
\tag{AR19.1}
\] We use the Bannai--Kobayashi (BK) convention: their Katz value
\(L_{p,\mathfrak f}(\rho_0)\) interpolates \(L(\rho_0^{-1},0)\), and
\(\rho_0\) has infinity type \((3,-2)\) in their convention. Castella's
infinity-type convention is opposite for the same ideal-valued
character, and his Katz function interpolates \(L(\xi,0)\). Thus the
relevant comparison is \(\xi=\rho^{-1}\). In particular \[
\rho^{-1}=\psi^3\mathbf N^{-2}=\chi\mathbf N^{-2},
\] where \(\chi\) is the weight-four Fermat character of AR-11.

\textbf{Proposition AR-19.} With compatible good-reduction CM periods,
the BK value \(L_{p,(4)}(\rho)\) is a 5-adic unit. Consequently the Katz
nonvanishing premise (AR-KATZ) holds: \[
\mathcal L_v(\chi\mathbf N^{-2})(0)\ne0.
\tag{AR19.2}
\] The statement concerns the Katz function itself and does not use the
ramified Rankin--Katz comparison factor in BDP §5.2.

\subsection{1. The exact fourth-order theta
coefficient}\label{the-exact-fourth-order-theta-coefficient}

Write \(E(\mathbf C)=\mathbf C/\Gamma\), \(\Gamma=4\Omega\mathbf Z[i]\),
using the differential \(\omega\), and put
\(A=\operatorname{area}(\mathbf C/\Gamma)/\pi\). The modified weight-two
Eisenstein constant vanishes for a square lattice: multiplication by
\(i\) fixes the lattice and multiplies that constant by \(i^{-2}=-1\).
Hence the reduced theta function is the Weierstrass sigma function with
derivative one at zero.

For a nonzero torsion point \(P=z_P\bmod\Gamma\), put \[
F(P)=\zeta(z_P)-\frac{\bar z_P}{A},\qquad
x=\wp(z_P),\qquad y=\frac{\wp'(z_P)}2.
\] The translated Kronecker function in BK equation (15), p.15, has the
expansion \[
\begin{aligned}
\Theta_{P,0}(z,w)&=\frac1w+F(z+P)
 +\frac{F(z+P)^2-\wp(z+P)}2w\\
&\quad +\frac{F(z+P)^3-3F(z+P)\wp(z+P)-\wp'(z+P)}6w^2+O(w^3).
\end{aligned}
\] where in this formula \(F(z+P)=\zeta(z+z_P)-\bar z_P/A\) is
holomorphic in \(z\); the conjugate of the variable \(z\) is not
subtracted. Its derivatives satisfy \[
F'=-x,\qquad x'=2y,\qquad y'=3x^2+1.
\] By BK Theorem 1.17, p.17, comparison of the coefficient of \(z^2w^2\)
gives \[
\mathcal E(P):=\frac{e^*_{2,3}(P,0)}{A^2}
 =-F(P)(2x^2+1)-y\bigl(F(P)^2+x\bigr).
\tag{AR19.3}
\] Indeed applying the derivation \(D(F)=-x\), \(D(x)=2y\),
\(D(y)=3x^2+1\) twice to \((F^3-3Fx-2y)/6\) gives the right side. The
two factorials give
\(\partial_z^2\partial_w^2\Theta_{P,0}(0,0)=2\mathcal E(P)\).

Choose a primitive \(4\)-torsion generator with \(2P=(i,0)\). Its two
conjugates over \(K\) can be written \[
P_\pm=\bigl(i(1\pm\sqrt2),\ (1\pm\sqrt2)(1-i)\bigr).
\tag{AR19.4}
\] The duplication formula verifies \(2P_\pm=(i,0)\); adding
\((-i,0)=i(2P_+)\) to \(P_+\) gives \(P_-\). Thus \(P_-=[1+2i]P_+\). The
reduced zeta function vanishes at a nonzero half-period, so its
duplication formula gives \[
F(P_\pm)=-\frac{3x(P_\pm)^2+1}{4y(P_\pm)}.
\] Substitute (AR19.4) into (AR19.3). Using only \(i^2=-1\) and
\((\sqrt2)^2=2\), one obtains the exact identities \[
\mathcal E(P_+)=\mathcal E(P_-)=-\frac{1+i}{4},
\qquad
\mathcal E(P_+)+\mathcal E(P_-)=-\frac{1+i}{2}.
\tag{AR19.5}
\] Changing a lattice generator by a Gaussian unit multiplies this sum
by a Gaussian unit and cannot affect the ensuing nonvanishing. Among the
two points in (AR19.4), at least one is represented by a lattice
generator divided by four: the primitive classes of
\(\mathbf Z[i]/(4)\), modulo global Gaussian units, are precisely its
two ray classes. Both points give the same coefficient, so that choice
is immaterial.

\subsection{2. The two ray classes and their measure
normalizations}\label{the-two-ray-classes-and-their-measure-normalizations}

The ray class group modulo \((4)\) has order two:
\(|(\mathbf Z[i]/4)^\times|=8\) and the four global units inject. Use
inverse ideal representatives \(1\) and \(\mathfrak a^{-1}\), with \[
\mathfrak a=(\beta),\qquad \beta=1+6i,\qquad N\mathfrak a=37.
\] Both are prime to \(5\mathfrak f\), and \(\beta\equiv1+2i\pmod4\).
For this curve \(\psi(\mathfrak a)=\beta\). This also follows from the
usual primary-generator formula \(\psi=\psi_{y^2=x^3-x}\mu\): \(-\beta\)
is primary and the quadratic character of \(K(\sqrt2)/K\) has value
\(-1\) at norm 37. The resulting isogeny scalar in BK §3.3 is \[
\Lambda(\mathfrak a)=\psi(\mathfrak a)=\beta.
\] The ray class field is \(K(\sqrt2)\). We may take the same equation
of \(E\) on that field; all its conjugates have the same differential
and period lattice. BK §3.3, pp.40--42, then gives \[
\Omega_p^{\sigma_{\mathfrak a}}
 =\frac{\Lambda(\mathfrak a)}{N\mathfrak a}\Omega_p,
\qquad \Omega_p\in W(\overline{\mathbf F}_5)^\times.
\tag{AR19.6}
\] The choices \(\alpha_0=1\in\mathfrak a^{-1}\cap P_K(\mathfrak f)\) in
their Definition 3.9 give translated points \(\Omega\) and
\(\Lambda(\mathfrak a)\Omega\). These are \(P_+,P_-\), up to the same
harmless choice of Gaussian unit.

For any such representative \(\mathfrak b^{-1}\), BK Definitions
3.8--3.9, equations (38)--(39), and the displayed formula following
Definition 3.9, pp.41--42, yield at \(\rho_0\) (with \(k=3,l=2\)) \[
L_{p,\mathfrak f}(\rho_0,\mathfrak b^{-1})
 =\Omega_p\rho_0(\mathfrak b)^{-1}(N\mathfrak b)^2(N\mathfrak f)^2
 \int_{\mathbf Z_5^\times\times\mathbf Z_5^\times}
       x^2y^2\,d\mu_{\Omega,0}(x,y;\mathfrak b^{-1}\Gamma).
\tag{AR19.7}
\] These measures are integral; this is Corollary 2.17, p.30, and
Definition 3.2, p.36, applied to prime-to-5 torsion on the good model.
By the lattice-scaling definition of the measure on p.41 and (AR19.6),
its moment over all of \(\mathbf Z_5^2\) is \[
\int_{\mathbf Z_5^2}x^2y^2\,d\mu_{\Omega,0}(x,y;\mathfrak b^{-1}\Gamma)
 =2\Lambda(\mathfrak b)
       (\Omega_p^{\sigma_{\mathfrak b}})^4
       \mathcal E\bigl(\Lambda(\mathfrak b)\Omega\bigr).
\tag{AR19.8}
\] Although the second torsion argument is zero, BK's regularization
subtracts a term depending only on the second formal variable. Two
derivatives in the first variable kill it, so it makes no contribution
to this moment.

The multiplier of \(2\Omega_p^5(N\mathfrak f)^2\mathcal E\) after
substituting (AR19.8) in (AR19.7) is exactly \[
\frac{\Lambda(\mathfrak b)^5}
 {\rho_0(\mathfrak b)(N\mathfrak b)^2}=1,
\tag{AR19.9}
\] since \(\Lambda=\psi\),
\(N\mathfrak b=\psi(\mathfrak b)\bar\psi(\mathfrak b)\), and
\(\rho_0=\psi^3/\bar\psi^2\).

\subsection{3. The unit congruence and the noncritical
value}\label{the-unit-congruence-and-the-noncritical-value}

On the complement of \(\mathbf Z_5^\times\times\mathbf Z_5^\times\) in
\(\mathbf Z_5^2\), the function \(x^2y^2\) is divisible by \(25\).
Integrality of the measure therefore permits replacing the unit-domain
moment in (AR19.7) by (AR19.8) modulo 5. Every exterior factor in
(AR19.7) is a unit. Since \(N\mathfrak f=16\), summing the two classes
and using (AR19.5) gives \[
L_{p,(4)}(\rho_0)
 \equiv2\Omega_p^5\,16^2
          \left(-\frac{1+i}{2}\right)
 =-256(1+i)\Omega_p^5\pmod5.
\tag{AR19.10}
\] This is a unit at either prime over 5: \(1+i\) has norm two, and
\(\Omega_p\) is a unit.

On the Galois group of the ray-class tower of conductor
\(4\cdot5^\infty\), both \(\psi\) and \(\bar\psi\), interpreted as
5-adic characters in the BK convention, take values in
\(\mathbf Z_5^\times\). Therefore \[
\frac{\rho_0}{\rho}
 =\left(\frac{\psi}{\bar\psi}\right)^4
 \equiv1\pmod5.
\tag{AR19.11}
\] Integrality of the Katz measure gives \[
L_{p,(4)}(\rho)\equiv L_{p,(4)}(\rho_0)\not\equiv0\pmod5.
\tag{AR19.12}
\] This proves the BK unit assertion without a numerical approximation.

Finally choose the \textbf{same compatible complex and 5-adic period
data} in BK and Castella, attached to the good equation and differential
fixed above. Compare BK Theorem 3.11, pp.44--46, with Castella Theorem
2.2.1, pp.7--8. Under \(\xi=\rho^{-1}\), their interpolation values
involve the same Hecke value, gamma factor, and the two Euler factors at
5. BK writes one factor explicitly and the other in
\(L_{\mathfrak f\bar{\mathfrak p}}\). In particular \(\rho_0^{-1}\) has
Castella infinity type \((3,-2)\), so that the displayed formulas match
literally with this period choice. More generally they match on the
critical characters in the component with \(k\equiv3\pmod4\),
\(l\equiv2\pmod4\); the BK sign \((-1)^{k+l-1}\) is then +1. These
critical characters are dense in that two-variable component. The
uniqueness of the integral Katz measure therefore identifies the two
functions there with this inverse-character convention. Both \(\rho\)
and \(\rho_0\) lie in the component by (AR19.11). This proves (AR19.2)
from (AR19.12). Changing the good integral period data multiplies the
evaluation at this fixed algebraic weight by a nonzero unit, so
nonvanishing is independent of that permitted choice; no assertion about
an arbitrary weight-dependent global analytic rescaling is needed.

\subsection{Scope}\label{scope}

AR-19 supplies the previously missing nonzero Katz value for the Fermat
character at 5. Together with a separately checked Euler-system and
duality argument at the actual conductor, it can supply a
Selmer-dimension bound. It does not give scalar-extended étale
Abel--Jacobi injectivity, does not identify a Selmer dimension with a
Chow dimension, and does not prove real-height factorization through a
p-adic Abel--Jacobi map.

\subsection{Primary locators and computational
check}\label{primary-locators-and-computational-check}

\begin{itemize}
\tightlist
\item
  Bannai--Kobayashi, \emph{Algebraic theta functions and the p-adic
  interpolation of Eisenstein--Kronecker numbers}, arXiv:math/0610163v4
  (11 December 2007), equation (15) p.15; Theorem 1.17 p.17; Corollary
  2.17 p.30; Proposition 3.1 and Definition 3.2 pp.35--36; §3.3
  pp.40--43, Definitions 3.8--3.9; Theorem 3.11 pp.44--46. Primary PDF
  directly read, including the period and Euler-factor formulas. Local
  source \path{work/h9/BK2007.pdf/.txt}.
\item
  Castella, \emph{Tamagawa number conjecture for CM modular forms and
  Rankin--Selberg convolutions}, 14 February 2025 manuscript, §2.1.1
  pp.5--6 (infinity types and reciprocity), §2.1.2 p.6 (periods),
  Theorem 2.2.1 pp.7--8 (the integral Katz function and uniqueness).
  This use is outside §5 and does not invoke its exact-conductor
  hypothesis.
\item
  BDP, \emph{p-adic L-functions and the coniveau filtration on Chow
  groups}, §5.2, for the character of \(y^2=x^3+x\), its conductor and
  the quadratic twist relation.
\item
  The rational algebra in (AR19.3)--(AR19.5) is also checked by the
  reproducible exact script \texttt{fermat\_katz\_coefficient.py}; no
  floating-point arithmetic is used.
\end{itemize}
% Proof-level conductor-scope repair for the next arithmetic gate.
% Run JAC1964-20261121; /root/h9_height; GPT-6 Astra Pro.
% Coordinate with AR19 (Katz value); this is H9-T15, not a Chow-rank theorem.

\subsection*{H9-T15. The conductor-general rank-one implication for the Fermat character}

\paragraph{Why a separate deduction is needed.}
Castella's \emph{Tamagawa number conjecture for CM modular forms and
Rankin--Selberg convolutions}, manuscript of 14 February 2025,
introduces in Section 5, printed p.~21, the standing hypothesis
$\mathfrak d_K\Vert\mathfrak c$. The displayed Theorem 5.2.1,
printed p.~23, does not repeat that hypothesis. It is therefore
inappropriate to disregard it merely by reading the theorem display.
For the Fermat character its conductor is $(1+i)^4$, while the
different of $\mathbf Q(i)$ is $(1+i)^2$; exact divisibility fails.
The argument below instead uses the arbitrary-conductor reciprocity
theorem of Section 4 and the Euler-system bound. It records exactly
the portion of the proof of Theorem 5.2.1 that does not use the
Section 5 good-pair or conductor hypothesis.

\begin{proposition}[H9-T15]
Let $K_0=\mathbf Q(i)$, let $p=5$, and let
$\chi=\psi^3\mu=(\psi')^3$ be the ideal-valued Hecke character of
the Fermat rank-two motive, in the notation of AR-11 and BDP
Section 5.2. Write $\lambda$ for this same character in Castella's
infinity-type convention, so that it has type $(-3,0)$ and central
character $\eta_{K_0}$. Put
\[
 \varphi=\lambda\mathbf N^{-2},\qquad
 V=V_\lambda(2).
\]
Thus $\varphi$ has type $(-1,2)$ and is self-dual in Castella's
Definition 2.2.2. Fix one embedding $K_0\hookrightarrow\mathbf Q_5$
and use the associated prime $v$ and its conjugate $\bar v$ as in
that source. Let $\Phi/\mathbf Q_5$ contain the character values
and let $T\subset V$ be a stable rank-one lattice over its integers
$\mathcal O$. One can take $\Phi=\mathbf Q_5$ in this example.

If
\[
                  \mathcal L_v(\varphi)(0)\ne0,              \tag{H9.56}
\]
then
\[
 \dim_\Phi\operatorname{Sel}_{\rm BK}(K_0,V)=1.             \tag{H9.57}
\]
This implication requires no exact-divisibility hypothesis at $2$.
It is a statement about this specified Selmer group; it does not
identify its dimension with a rational Chow rank.
\end{proposition}

\begin{proof}
We use the conventions of Castella throughout. In particular the
cyclotomic character has Hodge--Tate weight $+1$, and the two
Hodge--Tate weights of $V$ are $-1$ at $v$ and $2$ at $\bar v$.
Self-duality identifies its Tate dual $V^*(1)$ with the conjugate
rank-one representation $V^\tau$. Neither representation is
trivial globally or at either prime over $5$. Their local
Hodge--Tate weights are nonzero, which also proves
\[
 H^0(K_0,V)=H^0(K_0,V^*(1))=0,
 \qquad H^0((K_0)_w,V^*(1))=0\quad(w\mid5).                \tag{H9.58}
\]

\emph{The nonzero class and its local conditions.}
Castella's Theorem 4.2.1, printed p.~16, is stated for an arbitrary
self-dual Hecke character. Its elliptic-unit construction starts
with any prime-to-$p$ conductor ideal; it has no requirement that
the conductor have exact valuation equal to the different at a
ramified prime. Apply this construction to $\varphi$ and its
conjugate. Let
\[
 c\in H^1_{\rm Iw}(K_\infty,T),\qquad
 c_0\in H^1(K_0,T)
\]
be the normalized twisted elliptic-unit class and its trivial
anticyclotomic specialization. The specialization of the Coleman
maps, which interpolate the Bloch--Kato logarithm and dual
exponential in Theorem 4.2.1, gives the two consequences used in
the first half of Castella's printed p.~24 proof:
\[
 \mathcal L_v(\varphi)(0)\ne0
       \ \Longrightarrow\ \operatorname{loc}_{\bar v}(c_0)\ne0,
                                                               \tag{H9.59}
\]
and
\[
 \mathcal L_v(\varphi^\tau)(0)=0
       \ \Longrightarrow\ \operatorname{loc}_{v}(c_0)=0.
                                                               \tag{H9.60}
\]
These are the reciprocity and specialization statements of the
Section 4 theorem, not an invocation of the Section 5 conductor
hypothesis. Conjugating the two primes gives
$\mathcal L_{\bar v}(\varphi^\tau)(0)=\mathcal L_v(\varphi)(0)$
up to the compatible nonzero period unit, explaining the prime
$\bar v$ in (H9.59). The local weights $-1,2$ avoid the exceptional
weight-zero and weight-one cases; after inverting $5$ the relevant
logarithm or dual exponential identifies the one-dimensional local
space in question with its de Rham target.

The conjugate character has type $(2,-1)$, in the interpolation
range of Castella Theorem 2.2.1, printed p.~7. The central root
number of $\lambda$ is $-1$ by the Fermat character computation
in AR-11, so $L(\lambda,2)=L(\lambda^\tau,2)=0$.
Interpolation therefore gives
$\mathcal L_v(\varphi^\tau)(0)=0$, and (H9.60) applies.
Twisted elliptic units are unramified in the required rational
local condition away from $5$. At $5$, Lemma 3.4.2, printed p.~12,
identifies the Bloch--Kato local space with $0$ at $v$ and with
all of $H^1((K_0)_{\bar v},V)$ at $\bar v$. Consequently
\[
       0\ne c_0\otimes1
             \in\operatorname{Sel}_{\rm BK}(K_0,V).        \tag{H9.61}
\]

\emph{The Euler-system bound with the ramified prime retained.}
For clarity, the Euler system is taken over the ray-class tower
containing $K_0(\mathfrak c5^\infty)$ and all auxiliary ray-class
extensions at primes away from $6\mathfrak c5$, where
$\mathfrak c=(1+i)^4$. One starts with the elliptic units for
$\mathcal O(1)$ and twists by the rank-one character
$T\otimes\mathcal O(-1)$. This character factors through a finite
extension of the compositum of the $\mathbf Z_5$-extensions of
$K_0$, inside the displayed ray-class tower.

Rubin's twisting theorem constructs the full twisted Euler system
whose normalized base specialization is $c_0$: in the directly read author draft
\emph{Euler systems}, Chapter VI, Definition 3.1 and Theorem 3.5,
printed pp.~92--93, the new non-$p$ conductor is simply adjoined
to the excluded-prime ideal. Thus the prime above $2$ is among
the excluded bad primes. No constraint on its conductor exponent
is used. This is the same twisting input cited in Castella's
printed p.~24 proof as the published Rubin Theorem 6.4.1.
The harmless normalization of the elliptic units is the invertible
auxiliary-ideal factor in Castella Theorem 4.2.1: its proof explicitly
chooses $\mathfrak a$ prime to $6\mathfrak cp$ such that
$\sigma_{\mathfrak a}-\varphi^\tau(\sigma_{\mathfrak a})$
is invertible in the anticyclotomic Iwasawa algebra. More precisely,
apply Rubin to the \emph{unnormalized} full twisted elliptic-unit
Euler system. Its base class is the nonzero unit multiple of $c_0$
obtained by specializing that auxiliary-ideal factor. Thus one need
not invert the anticyclotomic normalization simultaneously on all
auxiliary ray extensions; its base specialization already supplies
the nonzero Euler-system class required below.

Rubin's Euler-system bound now gives
\[
           \operatorname{Sel}_{\rm str}(K_0,V^*(1))=0.
                                                               \tag{H9.62}
\]
The exact primary statement is Chapter II, Theorem 2.3, printed
p.~24 of that draft (published Theorem 2.2.3): a nonzero
Euler-system class for a nontrivial representation satisfying
$\operatorname{Hyp}(K_0,V)$ makes the strict dual torsion Selmer
group finite, and hence its rational space zero. Its hypotheses
hold for this rank-one representation: choose $\tau=1$ in the
rank-one-quotient condition, and irreducibility is automatic.
The same page's Remark 2.4 explicitly records this rank-one case.
The nontriviality condition follows also from (H9.58).

\emph{The dimension calculation.}
Let $\operatorname{Sel}_{\rm rel}(K_0,V)$ impose the Bloch--Kato
unramified condition away from $5$ and impose no condition at the
two primes over $5$. Poitou--Tate duality and the global Euler
characteristic give, using (H9.58),
\[
 \dim_\Phi\operatorname{Sel}_{\rm rel}(K_0,V)
    =d^-(V)+
       \dim_\Phi\operatorname{Sel}_{\rm str}(K_0,V^*(1)).    \tag{H9.63}
\]
Here $d^-(V)=1$, because $K_0$ has one complex place and $V$ has
rank one. A complete displayed derivation of (H9.63) is
Loeffler--Zerbes, \emph{Euler systems with local conditions},
arXiv:1710.04956v2, Proposition 2.8 and its proof, printed p.~3.
The only suppressed terms in (H9.63) are the global and local
$H^0$ terms explicitly checked in (H9.58).

By (H9.62), the relaxed space in (H9.63) has dimension one.
Its Bloch--Kato subspace contains the nonzero vector (H9.61).
It is therefore one-dimensional, which proves (H9.57).
\end{proof}

\paragraph{Scope of this check.}
This is a proof-level application of the arbitrary-conductor
reciprocity construction and rank-one Euler-system machinery.
It repairs the source-scope issue caused by the standing
Section 5 assumption; it does not claim that assumption was
explicitly withdrawn in the source. The hypothesis (H9.56)
remains a separate Katz-value calculation, assigned to AR19.
Neither a nonzero complex derivative alone nor the BDP
Rankin--Katz formula outside its stated ramification hypotheses
is used to prove it here. No scalar-extended Abel--Jacobi
injectivity, rational Chow-generation theorem, or positivity on
the full arithmetic complementary Chow group follows just from
(H9.57).

\paragraph{Primary sources directly inspected.}
Castella: \url{https://web.math.ucsb.edu/~castella/Katz.pdf},
Section 4.2 and printed p.~24 (the latter also visually inspected).
Rubin author draft:
\url{https://swc-math.github.io/aws/1999/99RubinES.pdf},
Chapter II pp.~23--24 and Chapter VI pp.~92--94.
Loeffler--Zerbes:
\url{https://discovery.ucl.ac.uk/id/eprint/10067153/1/1710.04956v2.pdf},
Theorem 2.5, Proposition 2.8 and its proof, printed pp.~2--3.
The general Euler-system theorems are used as established inputs;
they have not been reproved or formally verified in this run.
\part{Unresolved assertions and exploratory work}
\section{What remains to be proved}
The full primitive arithmetic height form on the complementary CM Chow
motive remains to be controlled. A formula for one special cycle does
not prove positivity for every cycle. Analytic derivative signs require
an applicable real self-dual $L$-function and a proved height identity
with its local terms and positive constants. The Bost slope inequality
requires an actual canonical height line and a proved norm estimate.
Neither missing input is silently supplied by purity or Dirichlet GRH.
For the specified Fermat class the height sign is now proved directly,
independently of any $L'$-height identity. The remaining Fermat gate is
positivity on its entire Griffiths quotient, together with vanishing
of its algebraically trivial subgroup if strict positivity on rational
Chow is required. The positive CM plane is known, but a spanning
theorem for the whole group has not been proved. The full abelian
endomorphism orbit cannot enlarge that plane.

The geometric H9 assertions cannot be completed in their original form:
the counterexamples above refute the general formula and the asserted
incidence domain. For abelian ambient varieties condition A is automatic.
The full geometric height problem reduces to an explicit primitive
intersection form on the axis. It closes for threefold fibres and for
geometric generic axes in every dimension covered by JAC-P09. These
statements extend through purely inseparable field extensions. For
finite separable covers the old subspace splits orthogonally from the
new norm-zero subspace; positivity on arbitrary new nonalgebraic cycles
is a remaining gate. The algebraically trivial primitive domain is
controlled after every finite extension.
\begin{thebibliography}{99}
\bibitem{Saito1979} H. Saito, \emph{Abelian varieties attached to cycles of intermediate dimension}, Nagoya Math. J. 75 (1979), 95--119. Definition 3.1; Proposition 3.4; Theorem 3.6; \S4.1; Proposition 4.4; Corollary 4.5; Theorem 4.6; Proposition 5.2. \url{https://doi.org/10.1017/S0027763000024855}.
\bibitem{Kleiman68} S. L. Kleiman, \emph{Algebraic cycles and the Weil conjectures}, in Dix expos\'es sur la cohomologie des sch\'emas (1968), 359--386, Theorem 2A11.
\bibitem{MilneAV} J. S. Milne, \emph{Abelian Varieties}, course notes, v2.00 (2008), Chapter I, Theorem 11.1; Chapter III, Theorem 6.8 and Summary 6.11. \url{https://www.jmilne.org/math/CourseNotes/AV.pdf}.
\bibitem{MumfordAV} D. Mumford, \emph{Abelian Varieties}, Oxford University Press, 1970. Chapter II, \S6, Application 1, p.~60; Chapter III, \S16, Index Theorem, p.~155.
\bibitem{ConradAV} B. Conrad, \emph{Abelian varieties}, notes by T. Feng (2015), Example 4.2.14, p.~38.
\bibitem{SGA7I} A. Grothendieck et al., \emph{Groupes de monodromie en g\'eom\'etrie alg\'ebrique, I}, SGA 7 I, Lecture Notes in Mathematics 288, Springer, 1972. Expos\'e IX, \S9.2, 10.3.5 (p.~443), Theorem 10.4 (p.~444), \S\S12.5 and 12.7.
\bibitem{SGA7II} P. Deligne and N. Katz, eds., \emph{Groupes de monodromie en g\'eom\'etrie alg\'ebrique, II}, SGA 7 II, Lecture Notes in Mathematics 340, Springer, 1973. Expos\'e XV, 3.3.3--3.3.6, pp.~191--192; XVII, \S2.2; XVIII, 5.3.2--5.3.5, pp.~287--288, Proposition 6.2, p.~319, Corollary 6.4, 6.7.2; XX, Theorem 4.1 and Remark 4.3, pp.~355--356.
\bibitem{Bloch1984L} S. Bloch, \emph{Algebraic cycles and values of L-functions}, J. reine angew. Math. 350 (1984), 94--108. Lemma 3.1, p.~99; Theorem 4.1, p.~104; Proposition 5.2, p.~107.
\end{thebibliography}
\end{document}
